\documentclass[11pt,letterpaper]{article}
\usepackage[T1]{fontenc}
\usepackage[utf8]{inputenc}
\usepackage{lmodern,microtype}
\usepackage[margin=.94in]{geometry}
\usepackage{amsmath,amssymb,amsthm,mathtools,mathrsfs,booktabs,array,longtable,enumitem}
\usepackage{xcolor,xurl}
\usepackage[colorlinks=true,linkcolor=blue!45!black,urlcolor=blue!45!black,citecolor=blue!45!black]{hyperref}
\hypersetup{pdftitle={A factorial-window bound for the irrationality exponent of zeta(5)},pdfauthor={Zeta Irrationality Project},pdfsubject={Ordinary mathematical proofs and exact finite certificates}}
\usepackage{fancyhdr}
\pagestyle{fancy}\fancyhf{}\fancyhead[L]{\small Zeta Irrationality Project}
\fancyhead[R]{\small Irrationality exponent of $\zeta(5)$}\fancyfoot[C]{\thepage}
\setlength{\headheight}{14pt}\setlength{\emergencystretch}{3em}\setlength{\parskip}{.22em}
\newtheorem{theorem}{Theorem}[section]\newtheorem{lemma}[theorem]{Lemma}
\newtheorem{proposition}[theorem]{Proposition}\newtheorem{corollary}[theorem]{Corollary}
\theoremstyle{definition}\newtheorem{definition}[theorem]{Definition}
\theoremstyle{remark}\newtheorem{remark}[theorem]{Remark}
\newcommand{\R}{\mathbb R}\newcommand{\C}{\mathbb C}\newcommand{\Q}{\mathbb Q}\newcommand{\Z}{\mathbb Z}
\newcommand{\dd}{\,\mathrm d}\newcommand{\J}{\mathcal J}\newcommand{\K}{\mathcal K}
\newcommand{\Lfun}{\mathcal L}\newcommand{\Bfun}{\mathcal B}\newcommand{\pos}[1]{(#1)_+}
\DeclareMathOperator{\arsinh}{arsinh}\DeclareMathOperator{\nullity}{nullity}
\DeclareMathOperator{\diag}{diag}\DeclareMathOperator{\rank}{rank}
% Generated from matching PASS certificates; do not edit numbers by hand.
\newcommand{\CertClaim}{87.95428}
\newcommand{\CertInteger}{88}
\newcommand{\CertSOne}{\frac{1179301}{10000000}}
\newcommand{\CertSTwo}{\frac{1286161}{10000000}}
\newcommand{\CertDelta}{\frac{5559831}{5000000}}
\newcommand{\CertLambda}{0.8713839}
\newcommand{\CertGap}{0.076728643701677087}
\newcommand{\CertGapExact}{\frac{76728643701677087}{1000000000000000000}}
\newcommand{\CertArithmeticExact}{\frac{5083839596948510913}{1000000000000000000}}
\newcommand{\CertArithmetic}{5.083839596948510913}
\newcommand{\CertAnalytic}{5.160568240650188}
\newcommand{\CertU}{\frac{1020742754711}{1000000000000}}
\newcommand{\CertAngularB}{\frac{39330907}{12500000000}}
\newcommand{\CertAlphaZero}{\frac{6908942461}{25000000000}}
\newcommand{\CertYcut}{\frac{350411541}{10000000000}}
\newcommand{\CertTcap}{\frac{9907255237}{1000000000000}}
\newcommand{\CertGamma}{7.80770196117103596}
\newcommand{\CertKappa}{4.475942079831855724}
\newcommand{\CertLinear}{0.287014328898465886}
\newcommand{\CertBeta}{-0.696205633310614304}
\newcommand{\CertMassBeta}{0.743814928429757219}
\newcommand{\CertMassGain}{0.731107191265849172}
\newcommand{\CertSaving}{0.003500391586353087}
\newcommand{\CertAppliedSavingExact}{\frac{3500391586353087}{1000000000000000000}}
\newcommand{\CertPhysicalBeta}{0.034901557955234868}
\newcommand{\CertMassFloor}{1.10466631717049132}
\newcommand{\CertBMinLower}{0.9366045667971780542655}
\newcommand{\CertMomentUpper}{0.0000068460160147188132}
\newcommand{\CertTLower}{0.0099072352312867641982}
\newcommand{\CertTUpper}{0.0099072352312867642239}
\newcommand{\CertILower}{0.0767286289451411252289}
\newcommand{\CertIUpper}{0.076728628945141125432903}
\newcommand{\CertISlack}{0.000000014756535961567097}
\newcommand{\CertOuterSlack}{0.0000000023079336972064}
\newcommand{\CertMuDetailed}{87.954278950885823705}
\newcommand{\CertMuLower}{87.954278950885823704}
\newcommand{\CertMuExact}{\frac{223288451270279412930934211711}{2538687758385978822649900000}}
\newcommand{\CertMuSimple}{87.954295992508049175}
\newcommand{\CertTransferC}{\frac{114413952945899249}{10000000000000000}}
\newcommand{\CertBaseCutoff}{\frac{1}{65536}}
\newcommand{\CertCells}{988\,462}
\newcommand{\CertCorrectionCells}{294\,826}
\newcommand{\CertArithmeticQmax}{16384}
\newcommand{\CertMixedQmax}{192}
\newcommand{\CertThirdQmax}{128}
\newcommand{\CertOddQmax}{128}
\newcommand{\CertSaturatedQmax}{128}
\newcommand{\CertCentralSecondQmax}{128}
\newcommand{\CertCrossedQmax}{128}
\newcommand{\CertCoreQmax}{512}
\newcommand{\CertConfluentQmax}{512}
\newcommand{\CertConfluentWitnesses}{300}
\newcommand{\CertUpperOddQmax}{128}
\newcommand{\CertUpperOddWitnesses}{6}
\newcommand{\CertOuterUpperWitnesses}{2}
\newcommand{\CertPaidLowerWitnesses}{2}
\newcommand{\CertOuterLowerWitnesses}{1}
\newcommand{\CertExtendedPanels}{4096}
\newcommand{\CertPanels}{128}
\newcommand{\CertCharges}{256}
\newcommand{\CertThresholdLower}{7.6873735423713731544532}
\newcommand{\CertThresholdUpper}{7.6873735423713731547347}
\newcommand{\CertResponseUpper}{0.000568010950097722320503}
\newcommand{\CertNextInteger}{87}
\newcommand{\CertNextIntegerGap}{0.000835384927827329434619}
\newcommand{\CertDynamicSaving}{0.000000014717352034885597}
\newcommand{\CertDynamicTangents}{1025}
\newcommand{\CertDynamicH}{0.69959185667459498}
\newcommand{\CertDynamicGuard}{10.929461912937890118}
\newcommand{\CertPreviousResponseMuUpper}{87.954295947256385102}
\newcommand{\CertCommonRootCutoff}{94}
\title{Factorial-window determinants and an upper bound for the irrationality exponent of $\zeta(5)$}
\author{Zeta Irrationality Project}\date{Revised 18 November 2026}
\begin{document}\maketitle
\begin{center}\small Ordinary mathematical proofs and exact finite certificates.\\
The quantitative extension in this report has not been verified in Lean.\end{center}
\begin{abstract}
For the irrationality exponent of $\zeta(5)$ we obtain
$\mu(\zeta(5))<87.897631683470$. The proof uses two factorial windows, a Hankel pencil
whose finite roots have multiplicity at most two, and prime-by-prime
bounds on the denominators of its determinant. Comparing four nested local
moment matrices determines when local row and column restrictions share a
single evaluation. Nonzero minors and upper rank bounds certify the resulting
denominator savings. We estimate the response of that determinant to perturbations of $\zeta(5)$. We bound the outer-pole contour and the omitted real residues on the same polynomial subspace. Green functions control the prescribed zeros at both sets of points. Rational interval
certificates give the constants in the final transfer. The bound gives
integer exponent $88$ and improves the preceding project bound
$87.925687270635509$.
\end{abstract}

\section{Result, provenance, and verification boundary}

For an irrational real number $x$, write $\mu(x)$ for the supremum of exponents
$M$ for which $|x-p/q|<q^{-M}$ has infinitely many coprime integer solutions
with $q>0$. Smaller upper bounds are stronger.

\begin{theorem}\label{thm:main}
The two-window determinant with
\[
 (s_1,s_2,\delta)=\left(\frac{73577}{625000},\frac{12904077}{100000000},\frac{27801871}{25000000}\right)
\]
satisfies the estimates needed to give
\[
 \boxed{\mu(\zeta(5))\le \frac{4463734983998220580418722476616000}{50783336234503933606439686370429} < 87.897631683469558.}
\]
Consequently $|\zeta(5)-p/q|>q^{-88}$ for every sufficiently large
positive integer $q$ and every integer $p$.
\end{theorem}

The numerical proof of Theorem~\ref{thm:main} appears in
Section~\ref{sec:current-cert}. The rational parameters in that section
specify the result on the title page. The earlier configuration giving
$87.954278950885824$ is retained in Section~\ref{sec:certificate} to
document the arithmetic and spectral comparisons. The intervening
$87.925687270635509$ certificate added the local refinements in
Sections~\ref{z87:structural}--\ref{z87:kernel-completion}.
The contour and residue estimates used by the final certificate are identified in Section~\ref{sec:current-cert}.

Niedbala Giraudin's arXiv:2611.01708 states the integer bound $171$, conditional
on specified inputs from Fauzan's preprint \cite{ng}. That paper motivated the
original comparison. The construction here uses the independent manuscript
\cite{independent} and its quantitative extension \cite{previous}; Fauzan's
irrationality theorem and arithmetic bounds are not premises.

\paragraph{Scope of verification.}
This is an ordinary mathematical proof relative to the positive-measure,
Stieltjes, local-arithmetic, and weighted-polynomial lemmas identified in the
preceding report. The new algebraic and analytic deductions are proved below.
The finite certificates use integer and rational arithmetic with directed
square-root and logarithm enclosures. Their replay does not replace the
asymptotic proofs. The new quantitative theorem has not been built or checked
in Lean. No explicit $q_0$, global optimality theorem, or publication-priority
claim is made. The original qualitative proof sources remain unchanged.

\section{The inherited determinant and the quantitative inputs}
\label{sec:setup}

Let $\xi=\zeta(5)$ and
\[
 C_m(t)=\prod_{a=0}^{m-1}(t+a^2),\qquad
 D_N(t)=\prod_{i=1}^{N}(t+i^2).
\]
The source functional $\J_X$ has rational polynomial moments
\[
 \J_X(t^k)=\frac{(-1)^k}{24}(2k+3)(2k+4)(2k+5)B_{2k+2}
\]
and simple-pole values
\[
 \J_X\!\left(\frac1{t+i^2}\right)
 =i^4X-i^4H_{i-1}^{(5)}-\frac1{2i}-\frac14.
\]
Here $B_1=-1/2$. At $X=\xi$, this functional is integration against the
strictly positive measure
\[
 \dd\mu(x^2)=\frac{(2\pi)^4}{12}x^5
                  \sum_{j\ge1}j^4e^{-2\pi jx}\dd x\qquad(x>0).
\]

Use an unbounded integral subsequence with $m_i=s_iN$ and $r=\delta N$ integers.
The factorial-window determinant is related by invertible, constant row and
column changes to the symmetric Hankel pencil
\begin{equation}\label{eq:H}
 H_N(X)=\left[\J_X\!\left(
  \frac{C_{m_1}(t)C_{m_2}(t)t^{j+k}}{D_N(t)}\right)
                  \right]_{0\le j,k<r}.
\end{equation}
Its value $G_N=H_N(\xi)$ is positive definite. Write $M_N=[X]H_N(X)$ and
\[
 T_N=G_N^{-1/2}M_NG_N^{-1/2}.
\]
After primitive rational normalization the determinant gives
$P_N\in\Z[X]$, and
\begin{equation}\label{eq:det-ratio}
 P_N(\xi+h)=P_N(\xi)\det(I+hT_N).
\end{equation}
In either parameter configuration the analytic estimate and arithmetic certificate give the following bounds. The displayed numerical constants here describe the reference configuration of Section~\ref{sec:certificate}. Section~\ref{sec:current-cert} substitutes the final values in the same formulas.
\begin{align}
 \deg P_N&=\lambda N+O(1), & \lambda&=\CertLambda,\label{eq:degree}\\
 \log|P_N(\xi)|&\le-\varepsilon N^2+o(N^2),
 &\varepsilon&=\CertGapExact>0.\label{eq:gap}
\end{align}
Both configurations have $r>N-m_2+1$, and the exact degree is
$N-m_2+1$. The arithmetic and analytic rates are respectively bounded by
\[
 C_{\rm ar}\le\CertArithmetic,\qquad C_{\rm an}\ge\CertAnalytic,
 \qquad\varepsilon=C_{\rm an}^{\rm low}-C_{\rm ar}^{\rm up}.
\]
The bundle includes the complete inherited exact checker, the new arithmetic
correction checker, and both results. Every finite stage can be replayed
with the supplied launcher.

% Report-ready proof appendix. Requires amsmath, amssymb, amsthm, booktabs.
% All labels are namespaced lyr:. No global macro definitions are introduced.
\section{General arithmetic layers}
\label{lyr:section}

This section proves the local divisibility bounds used in the new arithmetic
profile. All valuations of matrices with entries in $\mathbb Q[X]$ are
coefficientwise in $X$. The second-layer statements are uniform at every
prime satisfying their displayed hypotheses. The third-layer statements
also use finitely many fixed rational nonsingularity tests; their exact
modular witnesses are part of the arithmetic certificate. No numerical
Smith-form experiment is a hypothesis.

\subsection{Normalization, local classes, and Newton coefficients}
\label{lyr:setup}

Write
\[
 P_j(x)=\frac{\prod_{b=0}^{j-1}(x^2-b^2)}{(2j)!},\qquad
 b_j(t)=\frac{(-1)^j\prod_{b=0}^{j-1}(t+b^2)}{(2j)!},
 \qquad D_N(t)=\prod_{i=1}^N(t+i^2).
\]
Thus $P_j(x)=b_j(-x^2)$. At an odd prime the $P_j$ form a basis of
the even integer-valued polynomials: the identity
\[
 2P_j(x)=\binom{x+j}{2j}+\binom{x+j-1}{2j}
 \quad(j>0)
\]
shows integer-valuedness, and evaluation at $0,1,\ldots$ is triangular
with diagonal entries $1,1/2,1/2,\ldots$. The matrix under consideration is
\begin{equation}\label{lyr:matrix}
 A_{jk}(X)=(N!)^2\mathcal J_X
 \left(\frac{b_{m_1+j}(t)b_{m_2+k}(t)}{D_N(t)}\right),
 \qquad 0\le j,k<r,
\end{equation}
where $R_i=m_i+r$. The local calculation uses the identities
\[
 \mathcal J_X\!\left(\frac1{t+i^2}\right)
 =i^4X-i^4H_{i-1}^{(5)}-\frac1{2i}-\frac14,
\]
\begin{equation}\label{lyr:bernoulli-functional}
 \mathcal J_X(q(t))=\frac1{24}\mathcal B\,
 \frac{d^3}{dx^3}\bigl[x^5q(-x^2)\bigr],
 \qquad \mathcal B(x^j)=B_j,\quad B_1=-\frac12.
\end{equation}
Bernoulli distribution gives
\begin{equation}\label{lyr:distribution}
 \mathcal B(F(x))=\frac1p\sum_{c=0}^{p-1}
 \mathcal B_Y(F(c+pY)).
\end{equation}

Throughout the new local arguments, take $p>5$ and
\begin{equation}\label{lyr:good-prime}
 p^2>4\max(R_1,R_2)+20,\qquad p^2>2N+20.
\end{equation}
Write
\[
 N=qp+\sigma,\quad 0\le\sigma<p,\quad q\ge1,\qquad
 \epsilon=\left\lfloor\frac{2\sigma}{p}\right\rfloor,
 \quad n=\frac{p-1}{2},\quad h=2q+1+\epsilon.
\]
We use reflected classes $1\le c\le n$, with signed physical levels
\[
 K_c=\{a\in\mathbb Z:1\le|c+pa|\le N\}.
\]
The maximal set $M$, its complement $T$, and their sizes are
\begin{equation}\label{lyr:class-sets}
 \begin{array}{c|c|c|c}
 \epsilon&M&T&w=|M|\\\hline
 0&\{1,\ldots,\sigma\}&\{\sigma+1,\ldots,n\}&\sigma\\
 1&\{p-\sigma,\ldots,n\}&\{1,\ldots,p-\sigma-1\}&\sigma-n.
 \end{array}
 \qquad t=|T|=n-w.
\end{equation}
There are $h$ physical levels on $M$ and $h-1$ on $T$.
The nonmaximal level set is $\{-q,\ldots,q-1\}$ if $\epsilon=0$ and
$\{-q,\ldots,q\}$ if $\epsilon=1$.

For later use put
\[
 Z_{2a}(Y)=\prod_{v=-a}^{a-1}(Y-v),\qquad
 Z_{2a+1}(Y)=\prod_{v=-a}^{a}(Y-v),\qquad Z_0=1.
\]
The common-root polynomial of window $m_i$ in class $c$ is
$Z_{k_i(c)}$, where
\begin{equation}\label{lyr:root-count}
 k_i(c)=\#\{a\in\mathbb Z:|c+pa|<m_i\}.
\end{equation}
These root polynomials form one nested family. Those of even degree
are invariant under $S(Y)=-1-Y$.

\begin{lemma}[Newton coefficients and Bernoulli precision]
\label{lyr:newton}
If $g$ is integer-valued at $p$, $\deg g=L<p^2$, and
$E=\lfloor L/p\rfloor$, then
\[
 v_p([Y^j]g(c+pY))\ge\max(0,j-E).
\]
In particular Bernoulli evaluation after three derivatives can be
computed modulo $p$ or modulo $p^2$ by first reducing the local
polynomial to that precision. The retained local degrees are at most
$E$ or $E+1$, respectively.
\end{lemma}
\begin{proof}
Expand $g$ in the integral Newton basis $\binom{x}{a}$. Among the
numerator factors $c+pY-i$, $0\le i<a$, either
$\lfloor a/p\rfloor$ or one more has constant term divisible by $p$.
If this count is $b$, extracting those factors and taking a coefficient
of degree $j$ gives valuation at least
\[
 b-\lfloor a/p\rfloor+\max(0,j-b)
 \ge\max(0,j-E).
\]
For reduction modulo $p^2$, omitted terms with $j>E+1$ have
coefficient valuation at least two. If $j<p+2$, the Bernoulli number
$B_{j-3}$ is $p$-integral. If $j\ge p+2$, the coefficient valuation is
at least three because $E\le p-1$; the possible loss of one power of
$p$ in a Bernoulli denominator still leaves a multiple of $p^2$.
The modulo-$p$ argument is the same with one less retained order.
Terms of degree below three vanish under differentiation.
\end{proof}

\subsection{The universal local functional and complete features}
\label{lyr:features}

For $K=\{-q-\eta,\ldots,q\}$, $\eta\in\{0,1\}$, let
$L(Y)=\prod_{a\in K}(Y-a)$, and let $I_P$ denote the remainder of
$P$ modulo $L$. Define
\[
 H(a)=\begin{cases}
 \displaystyle\sum_{j=1}^a j^{-5},&a\ge0,\\[2pt]
 \displaystyle\sum_{j=1}^{-a-1}j^{-5},&a\le-1,
 \end{cases}
\]
with empty sums zero, and
\begin{equation}\label{lyr:universal-functional}
 \mathcal T_{q,\eta}(P)
 =\sum_{a\in K}\frac{H(a)P(a)}{L'(a)}
 +\frac1{24}\mathcal B\left[
 \left(\frac{P-I_P}{L}\right)'''\right].
\end{equation}
This is a rational functional; its reduction is used only when the
required denominators are units. Increasing the good-prime cutoff
for any fixed finite list of local degrees makes this automatic.
For $\eta=1$,
\begin{equation}\label{lyr:anti-invariant}
 \mathcal T_{q,1}(P\circ S)=-\mathcal T_{q,1}(P),
 \qquad S(Y)=-1-Y.
\end{equation}
Indeed $K$ is $S$-invariant, $L\circ S=L$, $H(Sa)=H(a)$, and
$L'(Sa)=-L'(a)$. The quotient term changes sign because of its three
derivatives and the identity
$\mathcal B(F(-1-Y))=\mathcal B(F(Y))$.

For endpoint $i$ put
\begin{equation}\label{lyr:degree-cuts}
 D_i=\left\lfloor\frac{2R_i-2}{p}\right\rfloor,
 \qquad J_i=\lceil D_i p/2\rceil,
 \qquad H_i=R_i-J_i.
\end{equation}
Assume $m_i\le J_i$ whenever the complete-feature statement is used.

\begin{lemma}[Complete low features and correlated top coefficients]
\label{lyr:complete-features}
The reduction of the lattice with basis
$P_0,\ldots,P_{J_i-1}$ is isomorphic to
\[
 \bigoplus_{c=1}^{n}\mathbb F_p[Y]_{\le D_i-1}
 \ \oplus\ \{\text{even polynomials of degree at most }D_i-1\},
\]
where the last factor is the zero class. The low window beginning at
$m_i$ has exactly the nonzero-class spaces
\[
 Z_{k_i(c)}\mathbb F_p[Y]_{\le D_i-1-k_i(c)}.
\]
Its zero-class conditions are the evaluations forced by the integers
$0,p,2p,\ldots$ below $m_i$.

The top $H_i$ vectors have leading local coefficients
\begin{equation}\label{lyr:lucas-top}
 \frac{P_h(c)}{D_i!}\quad(D_i\text{ even}),\qquad
 \frac{cQ_h(c^2)}{D_i!}\quad(D_i\text{ odd}),\quad0\le h<H_i,
\end{equation}
where $P_h(c)$ is a polynomial of degree $h$ in $c^2$, and
\[
 Q_h(X)=\frac{\prod_{v=1}^{h}(X-(v-1/2)^2)}{(2h+1)!}.
\]
Every displayed polynomial has nonzero leading coefficient.
Consequently, with
\begin{equation}\label{lyr:Ai-Ki}
 L_i=\sum_{c\in M}(D_i-k_i(c))_+,
 \qquad A_i=L_i+\min(H_i,w),\qquad K_i=(H_i-w)_+,
\end{equation}
the maximal restriction map has rank exactly $A_i$, and its kernel
has exactly $K_i$ top coefficient directions. Writing
$V_M(X)=\prod_{c\in M}(X-c^2)$, those remaining top coefficients
are contained in
\begin{equation}\label{lyr:kernel-top}
 c^{D_i\bmod2}V_M(c^2)Q(c^2),\qquad\deg Q<K_i.
\end{equation}
\end{lemma}
\begin{proof}
Lemma~\ref{lyr:newton} gives the local degree bounds. If every local
polynomial reduction vanishes, every integer value is divisible by
$p$, hence so are all coefficients in the triangular factorial basis.
The map is injective, and its target dimension is
$D_i n+\lceil D_i/2\rceil=J_i$, so it is an isomorphism.

The $m_i$ conditions at $0,1,\ldots,m_i-1$ are independent by the
triangular evaluation matrix. Under the isomorphism they become
precisely the distinct root evaluations in each class. This proves
both the root spaces and independence between classes.

Formula~\eqref{lyr:lucas-top} follows coefficientwise from the
central-binomial identity for $2P_j$ and Lucas reduction. In an even
band $j=bp+h$, the complete local factor is
$P_h(c)\binom{Y+b}{2b}$; in an odd band
$j=bp+(p+1)/2+h$, it is
$cQ_h(c^2)\binom{Y+b}{2b+1}$. All factorial denominators are units
in the indicated top band. Thus evaluation on any $w$ distinct
nonzero squared nodes has rank $\min(H_i,w)$. When its top
coefficient vanishes on $M$, the remaining restrictions there can be
cancelled independently with low features. This gives
\eqref{lyr:Ai-Ki}; polynomial divisibility at the squared nodes gives
\eqref{lyr:kernel-top}.
\end{proof}

\subsection{The general second local layer}
\label{lyr:second-layer}

For window polynomials $f,g$, let $I$ be the truncation of $fg$ in
the factorial basis through index $N$, and define the exact polynomial
\begin{equation}\label{lyr:quotient}
 Q_{fg}(x)=(N!)^2\frac{f(x)g(x)-I(x)}{D_N(-x^2)}.
\end{equation}
The truncation agrees with $fg$ at $0,\pm1,\ldots,\pm N$. In a
nonzero class the exact denominator factors as
\[
 D_N(-(c+pY)^2)=p^{h_c}L_c(Y)U_c(Y),
\]
with $L_c$ monic and $U_c$ a constant unit modulo $p$. At zero,
extracting $p^2$ from each physical multiple of $p$ shows that
$Q_{fg}(pY)$ is coefficientwise integral. It follows that
$p^{1+\epsilon}Q_{fg}$ is integer-valued. The same calculation with
\eqref{lyr:distribution} proves the leading formula
\[
 p^{5+\epsilon}A(X)\bmod p
 =\sum_{c\in M}\omega_c\,
 \mathcal T_{q,\epsilon}(f_cg_c),\qquad
 \omega_c=2c^5\frac{(N!/p^q)^2}{\overline U_c}\ne0.
\]
In particular every maximal restriction-kernel row or column is
zero in this leading reduction. The coefficient of $X$ in $A$ has
valuation at least $-\epsilon$, so all layers considered here are
constant in $X$.

\begin{theorem}[Uniform second layer]
\label{lyr:second-theorem}
On the two maximal restriction kernels, the reduction of
$p^{4+\epsilon}A(X)$ is the pullback of
\begin{equation}\label{lyr:second-form}
 \bigoplus_{c\in T}\omega_c\mathcal T_*(f_cg_c),\qquad
 \mathcal T_*=\begin{cases}
 \mathcal T_{q-1,1},&\epsilon=0,\\
 \mathcal T_{q,0},&\epsilon=1.
 \end{cases}
\end{equation}
If $B$ is any upper bound for the rank of this pullback, then
\begin{equation}\label{lyr:second-determinant}
 v_p(\det A(X))\ge-(3+\epsilon)r-A_1-A_2-B.
\end{equation}
\end{theorem}
\begin{proof}
On a maximal class both $f_c$ and $g_c$ vanish modulo $p$, so their
product is divisible by $p^2$. The reduction of $I$ has degree at
most $h-1$ and vanishes at the $h$ physical nodes, hence is zero.
Lemma~\ref{lyr:newton} then shows that $I/p$ modulo $p$ has degree
at most $h$ and is a scalar multiple of $L_c$. Thus
$p^\epsilon Q_{fg}$ is integral and constant modulo $p$ on $M$.

On $T$, $L_c$ has degree $h-1$. The reduction of $I$ agrees with
$f_cg_c$ at those nodes and has degree at most $h-1$, so it differs
from the interpolation remainder by a constant multiple of $L_c$.
Consequently
\[
 \overline{p^\epsilon Q_{fg}(c+pY)}
 =\frac{(N!/p^q)^2}{\overline U_c}
 \frac{f_cg_c-I_{f_cg_c}}{L_c}+\text{constant}.
\]
Distribution gives exactly
\[
 p^{4+\epsilon}\operatorname{Quot}(f,g)
 =\frac1{24}\sum_{c=0}^{p-1}\mathcal B_Y\partial_Y^3
 \bigl[(c+pY)^5p^\epsilon Q_{fg}(c+pY)\bigr].
\]
Lemma~\ref{lyr:newton} allows reduction before differentiation.
The maximal constants, interpolation constants, and zero class
vanish. Proper fractions in a maximal class gain $p^2$ from the
two kernel factors and vanish at this scale; nonmaximal fractions
supply exactly the residue term in \eqref{lyr:second-form}.
Physical multiples of $p$ are integral. The retained constant
quotient from $P_N$ has valuation at least $-\epsilon$ and also
vanishes here.

For the determinant, put $M_0=p^{5+\epsilon}A$ and choose constant
unimodular bases adapted to the restriction kernels. All kernel
rows and columns are divisible by $p$. In their common block,
$M_0/p$ has rank at most $B$ modulo $p$. Operations within the
kernel rows leave at least $r-A_1-B$ such rows divisible by $p^2$
in the kernel columns. Factoring from these rows and all
$r-A_2$ kernel columns gives
$v_p(\det M_0)\ge2r-A_1-A_2-B$. If the stated row count is negative,
the single column factorization already implies the same weaker
inequality. This proves \eqref{lyr:second-determinant}.
\end{proof}

\paragraph{Four feature cuts.}
For $u_i\in\{0,1\}$, let
$n_i(c;u_i)=(D_i+u_i-k_i(c))_+$; $u_i=0$ means low degree
$D_i-1$, and $u_i=1$ means full degree $D_i$. Let
$b_c(u_1,u_2)=\min(n_1,n_2)$, except that when $\epsilon=0$ and
both root counts are even it may be replaced by
\begin{equation}\label{lyr:parity-cap}
 \min(\lceil n_1/2\rceil,\lfloor n_2/2\rfloor)
 +\min(\lfloor n_1/2\rfloor,\lceil n_2/2\rceil).
\end{equation}
The latter bound follows from anti-invariance and the two centered
parity eigenspaces. All four bounds
\begin{equation}\label{lyr:four-cuts}
 B_{u_1,u_2}=\sum_{c\in T}b_c(u_1,u_2)
 +(1-u_1)K_1+(1-u_2)K_2
\end{equation}
are valid: a full feature enlargement absorbs that side's top
directions, while a low cut pays at most $K_i$ additional rows or
columns. Also cap $B$ by $r-A_1$ and $r-A_2$.

\subsection{Cauchy cancellation and retained roots}
\label{lyr:cauchy}

The nonmaximal weights have a useful exact form. Each complete
$p$-block contributes $-1$ to $\overline U_c$. The partial block gives
\begin{equation}\label{lyr:nonmax-units}
 \overline U_c=\begin{cases}
 (-1)^{q+\sigma}V_M(c^2),&\epsilon=0,\\
 (-1)^{q+n-1+w}c^{-1}V_M(c^2),&\epsilon=1.
 \end{cases}
\end{equation}
For the second identity, use
$D_\sigma(t)=D_n(t)D_M(t)$ modulo $p$ and
$D_n(t)=t^n-(-1)^n$. Thus, up to a common nonzero scalar,
\[
 \omega_c=\frac{c^{5+\epsilon}}{V_M(c^2)}.
\]
Let $V_T(X)=\prod_{c\in T}(X-c^2)$. Since the squared nonzero
classes are the roots of $X^n-1$,
\[
 V_M(c^2)V_T'(c^2)=n/c^2\qquad(c\in T).
\]
After the two factors $V_M$ in \eqref{lyr:kernel-top}, the effective
weight is $c^3/V_T'$ for $\epsilon=0$ and $c^4/V_T'$ for
$\epsilon=1$.

\begin{theorem}[Second-layer Cauchy bound]
\label{lyr:cauchy-theorem}
Suppose $D_2-D_1\equiv1-\epsilon\pmod2$. Choose
$0\le l_i\le D_i$ with $d=D_1-l_1=D_2-l_2$, and put
\[
 E_i=\sum_{c\in T}(l_i-k_i(c))_+.
\]
Then a valid second-layer rank bound is
\begin{equation}\label{lyr:cauchy-bound}
 B\le dt+(K_1+K_2-t+\nu_D)_++E_1+E_2,
\end{equation}
where $\nu_D=2$, except that $\nu_D=3$ when $\epsilon=1$ and both
endpoint degrees are odd.
\end{theorem}
\begin{proof}
Impose the roots $Z_{l_i}$, paying at most $E_i$ directions, and
enlarge the low spaces to the full $d$-dimensional root spaces.
Their local pairing is one fixed square matrix $T_0$, multiplied
by $\omega_c$ in each class. If $T_0$ is singular, the low rank is
at most $(d-1)t$, and adjoining the top directions gives the
stronger bound $(d-1)t+K_1+K_2$.

If $T_0$ is invertible, its Schur complement has entries equal to a
common scalar times inverse-Vandermonde sums of polynomials of
degree at most $h+j+\nu_D$, for $0\le h<K_1$, $0\le j<K_2$.
The parity hypothesis makes these polynomials in $c^2$.
The identity
\[
 \sum_{c\in T}\frac{(c^2)^v}{V_T'(c^2)}=0
 \qquad(0\le v\le t-2)
\]
leaves rank at most $(K_1+K_2-t+\nu_D)_+$. The empty low block is
included in this case. Restoring the root directions proves the
claim without any nonsingularity assumption.
\end{proof}

\begin{proposition}[Retaining one extra root]
\label{lyr:retained-root}
Assume $\epsilon=0$, the endpoint parities are opposite, and
$d=D_i-l_i$ is odd. Let $e$ denote the even-degree endpoint and put
\[
 C=\#\{c\in T:k_e(c)\ge l_e+1\}.
\]
If the fixed rational low matrix $T_0$ in the preceding proof has
invertible reduction, then
\begin{equation}\label{lyr:retained-bound}
 B\le dt-C+(K_1+K_2-t+2)_++E_1+E_2.
\end{equation}
\end{proposition}
\begin{proof}
Transpose if necessary so $e=1$. Write $u=Y+1/2$,
$l_1=2a+1$, $\gamma=a+1/2$, and
$Z_{l_1}=(u-\gamma)E(u^2)$, $Z_{l_2}=F(u^2)$.
For $d=2r+1$, use on both residual polynomial spaces the basis
\[
 1,u^2,\ldots,u^{2r};\qquad
 (u+\gamma),(u+\gamma)u^2,\ldots,(u+\gamma)u^{2r-2}.
\]
Anti-invariance makes the two blocks orthogonal. If
$\mathcal S(P)=\mathcal T_{q-1,1}(uE(u^2)F(u^2)P(u^2))$, their
matrices are the Hankel matrices of $\mathcal S$ of size $r+1$ and
of $(v-\gamma^2)\mathcal S$ of size $r$. Both are invertible.

Choose the residual top polynomial $(u+\gamma)u^{2r}$. It couples
only to the second block. Retaining the actual root $u=-\gamma$
replaces the first row block by
$(v-\gamma^2)\mathbb F_p[v]_{\le r-1}$, lowering its rank by one:
its first $r$ column pairings form the invertible second block.
The remaining column radical is orthogonal to the top polynomial,
so the top Schur scalar is unchanged. This removes exactly one
low rank on each of the $C$ classes. The Cauchy top bound remains
unchanged, proving \eqref{lyr:retained-bound}.
\end{proof}

A nonzero determinant modulo one witness prime proves that the
corresponding fixed rational determinant is nonzero. Only finitely
many actual primes then have bad reduction. For any fixed finite
list of bands bounded away from $\beta=0$, all relevant primes
$p\ge\beta_{\min}N$ eventually avoid that finite set. This is the
precise asymptotic justification for the retained-root certificates.

\subsection{Central roots and the invariant first jet}
\label{lyr:third-geometry}

We now restrict to $\epsilon=0$ and a fixed even degree $D$. Let
$m_i=ap+\tau_i$ with the same integer $a\ge0$ and
$0<\tau_1\le\tau_2<p$. Put
\[
 e=D-2a\ge2,\qquad z=e/2-1,\qquad W(Y)=Z_{2a+1}(Y).
\]
The maximal root polynomial is $W$ in either of the following
regimes. We distinguish the class sets $\mathsf A,\mathsf B,\mathsf C$
from their cardinalities $a_0,b_0,c_0$.
\begin{equation}\label{lyr:third-regimes}
\begin{array}{c|c|c}
 &\text{upper half}&\text{lower half}\\\hline
 \text{conditions}&p/2<\tau_1\le\tau_2<p,
       \ w\le p-\tau_2
     &0<\tau_1\le\tau_2<p/2,
       \ w\le\tau_1-1\\
 V&Z_{2a+2}&Z_{2a}\\
 \mathsf A\text{ roots}&(W,W)&(W,W)\\
 a_0&p-\tau_2-w&\tau_1-1-w\\
 \mathsf B\text{ roots}&(W,V)&(V,W)\\
 b_0&\tau_2-\tau_1&\tau_2-\tau_1\\
 \mathsf C\text{ roots}&(V,V)&(V,V)\\
 c_0&\tau_1-(p+1)/2&n+1-\tau_2.
\end{array}
\end{equation}
The three sets are the successive nonmaximal intervals in $c$, so
$a_0+b_0+c_0=n-w$.

Write $\mathcal T=\mathcal T_{q-1,1}$ and
$\mathcal T_{\max}=\mathcal T_{q,0}$. The following fixed rational
nonsingularity hypotheses will be used:
\begin{equation}\label{lyr:third-local-tests}
\begin{split}
 &\mathcal T_{\max}\text{ on }
   W\mathbb Q[Y]_{\le D-\deg W},\qquad
 \mathcal T\text{ on the same space};\\
 &\mathcal T\text{ on }
   V_0:=V\mathbb Q[Y]_{\le D-1-\deg V}.
\end{split}
\end{equation}
The first two matrices have size $e$. The last has size $e-2$ in
the upper half and $e$ in the lower half. At a finite prime these
matrices must have invertible reduction. Only the finitely
certified tuples $(q,D,a,\text{half})$ are used in the rate bound.

Anti-invariance and the last test produce a unique monic radical
polynomial $R$ of degree $D$ in the full $V$ space
$V\mathbb F_p[Y]_{\le D-\deg V}$. It is even in $Y+1/2$.
In the upper half this full $V$ space has dimension $e-1$ and is
contained in the full $W$ space; its radical belongs to both.
In the lower half it has dimension $e+1$ and contains the full
$W$ space; the latter and $V_0$ are complements to its radical.
The nonsingularity tests justify both containment arguments.

Let $\psi$ be the monic polynomial of degree $D+1$ with root
polynomial $W$ that is orthogonal to the full degree-$D$ maximal
$W$ space under $\mathcal T_{\max}$. Define
\begin{equation}\label{lyr:local-scalars}
 \nu=\mathcal T(YR^2),\qquad
 \theta=\mathcal T_{\max}(\psi^2),\qquad
 \kappa=\theta-\nu.
\end{equation}
The one-odd-band theorem below requires $\kappa\ne0$; the
both-odd theorem also requires $\theta\ne0$. No such extra test
is needed for the equal-even rank upper bound.

\begin{lemma}[Highest first jet]
\label{lyr:first-jet}
For a window polynomial of even endpoint degree $D$, let $P(c)$
be its polynomial Lucas top coefficient, of degree at most
$2H-2$ in $c$. If
\[
 f(c+pY)=f_{0,c}(Y)+p f_{1,c}(Y)\pmod {p^2},
\]
then
\begin{equation}\label{lyr:jet-identity}
 [Y^{D+1}]f_{1,c}(Y)=P'(c).
\end{equation}
The reduction $f_{0,c}$ has degree at most $D$ and the first jet
has degree at most $D+1$.
\end{lemma}
\begin{proof}
For a top factorial vector, clear the denominator valuation $D$.
Its numerator reduces to a unit times
$(x^p-x)^D S(x)$ with $\deg S<p$. At $x=c+t$, all degrees below
$p$ in this product are a unit times $(-t)^D S(c+t)$. The
coefficients of $t^D$ and $t^{D+1}$ are the same unit times
$(-1)^D S(c)$ and $(-1)^D S'(c)$. Substitute $t=pY$ and restore
the denominator. Lower-band vectors have valuation at least two
in their degree-$D+1$ coefficient by Lemma~\ref{lyr:newton}.
This proves the identity for arbitrary integral linear combinations.
In particular there is no factor $1/(D+1)$ in
\eqref{lyr:jet-identity}.
\end{proof}

\paragraph{First correction of the nonmaximal functional.}
On $T$, with $D_N(-(c+pY)^2)=p^{2q}L(Y)U_c(Y)$, the first
logarithmic coefficient is
\begin{equation}\label{lyr:logarithmic-jet}
 \frac{U_c'(Y)}{pU_c(Y)}
 \equiv 2c\frac{V_M'(c^2)}{V_M(c^2)}\pmod p.
\end{equation}
Every full $p$-block contributes zero: its off-class sum is $-1/c$
and its two class companions contribute $1/c$. The partial block
is $V_M(c^2)$ up to a constant sign and gives the displayed term.
Thus the multiplier $(c+pY)^5/U_c(Y)$ has relative first
$Y$-coefficient
\[
 \lambda_c=5/c-2cV_M'(c^2)/V_M(c^2).
\]
Keeping the proper fractions and polynomial quotient modulo $p^2$
gives the same universal functional applied to this multiplier
and the two local polynomial jets. To justify the extra precision,
$I$ has local degree at most $2q+1$ modulo $p^2$ on $T$.
Division by the degree-$2q$ polynomial $L$ leaves only a constant
or affine interpolation ambiguity; multiplication by the
first-order affine unit does not increase the surviving ambiguity
beyond degree one. Three derivatives kill it. Multiplying the
interpolation remainder by that affine unit introduces only a
constant polynomial quotient, also killed. The nonmultiple-$p$
harmonic terms and the other Stieltjes terms begin at order $p^5$
in this scaled expression. Bernoulli reduction modulo $p^2$ is
legitimate by Lemma~\ref{lyr:newton}.

Where the opposite leading vector is a local radical, every
lower-degree piece of a first jet pairs to zero. A monic degree
$D+1$ representative can be $YR$, or $(Y-a)R$ if an extra odd
root must be retained. These give the same scalar because
$\mathcal T(R^2)=0$. Consequently, on radical sections with top
coefficients $P_1,P_2$, the nonmaximal first correction is a
common scalar times
\begin{equation}\label{lyr:derivative-form}
 \nu\sum_{c\in T}\omega_c
 \left[P_1'P_2+P_1P_2'+\lambda_cP_1P_2\right](c)
 =C_p\nu\sum_{c\in T}H'(c),
 \quad H(c)=\frac{c^5P_1(c)P_2(c)}{V_M(c^2)},
\end{equation}
where $\omega_c=C_p c^5/V_M(c^2)$. This includes classes where
one leading polynomial vanishes but its derivative does not.

\begin{lemma}[Half-field differentiation]
\label{lyr:half-field}
If $H$ is an odd polynomial with no linear term and
$\deg H<2p-1$, then
\[
 \sum_{c=1}^{n}H'(c)=0\quad\text{in }\mathbb F_p.
\]
If both $P_i$ are divisible by $V_M(c^2)$, the polynomial in
\eqref{lyr:derivative-form} has no linear term, and at $c\in M$
\begin{equation}\label{lyr:max-derivative}
 H'(c)=\frac{c^4}{2V_M'(c^2)}P_1'(c)P_2'(c).
\end{equation}
\end{lemma}
\begin{proof}
Twice the half-field sum is the sum over $\mathbb F_p^\times$,
because $H'$ is even. Its only possible surviving exponents are
$0$ and $p-1$. The constant coefficient is absent, and the
coefficient of $c^{p-1}$ is $p[c^p]H=0$. Formula
\eqref{lyr:max-derivative} follows by differentiating the one
remaining factor of $V_M$ after cancellation.
\end{proof}

\paragraph{The maximal first Schur correction.}
For two maximal restriction-kernel lifts, the first local jets
have root $W$ and degree at most $D+1$. The maximal contribution
at the third scale is $\mathcal T_{\max}$ on those first jets.
Cross terms with an unrestricted maximal feature use the same
pairing. Eliminating the invertible maximal degree-$D$ block
therefore projects each first jet onto the one-dimensional
quotient represented by $\psi$. Its coefficient is $P_i'(c)$.
The nonmaximal order-$p$ cross term is zero for the radical
sections, since these are local radicals against every full
opposite root space.

Thus the maximal third contribution and the boundary term
\eqref{lyr:max-derivative} use the same $w$ pairs of jet
coordinates. Their combined scalar is $\kappa$, not a sum of two
independent rank bounds. On a maximal class, the same assertion
can be checked directly by applying the Newton bounds to $I/p$
and $I/p^2$: after division by its degree-$(2q+1)$ local
polynomial, the remaining ambiguities are constant and affine,
so again disappear under three derivatives. The zero class and
the retained constant quotient vanish at the required precision.

\subsection{The equal-even and one-odd third layers}
\label{lyr:third-theorems}

\begin{theorem}[Equal even endpoints]
\label{lyr:equal-even}
Assume $D_1=D_2=D$, the appropriate regime
\eqref{lyr:third-regimes}, and the local tests
\eqref{lyr:third-local-tests}. Put $H_i=R_i-Dp/2$ and suppose
\[
 H_1-w\ge a_0,\qquad H_2-w\ge a_0+b_0,
 \qquad H_1+H_2-w\le p-2.
\]
The first two inequalities are equivalent. Define
\begin{equation}\label{lyr:even-k}
 k_f=H_1-w-a_0=H_2-w-a_0-b_0
 =\begin{cases}
 r+m_1+m_2-(D/2+a+1)p,&\text{upper half},\\
 r-(D/2-a)p+1,&\text{lower half}.
 \end{cases}
\end{equation}
For $M=p^5A$, the first rank is $r_0=ew$, the second rank is
\begin{equation}\label{lyr:even-second-rank}
 r_1=\begin{cases}
 ea_0+(e-1)b_0+(e-2)c_0,&\text{upper half},\\
 e(a_0+b_0+c_0),&\text{lower half},
 \end{cases}
\end{equation}
and the third rank is at most $w$. Hence
\begin{equation}\label{lyr:even-valuation}
 v_p(\det M)\ge(r-ew)+(k_f+z)+(k_f+z-w)_+.
\end{equation}
\end{theorem}
\begin{proof}
The maximal feature images are the complete $e$-dimensional
$W$ spaces, so their pairing is invertible and has rank $ew$.
In the two restriction kernels, the low nonmaximal features are
independent. The top evaluations can be prescribed freely on
$\mathsf A$ on the row side and on
$\mathsf A\cup\mathsf B$ on the column side, by the displayed
width conditions.

In the upper half, pairing against low features first forces
multiples of $R$ on $\mathsf C$. Top tests can use $R$ on the
other classes without disturbing those pairings; they then test
the full $W$ spaces on $\mathsf A$ and the full $V$ spaces on
$\mathsf B$. Thus left radicals vanish on $\mathsf A$ and are
multiples of $R$ on $\mathsf B\cup\mathsf C$; right radicals
vanish on $\mathsf A\cup\mathsf B$ and are multiples of $R$
on $\mathsf C$. In the lower half, the full $V$ space modulo
its radical is nondegenerate, and both $V_0$ and the full $W$
space are complements. Their mutual pairing is nondegenerate.
The same low and top tests therefore give exactly the same
support description.

The complete nonzero-class radicals consequently have top
polynomials
\begin{equation}\label{lyr:radical-sections}
 P_1(c)=V_M(c^2)V_{\mathsf A}(c^2)Q_1(c^2),\qquad
 P_2(c)=V_M(c^2)V_{\mathsf A}(c^2)V_{\mathsf B}(c^2)Q_2(c^2),
 \quad\deg Q_i<k_f.
\end{equation}
The zero-class low space has dimension $D/2$, with $a+1$ root
evaluations imposed, leaving $z=D/2-a-1$ independent features.
Its top coefficient is already determined by $P_i$. Both second
kernels thus have dimension $k_f+z$, and elementary counting
gives \eqref{lyr:even-second-rank}.

The degree of $H$ in \eqref{lyr:derivative-form} is at most
$2(H_1+H_2-w)+1<2p-1$. Lemma~\ref{lyr:half-field} reduces the
nonmaximal derivative sum to maximal jet evaluations. The
maximal Schur correction uses those same coordinates, so the
third rank is at most $w$. Eliminating the second leading block
does not change this third form: cross blocks on second radicals
start one order higher, and their Schur product begins at the
next order. Counting the first three kernel dimensions proves
\eqref{lyr:even-valuation}.
\end{proof}

\begin{theorem}[One additional odd column band]
\label{lyr:one-odd}
Keep the central-root regime and local tests, but suppose
$D_1=D$ and $D_2=D+1$. Put
\[
 J_0=Dp/2,\quad J_1=((D+1)p+1)/2,\qquad
 H_1=R_1-J_0,\quad H_2=R_2-J_1,
\]
and put $k_f=H_1-w-a_0$. Assume
\[
 0\le H_2\le w,\qquad k_f\ge w,\qquad
 H_1+H_2-w\le n-2,\qquad\kappa\ne0.
\]
Let $k_c=k_f-H_2$. Then the first rank is $ew$, the second rank
is the expression \eqref{lyr:even-second-rank} plus $H_2$, and
the third rank is at most $w-H_2$. In particular
\begin{equation}\label{lyr:one-odd-valuation}
 v_p(\det(p^5A))\ge(r-ew)+(k_c+z)+(k_f+z-w).
\end{equation}
\end{theorem}
\begin{proof}
First truncate the column window at $J_1$. Its complete even
band has top width $n+1$. The preceding radical argument gives
left parameter dimension $k_f$ and right dimension
$k_c=n+1-w-a_0-b_0$, with $k_f-k_c=H_2$ exactly.

The extra odd top coefficients are $cS(c^2)$, $\deg S<H_2$.
After maximal degree-$D$ projection their local maximal
polynomial is that coefficient times $\psi$. For a left radical
$P_1=V_M V_{\mathsf A}Q$, the nonmaximal second pairing is
\[
 C_p\nu\sum_{c\in T}c^6V_{\mathsf A}(c^2)Q(c^2)S(c^2).
\]
The maximal weight is
$(C_p/2)c^4/V_M'(c^2)$; since
$P_1'(c)=2cV_M'(c^2)V_{\mathsf A}(c^2)Q(c^2)$ on $M$, its
first-jet contribution is the same sum over $M$ with coefficient
$C_p\theta$.

The summand is an even polynomial without constant term and
of degree at most $2(H_1+H_2-w)+2<p-1$. Its full half-field sum
vanishes. The total second constraint is therefore
\begin{equation}\label{lyr:odd-constraint}
 C_p\kappa\sum_{c\in M}
 c^6V_{\mathsf A}(c^2)Q(c^2)S(c^2).
\end{equation}
The $Q$ evaluations are arbitrary because $k_f\ge w$, and the
$S$ evaluation matrix has rank $H_2\le w$. All weights are
nonzero, so these are $H_2$ independent conditions. No odd top
direction survives in the full second right kernel.

Choose right radical lifts in the complete even column band.
The same third form as in Theorem~\ref{lyr:equal-even} depends
only on maximal derivative evaluations. Condition
\eqref{lyr:odd-constraint} leaves a row boundary space of
dimension $w-H_2$, proving the third rank bound. An alternative
lift using $p$ times an odd column changes the form by
\eqref{lyr:odd-constraint}, which already vanishes. The second
Schur correction again begins one order later. This proves
\eqref{lyr:one-odd-valuation}.
\end{proof}

\subsection{Two odd endpoint bands}
\label{lyr:both-odd-section}

The next theorem keeps both odd bands simultaneously. Its proof
uses a finite-dimensional boundary form; the unknown first
correction of the odd/odd block does not need to be computed.

\begin{theorem}[Boundary flags for two odd bands]
\label{lyr:both-odd}
Use a central-root regime from \eqref{lyr:third-regimes} and the
local tests \eqref{lyr:third-local-tests} at even degree $D$, but
assume $D_1=D_2=D+1$. Set
\[
 J=((D+1)p+1)/2,\qquad H_i=R_i-J,
\]
\begin{equation}\label{lyr:odd-k}
 k_f=\begin{cases}
 (D/2-a)p+1-r,&\text{upper half},\\
 (D/2+a+1)p+2-r-m_1-m_2,&\text{lower half}.
 \end{cases}
\end{equation}
Assume $\theta\ne0$, $\kappa\ne0$, and
\begin{equation}\label{lyr:odd-guards}
 w\ge3,\qquad0\le H_i\le w-3,
 \qquad k_f+\min(H_1,H_2)\ge w.
\end{equation}
Put $\ell=(H_1+H_2-w+3)_+$. The first rank of $M=p^5A$ is
$ew+\ell$, its second kernels have dimension $k_f+\ell+z$, and
its third rank is at most $w-H_1-H_2+2\ell$. Hence
\begin{equation}\label{lyr:both-odd-valuation}
 v_p(\det M)\ge
 r-(e+1)w+2k_f+2z+H_1+H_2-\ell.
\end{equation}
\end{theorem}
\begin{proof}
Truncate both windows at $J$. Their complete even bands have
width $n+1$. The maximal rank is $ew$, and the nonmaximal low
rank is $r_{1,0}$, the expression
\eqref{lyr:even-second-rank}. The even second left radical has
$k_f+H_2$ nonzero-class parameters and the right one has
$k_f+H_1$; each also has $z$ zero-class parameters. These are
exact counts. In the upper half, for example, the two parameter
counts are $\tau_2-n$ and $\tau_1-n$.

Their top polynomials still have the forms
\eqref{lyr:radical-sections}. By \eqref{lyr:odd-guards}, their
maximal derivative evaluations are arbitrary. Let
$F=\mathbb F_p^w$ with the nondegenerate form
\begin{equation}\label{lyr:boundary-form}
 B(x,y)=\sum_{c\in M}\frac{c^6}{V_M'(c^2)}x_cy_c.
\end{equation}
The odd bands give flags
\[
 S_i=\{(S(c^2))_{c\in M}:\deg S<H_i\}\subset F.
\]
For an even radical polynomial use boundary coordinates
$E_i(c)=P_i'(c)/c$. After eliminating the complete even maximal
pivots and applying one common unit rescaling, the relevant
blocks have the leading terms
\begin{equation}\label{lyr:boundary-normal-form}
 \begin{array}{c|cc}
 &O_2&E_2\\\hline
 O_1&\theta B+O(p)&p\kappa B+O(p^2)\\
 E_1&p\kappa B+O(p^2)&p^2\kappa B+O(p^3).
 \end{array}
\end{equation}
Here $O_i$ is restricted to $S_i$, while the $E_i$ coordinates
surject onto $F$. Additional parameters in the kernels of these
evaluation maps contribute zero to the displayed terms.

We justify all three entries simultaneously. The odd/odd leading
term is the maximal Schur scalar $\theta$ times the two odd
coefficients $cS_i(c^2)$; the maximal weight is proportional to
$c^4/(2V_M')$, giving \eqref{lyr:boundary-form}. The even/odd
term is the calculation \eqref{lyr:odd-constraint}, with
$\kappa=\theta-\nu$. Its full half-field degree bound follows
from the complete even width $n+1$ and $H_i\le w-3$.
The even/even third term is the maximal Schur contribution minus
the derivative sum, again $\kappa$; its degree guard follows
from $w\ge3$. The coordinates $E_i=P_i'/c$ absorb the factors
of two in these formulas. The remaining nonmaximal even low
block has order $p$ and rank $r_{1,0}$. Its coupling to an even
radical starts at $p^2$, and to an odd direction at $p$.
Eliminating it therefore changes only the odd/odd $O(p)$ term
and higher remainders in \eqref{lyr:boundary-normal-form}.

In monomial bases for the flags, $B|_{S_1\times S_2}$ has entries
\[
 \sum_{c\in M}\frac{(c^2)^{h+j+3}}{V_M'(c^2)}.
\]
They vanish for $h+j\le w-5$, and the anti-diagonal
$h+j=w-4$ consists of ones. The width guards imply that the
terminal nonzero square has size exactly $\ell$, at most each
$H_i$; it is anti-triangular with unit anti-diagonal. Thus this
flag pairing has rank $\ell$.

Choose $S_i=L_i\oplus N_i$, where $B$ pairs $L_1,L_2$
nondegenerately in dimension $\ell$ and
\[
 N_1=S_1\cap S_2^\perp,\qquad
 N_2=S_2\cap S_1^\perp.
\]
Then $\dim N_i=H_i-\ell$, and $N_1,N_2$ are $B$-orthogonal.
Eliminating the $\ell$ leading odd pivots changes the even
$p^2$ coefficient to
\begin{equation}\label{lyr:flag-schur}
 Q(x,y)=\kappa B(x,y)-\frac{\kappa^2}{\theta}
 B(x,L_2)B(L_1,L_2)^{-1}B(L_1,y).
\end{equation}
The remaining order-$p$ cross blocks are
$\kappa B(E_1,N_2)$ and $\kappa B(N_1,E_2)$.
They have full ranks $H_2-\ell$ and $H_1-\ell$ by
surjectivity and $\kappa\ne0$. An arbitrary order-$p$
$N_1/N_2$ block is cleared by these cross blocks and does not
add rank. The second kernels have no residual odd coordinates
and satisfy exactly
\[
 x\in N_2^\perp,\qquad y\in N_1^\perp.
\]
Their full dimensions are $k_f+\ell+z$.

On these spaces, the row radical of $Q$ contains $N_1$ and its
column radical contains $N_2$. The correction in
\eqref{lyr:flag-schur} vanishes on those flags by their defining
orthogonality, and the remaining $B$ term vanishes against the
opposite orthogonal space. Therefore
\[
 \rank Q\le\dim N_2^\perp-\dim N_1
 =w-H_1-H_2+2\ell.
\]
The unidentified odd/odd correction cannot change this third
rank: on the second kernels every coupling to a remaining odd
pivot starts at $p^2$, so its Schur product through an order-$p$
pivot starts at $p^3$. This also explains why no nonvanishing
assumption on $\nu$ is needed for the rank upper bound.

The first three kernel counts are
\[
 r-ew-\ell,\qquad k_f+\ell+z,\qquad
 k_f+z-w+H_1+H_2-\ell.
\]
The last is nonnegative under \eqref{lyr:odd-guards}, since it
is at least $\max(H_1,H_2)-\ell+z$. Their sum gives
\eqref{lyr:both-odd-valuation}.
\end{proof}

\begin{lemma}[Restricted second kernels]
\label{lyr:restricted-kernels}
Let $M$ be a square $p$-integral matrix, with successive Smith
nullities $n_1,n_2,n_3$. Suppose its left and right second
kernels contain subspaces of dimensions $d_1,d_2$, on which
the induced third form has rank at most $B$. Then
\[
 v_p(\det M)\ge n_1+d_1+d_2-B.
\]
\end{lemma}
\begin{proof}
The full third form has rank at most
$B+(n_2-d_1)+(n_2-d_2)$, by extending bases of the two chosen
subspaces. Hence $n_2+n_3\ge d_1+d_2-B$, and the first three
Smith nullities give the stated valuation bound.
\end{proof}

\begin{corollary}[Two odd bands without evaluation surjectivity]
\label{lyr:partial-flags}
The determinant bound \eqref{lyr:both-odd-valuation} remains
valid if the last guard in \eqref{lyr:odd-guards}, namely
$k_f+\min(H_1,H_2)\ge w$, is omitted. There is no replacement
sign condition on $k_f$. The exact second rank and third rank
statements of Theorem~\ref{lyr:both-odd} are not asserted in
this enlarged region.
\end{corollary}
\begin{proof}
The simultaneous normal form
\eqref{lyr:boundary-normal-form} still holds on the actual even
boundary images; its derivation uses the degree guards
$H_i\le w-3$, not their surjectivity. The odd flags still have
pairing rank $\ell$, so the first nullity is $r-ew-\ell$.
After eliminating their leading pivots, the even parameter
spaces have dimensions $k_f+H_2+z$ and $k_f+H_1+z$.
These complete-even dimensions are nonnegative from their
class counts.

Restrict the left even parameters to boundary values in
$N_2^\perp$ and the right ones to values in $N_1^\perp$.
The numbers of imposed linear equations are at most
$H_2-\ell$ and $H_1-\ell$. Thus the resulting even-only
subspaces of the true second kernels each have dimension at
least $K=k_f+\ell+z$. No lower bound on the ranks of the
cross blocks is required, and residual odd directions may
remain in the full second kernels.

The third restriction is the form
\eqref{lyr:flag-schur} on these actual boundary images. Its
ambient restriction to $N_2^\perp\times N_1^\perp$ has rank
at most $C=w-H_1-H_2+2\ell$, by the same radical argument as
before. Restricting its arguments further cannot increase
that rank. The unknown order-$p$ odd/odd block has no effect
on this third restriction: the chosen even-only vectors
annihilate every order-$p$ cross pairing, and changes to
their lifts affect only the next order.

If $K\ge0$, Lemma~\ref{lyr:restricted-kernels} gives
$v_p(\det M)\ge r-ew-\ell+2K-C$, exactly
\eqref{lyr:both-odd-valuation}. If $K<0$, then $2K-C\le0$,
and the first-layer bound alone implies this same, weaker
inequality. This covers every finite case.
\end{proof}

\begin{corollary}[One odd band without evaluation surjectivity]
\label{lyr:partial-one-odd}
The determinant bound \eqref{lyr:one-odd-valuation} remains
valid after omitting the guard $k_f\ge w$ in
Theorem~\ref{lyr:one-odd}. Retain $0\le H_2\le w$, the
full even/odd degree guard, and the same local hypotheses.
Only the determinant bound, rather than the exact second
rank, is asserted with this relaxation.
\end{corollary}
\begin{proof}
Before imposing the odd column constraints, the left and right
even radicals have dimensions $k_f+z$ and $k_c+z$, where
$k_c=k_f-H_2$. Impose the $H_2$ equations of orthogonality
to the odd boundary flag on the left even parameters. This
gives an even-only left second-kernel subspace of dimension
at least $k_c+z$; the right even radical is already a
second-kernel subspace of that dimension. The unchanged
degree guard identifies the odd equations with the boundary
pairing, independently of its rank.

On these subspaces the third form has rank at most $w-H_2$,
because its ambient left boundary values are orthogonal to
an independent flag of dimension $H_2$. If $k_c+z\ge0$,
Lemma~\ref{lyr:restricted-kernels} therefore gives
\[
 v_p(\det(p^5A))\ge
 r-ew+2(k_c+z)-(w-H_2),
\]
which is \eqref{lyr:one-odd-valuation}. If $k_c+z<0$, the
displayed addition to the first nullity is nonpositive, so
the first-layer bound gives the same conclusion.
\end{proof}

\begin{corollary}[Two-layer counts from the third-layer spaces]
\label{lyr:central-second-counts}
In a common central-root regime with the three local matrix
tests, the following determinant bounds are also valid:
\[
 \begin{array}{c|c}
 \text{endpoint bands}&v_p(\det(p^5A))\text{ lower bound}\\\hline
 \text{equal even}&r-ew+k_f+z\\
 \text{one odd column}&r-ew+k_f-H_2+z\\
 \text{both odd}&r-ew+k_f+z.
 \end{array}
\]
Use the corresponding definition of $k_f$ in each row.
These bounds do not need the third-layer degree guards,
boundary surjectivity, or scalar nonvanishing tests. The
odd widths need only belong to their actual endpoint bands.
\end{corollary}
\begin{proof}
In the equal-even case the first nullity is at least $r-ew$,
and the constructed second kernel has dimension at least
$k_f+z$. When $k_f<0$, the $z$ free zero-class directions
already give this weaker count.

In the one-odd case the left even radical has dimension
$k_f+z$. Testing all $H_2$ odd column directions imposes at
most $H_2$ equations, with no need to identify or rank their
pairing. The right complete-even radical has dimension
$k_f-H_2+z$. These even-only second-kernel subspaces give the
stated count.

For two odd bands let $\ell$ be the actual leading odd/odd
rank after the complete-even maximal pivots. The first
nullity is $r-ew-\ell$. The even radical dimensions are
$k_f+H_2+z$ and $k_f+H_1+z$. Imposing the equations from the
remaining $H_2-\ell$ and $H_1-\ell$ odd directions leaves
even-only second-kernel subspaces of dimension at least
$k_f+\ell+z$. This cancels $\ell$ in the sum of the first
two nullities. If the displayed second-kernel lower bound
is negative, the first-layer bound alone suffices.
\end{proof}

\paragraph{Independent finite diagnostic.}
The original rational matrix at
\[
 (p,N,m_1,m_2,r)=(251,1530,179,197,1707)
\]
has Smith counts
$0:339$, $1:1308$, $2:6$, and at least $3:54$.
Here $w=24$, $(H_1,H_2)=(3,21)$, $\ell=3$, $k_f=51$, and
$z=6$. All finite guards hold. The theorem predicts first rank
$339$, second kernel dimension $60$, and third rank at most
$6$; all three are attained. This is a diagnostic of the
independent implementation, not a premise of the theorem.

\subsection{Odd bands at the maximal evaluation width}
\label{lyr:saturated-section}

The degree guard in Theorem~\ref{lyr:both-odd} can be removed
when an odd band reaches the maximal evaluation width. The full
finite-field sum then contributes a highest-coefficient pairing.
Keeping that pairing gives a larger admissible region.

\begin{theorem}[Saturated odd bands]
\label{lyr:saturated}
Keep the central-root regime, local matrix tests, odd endpoint
degrees, and definitions $e,z,k_f,H_i$ of
Theorem~\ref{lyr:both-odd}. Assume
\begin{equation}\label{lyr:saturated-guards}
 w\ge3,\qquad k_f\ge3,\qquad 0\le H_i\le n,
 \qquad\theta\nu\kappa\ne0.
\end{equation}
Suppose either $H_1\le w-3\le H_2$, its transpose, or
$H_1,H_2\ge w-3$. Then
\begin{equation}\label{lyr:saturated-valuation}
 v_p(\det(p^5A))\ge r-ew+2k_f+2z-3.
\end{equation}
In the first case the first rank is $ew+H_1$, the second kernel
dimension is $k_f+H_1+z$, and the third rank is at most $H_1+3$.
In the doubly saturated case let $\ell$ be the leading odd/odd
rank and $s=\ell-(w-3)$. Then $0\le s\le3$, the second kernel
dimension is $k_f+\ell+z$, and the third rank is at most $w+s$.
\end{theorem}
\begin{proof}
Use $X=c^2$ and write a left complete-even radical as
$P_1(c)=V_M(c^2)V_A(c^2)Q(c^2)$. Its parameter dimension is
\[
 d_Q=k_f+H_2=n+1-w-|A|.
\]
For an odd column polynomial $S$, $\deg S<H_2$, the same local
calculation as in \eqref{lyr:odd-constraint} now gives, up to a
common nonzero scalar,
\begin{equation}\label{lyr:coefficient-constraint}
 U(Q,S)=\kappa\sum_{c\in M}c^6V_A(c^2)Q(c^2)S(c^2)
 +\nu n[X^{n-3}](V_AQS).
\end{equation}
Indeed, the full half-field polynomial $X^3V_AQS$ has zero
constant term and degree at most $n+H_2-w+2<2n$. The only
surviving power in its sum over all nonzero squares is $X^n$,
whose sum is $n$. This proves
\eqref{lyr:coefficient-constraint} without a small-width guard.
The transposed formula uses the corresponding forced-zero
polynomial on that side. The complete-even third form still
depends only on maximal derivative evaluations, because its
two even widths are $n+1$ and $w\ge3$.

Put $s_0=w-3$ and let $S_0=\mathbb F_p[X]_{<s_0}$. The boundary
form \eqref{lyr:boundary-form} is nonsingular on $S_0$: its
entries are zero below the anti-diagonal $h+j=w-4$, which
consists of ones. An empty $S_0$ is interpreted in the usual
way. On $Q=V_MT$, the boundary part of
\eqref{lyr:coefficient-constraint} vanishes, leaving
\begin{equation}\label{lyr:high-coefficient-pairing}
 U(V_MT,S)=\nu n[X^{n-3}](V_AV_MTS).
\end{equation}
Here $T$ has dimension $L=k_f+H_2-w$. If $H_2\ge s_0$, put
$d=H_2-s_0$. Since $k_f\ge3$, one has $L\ge d$.
The last $d$ monomials of $T$ pair with
$X^{s_0},\ldots,X^{H_2-1}$ in an anti-triangular matrix with
nonzero anti-diagonal: $V_AV_M$ is monic, and its degree fixes
the coefficient threshold in
\eqref{lyr:high-coefficient-pairing}. All columns in $S_0$
vanish. Thus the exact kernel, on the full odd polynomial
space, of pairing against all $Q$ divisible by $V_M$ is $S_0$.
In particular the high constraints can be solved while
preserving all maximal boundary evaluations.

First suppose $H_1\le s_0\le H_2$. The first odd flag is
contained in $S_0$, and the second contains $S_0$. Their
boundary pairing has rank $H_1$. Eliminate these $H_1$ leading
$\theta$ pivots, choosing the paired column directions inside
$S_0$. The remaining low column flag has dimension $s_0-H_1$.
Its constraints are purely the $\kappa$ boundary pairing.
The remaining high column flag has dimension $d=H_2-s_0$;
its constraints have rank $d$ on the boundary-zero even
parameters by \eqref{lyr:high-coefficient-pairing}.

The row boundary evaluation map is onto, because
$d_Q=k_f+H_2\ge w$. The low constraints therefore have their
full rank $s_0-H_1$, independently of the high constraints.
The added second rank is $H_2-H_1$, and both second kernels
have dimension $k_f+H_1+z$. The remaining row boundary space
has dimension $w-(s_0-H_1)=H_1+3$. All eliminated leading
odd columns were in $S_0$, so their even/odd pairings were
pure boundary pairings. The resulting third Schur correction
therefore has rank at most $H_1+3$ as well. There are no
remaining odd rows, and all remaining odd columns are consumed
by the full second cross block. Later Schur products start at
$p^3$, as in Theorem~\ref{lyr:both-odd}. The three kernel
counts are at least
\[
 r-ew-H_1,\qquad k_f+H_1+z,\qquad k_f+z-3,
\]
whose sum is \eqref{lyr:saturated-valuation}.

Now suppose both flags contain $S_0$. Their leading pairing
rank satisfies $w-3\le\ell\le w$. After its $\ell$ unit
pivots, each residual odd flag has zero intersection with
$S_0$, since the opposite flag contains this nondegenerate
space. By the exact kernel statement for
\eqref{lyr:high-coefficient-pairing}, all residual odd
directions pair injectively with boundary-zero even parameters
at order $p$. The second kernels thus have only even
coordinates and dimension $k_f+\ell+z$.

For the third rank, first eliminate the nondegenerate $S_0$
block, of size $w-3$. Its even/odd couplings are boundary
pairings, so the resulting third form depends on at most $w$
boundary coordinates. There are at most
$s=\ell-(w-3)\le3$ additional unit odd pivots. Their couplings
can involve the coefficient term in
\eqref{lyr:coefficient-constraint}, but their additional
third Schur correction has rank at most $s$. Consequently the
third rank is at most $w+s$. The order-$p$ cross eliminations
again affect only the subsequent order on the second kernels.
The three kernel counts give
\[
 \begin{split}
 v_p(\det(p^5A))
 &\ge (r-ew-\ell)+(k_f+\ell+z)
       +(k_f+\ell+z-w-s)\\
 &=r-ew+2k_f+2z-3.
 \end{split}
\]
The last kernel count is nonnegative, since it equals
$k_f+z-3$. This proves the remaining case.
\end{proof}

\paragraph{Independent saturated diagnostics.}
At $(p,N,m_1,m_2,r)=(101,626,72,80,700)$ all finite guards
hold, with $w=20$, $(H_1,H_2)=(14,22)$, $k_f=8$, $e=14$, and
$z=6$. The original matrix has first rank $294$, second kernel
dimension $28$, and third rank $17=H_1+3$. The doubly
saturated instance $(101,626,76,79,704)$ has
$(H_1,H_2)=(22,25)$, $k_f=4$, and $\ell=w=20$; its first
rank, second kernel dimension, and third rank are $300,30,23$.
Its first three Smith nullities sum to $441$, matching the
displayed lower bound. In particular the
finite correction by the extra $s\le3$ pivots is necessary.

\subsection{Adjacent root quotients with opposite residue halves}
\label{lyr:crossed-section}

The final family has different maximal root polynomials on its
two sides. A rectangular local pairing and a maximal first-jet
correction replace the common-root boundary form.

\begin{theorem}[Adjacent root quotients]
\label{lyr:crossed}
Assume $\epsilon=0$, $N=qp+w$, and both endpoint degrees
are $D+1$, with $D$ even. Put
\[
 J=((D+1)p+1)/2,\qquad H_i=R_i-J,\qquad0\le H_i\le n.
\]
Suppose
\[
 m_1=ap+\tau_1,\qquad m_2=(a+1)p+\tau_2,
 \qquad p/2<\tau_1<p,\quad0<\tau_2<p/2,
\]
and impose
\begin{equation}\label{lyr:crossed-guards}
 \tau_2\le p-\tau_1,\qquad
 3\le w\le\tau_2-1,\qquad H_2\ge w,
 \qquad z=D/2-a-1\ge1.
\end{equation}
The complete low windows are assumed available as in the
feature lemma, in particular $m_i\le Dp/2$. Define
\[
 l=D-2a-1=2z+1,\quad
 k_0=(D/2-a)p-r,\quad
 k_{\mathrm{col}}=(D+2)p/2-R_2.
\]
Thus $k_{\mathrm{col}}=n-H_2\ge0$ and
$k_0=k_{\mathrm{col}}+\tau_2$. Let
$W_-=Z_{2a+1}$, $V=Z_{2a+2}$, and $W_+=Z_{2a+3}$.
Require nonsingularity of the following fixed rational forms:
\begin{equation}\label{lyr:crossed-local}
 \begin{array}{c|c|c}
 \text{functional}&\text{row and column spaces}&\text{size}\\\hline
 \mathcal T_{q,0}
 &W_-\Q[Y]_{\le D-1-\deg W_-}\ ,\quad W_+\Q[Y]_{\le D+1-\deg W_+}&l\\
 \mathcal T_{q-1,1}
 &W_-\Q[Y]_{\le D-\deg W_-}\text{ on both sides}&l+1\\
 \mathcal T_{q-1,1}
 &V\Q[Y]_{\le D-1-\deg V}\text{ on both sides}&l-1.
 \end{array}
\end{equation}
For all sufficiently large good primes,
\begin{equation}\label{lyr:crossed-valuation}
 \begin{split}
 v_p(\det(p^5A))
 &\ge r-lw+2(k_0+z)-\tau_2-2\\
 &=r-lw+k_0+k_{\mathrm{col}}+2z-2.
 \end{split}
\end{equation}
\end{theorem}
\begin{proof}
Put $u=p-\tau_1$. The nonmaximal classes have the following
three root types:
\[
 \begin{array}{c|c|c}
 A&w<c<\tau_2&(W_-,W_+)\\
 B&\tau_2\le c\le u&(W_-,V)\\
 C&u<c\le n&(V,V).
 \end{array}
\]
Thus $|A|=\tau_2-1-w$. Every maximal class has roots
$(W_-,W_+)$. Its full column space has dimension $l$, and
the first local test pairs it nondegenerately with the row
space of degree at most $D-1$. The guard $H_2\ge w$ makes
all full maximal column features available. Hence the first
rank is $lw$, and the degree-$D$ row radical there is a
unique monic polynomial $L$ of degree $D$.

The third local test and anti-invariance give a unique monic,
reflection-even radical $R$ of degree $D$ in the full $V$
space. The second local test makes the full degree-$D$
$W_-$ space nondegenerate. On $A$ its left kernel against
the degree-$D$ $W_+$ space has dimension two, containing
$R$ and a direction $U$ of degree at most $D-1$. On $B$
the left kernel against the $V$ space is $R$, and on $C$
both radicals are $R$. The right kernels on $A,B$ are zero.
These statements follow from the indicated root-space
containments and the nonsingularity of the $W_-$ form.

Truncate both windows at $J$. Their even top polynomials
$P_i(c)$ have degree at most $n$ in $c^2$. A preliminary
left even section has local features $P_1L$ on $M$,
$P_1R+t_cU$ on $A$, and $P_1R$ on $B,C$, together with
$z$ free zero-class low features. Its dimension is
$n+1+|A|+z$. A preliminary right section vanishes on
$M,A,B$ and is $P_2R$ on $C$. Its top polynomial vanishes
at $c=1,\ldots,u$, leaving $n+1-u=\tau_1-n$ parameters;
it also has $z-1$ zero-class low parameters.

The maximal first Schur complement removes the entire first
jet of a right even section, since it has degree at most
$D+1$ and root $W_+$. The preliminary left sections
therefore annihilate every complete-even right first-kernel
direction at order $p$, and the preliminary right sections
annihilate every complete-even left direction. On the
nonmaximal classes these are exactly the local radical
statements above. The residual column odd directions have
dimension $H_2-w$, since their top coefficients vanish on
$M$. Imposing their equations leaves a true left
second-kernel subspace $U_1$ of dimension at least
\[
 n+1+|A|+z-(H_2-w)=k_0+z.
\]
Imposing at most $H_1$ row odd equations leaves a right
second-kernel subspace $U_2$ of dimension at least
\[
 \tau_1-n+z-1-H_1=k_0+z.
\]
No exact rank of either set of equations is needed.

Within $U_1$ impose $P_1(c)=0$ on $M$ and $t_c=0$ on
$A$. The codimension is at most $w+|A|=\tau_2-1$.
On the resulting subspace, both top polynomials are
divisible by $V_M(c^2)$, all nonmaximal features are
$P_iR$, and both maximal leading features vanish. Write
$P_1(c)=V_M(c^2)G_1(c^2)$, with $\deg G_1\le n-w$.

Choose the polynomial $S$, $\deg S<w$, satisfying
\begin{equation}\label{lyr:crossed-interpolant}
 S(c^2)=P_2'(c)/c\qquad(c\in M).
\end{equation}
Because $H_2\ge w$, a correction $h$ in the original
column window can have odd top coefficient $cS(c^2)$.
Its lower features can be chosen so that its leading
maximal feature equals the entire first jet of the right
section. Subtracting $ph$ is the allowed first Schur
correction through the full column maximal feature space.
The corrected right first jet is zero on $M$, so the
maximal contribution at the third scale is zero.

Choose the row pivots from the complete lower even band
ending at $Dp/2$. Their nonmaximal degrees are therefore at
most $D-1$ everywhere. Right $R$ sections are nonzero only
on $C$, where these row pivots have root $V$, so their
nonmaximal pairings vanish. The lower-degree part of the
column correction also pairs to zero with $R$: every
column root space contains $V$, and subtracting
$cS(c^2)YR$ leaves a root-$V$ polynomial of degree at most
$D$. If $\nu=\mathcal T_{q-1,1}(YR^2)$, the only new
nonmaximal third contribution is consequently
\[
 -C_p\nu\sum_{c\notin M}c^6G_1(c^2)S(c^2).
\]
No nonvanishing assumption on $\nu$ is required.

Before this correction the nonmaximal deformation is
\[
 C_p\nu\sum_{c\notin M}H'(c),\qquad
 H(c)=\frac{c^5P_1(c)P_2(c)}{V_M(c^2)}.
\]
This includes classes
where $P_2(c)=0$ but its first jet is nonzero. The fixed-root
orthogonality, first-jet identity, and modulo-$p^2$
Bernoulli reduction are the same as in
\eqref{lyr:derivative-form}. The interpolation ambiguities
and zero-class terms again vanish at this precision.

Since $w\ge3$, one has $\deg H\le2p-2w+3<2p-1$;
the full half-field derivative sum is zero. On $M$,
\eqref{lyr:crossed-interpolant} gives
\[
 H'(c)=\frac{c^4}{2V_M'(c^2)}P_1'(c)P_2'(c)
       =c^6G_1(c^2)S(c^2).
\]
The maximal boundary terms cancel exactly. The remaining
third form on the restricted row space is
\begin{equation}\label{lyr:crossed-three-coefficients}
 -C_p\nu\sum_{c=1}^n c^6G_1(c^2)S(c^2)
 =-C_p\nu n[X^{n-3}](G_1(X)S(X)).
\end{equation}
Indeed $X^3G_1S$ has zero constant term and degree at most
$n+2<2n$, so only its $X^n$ coefficient survives. The
coefficient in \eqref{lyr:crossed-three-coefficients}
depends on at most the three highest coefficients of
$G_1$, of degrees $n-w-2,n-w-1,n-w$. Its rank is at most
three. Restoring the at most $\tau_2-1$ exceptional row
directions gives third rank at most $\tau_2+2$ on
$U_1\times U_2$.

Finally apply Lemma~\ref{lyr:restricted-kernels} with first
nullity $r-lw$ and the two chosen dimensions at least
$k_0+z$. This gives \eqref{lyr:crossed-valuation}.
\end{proof}

\paragraph{Independent adjacent-quotient diagnostic.}
For $(p,N,m_1,m_2,r)=(101,814,88,111,885)$, all finite
guards hold. Here $w=6$, $\tau_2=10$, $(H_1,H_2)=(13,36)$,
$k_0=24$, $k_{\mathrm{col}}=14$, and $z=8$. The original matrix
has first rank $102=17w$, second kernel dimension $32$, and
third rank $12=\tau_2+2$. Removing the three exceptional
nonmaximal row features leaves rank nine; removing the six
maximal row features as well leaves rank exactly three.
The three fixed matrix tests were independently certified
for $(q,D,a)=(8,18,0),(16,38,1),(24,58,2)$.

\begin{corollary}[Mixed degrees across adjacent root quotients]
\label{lyr:crossed-mixed}
Keep the root, ordering, maximal, and local matrix hypotheses
of Theorem~\ref{lyr:crossed}, but let the actual row endpoint
degree be $D$ and the column degree be $D+1$. Put
$H_{\mathrm{even}}=R_1-Dp/2$, retain $H_2\ge w$, and define
$k_{\mathrm{even}}=\tau_1-n-1$. Then
\[
 \begin{split}
 v_p(\det(p^5A))&\ge r-lw+k_{\mathrm{even}}+z,\\
 v_p(\det(p^5A))&\ge
 r-lw+2(k_{\mathrm{even}}+z)-\tau_2-2.
 \end{split}
\]
\end{corollary}
\begin{proof}
The row even top width is now $H_{\mathrm{even}}\le n+1$,
and there are no odd row constraints. The preliminary left
dimension is $H_{\mathrm{even}}+|A|+z$. After imposing at
most $H_2-w$ odd column equations it is at least
\[
 H_{\mathrm{even}}+|A|+z-(H_2-w)
 =\tau_1-n-1+z=k_{\mathrm{even}}+z,
\]
using $H_{\mathrm{even}}-H_2=\tau_1-\tau_2-n$.
The preliminary right space already has this same dimension,
and there are no row odd equations to reduce it. The first
and second nullity counts give the first bound. The
exceptional row restrictions and first-jet correction are
unchanged; the row top degree is only smaller. Thus the
restricted third rank remains at most $\tau_2+2$, and
Lemma~\ref{lyr:restricted-kernels} gives the second bound.
\end{proof}

\begin{corollary}[Vanishing mixed cross constraint]
\label{lyr:crossed-mixed-strong}
In Corollary~\ref{lyr:crossed-mixed}, further suppose
\begin{equation}\label{lyr:crossed-mixed-strong-guards}
 H_{\mathrm{even}}\ge w,\qquad
 H_{\mathrm{even}}+H_2-w\le n-2.
\end{equation}
Then
\begin{equation}\label{lyr:crossed-mixed-strong-valuation}
 v_p(\det(p^5A))\ge
 r-lw+H_{\mathrm{even}}-w+k_{\mathrm{even}}+2z.
\end{equation}
\end{corollary}
\begin{proof}
Impose the maximal row zeros and remove the extra $A$
directions directly in the preliminary left even space.
The resulting space has dimension
$H_{\mathrm{even}}-w+z$, with nonmaximal features $P_1R$
and $V_M\mid P_1$. A residual odd column top coefficient
has the form $cV_M(c^2)T(c^2)$, $\deg T<H_2-w$. Its
second pairing is
\[
 C_p\nu\sum_{c=1}^n c^6P_1(c)T(c^2).
\]
The maximal summands vanish. As a polynomial in $X=c^2$,
the summand has zero constant term and degree at most
$H_{\mathrm{even}}+H_2-w+1\le n-1$. Thus the full sum
is zero. The chosen left space is already a true second
kernel, with no further loss from odd column equations.
The right preliminary space has dimension
$k_{\mathrm{even}}+z$ and is also a true second kernel,
since there are no odd row directions.

The third form is the coefficient in
\eqref{lyr:crossed-three-coefficients}. Here
$\deg(G_1S)\le H_{\mathrm{even}}-2$. The displayed degree
guard and $H_2\ge w$ imply
$H_{\mathrm{even}}\le n-2$, so that coefficient is zero.
Lemma~\ref{lyr:restricted-kernels}, with the two unequal
chosen dimensions, proves the claim.
\end{proof}

\paragraph{Independent mixed diagnostic.}
At $(p,N,m_1,m_2,r)=(101,814,88,111,860)$, one has
$H_{\mathrm{even}}=39$, $H_2=11$, $w=6$, $n=50$,
$k_{\mathrm{even}}=37$, and $z=8$. Both extra guards hold.
The first nullity is $758$, and the chosen dimensions are
$41$ and $45$, giving a valuation lower bound of $844$.
The original matrix
has first rank $102$, second kernel dimension $45$, and
third rank $4$, so its first three Smith nullities sum to
$758+45+41=844$.

\subsection{Passage to the arithmetic profile}
\label{lyr:rates}

Write $\beta=p/N$, $m_i/N\to s_i$, $r/N\to\delta$, and
$\rho_i=s_i+\delta$. On each open rational floor cell, use
\[
 D_i=\lfloor2\rho_i/\beta\rfloor,\qquad
 H_i=\rho_i-D_i\beta/2,\qquad
 w=1-q\beta-\epsilon\beta/2,\qquad t=\beta/2-w,
\]
where symbols such as $H_i,w,t$ now denote rates. Replace the
finite sums of class dimensions by their interval integrals.
The second-layer cost from \eqref{lyr:second-determinant} is
\begin{equation}\label{lyr:second-rate}
 E_2(\beta)=(3+\epsilon)\delta+A_1+A_2+B,
\end{equation}
where $B$ is the minimum of the four feature cuts, the Cauchy
bounds, and the finitely certified retained-root bounds.
The constants $\nu_D$ disappear in the rate.

For the central-root third layers let
$a=\lfloor s_1/\beta\rfloor=\lfloor s_2/\beta\rfloor$,
$e=D-2a$, and $w=1-q\beta$. The equal-even and one-odd
candidates are
\begin{equation}\label{lyr:even-mixed-rates}
 \begin{aligned}
 E_3&=4\delta+(e+1)w-2k
       &&\text{(equal even)},\\
 E_3&=4\delta+(e+1)w-2k+H_2
       &&\text{(one odd column band)},
 \end{aligned}
\end{equation}
with
\[
 k=\begin{cases}
 \delta+s_1+s_2-(D/2+a+1)\beta,&\text{upper half},\\
 \delta-(D/2-a)\beta,&\text{lower half}.
 \end{cases}
\]
For two odd bands, $H_i=\rho_i-(D+1)\beta/2$ and
\begin{equation}\label{lyr:both-odd-rate}
 E_3=4\delta+ew-2k+(w-H_1-H_2)_+,
 \qquad
 k=\begin{cases}
 (D/2-a)\beta-\delta,&\text{upper half},\\
 (D/2+a+1)\beta-\delta-s_1-s_2,&\text{lower half}.
 \end{cases}
\end{equation}
All residue-half, maximal-root, degree-band, top-width, and
boundary-evaluation guards in the finite theorems are enforced
in their limiting forms. By Corollary~\ref{lyr:partial-flags},
the two-odd determinant candidate needs only $H_i\le w$;
there is no additional boundary-surjectivity or sign guard
on $k$. In the one-odd candidate,
Corollary~\ref{lyr:partial-one-odd} removes $k\ge w$ while
retaining its even/odd degree guard. The saturated
extension applies when $\max(H_1,H_2)\ge w$ and $k>0$,
with the same residue and local matrix guards and the additional
scalar witness $\nu\ne0$. Its rate is
\begin{equation}\label{lyr:saturated-rate}
 E_3=4\delta+ew-2k.
\end{equation}
The additional candidates from
Corollary~\ref{lyr:central-second-counts} are
\begin{equation}\label{lyr:central-second-rates}
 \begin{aligned}
 E_2&=4\delta+ew-k&&\text{(equal even or both odd)},\\
 E_2&=4\delta+ew-k+H_2&&\text{(one odd column)}.
 \end{aligned}
\end{equation}
Their hypotheses consist of the common-root geometry and
the local matrix tests; the additional third-layer guards
are not imposed.

For Theorem~\ref{lyr:crossed}, retain its adjacent quotients,
opposite residue halves, and interval ordering. Write
$\tau_1=s_1-a\beta$, $\tau_2=s_2-(a+1)\beta$ and
$l=D-2a-1$. Its candidate is
\begin{equation}\label{lyr:crossed-rate}
 E_3=4\delta+lw-k_0-k_{\mathrm{col}},\qquad
 k_0=(D/2-a)\beta-\delta,\quad
 k_{\mathrm{col}}=(D+2)\beta/2-\rho_2.
\end{equation}
The guards include $\tau_2\le\beta-\tau_1$,
$w\le\tau_2$, and $H_2\ge w$, with $w=1-q\beta$.
Only the three fixed matrix witnesses in
\eqref{lyr:crossed-local} are needed for this family.
For its mixed-degree version let
$k_{\mathrm{even}}=s_1-(a+\tfrac12)\beta$. The two candidates
from Corollary~\ref{lyr:crossed-mixed} combine as
\begin{equation}\label{lyr:crossed-mixed-rate}
 E=4\delta+lw-k_{\mathrm{even}}
       -(k_{\mathrm{even}}-\tau_2)_+.
\end{equation}
Under the additional rate guard
$H_{\mathrm{even}}+H_2-w\le\beta/2$ and
$H_{\mathrm{even}}\ge w$,
Corollary~\ref{lyr:crossed-mixed-strong} supplies
\begin{equation}\label{lyr:crossed-mixed-strong-rate}
 E=4\delta+lw-(H_{\mathrm{even}}-w)-k_{\mathrm{even}},
 \qquad H_{\mathrm{even}}=\rho_1-D\beta/2.
\end{equation}

The new candidates are combined with the inherited arithmetic
profile by taking the pointwise minimum. Each candidate is
piecewise affine on an exact rational partition. Fixed finite
rounding offsets, zero-class dimensions, and the small-prime
cutoff contribute only the inherited lower-order error. At a
limiting equality in a guard, either omit the boundary primes or
use the conservative finite guard; the signed prime-summation
argument is unchanged. The previous tail bound is retained
outside the finitely corrected range.

For every retained-root or third-layer tuple actually used by
the exact integrator, the checker recomputes a nonzero modular
determinant of each required fixed rational matrix, and the
required nonzero scalar witnesses. A witness prime is not a
condition on the varying primes in the asymptotic sum: it proves
rational nonvanishing, which excludes only finitely many fixed
bad primes. This separates the ordinary asymptotic proof from
the finite exact calculations that certify its stated inputs.
\section{Further local arithmetic extensions}
\label{lyrx:section}

This section retains the notation and local precision hypotheses of
Section~\ref{lyr:section}. The new statements are determinant valuation
bounds obtained from selected second-kernel spaces. They do not extend
the earlier exact Smith-rank formulas beyond their stated hypotheses.

\subsection{An arbitrary odd column head in the common-root regime}

\begin{theorem}[Common mixed head without the old width guard]
\label{lyrx:common-mixed}
Keep the common-root geometry and the three local matrix tests
\eqref{lyr:third-local-tests}. Suppose the actual row and column endpoint
degrees are $D$ and $D+1$, respectively, with $D$ even. Put
\[
 n=(p-1)/2,\quad H_1=R_1-Dp/2,\quad
 H_2=R_2-((D+1)p+1)/2,
\]
\[
 e=D-2a,\quad z=e/2-1,\quad k_f=H_1-w-a_0,
 \quad k_c=k_f-H_2.
\]
Here $a_0$ is the size of the nonmaximal class set $A$ in
\eqref{lyr:third-regimes}. Assume $w\ge3$ and the actual width bounds
$1\le H_1\le n+1$, $0\le H_2\le n$. Set
\[
 d=(H_1+H_2-w-n+2)_+.
\]
Then, coefficientwise in $X$,
\begin{equation}\label{lyrx:common-mixed-finite}
 v_p(\det(p^5A))\ge r-ew+k_f+k_c-w-d+2z.
\end{equation}
No nonvanishing condition on $\theta$, $\nu$, or
$\kappa=\theta-\nu$ is needed.
\end{theorem}

\begin{proof}
The complete-even comparison spaces from the proof of
Theorem~\ref{lyr:one-odd} have dimensions $k_f+z$ on the left and
$k_c+z$ on the right. Their top polynomials have the forms
$P_1=V_MV_AQ$ and $P_2=V_MV_AV_BQ_2$, with $\deg Q<k_f$ and
$\deg Q_2<k_c$. The right space is a true second-kernel subspace,
because every actual row direction is even. The first rank is $ew$.

For an odd column top coefficient $cS(c^2)$, $\deg S<H_2$, the
full second constraint, up to the common nonzero factor $C_p$, is
\begin{equation}\label{lyrx:mixed-full-constraint}
 U(Q,S)=\kappa\sum_{c\in M}c^6V_A(c^2)Q(c^2)S(c^2)
       +\nu n[X^{n-3}](V_AQS).
\end{equation}
Indeed, adding and subtracting the maximal part of the nonmaximal
$\nu$ sum leaves the full half-field sum. Its polynomial in $X=c^2$
is $X^3V_AQS$, with zero constant term and degree at most
\[
 a_0+k_f+H_2+1=H_1+H_2-w+1\le2n+2-w<2n.
\]
Only the $X^n$ coefficient can survive. This proves
\eqref{lyrx:mixed-full-constraint} without the former degree guard.

Put $s=n-H_1+w-2\ge0$, $h=\min(H_2,s)$, and
$m=\min(w,h)$. For $\deg S<s$, the coefficient term in
\eqref{lyrx:mixed-full-constraint} vanishes, since
\[
 \deg(V_AQS)\le a_0+k_f+s-2=n-4.
\]
The low-head evaluations have rank $m$ at the distinct maximal
nodes. Impose their boundary-orthogonality equations on $Q$, and
then impose the actual equations from any complement of that low
head. The complement has dimension $H_2-h=d$. The resulting true
left second-kernel space has dimension at least
\[
 k_f-m-d+z,
\]
and its boundary image has dimension at most $w-m$. These are
upper bounds on numbers of equations; no evaluation surjectivity or
scalar nonvanishing is used.

The complete-even third form is a boundary form. Its degree guard
is automatic: $H_1+(n+1)-w\le2n-1=p-2$. It therefore has rank
at most $w-m$ on the chosen left space and the complete-even right
space. A lift change using an odd column changes this form by
\eqref{lyrx:mixed-full-constraint}, which vanishes on the chosen
left space. Subsequent second-pivot Schur products start at $p^3$.
Lemma~\ref{lyr:restricted-kernels} now gives
\[
 v_p(\det(p^5A))
 \ge r-ew+(k_f-m-d+z)+(k_c+z)-(w-m),
\]
which is \eqref{lyrx:common-mixed-finite}. A negative displayed
dimension lower bound may be replaced by the actual nonnegative
dimension, so it does not invalidate the inequality.
\end{proof}

In normalized variables, this supplies the candidate
\begin{equation}\label{lyrx:common-mixed-rate}
 E=4\delta+(e+1)w-2k+H_2
       +(H_1+H_2-w-\beta/2)_+.
\end{equation}
The old conditions $H_2\le w$ and $H_1+H_2-w\le\beta/2$ are
removed; all common-root and local matrix conditions remain.
For $H_2\ge w$, writing $k_c=k-H_2$, the same expression is
\[
 E=4\delta+ew-2k_c-\min(H_2-w,\beta/2-H_1).
\]

\subsection{Truncating a mixed head across adjacent quotients}

\begin{theorem}[Capped mixed crossed head]
\label{lyrx:crossed-capped}
Keep the mixed adjacent-quotient hypotheses of
Corollary~\ref{lyr:crossed-mixed}, including $H_2\ge w$, and put
$H=H_{\mathrm{even}}$. For the additional candidate assume $H\ge w$.
Define
\[
 A_* =\max\{0,\min(H-w,n-H_2-2)\}.
\]
Then
\begin{equation}\label{lyrx:crossed-capped-finite}
 v_p(\det(p^5A))\ge r-lw+k_{\mathrm{even}}+A_*+2z.
\end{equation}
The former guard $H+H_2-w\le n-2$ is not needed.
\end{theorem}

\begin{proof}
Use the same restricted spaces as in
Corollary~\ref{lyr:crossed-mixed-strong}, but choose the left top
polynomial to be
\[
 P_1(c)=V_M(c^2)G(c^2),\qquad\deg G<A_*.
\]
Remove all exceptional nonmaximal row directions as before. Since
$A_*\le H-w$, these choices fit in the actual row head. With the
independent low zero-class parameters, the left dimension is
$A_*+z$. The original right comparison space still has dimension
$k_{\mathrm{even}}+z$ and lies in the true second kernel.

A residual odd column head is $cV_M(c^2)T(c^2)$ with
$\deg T<H_2-w$. Its second pairing with the chosen left space is
the full half-field sum of $X^3V_MGT$. If $A_*>0$, this polynomial
has zero constant term and degree at most
\[
 A_*+H_2+1\le n-1.
\]
Its sum is zero. If $A_*=0$, the nonzero-class term vanishes
identically. Thus the chosen left space is a true second-kernel
subspace.

The established crossed cancellation gives the third restriction
as a scalar multiple of $[X^{n-3}](GS)$ with $\deg S<w$.
For $A_*>0$,
\[
 \deg(GS)\le A_*+w-2\le n-H_2+w-4\le n-4.
\]
It therefore vanishes as well. The low zero-class terms vanish at
the retained precision. Applying Lemma~\ref{lyr:restricted-kernels}
with first nullity $r-lw$ and third rank zero proves the claim.
\end{proof}

The limiting candidate is
\begin{equation}\label{lyrx:crossed-capped-rate}
 E=4\delta+lw-k_{\mathrm{even}}
       -\max\{0,\min(H-w,\beta/2-H_2)\}.
\end{equation}
It is combined with every old candidate. Each new minimum or
positive-part break belongs in the exact affine partition.

\subsection{Scalar-free odd-head valuation corollaries}

\begin{proposition}\label{lyrx:scalar-free}
The determinant bounds \eqref{lyr:both-odd-valuation} and
\eqref{lyr:saturated-valuation} retain their validity without the
nonvanishing conditions on $\theta$, $\nu$, and $\kappa$.
The saturated valuation bound also does not require $k_f\ge3$.
The three base local matrix tests and the remaining geometric and
width hypotheses are retained. The earlier exact-rank formulas
retain their original assumptions.
\end{proposition}

\begin{proof}
For the unsaturated case let $\ell=\operatorname{rank}B(S_1,S_2)$.
When $\theta=0$ and $\kappa\ne0$, the first nullity is $r-ew$.
Imposing all odd equations on even-only parameters gives chosen
second-kernel dimensions at least $k_f+z$ on both sides. Their
boundary images lie in $S_2^\perp$ and $S_1^\perp$. The restricted
boundary form has rank at most
\[
 (w-H_2)-(H_1-\ell)=w-H_1-H_2+\ell,
\]
because its row radical contains $S_1\cap S_2^\perp$. The
restricted-kernel lemma gives precisely
\eqref{lyr:both-odd-valuation}. If $\theta\ne0$ and $\kappa=0$,
the complete-even spaces have zero third restriction after the
leading odd pivots; their larger dimensions give a stronger bound.
If $\theta=\kappa=0$, then $\nu=0$ as well and the same conclusion
holds without leading odd pivots. The remaining branch is the
already proved partial-flag argument.

For saturated heads with $\theta\ne0$, impose all residual odd
equations on even-only parameters instead of asserting their full
rank. In the two-sided case the chosen dimensions are at least
$k_f+\ell+z$. Eliminating the nondegenerate degree-$<w-3$ odd
block leaves a third form through $w$ boundary coordinates; the
at most $\ell-(w-3)\le3$ extra first pivots add that many ranks.
Thus the third rank is at most $\ell+3$, independently of $k_f$.
In the one-sided case the residual low flag forces one boundary
image to have dimension at most $H_1+3$ when $\kappa\ne0$.
The selected dimensions are at least $k_f+H_1+z$. Both counts give
\eqref{lyr:saturated-valuation}. When $\kappa=0$, the low
even/odd terms and the raw third even form vanish; only the at
most three extra first pivots can contribute, giving a stronger
bound.

Finally suppose $\theta=0$ in the saturated case. Imposing all odd
equations gives even-only spaces of dimension at least $k_f+z$ on
both sides. If $\kappa\ne0$, the saturated opposite head contains
the degree-$<w-3$ flag, forcing one selected boundary image to have
dimension at most three. The third restriction has rank at most
three. If $\kappa=0$, also $\nu=0$ and the third restriction is
zero. The restricted-kernel lemma proves the same valuation bound.
In every branch, genuine second-kernel restrictions ensure that
later Schur products start at $p^3$.
\end{proof}

The removal of $k>0$ alone does not improve an envelope containing
the central two-layer candidate: for $k<0$,
$4\delta+ew-k<4\delta+ew-2k$. Scalar-free extensions are useful
only where a previously omitted fixed tuple satisfies the three
base matrix tests.

\subsection{Cauchy cancellation on a containment core}

This extension uses the actual nonmaximal root spaces. It has no
local matrix nonvanishing hypothesis.

\begin{theorem}\label{lyrx:containment-core}
Retain the finite notation and good-prime hypotheses of the general
second-layer theorem. Assume
$D_2-D_1\equiv1-\epsilon\pmod2$, and choose root degrees
$\ell_1,\ell_2$ such that, for every $c\in T$,
\[
 k_1(c)\le\ell_1,\qquad
 \ell_2\le k_2(c)\le D_2,\qquad
 d=D_1-\ell_1=D_2-\ell_2\ge0.
\]
Put
\[
 C=\{c\in T:k_1(c)=\ell_1,\ k_2(c)=\ell_2\},\qquad
 f=|T\setminus C|,\qquad
 L_2=\sum_{c\in T}(D_2-k_2(c)).
\]
Let $t=|T|$, $K_i=(H_i-w)_+$, and let $\nu=2$ except when
$\epsilon=1$ and both $D_i$ are odd, in which case $\nu=3$.
Then the second local rank satisfies
\begin{equation}\label{lyrx:containment-core-rank}
 B\le L_2+f+(K_1+K_2-t+\nu)_+.
\end{equation}
The transposed statement is also valid. The corresponding
coefficientwise determinant bound is
\[
 v_p(\det A(X))\ge -(3+\epsilon)r-A_1-A_2
                  -L_2-f-(K_1+K_2-t+\nu)_+.
\]
\end{theorem}

\begin{proof}
The nested root family gives the containments
\[
 U_1:=Z_{\ell_1}\mathbb F_p[Y]_{<d}\subseteq W_{1,c},
 \qquad
 W_{2,c}\subseteq U_2:=Z_{\ell_2}\mathbb F_p[Y]_{<d}.
\]
On $C$ both are equalities. Let $T_0$ be the square pairing of
$\mathscr T_*$ on $U_1\times U_2$, independent of $c$.
If $T_0$ is singular, the total low rank is at most $L_2-|C|$.
Adjoining the $K_1+K_2$ top directions gives
\[
 B\le L_2+f+K_1+K_2-t
   \le L_2+f+(K_1+K_2-t+\nu)_+.
\]
This covers every singular reduction; the empty matrix is
understood to be invertible.

Suppose $T_0$ is invertible. Every local low pairing then has full
column rank, since its row space contains $U_1$ and its column
space is a subspace of $U_2$. The total low rank is $L_2$.
Let $F=T\setminus C$. Restrict the column top polynomial to
vanish on the $f$ points $c^2$ with $c\in F$. If $K_2\le f$,
this removes all column top directions, and $B\le L_2+K_2$
already suffices. Otherwise write
\[
 Q_2=V_FQ'_2,\qquad \deg Q'_2<K_2-f.
\]
By independent low adjustments, the remaining top column lifts
have zero full local feature on $F$. Eliminate all local low
columns against low rows. Residual row low radicals on $F$ now
annihilate all remaining columns; on $C$ the low matrix is square
and invertible. Thus only a top Schur matrix remains.

All classes in $C$ have the same fixed root pair, and hence the
same local top Schur scalar. The residual global entries are a
common scalar times
\[
 \sum_{c\in C}\omega_cV_M(c^2)^2c^{\pi_1+\pi_2}
          V_F(c^2)P_i(c^2)Q_j(c^2),
 \quad i<K_1,\quad j<K_2-f,
\]
where $\deg P_i\le i$ and $\deg Q_j\le j$.
The exact weight identity and the parity hypothesis turn the
numerator into $(c^2)^\nu$. Since $V_T=V_FV_C$, the factor
$V_F$ cancels to leave
\[
 \sum_{c\in C}\frac{(c^2)^\nu P_i(c^2)Q_j(c^2)}{V_C'(c^2)}.
\]
For $C\ne\varnothing$, Lagrange interpolation makes this zero
when $i+j+\nu\le |C|-2$. Its rank is at most
\[
 (K_1+(K_2-f)-|C|+\nu)_+=(K_1+K_2-t+\nu)_+.
\]
For $C=\varnothing$ the Schur matrix is zero. Adding the $L_2$
low pivots and restoring the $f$ removed columns proves the
claim. No local Schur scalar has been assumed nonzero.
\end{proof}

For a nonempty core the only possible choice is
$\ell_1=\max_T k_1$, $\ell_2=\min_T k_2$. Empty-core choices
are dominated by the generic low-plus-top cap. With normalized
class lengths, the production candidate is
\begin{equation}\label{lyrx:containment-core-rate}
 E_{\mathrm{core}}=(3+\epsilon)\delta+A_1+A_2
 +\int_T(D_2-k_2(x))\,dx
 +|T\setminus C|+(K_1+K_2-t)_+.
\end{equation}
It is evaluated in both orientations and combined with every
inherited candidate. All terms are affine on the ordinary root
partition apart from the displayed positive part.

In particular, suppose $\epsilon=1$, $D_1=D_2=D$,
$m_1=ap+\tau_1$ with $(p+1)/2\le\tau_1<p$, and
$m_2=(a+1)p+\tau_2$ with $0\le\tau_2\le(p-1)/2$.
Take $\ell_1=\ell_2=2a+2$ and $d=D-2a-2\ge1$.
Writing $u=p-\tau_1$ and $v=\tau_2$, the exact class counts are
\[
 L_2=dt-\min\{t,(v-1)_+\},\qquad
 f=\min\{t,\max(u,v-1)\}.
\]
Both orderings of $u$ and $v$ are included. In the limiting
formula, the minus-one endpoint constants disappear.

\subsection{The leading-layer containment core}

The same core argument also applies before taking a maximal
restriction kernel.  Work on the maximal class set $M$ of size
$w$, with unrestricted top widths $H_i$, and assume
$D_2-D_1\equiv\epsilon\pmod2$.  Choose roots with
\[
 k_1(c)\le\ell_1,\qquad
 \ell_2\le k_2(c)\le D_2\quad(c\in M),\qquad
 d=D_1-\ell_1=D_2-\ell_2\ge0.
\]
Set
\[
 C=\{c\in M:k_1(c)=\ell_1,\ k_2(c)=\ell_2\},\qquad
 f=|M\setminus C|,\qquad L_2=\sum_{c\in M}(D_2-k_2(c)).
\]
Then the leading rank satisfies
\begin{equation}\label{lyrx:leading-core-rank}
 r_0\le L_2+f+(H_1+H_2-w+\nu)_+,
\end{equation}
where $\nu=2$ for $\epsilon=0$ and both $D_i$ even, and
$\nu=3$ in the other parity-compatible cases.  No matrix
nonvanishing hypothesis is needed.

Indeed, a singular fixed low core loses one rank on each class
of $C$, giving the bound by adjoining all top directions.  In
the invertible case every low map has full column rank.  Force
the column top polynomial to vanish on $M\setminus C$, eliminate
the low columns, and retain the common top Schur scalar on $C$.
The leading weights are proportional to $c^4/V_M'(c^2)$ when
$\epsilon=0$ and to $c^5/V_M'(c^2)$ when $\epsilon=1$.
The parity condition makes the numerator a power $(c^2)^\nu$.
The imposed zero factor cancels from $V_M'$ on $C$, and the
same Lagrange moment calculation proves
\eqref{lyrx:leading-core-rank}.

For a nonempty useful core choose
$\ell_1=\max_M k_1$, $\ell_2=\min_M k_2$.  The limiting
candidate, also evaluated after transposition, is
\begin{equation}\label{lyrx:leading-core-rate}
 E_{\mathrm{leading\ core}}=(4+\epsilon)\delta+L_2+f
                              +(H_1+H_2-w)_+.
\end{equation}
\section{Odd heads across a maximal root transition}
\label{sec:odd-max-transition}

This section extends the upper-half third-layer theorem into a region
where the maximal column root changes before the maximal row root.
The saturated statement has a deliberately conservative finite loss of
six.  All such constants disappear from the arithmetic rate at a fixed
local cutoff.  We first treat the upper transition and then give separately proved
lower-half families.

\subsection{Notation and fixed local certificates}

Let $\epsilon=0$, $N=qp+w$, $n=(p-1)/2$, and write
\[
 m_i=ap+\tau_i,\qquad p/2<\tau_1\le\tau_2<p,
 \qquad u_i=p-\tau_i.
\]
Assume both endpoint degrees are $D+1$, where $D$ is even, and put
\[
 e=D-2a\ge2,\qquad z=e/2-1,\qquad
 J=\frac{(D+1)p+1}{2},\qquad H_i=R_i-J,
 \qquad k_f=(D/2-a)p+1-r.
\]
The new geometric region is
\begin{equation}
 u_2<w\le u_1.
 \label{eq:odd-max-transition-region}
\end{equation}
For a class set $S$, let $V_S(X)=\prod_{c\in S}(X-c^2)$.  Set
\[
 \begin{aligned}
 A&=\{1,\ldots,u_2\},& B&=\{u_2+1,\ldots,w\},\\
 T_B&=\{w+1,\ldots,u_1\},&T_C&=\{u_1+1,\ldots,n\},
 \end{aligned}
 \qquad b=|B|=w-u_2.
\]
Thus $M=A\cup B$ is the maximal set.  Its root pairs are $(W,W)$
on $A$ and $(W,V)$ on $B$, where
\[
 W=Z_{2a+1},\qquad V=Z_{2a+2}.
\]
The nonmaximal pairs are $(W,V)$ on $T_B$ and $(V,V)$ on $T_C$.

Write $T_{\max}=\mathscr T_{q,0}$ and $T=\mathscr T_{q-1,1}$.
Retain the three central-root matrix hypotheses: $T_{\max}$ and $T$
are nonsingular on the full $W$ space of degree at most $D$, and $T$
is nonsingular on the $V$ space of degree at most $D-1$.  In addition,
require nonsingularity of the square cross pairing
\begin{equation}
 T_{\max}:\ 
 W\mathbb Q[Y]_{\le D-\deg W}\ \times\
 V\mathbb Q[Y]_{\le D+1-\deg V},
 \qquad\text{dimension }e\text{ on each side}.
 \label{eq:odd-max-cross-test}
\end{equation}
Let $R$ be the monic centered-even degree-$D$ radical of $T$ on the
full $V$ space.  Let $\psi$ be the monic degree-$(D+1)$ polynomial
with root $W$, orthogonal under $T_{\max}$ to the full degree-$D$
$W$ space.  Define
\[
 \nu=T(YR^2),\qquad \theta=T_{\max}(\psi^2),\qquad
 \kappa=\theta-\nu.
\]
We conservatively require $\theta\ne0$ and $\kappa\ne0$.  The
finite-prime statements use the same hypotheses after reduction,
together with the inherited good-prime precision conditions.
Nonzero rational certificates exclude only finitely many fixed primes.
No nonvanishing of $\nu$ is required.

\subsection{A saturated right odd head}

Assume
\begin{equation}
 w\ge3,\quad u_2\ge3,\quad 0\le H_i\le n,\quad
 H_2\ge w,\quad k_f\ge0.
 \label{eq:odd-max-saturated-guards}
\end{equation}
Then
\begin{equation}
 \boxed{\quad
 v_p\bigl(\det(p^5A)\bigr)
 \ge r-ew+2k_f+2z-6.
 \quad}
 \label{eq:odd-max-saturated-bound}
\end{equation}
Consequently the limiting candidate is
\begin{equation}
 \boxed{\quad E_3=4\delta+ew-2k,\qquad
 k=(D/2-a)\beta-\delta.\quad}
 \label{eq:odd-max-saturated-rate}
\end{equation}
Here class counts and head widths in the rate are normalized by $N$.
The explicit rate guards are
\[
 u_2<w\le u_1,\qquad H_2\ge w,\qquad
 0\le H_i\le\beta/2,\qquad k\ge0,
 \qquad u_i=\beta-\tau_i.
\]
The old condition $w\le\beta-\tau_2$ is therefore extended to
$w\le\beta-\tau_1$ on this saturated branch.

\paragraph{First pivots and complete-even second parameters.}
Each complete even band has top width $n+1$.  Its full permitted
leading features can be prescribed independently on $M$.  The odd
column head contains the $b$-dimensional interpolation space
\[
 \mathcal I_B=\{S:\deg S<w,\ S|_A=0\},
\]
whose evaluations on $B$ are arbitrary.  These odd columns fill the
one missing column feature at every class of $B$.
The cross test~\eqref{eq:odd-max-cross-test} and the old square test
therefore supply $ew$ first pivots, all using even row features.
The remaining odd columns have coefficient polynomial divisible by
$V_B$.  The odd row head retains its $H_1$ parameters, after its
leading $B$ features have been adjusted using even row features.

The remaining leading odd form is supported on $A$.  Up to a common
unit it is $\theta B_A$, where
\begin{equation}
 B_S(x,y)=\sum_{c\in S}\frac{c^6}{V_M'(c^2)}x_cy_c.
 \label{eq:odd-max-boundary}
\end{equation}
Since the column head contains $V_B\mathbb F_p[X]_{<u_2}$, its
values on $A$ are arbitrary.  Thus the additional first rank is
$\ell=\min(H_1,u_2)$, and the $\ell$ paired odd columns may all be
chosen with coefficient polynomials of degree less than $w$.
The first rank is $ew+\ell$.

Initially retain even vectors whose entire leading feature is zero
on $M$.  The nonmaximal radical argument gives top polynomials
\begin{equation}
 P_1(c)=V_M(c^2)Q_1(c^2),\qquad
 P_2(c)=V_M(c^2)V_{T_B}(c^2)Q_2(c^2),
 \label{eq:odd-max-even-parameters}
\end{equation}
where
\[
 \deg Q_1<d_1=n+1-w=k_f+H_2-b,\qquad
 \deg Q_2<d_2=n+1-u_1=k_f+H_1.
\]
The left local feature is $P_1R$ on $T_B\cup T_C$; the right
feature is zero on $T_B$ and $P_2R$ on $T_C$.  Each side also has
$z$ independent zero-class parameters.  These are even second
radicals against all even first-kernel directions.

After the $\ell$ leading odd pivots, impose all remaining odd
order-$p$ equations on these even-only spaces.  Their numbers are
at most $H_2-b-\ell$ and $H_1-\ell$, respectively.  Hence true
second-kernel subspaces have dimensions at least
\begin{equation}
 d_1+z-(H_2-b-\ell)=k_f+\ell+z,
 \qquad d_2+z-(H_1-\ell)=k_f+\ell+z.
 \label{eq:odd-max-second-dimensions}
\end{equation}
No exact rank or surjectivity of these equations is asserted.
The arbitrary order-$p$ odd/odd block is harmless because every
actual residual odd pairing of the selected even vectors is imposed.

\paragraph{The exceptional maximal pivots and their lift correction.}
Put $E_i(c)=P_i'(c)/c$, with the derivative taken in $c$.
The first local jet has leading degree-$(D+1)$ coefficient $P_i'(c)$.
For one common nonzero unit $C_p$, the nonmaximal and maximal weights
are
\[
 C_p\frac{c^5}{V_M(c^2)},\qquad
 \frac{C_p}{2}\frac{c^4}{V_M'(c^2)}.
\]
The inherited derivative-sum calculation gives the raw nonmaximal
third form
\begin{equation}
 C_p\nu\sum_{c\notin M}\frac{d}{dc}
 \left(\frac{c^5P_1(c)P_2(c)}{V_M(c^2)}\right)
 =-\frac{C_p\nu}{2}B_M(E_1,E_2).
 \label{eq:odd-max-raw-third}
\end{equation}
Its degree guard holds because both even widths are $n+1$ and
$w\ge3$.

For a fixed right even vector, let $S_B$ be the unique polynomial
of degree less than $w$ with
\begin{equation}
 S_B|_A=0,\qquad S_B(c^2)=E_2(c)\quad(c\in B).
 \label{eq:odd-max-interpolant}
\end{equation}
The odd coefficient $cS_B(c^2)$ equals $P_2'(c)$ on $B$.
The first column jet there has root $V$ and degree at most $D+1$;
it belongs to the column space in~\eqref{eq:odd-max-cross-test}.
Choose the lower even features of the corresponding first-pivot
combination so its leading feature on $B$ equals that entire jet,
and its leading feature on $A$ is zero.  Subtracting $p$ times this
combination from the right lift cancels its first jet on $B$.
The third Schur correction cancels the entire raw maximal term on
$B$.  Its additional nonmaximal contribution is
\begin{equation}
 -C_p\nu\sum_{c\notin M}c^6Q_1(c^2)S_B(c^2).
 \label{eq:odd-max-nonmax-correction}
\end{equation}
Only the odd leading coefficient contributes: all nonmaximal column
roots are $V$, and $R$ annihilates their degree-at-most-$D$ parts.

The exact half-field identity is
\begin{equation}
 \sum_{c=1}^{n}c^6Q_1(c^2)S_B(c^2)
 =n[X^{n-3}](Q_1S_B).
 \label{eq:odd-max-coefficient}
\end{equation}
Indeed $\deg(Q_1S_B)\le n-1$, so $X^3Q_1S_B$ has zero constant
coefficient and degree at most $n+2<2n$.  Only $X^n$ survives.
Moreover, $E_1(c)=2V_M'(c^2)Q_1(c^2)$ on $M$, whence
\[
 \sum_{c\in B}c^6Q_1(c^2)S_B(c^2)=\tfrac12B_B(E_1,E_2).
\]
Thus~\eqref{eq:odd-max-nonmax-correction} cancels precisely the
$B$ part of~\eqref{eq:odd-max-raw-third}.  The ordinary maximal
Schur contribution on $A$ has coefficient $C_p\theta/2$.
After the $ew$ first pivots, the full third form is
\begin{equation}
 \frac{C_p\kappa}{2}B_A(E_1,E_2)
 -C_p\nu n[X^{n-3}](Q_1S_B).
 \label{eq:odd-max-third-form}
\end{equation}
The coefficient term depends on at most the three highest
coefficients of $Q_1$: its contributing degrees are at least
$n-w-2$, while $\deg Q_1\le n-w$.

\paragraph{Remaining odd pivots and boundary constraints.}
For any residual odd column $S$, divisible by $V_B$, the even/odd
order-$p$ form is
\begin{equation}
 U(Q_1,S)=\frac{C_p\kappa}{2}B_A(E_1,S)
 +C_p\nu n[X^{n-3}](Q_1S).
 \label{eq:odd-max-even-odd}
\end{equation}
The half-field sum has only this one coefficient because
$H_2\le n$ and $w\ge3$.  The $\ell$ leading odd pivot columns
have degree less than $w$, so their row-side factors in the next
Schur correction use the same $A$-boundary values and the same
three high coefficients.  Hence the third form after all first
pivots still factors through these $u_2+3$ row coordinates.

The low column space $V_B\mathbb F_p[X]_{<u_2-3}$ has dimension
$u_2-3$.  Its intersection with the leading right odd kernel has
dimension at least $u_2-3-\ell$.  Each such polynomial has degree
less than $w-3$, so the coefficient term in
\eqref{eq:odd-max-even-odd} vanishes.  Since $\kappa\ne0$, the
residual second equations impose its boundary orthogonality.
Writing $S=V_BS'$ cancels $V_B$ in the denominator and leaves
\[
 \sum_{c\in A}\frac{c^6}{V_A'(c^2)}E_1(c)S'(c^2).
\]
Distinct nodes and nonzero weights show that a $j$-dimensional
subspace of these $S'$ gives $j$ independent constraints on the
ambient $A$-boundary space.  The selected left second-kernel image
therefore has dimension at most $\ell+3$ on $A$.  Including the
three high coefficients, the restricted third rank is at most
$\ell+6$.

The selected spaces are true second kernels.  Their couplings to
second pivots begin at order $p^2$, so further Schur products begin
at $p^3$ and do not alter this third form.  The restricted-kernel
valuation lemma now gives
\[
 v_p(\det(p^5A))
 \ge r-ew-\ell+2(k_f+\ell+z)-(\ell+6),
\]
which proves~\eqref{eq:odd-max-saturated-bound}.

\paragraph{Finite exceptions and passage to the rate.}
At a fixed $q$ cutoff there are only finitely many tuples $(q,D,a)$.
The fixed rational matrix and scalar tests, together with the
inherited precision cutoff, exclude finitely many fixed primes.
The $+1$ in $k_f$, the zero-class constant $2z$, and the loss six
are bounded on that tuple set.  Their prime sums divided by $N^2$
vanish.  Thus~\eqref{eq:odd-max-saturated-rate} is a valid additional
limiting candidate.  Finite Smith computations provide diagnostics
only; no sampled rank is used as a proof hypothesis.

\subsection{Unsaturated heads in the same transition region}
\label{subsec:odd-max-unsaturated}

Keep the geometry, fixed matrices, and scalar tests above.  Suppose
\begin{equation}
 w\ge3,\quad u_2\ge3,\quad k_f\ge0,\qquad
 0\le H_1\le u_2-3,\qquad b\le H_2\le w-3.
 \label{eq:odd-max-unsaturated-guards}
\end{equation}
The extra inequalities $H_2\ge b$ and $H_1\le u_2-3$ follow from
$H_1\ge0$, $H_2\le w-3$, and the geometry, since
\[
 H_2-H_1=u_1-u_2,\qquad H_2-b=H_1+u_1-w,\qquad
 u_2-H_1=u_1-H_2.
\]
Put $\ell=(H_1+H_2-w+3)_+$.  Then
\begin{equation}
 \boxed{\quad
 v_p(\det(p^5A))\ge
 r-(e+1)w+2k_f+2z+H_1+H_2-\ell.
 \quad}
 \label{eq:odd-max-unsaturated-bound}
\end{equation}
The limiting candidate is the inherited unsaturated expression
\begin{equation}
 \boxed{\quad E_3=4\delta+ew-2k+(w-H_1-H_2)_+.\quad}
 \label{eq:odd-max-unsaturated-rate}
\end{equation}
Its explicit rate guards are
\[
 u_2<w\le u_1,\qquad 0\le H_1\le u_2,\qquad
 w-u_2\le H_2\le w,\qquad k\ge0,
\]
together with the common upper-half, odd-endpoint, and local
certificate hypotheses above.  All geometric equalities may be
assigned to either adjacent cell; their endpoints have zero measure
in the arithmetic integral.
This statement does not require the old surjectivity condition
$k_f+\min(H_1,H_2)\ge w$.

\paragraph{A smaller interpolation space.}
Fill the $b$ missing maximal column features using odd polynomials
of degree less than $b$.  Their values on $B$ are arbitrary, and
$H_2\ge b$ makes them available.  On $A$, choose their lower even
features so the leading feature is $cS(c^2)\psi$.  This is
orthogonal to every even row feature there.  On $B$ their full
features span the degree-at-most-$(D+1)$ $V$ space.  The same cross
test therefore gives $ew$ first pivots using even row directions.
Residual odd columns have coefficient divisible by $V_B$ and
parameter dimension $H_2-b$.

For a right even vector, define $S_B$ by
\[
 \deg S_B<b,\qquad S_B(c^2)=E_2(c)\quad(c\in B),
 \qquad E_2^*=E_2-S_B.
\]
Thus $E_2^*$ vanishes on $B$; the values on $A$ are the new right
boundary coordinates.  In contrast to
\eqref{eq:odd-max-interpolant}, $S_B$ need not vanish on $A$.

An odd row may be lifted with leading feature $cS_1(c^2)\psi$ on
every maximal class.  Its pairing with each new $B$-pivot is
proportional to $B_M(S_1,S)$, where $\deg S<b$.  This is zero:
the numerator degree is at most
\[
 (H_1-1)+(b-1)+3\le w-2,
\]
and the Cauchy moment identity against $V_M'$ vanishes through
degree $w-2$.  These odd rows therefore already annihilate all
$ew$ first pivots.  In particular,
\begin{equation}
 B_M(S_1,S_B)=0.
 \label{eq:odd-max-low-cauchy-zero}
\end{equation}

\paragraph{The simultaneous boundary form.}
The complete-even parameters remain
\eqref{eq:odd-max-even-parameters}.  The $B$-pivot lift cancels the
entire first column jet on $B$; on $A$ it subtracts the odd
coefficient $cS_B$.  Hence the maximal contribution after the even
projection is $(C_p\theta/2)B_A(E_1,E_2-S_B)$.
Its additional nonmaximal contribution is
\[
 -C_p\nu\sum_{c\notin M}c^6Q_1(c^2)S_B(c^2)
 =-C_p\nu n[X^{n-3}](Q_1S_B)+\frac{C_p\nu}{2}B_M(E_1,S_B).
\]
Combining this with~\eqref{eq:odd-max-raw-third} gives
\[
 \frac{C_p\kappa}{2}B_A(E_1,E_2^*)
 -C_p\nu n[X^{n-3}](Q_1S_B).
\]
The coefficient term vanishes identically, since
\[
 \deg(Q_1S_B)\le n-w+b-1=n-u_2-1\le n-4.
\]
Thus the even/even third form is exactly
$(C_p\kappa/2)B_A(E_1,E_2^*)$.

For a residual odd column $S_2$, divisibility by $V_B$ makes its
boundary values zero on $B$, and $\deg S_2\le w-4$ kills the
full-field coefficient in~\eqref{eq:odd-max-even-odd}.  The
order-$p$ even/odd form is therefore
$(C_p\kappa/2)B_A(E_1,S_2)$.
The opposite order-$p$ form is initially
$(C_p\kappa/2)B_M(S_1,E_2)$.  Its coefficient term vanishes because
$\deg(P_2/V_M)\le n-w$ and $\deg S_1\le u_2-4\le w-4$.
Equation~\eqref{eq:odd-max-low-cauchy-zero}, together with
$E_2^*|_B=0$, rewrites it as
$(C_p\kappa/2)B_A(S_1,E_2^*)$.

The leading odd/odd form is $(C_p\theta/2)B_A(S_1,S_2)$.
After removal of the common unit $C_p/2$, the simultaneous normal
form is consequently
\[
 \begin{array}{c|cc}
 &O_2&E_2^*\\\hline
 O_1&\theta B_A+O(p)&p\kappa B_A+O(p^2)\\
 E_1&p\kappa B_A+O(p^2)&p^2\kappa B_A+O(p^3).
 \end{array}
\]
As before, nonmaximal even order-$p$ pivots change only the
unidentified odd/odd correction at order $p$, the displayed
mixed terms at order $p^2$, and the even/even term at order $p^3$.
Thus the displayed orders are unaffected.

\paragraph{Flags and restricted third rank.}
The ambient boundary space has dimension $u_2$, with nondegenerate
form $B_A$.  The row flag $U$ is the evaluation of
$\mathbb F_p[X]_{<H_1}$; the column flag $V$ is the evaluation of
$V_B\mathbb F_p[X]_{<H_2-b}$.  Their dimensions are $H_1$ and
$H_2-b$, each at most $u_2-3$.  Factoring $V_B$ cancels it from
$V_M'$; the pairing matrix is
\[
 \left(\sum_{c\in A}\frac{(c^2)^{h+j+3}}{V_A'(c^2)}\right).
\]
Its first nonzero anti-diagonal is $h+j=u_2-4$, with unit entries.
Its exact rank is therefore
\[
 \ell=(H_1+(H_2-b)-u_2+3)_+=(H_1+H_2-w+3)_+.
\]
The first rank is $ew+\ell$.  Choose paired leading flags
$L_1,L_2$ of dimension $\ell$, and put
\[
 N_1=U\cap V^\perp,\qquad N_2=V\cap U^\perp.
\]
Their dimensions are $H_1-\ell$ and $H_2-b-\ell$.
Imposing their actual mixed order-$p$ equations on the even spaces
gives true selected second kernels of dimension at least
$k_f+\ell+z$ on both sides, by
\eqref{eq:odd-max-second-dimensions}.  Their boundary values lie in
$N_2^\perp$ and $N_1^\perp$, respectively.

After the $\ell$ leading odd pivots the third form is
\[
 \kappa B_A-
 \frac{\kappa^2}{\theta}
 B_A(\,\cdot\,,L_2)
 B_A(L_1,L_2)^{-1}B_A(L_1,\,\cdot\,).
\]
On the selected boundary spaces it has row radical containing
$N_1$ and column radical containing $N_2$.  Its rank is at most
\[
 u_2-(H_2-b-\ell)-(H_1-\ell)
 =w-H_1-H_2+2\ell.
\]
This ambient bound remains valid without surjectivity of either
actual boundary map.  Second-pivot elimination again changes
only order $p^3$ and higher.  The restricted-kernel lemma yields
\[
 v_p(\det(p^5A))\ge
 r-ew-\ell+2(k_f+\ell+z)-(w-H_1-H_2+2\ell),
\]
which is~\eqref{eq:odd-max-unsaturated-bound}.  The finite-exception
and fixed-cutoff arguments above give
\eqref{eq:odd-max-unsaturated-rate}.

\subsection{Both saturated heads beyond the outer upper boundary}
\label{subsec:odd-max-outer}

A separate optional extension covers $w>u_1$.  Retain the common
upper-half, equal odd degree, and fixed local hypotheses above, and
add nonsingularity of $T_{\max}$ on
\[
 G^+=V\mathbb Q[Y]_{\le D+1-\deg V},
 \qquad\dim G^+=e.
\]
Assume $u_2\ge3$, $w\ge3$, $u_1<w\le n$,
$w\le H_1,H_2\le n$, and $k_f\ge0$.  We deliberately set all
free zero-class parameters to zero in this optional theorem.  Then
\begin{equation}
 \boxed{\quad v_p(\det(p^5A))\ge r-ew+2k_f-12.\quad}
 \label{eq:odd-max-outer-finite}
\end{equation}
The rate is unchanged:
\begin{equation}
 \boxed{\quad E_3=4\delta+ew-2k,\quad
 u_1<w\le\beta/2,\quad w\le H_1,H_2\le\beta/2,
 \quad k\ge0.\quad}
 \label{eq:odd-max-outer-rate}
\end{equation}
All common quotient, parity, and local certificate guards remain.
This is an upper-half statement only.

\paragraph{A local rank-one projection.}
The defining rational functions give the exact Christoffel identity
\[
 T(P)=T_{\max}((Y-q)P):
\]
the maximal denominator is $(Y-q)$ times the nonmaximal one, so
both quotient and residue expressions agree.  Let
$G=V\mathbb F_p[Y]_{\le D-\deg V}$ and
$\psi_V=(Y-q)R$.  Then
\[
 T_{\max}(\psi_V F)=0\quad(F\in G),\qquad
 T_{\max}(\psi_V^2)=\nu.
\]
For the top functional $t(F)=[Y^{D+1}]F$, this implies
$T_{\max}(\psi_V,F)=\nu t(F)$ for all $F\in G^+$.
Nonsingularity on $G^+$ implies $\nu\ne0$.  Consequently
\begin{equation}
 K=T_{\max}-\nu\,t\otimes t
 \qquad\Longrightarrow\qquad
 K T_{\max}^{-1}K=K\quad\text{on }G^+.
 \label{eq:odd-max-projection-identity}
\end{equation}
This removes precisely the $\psi_V$ summand.

\paragraph{Pivots and selected kernels.}
Now split $M$ into
\[
 A=\{1,\ldots,u_2\},\quad
 B=\{u_2+1,\ldots,u_1\},\quad
 C=\{u_1+1,\ldots,w\}.
\]
Their maximal root pairs are $W/W$, $W/V$, and $V/V$.
Use odd row interpolation on $C$ and odd column interpolation on
$B\cup C$, with polynomials of degree less than $w$ and zero
values on the other maximal nodes.  The saturated heads supply
these spaces.  Together with the complete even bands they give
$ew$ first pivots with local spaces
\[
 F\times F\text{ on }A,\qquad
 F\times G^+\text{ on }B,\qquad
 G^+\times G^+\text{ on }C,
\]
where $F$ is the full degree-$D$ $W$ space.
Residual odd row heads vanish on $C$, and column heads vanish on
$B\cup C$.  Both can prescribe arbitrary values on $A$, so the
remaining first rank is $\ell=u_2$.  These additional odd pivots
may also be chosen with degree less than $w$ and supported on $A$.

All nonmaximal roots are now $V/V$.  The selected even parameters
are $P_i(c)=V_M(c^2)Q_i(c^2)$, with
$\deg Q_i<n+1-w$.  Their local features are $P_iR$ outside $M$
and zero on $M$.  Set the free zero-class coordinates to zero.
After all first pivots, impose the actual remaining odd equations.
Their counts are at most $H_2-w$ on the left and
$H_1-(w-u_1)-u_2$ on the right.  The identities
$k_f+H_i=n+1-u_i$ show that both selected true second kernels have
dimension at least $k_f+u_2$.

\paragraph{First-jet coefficient errors have bounded rank.}
The raw nonmaximal derivative sum converts the uncorrected third
form into the weighted sum of the local $K(f_1,g_1)$, with $K$ as
in~\eqref{eq:odd-max-projection-identity}.
For an initial pivot with odd coefficient $S$, $\deg S<w$, its
nonmaximal order-$p$ coupling to a selected even row is
\[
 C_p\nu\sum_{c\notin M}c^6Q_1(c^2)S(c^2)
 =C_p\nu n[X^{n-3}](Q_1S)-\frac{C_p\nu}{2}B_M(E_1,S).
\]
The analogous right coupling has the same formula with $Q_2$.
The coefficient term uses at most the three highest coefficients
of $Q_i$.  Thus each first-jet coupling is its weighted local
$K$ coupling plus a matrix of rank at most three.

First omit these two bounded-rank remainders.  On $A$, removing
$F\times F$ leaves the usual $\kappa$ top pairing.  On $B$, the
right first jet belongs to the full pivot space $G^+$, and
$t|_F=0$; the cross inverse cancels $K(f_1,g_1)$ exactly.
On $C$, identity~\eqref{eq:odd-max-projection-identity} makes the
third Schur term zero.  Restoring the two remainder matrices
changes the third form by rank at most six: for a pivot matrix
$P$, the difference between
$(X+R_1)P^{-1}(Y+R_2)$ and $XP^{-1}Y$ is
\[
 R_1P^{-1}(Y+R_2)+XP^{-1}R_2,
 \qquad\operatorname{rank}R_i\le3.
\]
After the initial $ew$ pivots, the third form is therefore the
$A$ boundary form plus a form of rank at most six.

The $u_2$ remaining leading odd pivots are supported on $A$ with
degree-less-than-$w$ coefficient polynomials.  Their mixed
order-$p$ pairings are again the pure $\kappa$ boundary pairings
plus rank-at-most-three coefficient remainders on each side.
Their pure Schur correction is supported on the $u_2$ boundary
coordinates; restoring its errors adds rank at most six.  Thus
all first-pivot corrections together give third rank at most
$u_2+12$.  The chosen spaces are true second kernels, so second
pivot corrections start only at order $p^3$.  The valuation lemma
now gives
\[
 r-ew-u_2+2(k_f+u_2)-(u_2+12),
\]
proving~\eqref{eq:odd-max-outer-finite}.  The omitted zero-class
directions, bounded loss twelve, and finite exceptional primes
have no effect on~\eqref{eq:odd-max-outer-rate}.

\subsection{A paid one-sided lower transition}
\label{subsec:odd-lower-paid}

This is a separate lower-half result.  Let
\[
 m_i=ap+\tau_i,\qquad 0<\tau_1\le\tau_2<p/2,\qquad
 v_i=\tau_i-1,\qquad v_1<w\le v_2.
\]
Both endpoint degrees are $D+1$, with $D$ even.  Set
\[
 e=D-2a\ge2,\quad H_i=R_i-\frac{(D+1)p+1}{2},\quad
 k_f=(D/2+a+1)p+2-r-m_1-m_2,\quad b=w-v_1.
\]
Assume $v_1\ge3$, $w\ge3$, $0\le H_i\le n$, and $k_f\ge0$.
Here $W=Z_{2a+1}$ and $V=Z_{2a}$, so the $W$ spaces are contained
in the corresponding $V$ spaces.  Retain the three lower
central-root tests and $\theta,\kappa\ne0$, with $R,\nu,\psi$
defined using this lower $V$.  The extra fixed test is
\begin{equation}
 T_{\max}:\quad
 V\mathbb Q[Y]_{\le D-\deg V}\ \times\
 W\mathbb Q[Y]_{\le D+1-\deg W},
 \qquad\text{dimension }e+1\text{ on each side}.
 \label{eq:odd-lower-cross-test}
\end{equation}
Set every free zero-class coordinate to zero.  If $H_2\ge w$, then
\begin{equation}
 v_p(\det(p^5A))\ge r-ew-b+2k_f-6.
 \label{eq:odd-lower-saturated-finite}
\end{equation}
If $b\le H_2\le w-3$ and $0\le H_1\le v_1-3$, put
$\ell=(H_1+H_2-w+3)_+$.  Then
\begin{equation}
 v_p(\det(p^5A))\ge
 r-(e+1)w-b+2k_f+H_1+H_2-\ell.
 \label{eq:odd-lower-unsaturated-finite}
\end{equation}
The two extra unsaturated inequalities are automatic from the
geometry and $H_2\le w-3$, because
\[
 H_2-H_1=v_2-v_1,\qquad H_2-b=H_1+v_2-w,\qquad
 v_1-H_1=v_2-H_2.
\]

The common limiting candidate is
\begin{equation}
 \boxed{\quad
 E_3=4\delta+ew-2k+(w-\tau_1)+(w-H_1-H_2)_+,
 \qquad k=(D/2+a+1)\beta-\delta-s_1-s_2.
 \quad}
 \label{eq:odd-lower-paid-rate}
\end{equation}
Its explicit rate guards are the common lower-half quotient and
equal odd endpoints, $\tau_1<w\le\tau_2$, $k\ge0$, and either
$H_2\ge w$, or
\[
 w-\tau_1\le H_2\le w,\qquad 0\le H_1\le\tau_1.
\]
The fixed local tests above are mandatory.  No claim about
$w>\tau_2$ is included in this subsection.

\paragraph{Geometry, pivots, and counts.}
The maximal roots are $W/W$ on $A=\{1,\ldots,v_1\}$ and $V/W$
on $B=\{v_1+1,\ldots,w\}$.  The nonmaximal roots are $V/W$ on
$T_B=\{w+1,\ldots,v_2\}$ and $V/V$ thereafter.
On $B$ the complete even row space has dimension $e+1$, whereas
the even column space has dimension $e$.  Fill its missing column
feature using $b$ odd column interpolants.  Use degree-less-than-$w$
interpolants vanishing on $A$ in the saturated case, or
 degree-less-than-$b$ interpolants in the unsaturated case.
Test~\eqref{eq:odd-lower-cross-test} supplies $e+1$ pivots at each
class of $B$, while the usual square test supplies $e$ on $A$.
Thus the initial first rank is $ew+b$, and all these pivots use
even row directions.

Residual odd columns are divisible by $V_B$.  The leading odd
boundary rank is $\ell=\min(H_1,v_1)$ in the saturated case and
the displayed $\ell$ in the unsaturated case.  The nonmaximal
even radical parameters have
\[
 P_1=V_MQ_1,\qquad P_2=V_MV_{T_B}Q_2,
\]
with the arguments $c^2$ understood, and
\[
 \deg Q_1<d_1=n+1-w=k_f+H_2-b,\qquad
 \deg Q_2<d_2=n+1-v_2=k_f+H_1.
\]
The left local feature is $P_1R$ throughout the nonmaximal set;
the right feature is zero on $T_B$ and $P_2R$ thereafter.
After all leading odd pivots, imposing every residual odd equation
leaves true selected second kernels of dimension at least
\[
 d_1-(H_2-b-\ell)=k_f+\ell,
 \qquad d_2-(H_1-\ell)=k_f+\ell.
\]
No zero-class contribution or exact equation rank is used.

\paragraph{The one-sided correction in the reversed containment.}
The weights and derivative sum are unchanged.  On $B$, the first
column jet has root $W$ and degree at most $D+1$, so it lies in the
entire column pivot space of~\eqref{eq:odd-lower-cross-test}.
The same first-pivot lift removes its whole maximal contribution.
All first row pivots are even, so there is no additional odd-row
nonmaximal correction at this stage.

The nonmaximal coupling to the left radical remains
\[
 C_p\nu\sum_T c^6Q_1S.
\]
On $T_B$ the interpolant's root is $W$,
which is contained in the full $V$ space in this lower geometry;
on the remaining classes its root is $V$.  Thus $R$ annihilates
all degree-at-most-$D$ parts in both cases.  Conversely, the right
radical is zero on $T_B$, and its opposite row features have root
$V$ wherever it is nonzero.  This verifies both containment steps
explicitly.

All interpolation and coefficient identities of the accepted
one-sided proofs now apply.  In the saturated case the third form
factors through the $A$ boundary and three high $Q_1$ coefficients;
the low residual flag $V_B\mathbb F_p[X]_{<v_1-3}$ gives restricted
third rank at most $\ell+6$.  In the unsaturated case,
$\deg S_B<b$ makes the coefficient term zero and
$B_M(S_1,S_B)=0$.  The row odd lift $cS_1\psi$ is permitted even
on $B$, since $\psi$ has root $W\subset V$.  The transformed
right boundary $E_2-S_B$ therefore gives exactly the same four-block
$\theta,\kappa,\kappa$ normal form on $A$, with flag widths $H_1$
and $H_2-b$.  Its restricted third rank is
$w-H_1-H_2+2\ell$.

The first nullity is now $r-ew-b-\ell$.  Combining it with the
two selected kernel dimensions proves
\eqref{eq:odd-lower-saturated-finite} and
\eqref{eq:odd-lower-unsaturated-finite}.  Further second-pivot
corrections start at $p^3$.  Bounded finite losses, omitted
zero-class directions, and fixed exceptional primes disappear in
\eqref{eq:odd-lower-paid-rate}.

\subsection{Cauchy cancellation beyond the outer lower boundary}
\label{subsec:odd-lower-outer-cauchy}

This final lower-half family uses equal coefficients in all four
boundary blocks.  Retain the notation and fixed tests of
Subsection~\ref{subsec:odd-lower-paid}, and add nonsingularity of
$T_{\max}$ on the full $V$ space of degree at most $D$, of dimension
$e+1$.  Assume
\begin{equation}
 3\le v_1\le v_2<w\le n,\qquad k_f\ge0,\qquad
 0\le H_i\le n,\qquad H_1+H_2\le w-3.
 \label{eq:odd-lower-outer-guards}
\end{equation}
Every free zero-class coordinate is set to zero.  Then
\begin{equation}
 \boxed{\quad
 v_p(\det(p^5A))\ge r-ew-w+2(n+1-w).
 \quad}
 \label{eq:odd-lower-outer-bound}
\end{equation}
The limiting candidate is
\begin{equation}
 \boxed{\quad E_3=4\delta+(e+3)w-\beta.\quad}
 \label{eq:odd-lower-outer-rate}
\end{equation}
Its explicit rate guards are the common lower-half quotient, equal
odd endpoints, all stated local tests, and
\[
 \tau_2<w\le\beta/2,\qquad 0\le H_i\le\beta/2,
 \qquad H_1+H_2\le w,\qquad k\ge0,
 \quad k=(D/2+a+1)\beta-\delta-s_1-s_2.
\]
On this region the cost equals
\[
 4\delta+ew-2k+(w-H_1-H_2)+(w-\tau_1)+(w-\tau_2).
\]
There is no additional lower bound on $H_1$.

\paragraph{Initial pivots and two maximal odd scalars.}
Split $M$ into
\[
 A=\{1,\ldots,v_1\},\quad
 B=\{v_1+1,\ldots,v_2\},\quad
 C=\{v_2+1,\ldots,w\}.
\]
Put $a_0=v_1$, $b=v_2-v_1$, $c_0=w-v_2$, and $b_1=w-v_1$.
Their roots are $W/W$, $V/W$, and $V/V$.  All nonmaximal roots
are $V/V$.  On $A$ use the even $W$ square block of dimension $e$;
on $B$ use the cross block~\eqref{eq:odd-lower-cross-test} of
dimension $e+1$, filling its missing column feature with odd
polynomials of degree less than $b$; on $C$ use the new even $V$
square block of dimension $e+1$.  The initial first rank is
$ew+b_1$, and every initial row pivot is even.

The exact identity $T(P)=T_{\max}((Y-q)P)$ shows that
$\psi_V=(Y-q)R$ is orthogonal to the full degree-$D$ $V$ space
and has maximal norm $\nu$.  Choose the lower features of the
$B$ odd pivot columns so their leading feature is $cS(c^2)\psi$
on $A$ and $cS(c^2)\psi_V$ on $C$.  These pair to zero with the
even row pivot spaces outside $B$.  Thus the leading pivot matrix
is block diagonal, although these odd coefficients are global.

Residual odd column polynomials have $S_2=V_BS_2'$.  Since
$H_2-H_1=v_2-v_1=b$, both remaining heads have parameter dimension
$H_1$: $\deg S_1<H_1$ and $\deg S_2'<H_1$.
Their leading form, up to the common unit $C_p/2$, is
\begin{equation}
 \theta B_A(S_1,S_2)+\nu B_C(S_1,S_2).
 \label{eq:odd-lower-outer-two-scalars}
\end{equation}
A row correction through the $B$ even pivots is supported on $B$
and does not change this expression, because $S_2|_B=0$.

\paragraph{The leading Cauchy zero.}
Factoring $V_B$ gives
\[
 B_{A\cup C}(S_1,S_2)
 =\sum_{c\in A\cup C}
 \frac{(c^2)^3S_1(c^2)S_2'(c^2)}{V_{A\cup C}'(c^2)}.
\]
The numerator has degree at most $H_1+H_2-b+1\le w-b-2$.
The Cauchy identity over the $w-b$ nodes makes it zero.
Therefore~\eqref{eq:odd-lower-outer-two-scalars} is exactly
$\kappa B_A(S_1,S_2)$.

Let $U$ and $V_0$ be the two evaluation flags on $A$.  Each has
dimension $h=\min(H_1,v_1)$, since $V_B$ is a unit on $A$.
Let $\ell=\operatorname{rank}(B_A|_{U\times V_0})$; no explicit
formula for this rank is needed.  The full first rank is
$ew+b_1+\ell$.

\paragraph{The even and mixed forms.}
The complete-even radical parameters are
\[
 P_i(c)=V_M(c^2)Q_i(c^2),\qquad \deg Q_i<d=n+1-w,
 \qquad E_i(c)=P_i'(c)/c.
\]
Choose the degree-less-than-$b$ polynomial $S_B$ with $S_B=E_2$
on $B$, and put $E_2^*=E_2-S_B$.  Its pivot lift cancels the
whole first column jet on $B$, subtracts $cS_B\psi$ on $A$, and
subtracts $cS_B\psi_V$ on $C$.
The nonmaximal correction is $-C_p\nu\sum_Tc^6Q_1S_B$.  Its
coefficient remainder vanishes because
\[
 \deg(Q_1S_B)\le n-w+b-1=n-(v_1+c_0)-1\le n-4.
\]
Combining it with the derivative boundary leaves
$(C_p\kappa/2)B_A(E_1,E_2^*)$ as the even/even third form:
the $C$ maximal scalar $\nu$ cancels its nonmaximal scalar, and
the $B$ first jet has been completely pivoted.

A residual odd column vanishes on $B$.  Its coefficient remainder
is zero since $H_2\le w-3$, so the even/odd order-$p$ form is
$(C_p\kappa/2)B_A(E_1,S_2)$.
For the opposite form, initially lift an odd row with $cS_1\psi$
on $A\cup B$ and $cS_1\psi_V$ on $C$.  Its raw maximal and
nonmaximal contributions combine to
$(C_p\kappa/2)B_{A\cup B}(S_1,E_2)$; the coefficient term is zero
because $H_1\le w-3$.  Its leading pairing with the $B$ pivot
combination for $S_B$ is
\[
 \frac{C_p\theta}{2}B_{A\cup B}(S_1,S_B)
 +\frac{C_p\nu}{2}B_C(S_1,S_B).
\]
Subtracting this full first Schur correction gives
\[
 \frac{C_p\kappa}{2}B_A(S_1,E_2^*)
 -\frac{C_p\nu}{2}B_M(S_1,S_B).
\]
The last term is zero by the Cauchy identity: $\deg S_B<b$ and
$H_1+b\le H_1+H_2\le w-3$.
After removing the common unit, the simultaneous form is therefore
\begin{equation}
 OO=\kappa B_A+O(p),\quad
 EO=OE=p\kappa B_A+O(p^2),\quad
 EE=p^2\kappa B_A+O(p^3).
 \label{eq:odd-lower-equal-coefficients}
\end{equation}
The usual nonmaximal even pivots change only higher remainders at
the displayed orders.

\paragraph{An equal-coefficient flag lemma.}
Choose paired flags $L_1,L_2$ of dimension $\ell$, and set
\[
 N_1=U\cap V_0^\perp,\qquad N_2=V_0\cap U^\perp.
\]
Each has dimension $h-\ell$.  Polynomial head directions vanishing
on $A$ give zero mixed pairings.  Thus the residual odd equations
have ranks at most $h-\ell$, and there are true selected even
second kernels of dimension at least $d-h+\ell$ on both sides.
Their boundaries lie in $N_2^\perp$ and $N_1^\perp$.

All three coefficients in~\eqref{eq:odd-lower-equal-coefficients}
are the same.  The third Schur form, up to its nonzero scalar, is
\[
 Q=B_A-B_A(\,\cdot\,,L_2)
 B_A(L_1,L_2)^{-1}B_A(L_1,\,\cdot\,).
\]
It annihilates $L_1$ on the left and $L_2$ on the right against the
entire ambient boundary space.  On the constrained product it also
annihilates $N_1$ and $N_2$.  Thus its row radical contains all of
$U=L_1+N_1$, of dimension $h$, within the domain $N_2^\perp$ of
dimension $v_1-h+\ell$.  Consequently
\[
 \operatorname{rank}Q\le v_1-2h+\ell.
\]
The right side is nonnegative by
$\ell\ge2h-v_1$, the elementary dimension inequality for two
$h$-dimensional subspaces of a nondegenerate $v_1$-dimensional
bilinear space.  Actual boundary maps need not be surjective.

The selected even vectors are true second kernels; the unknown
residual odd/odd term and further second pivots affect only order
$p^3$ and higher.  The valuation lemma gives
\[
 \begin{aligned}
 v_p(\det(p^5A))
 &\ge r-ew-b_1-\ell+2(d-h+\ell)-(v_1-2h+\ell)\\
 &=r-ew-w+2d,
 \end{aligned}
\]
proving~\eqref{eq:odd-lower-outer-bound}.  A negative displayed
lower bound for a selected dimension may be replaced by its actual
nonnegative dimension without weakening the conclusion.  Bounded
rounding terms, omitted zero-class directions, and fixed exceptional
primes disappear in~\eqref{eq:odd-lower-outer-rate}.
\section{A confluent second arithmetic layer}
\label{conf:section}

The feature-vanishing second layer discards every maximal local
feature.  In the epsilon-one regime this can discard genuine local
radicals of the first form.  Retaining these radicals gives a
second form whose top block is a sum of residues at simple and
double nodes.  The cancellation is uniform in the prime; finite
Smith computations are diagnostics only.

\subsection{Notation and the local matrix hypothesis}

Use the original window matrix $A(X)$ and write $R_i=m_i+r$.
Let $p>5$ be a good prime, put
\[
 N=qp+\sigma,\qquad n=(p-1)/2,
 \qquad \epsilon=1,\qquad
 w=\sigma-n,\qquad t=n-w.
\]
The maximal classes are $M=\{t+1,\ldots,n\}$ and the
nonmaximal classes are $T=\{1,\ldots,t\}$.  Retain the
good-prime feature and Newton--Bernoulli precision hypotheses.
For each fixed local tuple, the good-prime cutoff may be enlarged
so that all retained local degrees are less than $p-1$.
Assume $0\le m_i\le J_i$, where
\[
 D_i=\left\lfloor\frac{2R_i-2}{p}\right\rfloor,
 \qquad J_i=\lceil D_ip/2\rceil,
 \qquad H_i=R_i-J_i.
\]

Initially suppose that both $D_i$ are even.  Choose even integers
$\ell_i$ with
\begin{equation}\label{conf:equal-residual}
 d=D_1-\ell_1=D_2-\ell_2\ge0.
\end{equation}
Thus $d$ is even.  Put
\[
 V_i=Z_{\ell_i},\qquad
 U_i=V_i\mathbb Q[Y]_{<d},\qquad
 \Phi=\mathscr T_{q,1},\qquad \Psi=\mathscr T_{q,0}.
\]
The only local matrix hypothesis is that the $d\times d$ pairing
of $\Phi$ on $U_1\times U_2$ is nonsingular over $\mathbb Q$.
Primes dividing its denominators or determinant are excluded.
The empty matrix is invertible by convention.

\begin{lemma}[Parity form of the witness]
\label{conf:witness}
Let $d=2h$ and put $u=Y+1/2$.  In centered monomial bases, the
low matrix is, after grouping even and odd powers,
\[
 \begin{pmatrix}0&B\\B&0\end{pmatrix},\qquad
 B_{ab}=\Phi(V_1V_2u^{2a+2b+1}),\quad 0\le a,b<h.
\]
Its determinant is $(-1)^h\det(B)^2$.  A nonsingular reduction
of $B$ at one auxiliary prime whose denominators are units
therefore proves the required rational nonsingularity.
\end{lemma}

\begin{proof}
The even root polynomials $V_i$ are even in $u$, whereas $\Phi$
is anti-invariant under $u\mapsto-u$.  Its moments of centered
even polynomials vanish.  The two off-diagonal blocks have the
displayed entries and are equal and symmetric.  The determinant
formula follows by exchanging the two column blocks.  Reduction
at an admissible auxiliary prime preserves any rational zero
determinant, so a nonzero modular determinant proves that the
rational determinant is nonzero.
\end{proof}

This is a finite rational hypothesis for each fixed tuple.  It
does not assert that the determinant is a unit at every prime.
With a fixed production cutoff in $q$, only finitely many tuples
and finitely many exceptional reduction primes arise.

\subsection{Radicals and a Christoffel identity}

Write
\[
 \widetilde U_i=V_i\mathbb F_p[Y]_{\le d}.
\]
Each full space has one more even centered direction than odd
centered directions.  The nonsingular low pairing implies that
there are unique monic, centered-even polynomials $R_i$ of degree
$D_i$ satisfying
\begin{equation}\label{conf:radicals}
 \Phi(R_1F_2)=0\quad(F_2\in\widetilde U_2),\qquad
 \Phi(F_1R_2)=0\quad(F_1\in\widetilde U_1).
\end{equation}
Indeed the nonsingular off-diagonal parity blocks determine the
lower even coefficients of $R_i$; pairings with every even
direction vanish automatically.  The full pairing has rank $d$,
with the indicated one-dimensional left and right radicals.

The maximal level polynomial has nodes $-q-1,\ldots,q$ and the
nonmaximal one has nodes $-q,\ldots,q$.  Thus
\[
 L_\Phi(Y)=(Y+q+1)L_\Psi(Y).
\]
The residue and polynomial-quotient definition of the universal
functional gives the exact identity
\begin{equation}\label{conf:christoffel}
 \Psi(P)=\Phi((Y+q+1)P).
\end{equation}
At the added node $-q-1$ the numerator vanishes; at every other
node the added factor cancels from the numerator and denominator
derivative.  Polynomial division gives the same quotient on both
sides.  This proves the identity over $\mathbb Q$ before
reduction.

Since $(Y+q+1)F_i\in\widetilde U_i$ for $F_i\in U_i$,
\eqref{conf:radicals}--\eqref{conf:christoffel} imply
\begin{equation}\label{conf:cross-orthogonality}
 \Psi(R_1F_2)=\Psi(F_1R_2)=0,\qquad
 \Psi(R_1R_2)=\nu,
 \qquad \nu:=\Phi(YR_1R_2).
\end{equation}
No assumption that $\nu\ne0$ is made.  In particular, no
nonsingularity test for the nonmaximal low matrix is needed.

\subsection{Crossed flags give true first kernels}

Let $k_i(c)$ be the actual lower-root degree of side $i$ at class
$c$.  On the maximal set define
\begin{align}
 F_1&=\{c\in M:k_1(c)>\ell_1\text{ or }k_2(c)<\ell_2\},
 \nonumber\\
 F_2&=\{c\in M:k_2(c)>\ell_2\text{ or }k_1(c)<\ell_1\}.
 \label{conf:flags}
\end{align}
Put
\[
 f_i=|F_i|,\qquad c_i=\min(H_i,f_i),\qquad
 a_i=(H_i-f_i)_+,
\]
and
\begin{equation}\label{conf:costs}
 L_i^M=\sum_{c\in M}(D_i-k_i(c))_+,
 \qquad E_i^T=\sum_{c\in T}(\ell_i-k_i(c))_+.
\end{equation}
The sets in \eqref{conf:flags} are unions: a class satisfying both
conditions is counted once.

The available top coefficient on side $i$ is
\[
 P_i(c)=Q_i(c^2),\qquad \deg Q_i<H_i.
\]
Impose $P_i(c)=0$ on $F_i$ and prescribe the maximal leading
feature to be $P_i(c)R_i(Y)$.  If $k_i(c)>\ell_i$, the flag
forces the zero feature.  Otherwise $Z_{k_i(c)}$ divides $R_i$,
so the feature has all actual roots.  Independent maximal low
features realize the prescription.

This is a true first kernel against the entire opposite original
window.  If $k_j(c)<\ell_j$, the opposite missing-root flag forces
$P_i(c)=0$.  Otherwise the opposite full feature space is a
subspace of $\widetilde U_j$, against which $R_i$ is a radical.
The top flags cost $c_i$ directions, and prescribing the lower
maximal coefficients costs $L_i^M$.  Imposing the missing
$\ell_i$ roots on $T$ costs at most $E_i^T$ further directions.
Thus there are chosen true first-kernel subspaces of dimensions
at least
\begin{equation}\label{conf:kernel-dimensions}
 b_i=r-L_i^M-c_i-E_i^T.
\end{equation}

After imposing those roots on $T$, enlarge all nonmaximal lower
features to $U_i$.  Every feature is then expressible as
\[
 P_i(c)R_i(Y)+G_{i,c}(Y),\qquad G_{i,c}\in U_i.
\]
This enlargement can only increase the second rank.  By
\eqref{conf:cross-orthogonality}, its low part is orthogonal to
the top directions and has total rank at most $dt$.

The opposite missing-root flags are essential.  Imposing extra
roots only on a selected side does not make the opposite radical
orthogonal to the larger unrestricted actual space.

\subsection{The maximal first jet}

Choose $p$-integral lifts of the selected vectors.  On a maximal
class write
\[
 f_i(c+pY)=P_i(c)R_i(Y)+p f_{i,1,c}(Y)\pmod {p^2}.
\]
The strengthened Newton--Lucas coefficient identity gives
\begin{equation}\label{conf:first-jet}
 [Y^{D_i+1}]f_{i,1,c}=P_i'(c).
\end{equation}
There is no factor $1/(D_i+1)$.  After clearing the factorial
denominator of a top basis vector, its numerator reduces to
$(x^p-x)^{D_i}S(x)$ with $\deg S<p$.  Comparing the coefficients
of $(x-c)^{D_i}$ and $(x-c)^{D_i+1}$, then substituting $x=c+pY$,
proves \eqref{conf:first-jet}.  Lower-band vectors do not
contribute to the latter coefficient at this precision.

If the opposite leading coefficient $P_j(c)$ is nonzero, then
$c\notin F_j$, hence $k_i(c)\ge\ell_i$.  The first jet on side
$i$ is therefore divisible by $V_i$.  Subtracting
$P_i'(c)YR_i$ leaves a polynomial in $\widetilde U_i$, whose
pairing with the opposite radical vanishes.  If $P_j(c)=0$, the
corresponding first-jet term is zero without a root argument.
Consequently the surviving first-order terms on every maximal
class are governed by the same scalar $\nu$.  This includes the
case in which exactly one $P_i(c)$ vanishes: its derivative must
still be retained.

The first Schur complement does not change this order-$p$ form.
Both cross blocks from true first-kernel vectors to the leading
nonsingular block are divisible by $p$, so their Schur product
starts at $p^2$.

\subsection{Exact weights and first-order precision}

Let
\[
 V_M(Z)=\prod_{c\in M}(Z-c^2),\qquad
 V_T(Z)=\prod_{c\in T}(Z-c^2),\qquad
 V_MV_T=Z^n-1.
\]
There is a nonzero constant $C$, independent of $c$, such that
\begin{equation}\label{conf:weights}
 \omega_{M,c}=C\frac{c^5}{V_M'(c^2)},\qquad
 \omega_{T,c}=2C\frac{c^6}{V_M(c^2)}.
\end{equation}
For completeness set $S=(-1)^{q+n+w+1}$.  Factoring with monic
level polynomials gives reduced denominator units
\[
 U_{M,c}=2S V_M'(c^2),\qquad
 U_{T,c}=S c^{-1}V_M(c^2),\qquad
 C=(N!/p^q)^2/S.
\]
Each complete $p$-block contributes $-1$ to the unit.  The partial
block is, up to a constant sign, $(Z^n-1)V_M(Z)$, with double
nodes on $M$ and simple nodes on $T$.  This proves the formulas,
including the factor $2$ between the two weights.

At $c\in M$ put $x=c^2$ and $M_c(Z)=V_M(Z)/(Z-x)$.
The coefficient of $pY$ in the logarithm of the denominator unit
is
\[
 \frac1c+2c\left(
 \frac{V_T'}{V_T}+2\frac{M_c'}{M_c}\right)(x).
\]
Indeed, after dividing out the two physical class factors, the
partial block becomes, up to a constant,
$(z+c)^2V_T(z^2)M_c(z^2)^2$; complete blocks have zero first
logarithmic derivative modulo $p$.  Since
\[
 \left(\frac{V_T'}{V_T}+\frac{M_c'}{M_c}\right)(x)
 =\frac{n-1}{2x},\qquad
 \frac{M_c'(x)}{M_c(x)}=\frac{V_M''(x)}{2V_M'(x)},
\]
and $n=-1/2$ in $\mathbb F_p$, the multiplier
$(c+pY)^5/U_c(Y)$ has first logarithmic coefficient
\begin{equation}\label{conf:lambda}
 \lambda_c=\frac{11}{2c}
              -c\frac{V_M''(c^2)}{V_M'(c^2)}.
\end{equation}

Here is the precise first-order expansion.  Write
$f_i=f_{i,0}+pf_{i,1}$ on a maximal class.  Modulo $p^2$, its
contribution to $p^6A$ is
\begin{align}
 \omega_{M,c}\bigl\{&\Phi(f_{1,0}f_{2,0})
 +p\Phi(f_{1,1}f_{2,0}+f_{1,0}f_{2,1}
              +\lambda_cYf_{1,0}f_{2,0})\bigr\},
 \label{conf:precision-max}
\end{align}
up to a first-order unit-constant correction times
$\Phi(f_{1,0}f_{2,0})$, which vanishes on the selected radicals.
The corresponding nonmaximal contribution is
\begin{equation}\label{conf:precision-nonmax}
 p\omega_{T,c}\Psi(f_{1,0}f_{2,0}).
\end{equation}

The central factorial interpolation truncation has local degree
at most $2q+2$ modulo $p^2$, by the strengthened Newton bound.
The maximal level polynomial has degree $2q+2$, so dividing the
truncation by it leaves only a constant quotient ambiguity.
Multiplication by the affine first-order unit makes this at most
affine, and three derivatives annihilate it.  Multiplying an
interpolation remainder by that affine unit introduces only a
constant quotient, again annihilated.  On $T$ only the leading
functional is needed because of the additional factor $p$.

The nonmultiple-$p$ harmonic terms and the remaining Stieltjes
terms start at order $p^5$ in the scaled local expression and do
not enter \eqref{conf:precision-max}--\eqref{conf:precision-nonmax}.
The inherited strengthened Newton--Bernoulli argument permits
coefficient reduction modulo $p^2$ before the three derivatives
and Bernoulli evaluation.  Alternatively, each fixed local tuple
admits a larger good-prime cutoff placing every retained local
degree below $p-1$.  At zero the quotient is $p$-integral after
factoring the physical multiples of $p$, and the factor $x^5$
annihilates the contribution at this precision; the zero-class
proper fractions have the same extra valuation.  The retained
constant quotient and the affine $X$-slope start beyond $p^2$ in
$p^6A$.

\subsection{Confluent residues and the finite rank cap}

By the first-jet identities and \eqref{conf:cross-orthogonality},
the second top form is $C\nu$ times
\begin{align}
 &\sum_{c\in M}\frac{c^5}{V_M'(c^2)}
 \bigl(P_1'(c)P_2(c)+P_1(c)P_2'(c)
                    +\lambda_cP_1(c)P_2(c)\bigr)
 \nonumber\\
 &\hspace{35mm}
 +2\sum_{c\in T}\frac{c^6P_1(c)P_2(c)}{V_M(c^2)}.
 \label{conf:top-form}
\end{align}
Write $F(Z)=Q_1(Z)Q_2(Z)$ and consider
\begin{equation}\label{conf:residue-function}
 \mathcal R_F(Z)=\frac{Z^2F(Z)}{(Z^n-1)V_M(Z)}.
\end{equation}
At a simple node $x=c^2$, $c\in T$, its residue is
\[
 \frac{x^3F(x)}{nV_M(x)}=-2\frac{x^3F(x)}{V_M(x)}.
\]
At a double node $x=c^2$, $c\in M$, its residue is
\begin{equation}\label{conf:double-residue}
 \frac1n\left[
 \frac{x^3F'(x)}{V_M'(x)}
 +\frac{5-n}{2}\frac{x^2F(x)}{V_M'(x)}
 -\frac12\frac{x^3V_M''(x)F(x)}{V_M'(x)^2}\right].
\end{equation}
Using $n=-1/2$ and $P_i'(c)=2cQ_i'(c^2)$, one sees that
\eqref{conf:top-form} equals
\[
 -\sum_{c=1}^n\operatorname{Res}_{Z=c^2}\mathcal R_F(Z).
\]
This verifies the relative sign and constant at every simple and
double node.

The denominator has degree $n+w$.  Its finite residue sum
vanishes when the numerator degree is at most $n+w-2$.
In monomial top bases the entry indexed by $i<H_1$, $j<H_2$
therefore vanishes if $i+j+2\le n+w-2$.  The full top rank is
at most $(H_1+H_2-n-w+2)_+$.  Restricting to the flagged top
spaces cannot increase rank, and their dimensions are $a_i$.
Consequently
\begin{equation}\label{conf:top-rank}
 b_{\rm top}=\min\{a_1,a_2,(H_1+H_2-n-w+2)_+\}
\end{equation}
is a valid finite top-rank cap.  The total second rank on the
selected first kernels is at most $dt+b_{\rm top}$.

\begin{lemma}[Restricted first kernels]
\label{conf:restricted-first}
Let an integral square matrix $M_0$ have chosen left and right
first-kernel subspaces of dimensions $b_1,b_2$.  If the reduction
of $M_0/p$ on their common block has rank at most $B$, then
\[
 v_p(\det M_0)\ge b_1+b_2-B.
\]
For a rectangular matrix, if the chosen first-kernel
codimensions are $a_1,a_2$, its image length over
$\mathbb Z/p^2\mathbb Z$ is at most $a_1+a_2+B$.
\end{lemma}

\begin{proof}
In the square case divide the $b_1$ selected rows by $p$.
On the $b_2$ selected columns the new reduction has rank at most
$B$: outside those rows the entries remain divisible by $p$,
and inside them it is the stated common block.  At least
$b_2-B$ additional independent columns can therefore be divided
by $p$ after integral column operations.

For the rectangular statement, let $r_0$ be the first rank and
$r_1$ the second rank on the complete first kernels.  The image
length is $2r_0+r_1$.  Restricting the two complete first kernels
loses codimensions $a_1-r_0$ and $a_2-r_0$, so
$r_1\le B+a_1+a_2-2r_0$.  This gives the assertion.  Image
length cannot increase under deletion of rows or columns,
since inclusion of the retained domain and projection onto the
retained codomain cannot increase module length.
\end{proof}

\begin{theorem}[Confluent second-layer valuation]
\label{conf:valuation}
Under the hypotheses above, with no constant-maximal-root
assumption, one has coefficientwise in $X$
\begin{equation}\label{conf:finite-valuation}
 v_p(\det(p^6A(X)))\ge
 2r-L_1^M-L_2^M-c_1-c_2-E_1^T-E_2^T-dt-b_{\rm top}.
\end{equation}
\end{theorem}

\begin{proof}
Apply Lemma~\ref{conf:restricted-first} using
\eqref{conf:kernel-dimensions} and the second rank bound
$dt+b_{\rm top}$.  All leading forms are constant in $X$,
so the basis changes may be chosen constant and the estimate
holds coefficientwise.
\end{proof}

If $k_i(c)=\ell_i$ throughout $M$, then $F_i$ is empty and
$L_i^M=dw$.  Thus the theorem includes
\begin{equation}\label{conf:constant-core}
 v_p(\det(p^6A(X)))\ge
 2r-2dw-E_1^T-E_2^T-dt
 -(H_1+H_2-n-w+2)_+.
\end{equation}

\subsection{A simpler missing-root corollary}

For $d\ge1$, hence $d\ge2$ in the even setting, define
\[
 E_i^M=\sum_{c\in M}(\ell_i-k_i(c))_+.
\]
Class by class,
\[
 (D_i-k_i(c))_+
 \le d+(\ell_i-k_i(c))_+
               -\mathbf1_{k_i(c)>\ell_i}.
\]
Also the union definition of the flags gives
\[
 f_1+f_2\le
 \sum_i\#\{c\in M:k_i(c)>\ell_i\}
 +E_1^M+E_2^M.
\]
Hence
\[
 L_1^M+L_2^M+c_1+c_2
 \le L_1^M+L_2^M+f_1+f_2
 \le2dw+2E_1^M+2E_2^M.
\]
The simpler uniformly valid valuation bound is therefore
\begin{equation}\label{conf:missing-root-valuation}
 \begin{split}
 v_p(\det(p^6A(X)))\ge{}&
 2r-2dw-dt-E_1^T-E_2^T-2E_1^M-2E_2^M\\
 &-(H_1+H_2-n-w+2)_+.
 \end{split}
\end{equation}

For natural upper-half roots, the extra cost beyond the
constant-maximal-root regime is twice the number of missing
maximal roots on each side.  For natural lower-half roots,
there are no missing-root costs; actual extra maximal roots
lower the low-feature count by exactly the top-flag cost.
Thus the lower-half maximal-root guard disappears without a
new cost.

\subsection{Limiting profiles and finite certification}

On a rational $\beta=p/N$ cell, use normalized class lengths and
\[
 w=1-(q+\tfrac12)\beta,\quad t=\beta/2-w,
 \quad H_i=s_i+\delta-D_i\beta/2.
\]
The exact-flag candidate is
\begin{equation}\label{conf:rate}
 \begin{split}
 E_{\rm conf}={}&4\delta+L_1^M+L_2^M
 +\min(H_1,f_1)+\min(H_2,f_2)+E_1^T+E_2^T\\
 &+dt+(H_1+H_2-\beta/2-w)_+.
 \end{split}
\end{equation}
The top-dimension caps need not be included separately in this
limit.  Since $H_i\le\beta/2$ and $f_i\le w$, the positive
Hermite terminal is at most both $H_1-f_1$ and $H_2-f_2$.
The additional switches $H_i=f_i$ belong in the exact affine
partition.

A simpler candidate is
\begin{equation}\label{conf:missing-root-rate}
 E_{\rm miss}=4\delta+d(\beta/2+w)+E_1^T+E_2^T
 +2E_1^M+2E_2^M+(H_1+H_2-\beta/2-w)_+.
\end{equation}
The even root degrees may be artificial uniform choices.
For \eqref{conf:missing-root-rate}, the root objective is convex
and piecewise linear: $d$ is affine in the chosen root degree
and the other terms are sums of positive parts.  Endpoints and
root breakpoints, rounded to admissible even degrees, suffice.
The exact-flag candidate can instead be evaluated on the finite
class-type partition.

Use only fixed tuples whose rational low determinant has been
certified by Lemma~\ref{conf:witness}.  The finite exceptional
primes and finite endpoint rounding constants are negligible
in the prime-integral limit.  Beyond the finite local cutoff,
retain the inherited tail bound.  The new candidate is combined
by taking the minimum with the complete inherited envelope.

\subsection{Finite diagnostics}

For $p=101$, $N=675$, $m_1=m_2=82$, and $r=753$, one has
$D=16$, $\ell=2$, $d=14$, $w=19$, $t=31$,
$E_1^T=E_2^T=19$, and $H_1=H_2=27$.  Formula
\eqref{conf:constant-core} gives $502$, equal to the observed
first-two-nullity sum $487+15$.  At $p=151$, $N=1009$,
$m_1=m_2=123$, $r=1126$, it gives $754=734+20$.

Beyond the old maximal-root guard, take $p=61$, $N=418$,
$m_1=m_2=51$, $r=466$.  Here $d=14$, $w=22$, $t=8$,
$H_1=H_2=29$, $E_i^T=8$, and $E_i^M=2$.  The terminal is
$8$, so \eqref{conf:missing-root-valuation} gives
\[
 2\cdot466-2\cdot14\cdot22-14\cdot8-16-8-8=172,
\]
equal to the observed sum $154+18$.  These matches diagnose
the source of the previous growing deficit.  The uniform
theorem follows from the proof above, not from these examples.

\subsection{Odd endpoint heads and arbitrary parities}
\label{conf:all-parities}

The even theorem extends to arbitrary endpoint parities by using
even-only subspaces and imposing orthogonality to the deleted odd
heads.  This constructs true first kernels of the original matrix;
it does not presume a heuristic cost for restoring deleted rows.

Let $D_i,H_i$ now denote the original endpoint data and put
\[
 \pi_i=D_i\bmod2,\qquad D_i^e=D_i-\pi_i,
 \qquad O_i=\pi_iH_i,
\]
\begin{equation}\label{conf:even-truncation}
 h_i=\begin{cases}H_i,&\pi_i=0,\\n+1,&\pi_i=1.
 \end{cases}
\end{equation}
On an odd side delete its $O_i=H_i$ highest basis vectors,
leaving upper endpoint $(D_ip+1)/2$.  This truncated window has
degree $D_i^e=D_i-1$ and even-band top width $n+1$.  Its dimension
is $r-O_i$; the two truncated dimensions may differ.  Retain the
feature guard $m_i\le D_i^ep/2$.

Choose even roots $\ell_i$ with equal residual
$d=D_1^e-\ell_1=D_2^e-\ell_2$, and require the same low
$\Phi$ nonsingularity test.  Define $F_i,f_i$ by
\eqref{conf:flags} using the unchanged actual lower-root counts,
and set
\[
 c_i=\min(h_i,f_i),\qquad a_i=(h_i-f_i)_+,
\]
\[
 L_i^M=\sum_{c\in M}(D_i^e-k_i(c))_+,
 \qquad E_i^T=\sum_{c\in T}(\ell_i-k_i(c))_+.
\]

The preceding construction gives first kernels against the two
truncated windows.  To obtain first kernels against the entire
original matrix, the selected vector on side $i$ must also
annihilate the $O_j$ deleted odd-head vectors on side $j$.
An odd head has leading local coefficient $cS_j(c^2)$, with
$\deg S_j<O_j$.  Where $P_i(c)\ne0$, the crossed flags give
$k_j(c)\ge\ell_j$.  Subtracting
$cS_j(c^2)YR_j$ from its local feature leaves a polynomial of
degree at most $D_j^e$ in the full $\ell_j$-root space.
Thus the complete leading pairing with the deleted head is
\begin{equation}\label{conf:odd-head-pairing}
 C\nu\sum_{c\in M}
 \frac{c^6Q_i(c^2)S_j(c^2)}{V_M'(c^2)}.
\end{equation}
Only the selected even top polynomial is involved; nonmaximal
low features remain free.

If $h_i\le f_i$, the selected top polynomial is zero.  Otherwise
write $Q_i=V_{F_i}Q_i'$, with $\deg Q_i'<a_i$.  On
$C_i=M\setminus F_i$ the factor $V_{F_i}$ cancels from the
denominator derivative, and \eqref{conf:odd-head-pairing}
becomes
\[
 C\nu\sum_{c\in C_i}
 \frac{(c^2)^3Q_i'(c^2)S_j(c^2)}{V_{C_i}'(c^2)}.
\]
Lagrange interpolation therefore bounds its rank by
\begin{equation}\label{conf:odd-constraint-rank}
 \eta_i=
 \min\{O_j,a_i,w-f_i,(h_i+O_j-w+3)_+\}.
\end{equation}
The cap $w-f_i$ also covers an empty class set.  No nonzero
assumption on $\nu$ is needed.

After imposing these equations, the chosen spaces are true
first kernels of the whole original matrix, of dimensions at
least
\begin{equation}\label{conf:odd-kernel-dimensions}
 r-O_i-L_i^M-c_i-E_i^T-\eta_i.
\end{equation}
Every selected vector lies in its even truncated window, so its
first-jet degree is $D_i^e+1$.  The crossed-flag first-jet proof,
Christoffel identity and confluent residue calculation are
unchanged.  Additional odd-head equations only restrict the top
polynomial spaces.  The nonmaximal low rank remains at most $dt$,
and a valid top-rank cap is
\begin{equation}\label{conf:odd-top-rank}
 b_{
 \rm top}=\min\{a_1,a_2,(h_1+h_2-n-w+2)_+\}.
\end{equation}
We do not subtract $\eta_i$ from these dimensions: an upper bound
on the constraint rank would not justify that subtraction.

\begin{theorem}[All-parity confluent valuation]
\label{conf:all-parity-theorem}
Under these hypotheses, coefficientwise in $X$,
\begin{equation}\label{conf:all-parity-finite}
 \begin{split}
 v_p(\det(p^6A(X)))\ge{}&2r
 -\sum_{i=1}^2(O_i+L_i^M+c_i+E_i^T+\eta_i)\\
 &-dt-b_{\rm top}.
 \end{split}
\end{equation}
\end{theorem}

\begin{proof}
Apply Lemma~\ref{conf:restricted-first} to
\eqref{conf:odd-kernel-dimensions} and the second-rank cap
$dt+b_{\rm top}$.  The selected vectors annihilate the original
odd heads by \eqref{conf:odd-head-pairing}, so the first Schur
product still starts at $p^2$ and does not change the order-$p$
form used in the argument.
\end{proof}

In the limiting formula, put $h_i=\beta/2$ on an odd side and
$h_i=H_i$ on an even side, with $O_i=H_i$ or $0$ respectively.
Define
\[
 \eta_i=
 \min\{O_j,(h_i-f_i)_+,w-f_i,(h_i+O_j-w)_+\}.
\]
The valid all-parity cost is
\begin{equation}\label{conf:all-parity-rate}
 \begin{split}
 E_{\rm all}={}&4\delta+
 \sum_{i=1}^2\bigl(O_i+L_i^M+\min(h_i,f_i)+E_i^T+\eta_i\bigr)\\
 &+dt+(h_1+h_2-\beta/2-w)_+.
 \end{split}
\end{equation}
As before, the residual-top dimension caps are redundant in the
limiting terminal.  All min/max switches are inserted into the
exact affine partition.  The coarser version adds
$2(O_1+O_2)$ to the even-truncated cost, since
$\eta_1+\eta_2\le O_1+O_2$.

The sharper constraint cap can matter in mixed-parity cells:
if $h_{\rm even}+O_{\rm odd}\le w$, its limiting leading
constraint rank is zero.  For two odd endpoints,
$h_i=\beta/2$, so this particular positive-part cap is
redundant, though the other caps may still apply.

\subsection{Optional ceiling enlargement}

An odd endpoint can alternatively be enlarged to
\[
 R_i^+=(D_i+1)p/2+1,
\]
which has even degree $D_i+1$ and top width $1$.
The original matrix is obtained by deleting the added rows or
columns of the enlarged rectangular window matrix.
Lemma~\ref{conf:restricted-first} bounds the enlarged image
length over $\mathbb Z/p^2\mathbb Z$, and that length cannot
increase on deletion.  The original $r\times r$ determinant
therefore has first-two valuation at least $2r$ minus the same
image-length bound.  Its rate retains the original prefix
$4\delta$, with no enlargement penalty.

Floor and ceiling choices may be made independently on odd sides.
For a floored side use $D_i^e=D_i-1$, $h_i=n+1$, and its
deleted original head $O_i=H_i$.  For a ceiling side use
$D_i^e=D_i+1$, $h_i=1$, and $O_i=0$; the limiting top width is
zero.  The equal-residual and low-matrix hypotheses, as well as
the feature and precision guards for the enlarged windows, must
hold for the chosen effective degrees.  These are additional
valid candidates; they need not improve the final envelope.

For a finite odd-head diagnostic, take $p=61$, $N=426$,
$m_1=m_2=52$, $r=475$.  The original degrees are $17$ and
$O_i=8$.  The even truncations have degree $16$, $h_i=31$,
$\ell_i=2$, $d=14$, $w=30$ and $t=0$.  Here
$L_i^M=429$, $f_i=c_i=9$, $\eta_i=8$, and $b_{\rm top}=4$.
Theorem~\ref{conf:all-parity-theorem} gives
\[
 2\cdot475-2(8+429+9+8)-4=38,
\]
equal to the observed first-two-nullity sum $21+17$.

\section{Further arithmetic refinements}
\label{z87:structural}

The refinements in this section preserve the original determinant family.
The first keeps the actual nonmaximal root spaces in the confluent second
layer.  The second uses enlargement to a centered mixed core, and applies
without a local nonsingularity test.  The last extends the outer lower
transition by retaining a low-rank head correction.  Every statement is a
uniform finite-prime estimate; sampled Smith forms are diagnostics only.

\subsection{The actual nonmaximal spaces in the confluent layer}
\label{z87:actual-nonmax}

Retain the hypotheses and notation of the all-parity confluent theorem.
Thus $\epsilon=1$, $E_i=D_i-(D_i\bmod2)$ are the even truncated
endpoint degrees, $O_i$ are the deleted odd-head dimensions, and $h_i$
are the even top widths.  Choose even $\ell_i\ge0$ with
\[
 d=E_1-\ell_1=E_2-\ell_2\ge2,
\]
and assume the same rational low pairing of
$\Phi=\mathscr T_{q,1}$ on
$U_i=Z_{\ell_i}\mathbb F_p[Y]_{<d}$ is nonsingular after reduction.
Let $R_i$ be its monic centered-even radicals, put
$\Psi=\mathscr T_{q,0}$, and set $\nu=\Phi(YR_1R_2)$.
No condition $\nu\ne0$ is imposed.

Keep the crossed maximal flags $F_i$, their sizes $f_i$, the costs
$c_i=\min(h_i,f_i)$ and $L_i^M$, and the odd-head constraint bounds
\[
 a_i=(h_i-f_i)_+,\qquad
 \eta_i=\min\{O_j,a_i,w-f_i,(h_i+O_j-w+3)_+\}.
\]
On a nonmaximal class define the actual full even-truncated feature
space and its dimension by
\[
 W_{i,c}=Z_{k_i(c)}\mathbb F_p[Y]_{\le E_i-k_i(c)},\qquad
 g_i(c)=(E_i-k_i(c)+1)_+.
\]
An empty polynomial space has dimension zero.  Put
\[
 e_i(c)=\min\{g_i(c),(\ell_i-k_i(c))_+\},
\]
\begin{equation}\label{z87:local-K-cap}
 B_c=\min\{d+e_1(c)+e_2(c),g_1(c),g_2(c)\},\qquad
 B_T=\sum_{c\in T}B_c,
\end{equation}
and
\[
 B_{\mathrm{top}}=
 \min\{a_1,a_2,(h_1+h_2-n-w+2)_+\}.
\]

\begin{theorem}[Confluent valuation with actual nonmaximal roots]
\label{z87:actual-confluent}
With these hypotheses, coefficientwise in $X$,
\begin{equation}\label{z87:actual-confluent-finite}
 v_p\bigl(\det(p^6A(X))\bigr)
 \ge 2r-\sum_{i=1}^2(O_i+L_i^M+c_i+\eta_i)
        -B_T-B_{\mathrm{top}}.
\end{equation}
No additional fixed local matrix test is needed.
\end{theorem}

\begin{proof}
Construct the same maximal features $P_i(c)R_i(Y)$, with
$P_i(c)=Q_i(c^2)$, $\deg Q_i<h_i$, and $P_i|_{F_i}=0$.
The crossed flags make these true first kernels against the two
truncated windows.  The leading form is supported on $M$, so this
construction places no extra condition on a nonmaximal feature.
Restoring the opposite odd heads imposes at most $\eta_i$ equations,
exactly as in the preceding all-parity argument.  We obtain true first
kernels of the entire original matrix of dimensions at least
\begin{equation}\label{z87:actual-kernel-dims}
 b_i=r-O_i-L_i^M-c_i-\eta_i.
\end{equation}
In particular, the old missing-root payment on $T$ is not subtracted.

For arbitrary polynomials of degrees at most $E_i$, define
\[
 K(F_1,F_2)=\Psi(F_1F_2)
       -\nu[Y^{E_1}]F_1\,[Y^{E_2}]F_2.
\]
On each nonmaximal class the actual form is exactly its $K$ remainder
plus the scalar top term $\nu P_1(c)P_2(c)$.  Combining these top terms
with the maximal first jets gives the unchanged confluent residue form
associated with
\[
 \frac{Z^2Q_1(Z)Q_2(Z)}{(Z^n-1)V_M(Z)}.
\]
Its rank is at most $B_{\mathrm{top}}$.  This algebraic splitting does
not require an actual feature to have the canonical roots $Z_{\ell_i}$.
The first Schur correction still starts at $p^2$, since the spaces in
\eqref{z87:actual-kernel-dims} are true first kernels; it does not change
the order-$p$ form under consideration.

On the canonical full spaces
$\widetilde U_i=Z_{\ell_i}\mathbb F_p[Y]_{\le d}$ write
$F_i=t_iR_i+G_i$, $G_i\in U_i$.  Christoffel orthogonality gives
\[
 \Psi(R_1G_2)=\Psi(G_1R_2)=0,\qquad
 \Psi(R_1R_2)=\nu,
\]
so $K(F_1,F_2)=\Psi(G_1G_2)$ and its rank there is at most $d$.
The root polynomials form a nested family.  Consequently
$W_{i,c}\cap\widetilde U_i$ has codimension $e_i(c)$ in $W_{i,c}$.
Restricting $K$ to the intersections has rank at most $d$; restoring
$e_1(c)$ rows and $e_2(c)$ columns raises rank by at most their sum.
The independent dimension bounds $g_i(c)$ prove
\eqref{z87:local-K-cap}.

The direct sum of these actual local remainders has rank at most $B_T$.
Restriction to the global feature images cannot increase rank.  Adding
the confluent top form gives total second rank at most
$B_T+B_{\mathrm{top}}$.  The restricted-first-kernel determinant lemma,
applied to \eqref{z87:actual-kernel-dims}, proves the result.  All retained
forms are constant in $X$, so the estimate is coefficientwise.
\end{proof}

On a rational $\beta=p/N$ cell, class counts become normalized lengths.
The new arithmetic candidate is
\begin{equation}\label{z87:actual-confluent-rate}
 \begin{split}
 E_{\mathrm{actual}}={}&4\delta+
 \sum_i\bigl(O_i+L_i^M+\min(h_i,f_i)+\eta_i\bigr)\\
 &+\sum_{\text{root types on }T}
  |T_{k_1,k_2}|\min\{d+(\ell_1-k_1)_++(\ell_2-k_2)_+,
                        (E_1-k_1+1)_+,(E_2-k_2+1)_+\}\\
 &+(h_1+h_2-\beta/2-w)_+.
 \end{split}
\end{equation}
The $+1$ in a local feature dimension is retained: there are a positive
proportion of such classes.  The bounded offsets in the global Cauchy
term disappear in the usual fixed-cutoff prime-integral limit.
The old candidate is never smaller solely because of the replaced
nonmaximal term, since
$B_c\le d+(\ell_1-k_1(c))_++(\ell_2-k_2(c))_+$.
Optional ceiling enlargements remain valid by monotonicity of image
length over $\mathbb Z/p^2\mathbb Z$.

\subsection{A centered mixed core by enlargement}
\label{z87:centered-core}

Let $\Phi=\mathscr T_{q,1}$ be an anti-invariant local functional.
Choose even $E_i,\ell_i\ge0$ with
\[
 d=E_1-\ell_1=E_2-\ell_2\ge2.
\]
The row endpoint has odd degree $D_1=E_1+1$ and the column endpoint
has even degree $D_2=E_2$.  On a class set $S$, suppose only that
\begin{equation}\label{z87:centered-root-guard}
 k_i(c)\ge\ell_i\qquad(c\in S,\ i=1,2).
\end{equation}
The odd and even correlated top coefficients are $cP(c^2)$ and
$Q(c^2)$, where $\deg P<h_1$, $\deg Q<h_2$; fixed common vanishing
factors may be included as in the applications below.  Suppose the
weighted top product becomes a common scalar times
\[
 \frac{(c^2)^\gamma P(c^2)Q(c^2)}{V_S'(c^2)}
\]
for a fixed nonnegative integer $\gamma$.

\begin{theorem}[Unconditional centered mixed rank bound]
\label{z87:centered-rank}
The entire class form has rank at most
\begin{equation}\label{z87:centered-rank-bound}
 d|S|+(h_1+h_2-|S|+\gamma)_+.
\end{equation}
The assertion is unchanged by transposing the two sides and requires
no local nonsingularity test.
\end{theorem}

\begin{proof}
Enlarge the actual root spaces to uniform roots $Z_{\ell_i}$, preserving
the original correlated top polynomials.  This is an enlargement because
of \eqref{z87:centered-root-guard}; any actual feature is its top
coefficient times a fixed monic canonical top polynomial, plus a
canonical low feature.  The enlarged low row and column spaces are
\[
 V_1=Z_{\ell_1}\mathbb F_p[Y]_{\le d},\qquad
 U_2=Z_{\ell_2}\mathbb F_p[Y]_{<d},
\]
of dimensions $d+1$ and $d$.  Rank cannot increase on restricting the
resulting canonical global pairing back to the actual feature images.

Consider the square pairing on
$U_i=Z_{\ell_i}\mathbb F_p[Y]_{<d}$.  In centered powers
$u=Y+1/2$, its parity blocks are
\[
 \begin{pmatrix}0&B\\B&0\end{pmatrix},
\]
where $B$ is square of size $d/2$.  If it is singular, the square
pairing has rank at most $d-2$.  Adding the extra even row direction
raises the local low rank by at most one.  Adjoining all global top
heads therefore bounds the full rank by
\[
 (d-1)|S|+h_1+h_2
 \le d|S|+(h_1+h_2-|S|+\gamma)_+.
\]
This includes singular reductions at every good prime.

If the square pairing is nonsingular, its canonical full spaces
$V_i=Z_{\ell_i}\mathbb F_p[Y]_{\le d}$ have monic centered-even
radicals $R_i$ of degree $E_i$.  In particular $R_2$ annihilates all
row low features, and the row low radical $R_1$ annihilates the entire
column space $V_2$.  Choose the column top polynomial to be $R_2$ and
the row top polynomial to be $YR_1$.  After $d$ local low pivots at
each class, the only remaining scalar is
\[
 \nu_0=\Phi(YR_1R_2),
\]
the same at every class.  No nonzero assumption on $\nu_0$ is needed.
The residual global form is a scalar multiple of the displayed Cauchy
moment matrix.  Its entries vanish when
$i+j+\gamma\le |S|-2$, so its rank is at most the positive part in
\eqref{z87:centered-rank-bound}.  Adding the $d|S|$ low pivots proves
the theorem.
\end{proof}

For the leading form with $\epsilon=1$, take $S=M$ and $h_i=H_i$.
The maximal weight has the factor $c^5/V_M'(c^2)$, and the odd/even
top pair adds one factor $c$.  Thus $\gamma=3$, and
\[
 v_p\bigl(\det(p^6A)\bigr)
 \ge r-dw-(H_1+H_2-w+3)_+.
\]
The limiting candidate is
\begin{equation}\label{z87:centered-leading-rate}
 E_{\mathrm{lead}}=5\delta+dw+(H_1+H_2-w)_+.
\end{equation}

For the second form with $\epsilon=0$, take $S=T$ after killing both
maximal feature maps.  The local functional is
$\mathscr T_{q-1,1}$ and the residual top dimensions are
$K_i=(H_i-w)_+$.  Since
\[
 \omega_{T,c}V_M(c^2)^2=C'\frac{c^3}{V_T'(c^2)},
\]
the odd/even pair gives $\gamma=2$.  With the original maximal
restriction ranks $A_i$, the finite consequence is
\[
 v_p\bigl(\det(p^5A)\bigr)
 \ge2r-A_1-A_2-dt-(K_1+K_2-t+2)_+,
\]
and the limiting candidate is
\begin{equation}\label{z87:centered-second-rate}
 E_{\mathrm{second}}=3\delta+A_1+A_2+dt+(K_1+K_2-t)_+.
\end{equation}
In both applications use the smallest admissible even $d\ge2$, since
its coefficient is a nonnegative class length.  The respective root
guard is imposed on $M$ or $T$.  These statements have no fixed local
matrix witnesses, and hence no matrix-certificate cutoff.

\subsection{An outer lower transition with a head correction}
\label{z87:outer-lower-extension}

Retain the lower-half notation and all fixed local tests from the outer
lower Cauchy theorem: $\epsilon=0$, $N=qp+w$, $n=(p-1)/2$,
$m_i=ap+\tau_i$, $v_i=\tau_i-1$, equal odd endpoint degrees $D+1$,
\[
 e=D-2a,\quad
 H_i=R_i-\frac{(D+1)p+1}{2},\quad
 k_f=(D/2+a+1)p+2-r-m_1-m_2.
\]
Set all free zero-class parameters to zero.  Suppose
\[
 3\le v_1\le v_2<w\le n,\qquad k_f\ge0,\qquad
 0\le H_1\le H_2\le w-3.
\]
The head-sum restriction is replaced by
\[
 j=(H_1+H_2-w+3)_+.
\]

\begin{theorem}[Outer lower rank correction]
\label{z87:outer-lower}
Under these hypotheses,
\begin{equation}\label{z87:outer-lower-finite}
 v_p\bigl(\det(p^5A)\bigr)
 \ge r-ew-w+2(n+1-w)-j.
\end{equation}
\end{theorem}

\begin{proof}
Keep the class partition $M=A\cup B\cup C$, with
$A=\{1,\ldots,v_1\}$, $B=\{v_1+1,\ldots,v_2\}$, and
$C=\{v_2+1,\ldots,w\}$.  Put $a_0=v_1$, $b=v_2-v_1$,
and $b_1=w-v_1$.  The inherited initial construction gives
$ew+b_1$ first pivots.  Residual odd column polynomials have
$S_2=V_BS_2'$, and $H_2-H_1=b$, so both odd parameter spaces have
dimension $H_1$.

Their leading form, after removal of one common unit, is
\[
 \theta B_A(S_1,S_2)+\nu B_C(S_1,S_2)
 =\kappa B_A(S_1,S_2)+\nu J(S_1,S_2'),
\]
where $\kappa=\theta-\nu\ne0$ and
\[
 J(S_1,S_2')=
 \sum_{c\in A\cup C}
 \frac{(c^2)^3 S_1(c^2)S_2'(c^2)}{V_{A\cup C}'(c^2)}.
\]
There are $w-b$ nodes.  The Cauchy anti-diagonal vanishing gives
\[
 \operatorname{rank}J\le
 (2H_1-(w-b)+3)_+=j.
\]
This is precisely the term previously set to zero by the head-sum guard.

The complete-even parameter spaces still have dimension
$d_0=n+1-w$.  Their boundary coordinates are $E_1,E_2^*=E_2-S_B$,
where $S_B$ has degree less than $b$ and interpolates $E_2$ on $B$.
Every mixed coefficient remainder vanishes under $H_2\le w-3$.
In particular
$B_M(S_1,S_B)=0$, since
$\deg(S_1S_B)+3\le H_1+b+1=H_2+1\le w-2$.
The nonmaximal correction to the even/even form also vanishes, since
\[
 \deg(Q_1S_B)\le n-w+b-1\le n-4.
\]
Thus the simultaneous boundary blocks retain the exact leading orders
\[
 OO=\kappa B_A+\nu J+O(p),\quad
 OE=EO=p\kappa B_A+O(p^2),\quad
 EE=p^2\kappa B_A+O(p^3).
\]
Only the odd/odd block has acquired a rank-at-most-$j$ correction.

Here is the needed linear algebra.  For a nondegenerate pairing $B$ on
an $a_0$-dimensional boundary space, arbitrary flag maps $U,V$, and a
matrix $T$ of rank at most $j$, the augmented matrix
\[
 F=\begin{pmatrix}
 U^{\mathsf T}BV+T&U^{\mathsf T}B\\
 BV&B
 \end{pmatrix}
\]
has rank $a_0+\operatorname{rank}T\le a_0+j$, by taking the Schur
complement of $B$.  If the upper-left block has rank $\ell$, remove
its pivots and let $c_L,c_R$ be the remaining two mixed ranks.
The common restricted bottom-right Schur form $Q$ then satisfies
\[
 \ell+c_L+c_R+\operatorname{rank}Q
 =\operatorname{rank}F\le a_0+j.
\]
This is the elementary rank decomposition after first eliminating the
upper-left pivots and then the two residual mixed blocks.  Neither flag
map nor actual boundary map need be surjective.

Apply this identity with $B=\kappa B_A$ and $T=\nu J$.
The true selected even second kernels have dimensions at least
$d_0-c_L$ and $d_0-c_R$, after imposing every actual residual odd
order-$p$ equation.  Their third pairing is a restriction of $Q$.
The first nullity is $r-ew-b_1-\ell$.  Consequently
\[
 \begin{aligned}
 v_p\bigl(\det(p^5A)\bigr)
 &\ge r-ew-b_1-\ell+2d_0-c_L-c_R-\operatorname{rank}Q\\
 &\ge r-ew-b_1+2d_0-a_0-j
  =r-ew-w+2d_0-j.
 \end{aligned}
\]
The selected spaces are true second kernels, so subsequent second-pivot
corrections start at $p^3$.  This proves
\eqref{z87:outer-lower-finite}.
\end{proof}

The corresponding limiting candidate is
\begin{equation}\label{z87:outer-lower-rate}
 E_{\mathrm{outer}}=
 4\delta+(e+3)w-\beta+(H_1+H_2-w)_+.
\end{equation}
Equivalently, with
$k=(D/2+a+1)\beta-\delta-s_1-s_2$, it is
\[
 4\delta+ew-2k+(w-\tau_1)+(w-\tau_2)+(w-H_1-H_2)_+.
\]
The rate guards retain the common lower-half quotient, equal odd degrees,
all inherited local tests, $\tau_2<w\le\beta/2$, $k\ge0$, and
$0\le H_1\le H_2\le w$.  Endpoints are treated by adjacent open-cell
limits; the finite theorem keeps its stated three-unit buffer.

\section{Unequal endpoint degrees with a common effective dimension}
\label{z87:unequal-central}

The central-root third layer does not require the two even endpoint
degrees, or their two integral root levels, to agree.  The necessary
equality is their effective dimension.  We give the two-odd-band
version, including the saturated case, because it detects a new family
of growing third Smith nullities.  Every local nonsingularity condition
below is a condition on a fixed rational cross Gram matrix.  Numerical
Smith computations are not hypotheses.

\subsection{Geometry and the three cross Gram conditions}

Fix $q\ge1$, and work at an odd prime $p$ for which
$N=qp+w$, $1\le w\le n=(p-1)/2$, so $\epsilon=0$.
Retain the good-prime and complete-feature hypotheses
\eqref{lyr:good-prime}, in particular $m_i\le E_ip/2$,
and enlarge the fixed good-prime cutoff so that all local
rational denominators and degrees below have valid reduction.
Let the two windows have lengths $r$, initial indices
$m_i=a_ip+\tau_i$, and right endpoints $R_i=m_i+r$.
Assume their actual endpoint degrees are odd:
\[
 D_i=\left\lfloor\frac{2R_i-2}{p}\right\rfloor=E_i+1,
 \qquad E_i\ \hbox{even},\qquad
 e=E_1-2a_1=E_2-2a_2\ge2.
\]
Write
\[
 J_i=\frac{(E_i+1)p+1}{2},\qquad H_i=R_i-J_i,
 \qquad z=e/2-1,
\]
and, after transposing the two windows if necessary, suppose
$0<\tau_1\le\tau_2<p$.  This transposition does not assume
that the original initial indices have the same ordering as their
residues.  The identity
\begin{equation}\label{z87:unequal-head-relation}
 H_i=r+\tau_i-\frac{(e+1)p+1}{2}
\end{equation}
implies $H_2-H_1=\tau_2-\tau_1$.

As usual let
\[
 Z_{2a}(Y)=\prod_{v=-a}^{a-1}(Y-v),\qquad
 Z_{2a+1}(Y)=\prod_{v=-a}^{a}(Y-v),\qquad W_i=Z_{2a_i+1}.
\]
There are two central regimes.  In the upper regime assume
$p/2<\tau_1\le\tau_2<p$ and $w\le p-\tau_2$, and put
$V_i=Z_{2a_i+2}$.  In the lower regime assume
$0<\tau_1\le\tau_2<p/2$ and $w\le\tau_1-1$, and put
$V_i=Z_{2a_i}$.  On the maximal classes $M=\{1,\ldots,w\}$
the physical roots are $(W_1,W_2)$ in both regimes.  The
successive nonmaximal intervals $\mathsf A,\mathsf B,\mathsf C$
have the following root pairs and sizes:
\[
\begin{array}{c|cc}
 &\text{upper}&\text{lower}\\\hline
\mathsf A&(W_1,W_2)&(W_1,W_2)\\
|\mathsf A|&p-\tau_2-w&\tau_1-1-w\\
\mathsf B&(W_1,V_2)&(V_1,W_2)\\
|\mathsf B|&\tau_2-\tau_1&\tau_2-\tau_1\\
\mathsf C&(V_1,V_2)&(V_1,V_2)\\
|\mathsf C|&\tau_1-(p+1)/2&n+1-\tau_2.
\end{array}
\]
In particular, their sizes sum to $n-w$.

Put $\Psi=\mathscr T_{q-1,1}$ and $\Phi=\mathscr T_{q,0}$.
Define
\[
 \mathcal W_i=W_i\mathbb Q[Y]_{\le E_i-\deg W_i},\quad
 \mathcal V_i=V_i\mathbb Q[Y]_{\le E_i-\deg V_i},\quad
 \mathcal V_i^0=V_i\mathbb Q[Y]_{\le E_i-1-\deg V_i}.
\]
The three required cross Gram pairings are
\begin{equation}\label{z87:unequal-cross-tests}
 \Phi\big|_{\mathcal W_1\times\mathcal W_2},\qquad
 \Psi\big|_{\mathcal W_1\times\mathcal W_2},\qquad
 \Psi\big|_{\mathcal V_1^0\times\mathcal V_2^0}.
\end{equation}
Assume all three are nonsingular over $\mathbb Q$ and have
invertible reduction at $p$.  Their dimensions are respectively
$e,e,e-2$ in the upper regime and $e,e,e$ in the lower regime.
The determinant of the empty matrix is one.

The anti-invariance of $\Psi$ under $Y\mapsto-1-Y$, and the
last test, give unique monic left and right radical polynomials
$\mathcal R_i\in\mathcal V_i$ of degree $E_i$.  Each is even
in $Y+1/2$.  Indeed, the even root factors $V_i$ are invariant
under that reflection.  In centered monomial bases the full
cross pairing has zero equal-parity blocks.  The invertible
even-dimensional low matrix forces its two off-diagonal blocks
to be nonsingular, while adjoining the one extra even coordinate
on each side leaves exactly one left and one right radical.

Let $\psi_i$ be the monic degree-$E_i+1$ polynomial with roots
$W_i$ orthogonal to the opposite $\mathcal W_j$ for $\Phi$.
Existence and uniqueness follow from the first cross Gram test.
Define
\begin{equation}\label{z87:unequal-scalars}
 \nu=\Psi(Y\mathcal R_1\mathcal R_2),\qquad
 \theta=\Phi(\psi_1\psi_2),\qquad \kappa=\theta-\nu.
\end{equation}
The finite theorem only applies when the scalars explicitly
required below are units at $p$.  Their nonzero rational values
exclude at most finitely many additional primes.

\subsection{Even radical sections and their boundary form}

Truncate each window at $J_i$, thereby retaining its complete
even degree-$E_i$ band of top width $n+1$.  Define
\begin{equation}\label{z87:unequal-k}
 k_f=\begin{cases}
 (e/2)p+1-r,&\text{upper},\\
 (e/2+1)p+2-r-\tau_1-\tau_2,&\text{lower}.
 \end{cases}
\end{equation}
For any class set $S$, write $V_S(X)=\prod_{c\in S}(X-c^2)$.
The nonzero-class even second radicals have sections
\begin{equation}\label{z87:unequal-sections}
 P_1(c)=V_M(c^2)V_{\mathsf A}(c^2)Q_1(c^2),\qquad
 P_2(c)=V_M(c^2)V_{\mathsf A}(c^2)V_{\mathsf B}(c^2)Q_2(c^2),
\end{equation}
with parameter dimensions
\begin{equation}\label{z87:unequal-parameters}
 \dim Q_1=k_f+H_2,\qquad \dim Q_2=k_f+H_1.
\end{equation}
These numbers are nonnegative because they count actual complete
even top polynomials.  Each side also has $z$ free zero-class
parameters.  On its support a section has local polynomial
$P_i(c)\mathcal R_i(Y)$.

Here is the justification, including the points at which unequal
root levels could otherwise matter.  In the upper regime
$\mathcal V_i\subset\mathcal W_i$ has codimension one.
The full $\Psi$ pairing on $\mathcal W_1\times\mathcal W_2$
is nondegenerate.  Its restriction to
$\mathcal W_1\times\mathcal V_2$ therefore has full column
rank and its left radical is exactly $\mathcal R_1$; this
polynomial lies in $\mathcal V_1\subset\mathcal W_1$ and
annihilates $\mathcal V_2$.  The transposed assertion is
identical.  On $\mathcal V_1\times\mathcal V_2$ the radicals
are precisely $\mathcal R_1,\mathcal R_2$.  In the lower
regime $\mathcal W_i$ and $\mathcal V_i^0$ are both complements
to $\mathcal R_i$ in $\mathcal V_i$: otherwise their assumed
nonsingularity would fail.  Any left and right complements pair
nondegenerately in the quotient by the two radicals.  Consequently
the usual independent low features and correlated top tests give
exactly the supports in \eqref{z87:unequal-sections} in either
regime.  The complete widths $n+1$ make those tests available.
Subtracting the forced-zero factors gives
\eqref{z87:unequal-parameters}; for example, the upper dimensions
are $\tau_2-n$ and $\tau_1-n$.  The zero-class count is
$E_i/2-a_i-1=z$ on either side.

\paragraph{One-sided normalization.}
Use the literal local expansion of the integer-valued window
polynomial, $f_i(c+pY)=f_{i,0,c}(Y)+p f_{i,1,c}(Y)$, and
define $P_i(c)=[Y^{E_i}]f_{i,0,c}$.  No class-dependent scalar
is divided out of this coefficient.  The Lucas formula
\eqref{lyr:lucas-top} contains a factor $1/E_i!$, but this is
constant in $c$ and merely changes the basis of the polynomial
parameter space.  The same observation applies to the odd
coefficient $cS_i(c^2)$ and its constant factorial denominator.
Before either side's feature is inserted, the denominator
factorization and the universal-functional calculation give
the local weight
\[
 2c^5\frac{(N!/p^q)^2}{\overline U_c},\qquad
 D_N(-(c+pY)^2)=p^{h_c}L_c(Y)U_c(Y).
\]
This depends only on $N,p,c$, and not on $E_i,a_i$ or the chosen
window vectors.  In the present $\epsilon=0$ geometry it gives
$C_pc^5/V_M(c^2)$ on the nonmaximal classes and
$(C_p/2)c^4/V_M'(c^2)$ on the maximal classes, with the same
nonzero $C_p$.  Hence unequal endpoints introduce neither a
class-dependent ratio nor a different relative normalization
between the two strata.

The first local jet identity is separate on the two sides:
$[Y^{E_i+1}]f_{i,1,c}=P_i'(c)$.  The fixed Christoffel identity
$\Psi(F)=\Phi((Y-q)F)$ is unchanged.  Thus every calculation of
the central derivative and maximal Schur correction holds with
$\mathcal R_1\mathcal R_2$ in place of the square of one radical,
and $\psi_1\psi_2$ in place of the square of one maximal jet.
Lower jet terms vanish against the appropriate opposite radical;
on a mixed class the side with a zero leading section retains
the physical $V_i$ roots, so this statement also applies there.
On a class where both leading sections vanish there is no
nonmaximal first-order cross term.  In particular the nonmaximal
derivative form is
\[
 C_p\nu\sum_{c\notin M}\frac{d}{dc}
 \left(\frac{c^5P_1(c)P_2(c)}{V_M(c^2)}\right),
\]
while the maximal Schur term has coefficient $\theta$ on the
same pairs of derivative coordinates.  The zero class vanishes
at these orders by the same Newton divisibility bounds; its
$z$ independent directions contribute to both second kernels.

On $F=\mathbb F_p^w$ set
\[
 B(x,y)=\sum_{c\in M}\frac{c^6}{V_M'(c^2)}x_cy_c,
 \qquad S_i=\{(S(c^2))_{c\in M}:\deg S<H_i\}.
\]
For the even radical boundary use $E_i(c)=P_i'(c)/c$.
If $H_i\le w-3$, the unchanged finite-field degree bounds give
the simultaneous normal form
\begin{equation}\label{z87:unequal-normal}
\begin{array}{c|cc}
 &O_2&E_2\\\hline
O_1&\theta B+O(p)&p\kappa B+O(p^2)\\
E_1&p\kappa B+O(p^2)&p^2\kappa B+O(p^3).
\end{array}
\end{equation}
The $O_i$ coordinates range over $S_i$; the even boundary images
need not be surjective.  Eliminating the remaining nonmaximal
order-$p$ low block changes only the unidentified odd/odd
$O(p)$ term and subsequent remainders.  Its coupling to an even
radical starts at $p^2$, so it does not alter the displayed
even/even third form.

\subsection{The finite third-layer estimates}

\begin{theorem}[Unequal central odd endpoints]\label{z87:unequal-third}
Under the geometry and cross Gram hypotheses above, the following
bounds hold coefficientwise in $X$.

(i) Suppose $w\ge3$, $0\le H_i\le w-3$, and
$\theta\kappa\ne0$ after reduction.  Set
$\ell=(H_1+H_2-w+3)_+$.  Then
\begin{equation}\label{z87:unequal-unsaturated}
 v_p\bigl(\det(p^5A(X))\bigr)
 \ge r-(e+1)w+2k_f+2z+H_1+H_2-\ell.
\end{equation}
No sign condition on $k_f$ or evaluation-surjectivity condition
is required for this determinant inequality.

(ii) Suppose $w\ge3$, $k_f\ge3$, $0\le H_i\le n$, and
$\theta\nu\kappa\ne0$ after reduction.  If at least one odd
band has width $H_i\ge w-3$, then
\begin{equation}\label{z87:unequal-saturated}
 v_p\bigl(\det(p^5A(X))\bigr)
 \ge r-ew+2k_f+2z-3.
\end{equation}

(iii) Without scalar nonvanishing or third-layer degree guards,
the two-layer estimate
\begin{equation}\label{z87:unequal-second}
 v_p\bigl(\det(p^5A(X))\bigr)\ge r-ew+k_f+z
\end{equation}
holds whenever the same central geometry and three cross Gram
conditions hold.
\end{theorem}

\begin{proof}
For (i), the odd/odd boundary flag pairing has rank exactly
$\ell$: its entries are sums of $X^{h+j+3}/V_M'(X)$, whose
first nonzero anti-diagonal has index $h+j=w-4$ and entries one.
Eliminate these $\ell$ unit pivots.  Write
$S_i=L_i\oplus N_i$, where $\dim L_i=\ell$ and
$N_1=S_1\cap S_2^\perp$, $N_2=S_2\cap S_1^\perp$.
The even-only left and right subspaces annihilating the remaining
odd order-$p$ cross blocks each have dimension at least
$K=k_f+\ell+z$, by \eqref{z87:unequal-parameters}.
These are subspaces of the true second kernels, regardless of
the ranks of their boundary evaluation maps.  The induced third
form on them is
\[
 Q(x,y)=\kappa B(x,y)-\frac{\kappa^2}{\theta}
 B(x,L_2)B(L_1,L_2)^{-1}B(L_1,y),
 \quad x\in N_2^\perp,\quad y\in N_1^\perp.
\]
Its left and right radicals contain $N_1,N_2$, so its rank
is at most $C=w-H_1-H_2+2\ell$.  The restricted-second-kernel
lemma gives the valuation at least
$r-ew-\ell+2K-C$, which is \eqref{z87:unequal-unsaturated}.
If $K<0$, then $2K-C\le0$, and the first nullity alone gives
the same inequality.  The unknown odd/odd next coefficient
cannot affect this restriction: all its remaining couplings
to a chosen even second-kernel vector begin at order $p^2$,
and their Schur products through order-$p$ pivots begin at $p^3$.

For (ii), set $s_0=w-3$ and $S_0=\mathbb F_p[X]_{<s_0}$.
The restriction of $B$ to $S_0$ is nonsingular by the same
anti-diagonal identity.  The full second constraint on a left
even parameter $Q$ and an odd column polynomial $S$ is
\begin{equation}\label{z87:unequal-high-constraint}
 U(Q,S)=\kappa\sum_{c\in M}c^6V_{\mathsf A}(c^2)Q(c^2)S(c^2)
 +\nu n[X^{n-3}](V_{\mathsf A}QS),
\end{equation}
up to one common nonzero scalar.  The transposed formula uses
the forced-zero polynomial on that side.  This follows by
summing $X^3V_{\mathsf A}QS$ over the nonzero squares.  Its
degree is at most $n+H_2-w+2<2n$, and its constant coefficient
vanishes, so only its coefficient of $X^n$ survives.  Unequal
$E_i$ and $a_i$ do not enter this degree calculation.

On boundary-zero parameters $Q=V_MT$ the first term vanishes.
If $H_2\ge s_0$, put $d=H_2-s_0$.  Then $T$ has dimension
$k_f+H_2-w\ge d$ because $k_f\ge3$.  Its last $d$ monomials
pair with $X^{s_0},\ldots,X^{H_2-1}$ in an anti-triangular
matrix with nonzero anti-diagonal.  All columns in $S_0$ vanish.
Thus the exact odd kernel against boundary-zero even parameters
is $S_0$.  High odd constraints can be solved independently of
all prescribed maximal boundary values.

If $H_1\le s_0\le H_2$, the leading odd rank is $H_1$.
Choose its pivots inside $S_0$.  The residual low odd constraints
have rank $s_0-H_1$ on the boundary, and the high constraints
have rank $H_2-s_0$ on boundary-zero even parameters.  Both
second kernels therefore have dimension $k_f+H_1+z$; no odd
coordinates remain.  The third restriction has rank at most
$H_1+3$, the remaining left boundary dimension, because the
leading odd pivots used only pure boundary constraints.
The first three nullities are at least
\[
 r-ew-H_1,\qquad k_f+H_1+z,\qquad k_f+z-3,
\]
and give \eqref{z87:unequal-saturated}.  The transpose is identical.

For clarity, when $H_i>w$ the odd space throughout this argument
is the full polynomial space $\mathbb F_p[X]_{<H_i}$, not its
quotient by the maximal evaluation kernel.  The leading form
factors through its evaluation image, but
\eqref{z87:unequal-high-constraint} does not.  Its exact kernel
against boundary-zero even parameters is the polynomial space
$S_0$ itself.  Thus every evaluation-kernel direction outside
$S_0$ is included in the high constraints just proved.

If both odd polynomial spaces contain $S_0$, let their leading rank be
$\ell=w-3+s$, $0\le s\le3$.  Each residual odd direction
pairs injectively with boundary-zero even parameters, so both
second kernels have dimension $k_f+\ell+z$ and contain only
even coordinates.  Eliminating the $S_0$ pivots leaves a third
form on at most $w$ boundary coordinates.  Each of the further
$s$ leading odd pivots adds rank at most one, so the third rank
is at most $w+s$.  Their nullity sum is
\[
 (r-ew-\ell)+(k_f+\ell+z)+(k_f+\ell+z-w-s)
 =r-ew+2k_f+2z-3.
\]
Again order-$p$ odd cross eliminations only affect the next
order on the true second kernels.  This proves (ii).

Finally, let $\ell$ be the actual leading odd rank, with no
assumption on its value.  The first nullity is $r-ew-\ell$.
The even radical dimensions in \eqref{z87:unequal-parameters},
including $z$, lose at most $H_2-\ell$ and $H_1-\ell$
dimensions when tested against the residual odd directions.
This gives true second-kernel subspaces of dimension at least
$k_f+\ell+z$.  Adding the first and second nullities cancels
$\ell$ and proves (iii); a negative lower bound for the second
dimension is harmless.  All displayed forms at these orders
are constant in $X$, so the inequalities are coefficientwise.
\end{proof}

\paragraph{Normalized rates and uniformity.}
For $\beta=p/N$, write all interval and width symbols in their
normalized sense, and put $\sigma_i=s_i-a_i\beta$.  The new
odd candidates are
\[
 \begin{array}{c|c}
\text{regime}&\text{arithmetic exponent candidate}\\\hline
\text{unsaturated}&4\delta+(e+1)w-2k-H_1-H_2+(H_1+H_2-w)_+\\
\text{saturated}&4\delta+ew-2k\\
\text{two layers}&4\delta+ew-k,
 \end{array}
\]
where
\[
 k=\begin{cases}(e/2)\beta-\delta,&\text{upper},\\
 (e/2+1)\beta-\delta-\sigma_1-\sigma_2,&\text{lower}.
 \end{cases}
\]
Use only the corresponding normalized guards, with $k\ge0$
for the saturated rate.  A fixed finite collection of rational
cells has finitely many tuples $(q,E_1,E_2,a_1,a_2,\mathrm{half})$.
Each tuple has a fixed finite exceptional prime set from the
three determinants and the required scalars.  On cells bounded
away from $\beta=0$, these primes eventually disappear.  The
integer floors and the constants $z,3$ change each finite-prime
valuation by $O(1)$ uniformly on that collection.  Their total
weighted prime contribution is $O(N)$, with the same limiting
prime-integral argument as for the original central-root family.
No sampled matrix rank is used to justify this uniformity.

\paragraph{A finite check that motivated the extension.}
For $(q,E_1,a_1,E_2,a_2)=(4,8,0,10,1)$ in the upper regime,
the three cross Gram determinants and $\theta\nu\kappa$ are
nonzero over $\mathbb Q$.  Exact rational witnesses are recorded
separately.  At
\[
\begin{array}{c|rrrr|rrr|r}
p&N&m_1&m_2&r&n_1&n_2&n_3&k_f\\\hline
61&247&48&103&234&209&15&11&11\\
101&409&79&171&387&345&23&18&18\\
151&611&119&255&579&520&32&26&26
\end{array}
\]
the original-matrix Smith computations attain
\eqref{z87:unequal-saturated} exactly in the sum of their first
three nullities, namely $235,386,578$.  These checks illustrate
the theorem and are independent of its proof.

\section{Retaining maximal column radicals in a mixed endpoint chamber}
\label{z87:mixed-retained}

This is a second-layer refinement for $\epsilon=0$ with an even row
endpoint and an odd column endpoint.  Its new columns are true first
kernels, and their next pairing with the selected row first kernel
vanishes by a complete half-field sum.  Thus their dimensions may be
retained without increasing the second rank bound.

\subsection{Hypotheses and fixed rational tests}

Write $N=qp+w$, $n=(p-1)/2$, $M=\{1,\ldots,w\}$, and retain the
good-prime and complete-feature hypotheses of the general arithmetic
layer.  Suppose
\[
 m_i=a_ip+\tau_i,\quad
 D_1=E_1\text{ even},\quad D_2=E_2+1\text{ odd},\quad
 E_1-2a_1=E_2-2a_2=e\ge2,
\]
and impose the geometric guards
\begin{equation}\label{z87:mixed-retained-geometry}
 0<\tau_1<p/2<\tau_2<p,\qquad
 3\le w,\qquad \tau_1\le w\le u_2=p-\tau_2.
\end{equation}
Let
\[
 A=\{1,\ldots,\tau_1-1\},\quad B=M\setminus A,
 \quad C=\{w+1,\ldots,u_2\},\quad D=\{u_2+1,\ldots,n\},
\]
and set $a_0=|A|=\tau_1-1$, $d_0=|D|=n-u_2$.
The physical roots are
\[
\begin{array}{c|cccc}
\text{class set}&A&B&C&D\\\hline
\text{row}&W_1&V_1^-&V_1^-&V_1^-\\
\text{column}&W_2&W_2&W_2&V_2^+
\end{array}
\]
where $W_i=Z_{2a_i+1}$, $V_1^-=Z_{2a_1}$, and
$V_2^+=Z_{2a_2+2}$.

The actual top widths are
\[
 H_1=R_1-E_1p/2,\qquad
 H_2=R_2-((E_2+1)p+1)/2.
\]
For the explicit bound below assume $0\le H_1,H_2\le w$.
Write $\Phi=\mathscr T_{q,0}$ and $\Psi=\mathscr T_{q-1,1}$.
Require nonsingularity, over $\mathbb Q$ and after reduction at $p$,
of the following three square cross Gram matrices:
\begin{equation}\label{z87:mixed-retained-tests}
\begin{array}{c|cc|c}
\text{functional}&\text{row space}&\text{column space}&\text{size}\\\hline
\Phi&W_1\mathbb Q[Y]_{\le E_1-1-\deg W_1}
     &W_2\mathbb Q[Y]_{\le E_2-1-\deg W_2}&e-1\\
\Phi&W_1\mathbb Q[Y]_{\le E_1-\deg W_1}
     &W_2\mathbb Q[Y]_{\le E_2-\deg W_2}&e\\
\Psi&V_1^-\mathbb Q[Y]_{\le E_1-3-\deg V_1^-}
     &V_2^+\mathbb Q[Y]_{\le E_2-1-\deg V_2^+}&e-2.
\end{array}
\end{equation}
Empty determinants are one.  The first test produces monic
$\chi_i$ of degree $E_i$, with root factor $W_i$, orthogonal
to the opposite degree-$E_j-1$ space for $\Phi$.  The second
test implies
\[
 \lambda=\Phi(\chi_1\chi_2)\ne0.
\]
The last test and centered anti-invariance produce the monic
left and right radicals $\mathcal R_1,\mathcal R_2$ of degrees
$E_1-2,E_2$ with roots $V_1^-,V_2^+$, respectively.  They
annihilate the full opposite spaces with those endpoint degrees.
Require in addition
\begin{equation}\label{z87:mixed-retained-nu}
 \nu=\Psi(Y\mathcal R_1\mathcal R_2)\ne0
\end{equation}
over $\mathbb Q$ and after reduction at $p$.
All these are fixed rational conditions on $(q,E_1,E_2,a_1,a_2)$.

\subsection{The extra columns and the determinant bound}

Put
\[
 A_1=ew-a_0+H_1,\qquad A_2=ew+H_2,\qquad
 B_T=e(n-w)-d_0,\qquad g=(a_0-H_1-2)_+.
\]
The quantities $A_i$ are the exact maximal restriction ranks.
The two ordinary maximal restriction kernels have dimensions
$r-A_i$.  Their top coefficients vanish identically because
$H_i\le w$, and the nonmaximal second pairing has rank at most
$B_T$: it has local dimensions $e,e$ on $C$ and $e,e-1$ on $D$.

\begin{theorem}[Mixed retained maximal columns]
\label{z87:mixed-retained-theorem}
Under the displayed hypotheses,
\begin{equation}\label{z87:mixed-retained-finite}
 v_p\bigl(\det(p^5A(X))\bigr)
 \ge 2r-A_1-A_2-B_T+g
\end{equation}
coefficientwise in $X$.
\end{theorem}

\begin{proof}
Let $U_1,U_2$ be the ordinary maximal restriction kernels.  Since
the row even top coefficient vanishes, an integral lift of a vector
in $U_1$ may be chosen entirely in the row band truncated at
$E_1p/2$.  Its endpoint degree is $E_1-1$, and its complete odd
top band has width $n$.  Its leading coefficient therefore has
the form
\begin{equation}\label{z87:mixed-row-odd-top}
 T(c)=cV_M(c^2)Q(c^2),\qquad \deg Q<n-w.
\end{equation}
On $T=M^c$ the row root factor is always $V_1^-$.

Consider the polynomial space of $S(X)$ with
\[
 \deg S\le w-3,\qquad S|_B=0.
\]
It has dimension $(a_0-2)_+$.  Impose the at most $H_1$
linear equations
\begin{equation}\label{z87:mixed-leading-constraints}
 \sum_{c\in A}\frac{c^4P(c^2)S(c^2)}{V_M'(c^2)}=0
 \qquad(\deg P<H_1).
\end{equation}
The remaining space has dimension at least $g$.

For such an $S$, choose a column vector $v_S$ with zero odd
top coefficient and zero zero-class feature, and with independent
low local features
\begin{equation}\label{z87:mixed-column-completion}
 f_{2,0,c}(Y)=
 \begin{cases}
 S(c^2)\chi_2(Y),&c\in A,\\
 0,&c\in B,\\
 (\lambda/\nu)S(c^2)\mathcal R_2(Y),&c\in C\cup D.
 \end{cases}
\end{equation}
These features belong to the actual column root spaces:
$\mathcal R_2$ has roots $V_2^+$, which include $W_2$ on $C$.
They have degree at most $E_2$, so completeness of the column
low features constructs an actual window vector.

The vector $v_S$ is a true first column kernel.  On $A$, its
feature annihilates every row low feature by the definition of
$\chi_2$.  A row top feature with coefficient $P(c^2)$ pairs
with it by $\lambda P(c^2)S(c^2)$, because its difference from
$P(c^2)\chi_1$ has degree at most $E_1-1$.  The universal
maximal weight is $(C_p/2)c^4/V_M'(c^2)$, so
\eqref{z87:mixed-leading-constraints} kills all remaining top
pairings.  On $B$ its feature is zero.  Thus it annihilates the
whole original row window in the leading form.

Next evaluate the order-$p$ form of $p^5A$ between $u\in U_1$
and $v_S$.  On a maximal class the row leading feature is zero.
Consequently the only term is the universal functional on its
first jet and the column leading feature; no derivative of the
unit multiplier remains.  The highest first-jet identity is valid
also for an odd endpoint:
\[
 [Y^{E_1}]f_{1,1,c}=T'(c).
\]
To verify this directly, clear the fixed denominator valuation
of a factorial basis vector.  Its numerator modulo $p$ is a
constant times $(x^p-x)^{E_1-1}F(x)$ with $\deg F<p$.
At $x=c+t$, its coefficients of $t^{E_1-1}$ and $t^{E_1}$
are the same constant times $F(c)$ and $F'(c)$, including the
common sign.  Substituting $t=pY$ proves the identity.  Lower
bands have at least two extra powers of $p$ at degree $E_1$.
All factorial factors are constants in $c$; the expansion is
the literal local polynomial, without a class-dependent rescaling.

On $A$, the first jet retains the physical roots $W_1$ and
$\chi_2$ kills its terms of degree at most $E_1-1$.  At a
maximal node, \eqref{z87:mixed-row-odd-top} gives
$T'(c)=2c^2V_M'(c^2)Q(c^2)$.  The entire maximal second
contribution is therefore
\[
 C_p\lambda\sum_{c\in A}c^6Q(c^2)S(c^2).
\]
On a nonmaximal class, every row feature of degree at most
$E_1-2$ is annihilated by $\mathcal R_2$.  A monic representative
of its odd top is $Y\mathcal R_1$, so the local scalar on that
top is $\nu$.  The weight is $C_pc^5/V_M(c^2)$ with the same
$C_p$, independently of either endpoint degree.  Using
\eqref{z87:mixed-column-completion}, the nonmaximal contribution is
\[
 C_p\lambda\sum_{c\in C\cup D}c^6Q(c^2)S(c^2).
\]
The zero class contributes nothing at this scale by the usual
Newton and Bernoulli bounds (and its chosen column feature is zero).
Since $S$ vanishes on $B$, the two contributions sum to
\begin{equation}\label{z87:mixed-half-field-zero}
 C_p\lambda\sum_{c=1}^{n}(c^2)^3Q(c^2)S(c^2)=0.
\end{equation}
Indeed the polynomial $X^3QS$ has zero constant coefficient and
degree at most $3+(n-w-1)+(w-3)=n-1$.  Each of its monomials
has zero sum over the nonzero squares.  If $Q=0$, the assertion
is immediate.

Thus every $v_S$ has zero second pairing with all of $U_1$.
Its maximal feature is nonzero unless $S=0$, since $\deg S<w$
and $S$ already vanishes on $B$.  The $g$ constructed directions
are therefore independent modulo $U_2$.  Enlarge the true first
column kernel to $\widetilde U_2=U_2+\{v_S\}$, of dimension
at least $r-A_2+g$.  The second form on
$U_1\times\widetilde U_2$ has the same rank bound $B_T$ as on
$U_1\times U_2$.  The restricted-first-kernel determinant lemma
now gives \eqref{z87:mixed-retained-finite}.  The first Schur
correction starts at $p^2$ because both chosen spaces are true
first kernels; it does not change the order-$p$ form.  These
orders are constant in $X$, so the estimate is coefficientwise.
\end{proof}

\paragraph{Normalized candidate.}
Set $k_0=(e/2)p-r$.  The actual width identities give
\[
 a_0-H_1=k_0-1,\qquad d_0=k_0+H_2,
 \qquad A_1+A_2+B_T=e(n+w)-2k_0+1.
\]
Consequently the extra dimension is $(k_0-3)_+$.  On a fixed
rational rate cell with $k=e\beta/2-\delta>0$, the new
arithmetic candidate simplifies to
\begin{equation}\label{z87:mixed-retained-rate}
 E_{\mathrm{mixed-retained}}=6\delta+ew-e\beta.
\end{equation}
It improves the corresponding complete-feature candidate by
$k$.  Use the geometric and width guards, their transpose, and
the fixed local certificates.  All bounded offsets affect the
prime-weighted sum by $O(N)$, and fixed nonzero rational tests
exclude only finitely many primes on each cell.

\paragraph{Independent diagnostic.}
For $(q,E_1,E_2,a_1,a_2)=(3,8,8,1,1)$, the three rational
determinants are
\[
 \frac{312402369935}{1990656},\qquad
 -\frac{186500250298574837}{15479341056},\qquad
 \frac{133225}{1024},
 \quad\text{and}\quad \nu=-\frac{90487}{315360}.
\]
At $(p,N,m_1,m_2,r)=(151,507,188,242,443)$ one has
$a_0=36$, $H_1=27$, $H_2=5$, $g=7$.  The original
complete-feature lower bound for the first two nullities is
$131$, the new bound is $138$, and the actual sum is $140$.
The difference of two is bounded; the recovered growing term
is exactly the observed maximal-radical freedom $a_0-H_1=9$.

\section{A certificate-driven completion of one-sided maximal kernels}
\label{z87:kernel-completion}

The following construction retains correlated top directions in a true
first kernel, while completing their nonmaximal features so that the
next pairing has a global coefficient-rank bound.  It applies to the
original determinant family at $\epsilon=1$ and does not depend on an
experimental Smith-form identity.

\subsection{Fixed local feasibility conditions}

Retain the general good-prime and complete-feature hypotheses.  Put
$n=(p-1)/2$, let $M$ be the maximal classes, $w=|M|$, and let
$T$ be their nonempty complement.  The actual endpoint degrees and
top widths are $D_i,H_i$, with $\alpha_i=D_i\bmod2$.
Write $m_i=a_ip+\tau_i$, $0\le\tau_i<p$, and impose the
conservative availability guard $D_i-2a_i\ge2$ on both sides.
This ensures that the zero-class low features can meet all
physical root evaluations for every prescribed correlated top.
Write
\[
 \Phi=\mathscr T_{q,1},\qquad \Psi=\mathscr T_{q,0},\qquad
 k_R=\max_{c\in T}k_1(c),\qquad k_C=\min_{c\in T}k_2(c).
\]
Assume $k_R\le D_1$, $k_C\le D_2$, and that a fixed rational
monic polynomial $\chi$ of degree $D_1$ satisfies
\begin{equation}\label{z87:completion-nonmax-tests}
 Z_{k_R}\mid\chi,\qquad
 \Psi(\chi Z_{k_C}Y^j)=0\quad(0\le j<D_2-k_C).
\end{equation}
Define the rational scalar
\begin{equation}\label{z87:completion-theta}
 \theta=\Psi(\chi Z_{k_C}Y^{D_2-k_C}).
\end{equation}
Thus $\chi$ is allowed on every actual nonmaximal row class,
and its pairing with any actual nonmaximal column feature $F$
is exactly $\theta[Y^{D_2}]F$.

Select a collection $\mathcal G$ of maximal root types $(a,b)$.
For each selected type assume $a\le D_1$, $b\le D_2+1$, and
the existence of a fixed rational monic degree-$D_1$ polynomial
$G_{a,b}$ satisfying
\begin{equation}\label{z87:completion-max-tests}
 \begin{split}
 &Z_a\mid G_{a,b},\\
 &\Phi(G_{a,b}Z_bY^j)=0\quad(0\le j\le D_2-b),\\
 &\Phi(G_{a,b}Z_bY^{D_2+1-b})=\theta.
 \end{split}
\end{equation}
When $b=D_2+1$, the middle set of equations is empty.
All coefficients and denominators in these fixed rational systems
must reduce at $p$, and the stated identities must remain valid
after reduction.  These conditions exclude only finitely many
primes for any fixed certified tuple.

Call a maximal class good if $(k_1(c),k_2(c))\in\mathcal G$.
Let $F\subseteq M$ be the remaining bad classes and $f=|F|$.
No assertion is made that unselected types are infeasible: they
are simply omitted from this particular construction.

\paragraph{A simple determinant-only sufficient test.}
One need not print large rational polynomial coefficients when
the defining systems are square.  For example, suppose
$D_1-k_R=D_2-k_C=h$.  The $h$ unknown low coefficients of
$\chi=Z_{k_R}(Y^h+\cdots)$ satisfy the square cross Gram
system in \eqref{z87:completion-nonmax-tests}.  A nonzero
rational determinant gives a unique $\chi$ and hence a fixed
$\theta$.  If a proposed maximal type is
$(a,b)=(k_R-1,k_C+1)$, then the annihilation equations and the
jet equation in \eqref{z87:completion-max-tests} form a square
system of size $h+1$ for the low coefficients of $G_{a,b}$.
Nonsingularity of this second system realizes every rational
right side, in particular the already defined $\theta$.
Its rows and columns are the $\Phi$ cross Gram spaces of
degrees at most $D_1-1$ and $D_2+1$, with roots $Z_a,Z_b$.
This is an ordinary fixed rational nonsingularity certificate;
modular nonzero determinants certify it exactly.

\subsection{Global completion and its rank}

Let $U_i$ be the usual maximal restriction kernels, with
dimensions $r-A_i$, where
\[
 A_i=\sum_{c\in M}(D_i-k_i(c))_++\min(H_i,w).
\]
Put
\begin{equation}\label{z87:completion-dimensions}
 K_2=(H_2-w)_+,\qquad
 g=(H_1-f)_+-(H_1-w)_+.
\end{equation}
The number $g$ is nonnegative because $f\le w$.
Define a rank bound $R$ by
\begin{equation}\label{z87:completion-rank}
 R=\begin{cases}
 0,&\theta=0,\\
 \min\{g,K_2,(H_1+K_2-n+2+\alpha_1)_+\},
       &\theta\ne0,\ \alpha_1=\alpha_2,\\
 \min\{g,K_2\},&\theta\ne0,\ \alpha_1\ne\alpha_2.
 \end{cases}
\end{equation}
For a uniform bound when the rational scalar is not known to
vanish, always use the second or third line.  Its reduction
vanishing at an occasional prime can only improve the rank.

\begin{theorem}[One-sided maximal kernel completion]
\label{z87:completion-theorem}
Under the fixed local conditions above, there is a true first
row-kernel space $\widetilde U_1\supseteq U_1$ of dimension
$r-A_1+g$ such that the added order-$p$ pairing with $U_2$
has rank at most $R$.
Consequently, if $B$ bounds the second pairing on
$U_1\times U_2$, then coefficientwise in $X$,
\begin{equation}\label{z87:completion-valuation}
 v_p\bigl(\det(p^6A(X))\bigr)
 \ge 2r-A_1-A_2-B+g-R.
\end{equation}
The same conclusion with the same improvement $g-R$ applies
to a second-layer bound obtained by further restricting either
of the two ordinary maximal restriction kernels and charging
the lost dimensions.  The transposed construction is also valid.
\end{theorem}

\begin{proof}
Choose a polynomial $P(X)$ of degree less than $H_1$ that
vanishes on $F$, and prescribe the row leading local features
\begin{equation}\label{z87:completion-row-features}
 f_{1,0,c}(Y)=
 \begin{cases}
 c^{\alpha_1}P(c^2)G_{k_1(c),k_2(c)}(Y),&c\in M\setminus F,\\
 0,&c\in F,\\
 c^{\alpha_1}P(c^2)\chi(Y),&c\in T.
 \end{cases}
\end{equation}
Every displayed polynomial has the actual row roots and the
same globally correlated monic top coefficient
$c^{\alpha_1}P(c^2)$.  Complete low features therefore give
an actual row window vector.  Choose any compatible zero-class
low features; the zero class does not contribute at the two
scales under consideration.
More explicitly, its at most $a_1+1$ independent even root
evaluations can be solved by the zero-class low space, whose
dimension is $\lceil D_1/2\rceil\ge a_1+1$ under the
availability guard.  Thus a forced zero-class top coefficient
causes no extra condition on $P$; it is not being set to zero
independently of its Lucas value.

These are true first row kernels: at each good maximal class,
\eqref{z87:completion-max-tests} annihilates the entire actual
full column space, and the row feature is zero on the bad
classes.  The space of allowed $P$ has dimension $(H_1-f)_+$.
Its intersection with the ordinary restriction kernel consists
exactly of the polynomials divisible by $V_M$, of dimension
$(H_1-w)_+$.  Indeed each $G_{a,b}$ is monic and nonzero, and
$c\ne0$ on every class.  A complement therefore adds exactly
$g$ independent first-kernel row directions.

For a column vector in $U_2$, its leading top coefficient is
\begin{equation}\label{z87:completion-column-top}
 T_2(c)=c^{\alpha_2}V_M(c^2)Q(c^2),\qquad \deg Q<K_2.
\end{equation}
On a maximal class its leading feature vanishes, so its first
jet retains the actual roots $Z_{k_2(c)}$.  The literal highest
first-jet identity, valid for either endpoint parity, is
$[Y^{D_2+1}]f_{2,1,c}=T_2'(c)$.  After subtracting this
coefficient times the monic polynomial
$Z_{k_2(c)}Y^{D_2+1-k_2(c)}$, its remaining first jet has
degree at most $D_2$ and is annihilated by $G_{k_1(c),k_2(c)}$.
The local pair at the next scale is consequently
$\theta c^{\alpha_1}P(c^2)T_2'(c)$ on a good class.
There is no unit-derivative term because the column leading
feature is zero.  The row first jet also multiplies that zero
feature and contributes nothing.

Use the exact universal weights
\[
 \omega_{M,c}=C_p\frac{c^5}{V_M'(c^2)},\qquad
 \omega_{T,c}=2C_p\frac{c^6}{V_M(c^2)}.
\]
They are independent of either endpoint degree and root level.
At a maximal node,
$T_2'(c)=2c^{\alpha_2+1}V_M'(c^2)Q(c^2)$.
Hence the maximal order-$p$ contribution is
\[
 2C_p\theta\sum_{c\in M\setminus F}
 c^{6+\alpha_1+\alpha_2}P(c^2)Q(c^2).
\]
On $T$, condition \eqref{z87:completion-nonmax-tests} kills
all lower column coefficients.  Its scalar on the top is
$\theta$, so \eqref{z87:completion-column-top} gives the same
expression summed over $T$.  Since $P$ vanishes on $F$, the
total added pairing is exactly
\begin{equation}\label{z87:completion-full-sum}
 2C_p\theta\sum_{c=1}^{n}
 c^{6+\alpha_1+\alpha_2}P(c^2)Q(c^2).
\end{equation}
The usual Newton, quotient-interpolation, and Bernoulli
precision estimates justify these reductions exactly as in the
general second layer.  On a maximal class only one local factor
has been divided by $p$; its product has the same interpolation
remainder, and the constant quotient ambiguity disappears under
three derivatives.  The first Schur correction starts at $p^2$
because both spaces are true first kernels, so it cannot alter
the displayed order-$p$ form.

If $\theta=0$, this form is zero.  If the parities agree, put
$\gamma=3+\alpha_1$.  It is a constant times the half-field
sum of $X^\gamma P(X)Q(X)$.  The monomial entry indexed by
$i,j$ vanishes whenever $\gamma+i+j<n$, since its exponent
is positive and less than the order $n$ of the nonzero squares.
Thus the full coefficient matrix has rank at most
$(H_1+K_2-n+\gamma-1)_+$, capped by its two dimensions.
Restricting to the chosen $g$-dimensional complement gives
the second line of \eqref{z87:completion-rank}.  No upper bound
on the total polynomial degree is needed: higher multiples of
$n$ cannot create an entry below that same zero region.
With opposite parities the dimension bound $\min(g,K_2)$
always holds.  This proves all three cases of the rank bound.

Adding these rows to the existing pairing raises its second
rank by at most $R$.  The restricted-first-kernel determinant
lemma now yields \eqref{z87:completion-valuation}.  If either
original kernel had already been further restricted, the new
directions remain independent modulo that smaller row space,
and the added pairing has no larger rank on the smaller column
space.  The same improvement therefore holds.  All forms at
these orders are constant in $X$, proving the coefficientwise
assertion.
\end{proof}

\paragraph{Rate form and certification scope.}
On a fixed rational $\beta=p/N$ cell use normalized $H_i,w,f$
and put
\[
 K_2=(H_2-w)_+,\qquad g=(H_1-f)_+-(H_1-w)_+.
\]
\[ 
 \quad R_{\mathrm{rate}}=
 \begin{cases}
 \min\{g,K_2,(H_1+K_2-\beta/2)_+\},&D_1\equiv D_2\pmod2,\\
 \min\{g,K_2\},&D_1\not\equiv D_2\pmod2.
 \end{cases}
\]
Use $R_{\mathrm{rate}}=0$ if $\theta=0$ has been proved over
$\mathbb Q$.  Subtract $g-R_{\mathrm{rate}}$ from a compatible
feature-restriction second-layer candidate, not from an unrelated
minimum of arithmetic estimates.  A fixed rational cutoff has
finitely many root-type tuples and exceptional primes.  Bounded
floor and coefficient offsets contribute $O(N)$ to the weighted
prime sum, so the usual prime-integral limit applies uniformly.

For the active floating signatures $(q,D_1,D_2,k_R,k_C)
=(19,48,48,5,5)$ and $(28,70,70,7,7)$, the selected maximal
types are $(4,6)$ and $(6,8)$ respectively.  Both defining
systems are square, of sizes $43,44$ and $63,64$.  Their
ordinary rational nonsingularity can be certified by modular
determinants; no modular equality of two unknown scalar values
is being promoted to a rational identity.

\subsection{The corresponding epsilon-zero completion}
\label{z87:completion-epsilon-zero}

There is a direct version for $\epsilon=0$.  Keep the local
feasibility conditions and all availability guards, but now use
$M=\{1,\ldots,w\}$, $w=N-qp$, and
\[
 \Phi=\mathscr T_{q,0},\qquad \Psi=\mathscr T_{q-1,1}.
\]
These are the functionals on the physical sets with $2q+1$
and $2q$ nodes respectively; there is no deletion of a physical
node after the first elimination.  Define $g,K_2$ as before.
The rank bound is now
\begin{equation}\label{z87:completion-epsilon-zero-rank}
 R_0=\begin{cases}
 0,&\theta=0,\\
 \min\{g,K_2,(H_1+K_2-n+2)_+\},&\alpha_1+\alpha_2=1,\\
 \min\{g,K_2\},&\alpha_1=\alpha_2.
 \end{cases}
\end{equation}
As before, if a rational zero scalar has not been proved, use
the corresponding nonzero-scalar row even if its reduction
occasionally vanishes.  The finite valuation becomes
\[
 v_p\bigl(\det(p^5A(X))\bigr)
 \ge 2r-A_1-A_2-B+g-R_0.
\]

To prove this, construct exactly the same actual first-kernel
rows and pair them with the ordinary column restriction kernel.
The only change is the universal weight pair
\[
 \omega_{M,c}=(C_p/2)\frac{c^4}{V_M'(c^2)},\qquad
 \omega_{T,c}=C_p\frac{c^5}{V_M(c^2)}.
\]
Together with
$T_2'(c)=2c^{\alpha_2+1}V_M'(c^2)Q(c^2)$, the added
second form is
\[
 C_p\theta\sum_{c=1}^{n}
 c^{5+\alpha_1+\alpha_2}P(c^2)Q(c^2).
\]
Opposite endpoint parities make this the half-field polynomial
sum $X^3PQ$, giving the terminal rank in
\eqref{z87:completion-epsilon-zero-rank}.  Equal parities leave
an odd power of $c$, for which the dimension bound always
holds.  The same dimension, Schur, and uniformity arguments
complete the proof.  In the normalized opposite-parity case
the terminal rank is $\min\{g,K_2,(H_1+K_2-\beta/2)_+\}$.

\subsection{A common polynomial for both local functionals}
\label{sec:single-orthogonal}

Let $\Phi$ be the functional on the maximal physical node set and
$\Psi$ the functional on the nonmaximal set. The defining quotient and
residue formulas give, over $\mathbb Q$,
\[
 \Psi(F)=\Phi((Y-\ell)F),\qquad
 \ell=\begin{cases}q,&\epsilon=0,\\-q-1,&\epsilon=1.\end{cases}
\]
Indeed the maximal denominator has the additional linear factor
$Y-\ell$. It cancels in the quotient and at all common nodes, while the
residue at the added node vanishes. This extends the use of
\eqref{conf:christoffel} to a simultaneous completion.

\begin{proposition}\label{prop:common-christoffel}
Let $D_r,D_c$ be the endpoint degrees in either orientation.
Choose integers $U,V$ such that every nonmaximal class has
$k_r\le U$, $k_c\ge V$. Select all maximal classes satisfying these
same inequalities. Write $h=D_r-U$ and $b=D_c-V\ge0$.
For $s=b+1$ or $s=b+2$, suppose $h\ge s$ and
\[
 \det\bigl(\Phi(Z_U Z_V Y^{i+j})\bigr)_{0\le i,j<s}\ne0.
\]
The hypotheses for the one-sided kernel completion hold on these
selected classes. If $s=b+2$, the common next coefficient is zero.
\end{proposition}
\begin{proof}
Solve the displayed nonsingular system for the $s$ low coefficients of
a monic $G=Z_U(Y^h+\sum_{j=0}^{s-1}a_jY^j)$ so that
$\Phi(GZ_VY^i)=0$ for $0\le i<s$. Because the $Z_j$ form a nested
monic sequence, $G$ meets every prescribed row root. Each column of
degree at most $D_c$ having the actual column roots is of the form
$Z_V Q$, $\deg Q\le b$, and is annihilated by $G$ under $\Phi$.
For an admissible column of degree $D_c+1$ the result under $\Phi$
is its leading coefficient times
$\theta=\Phi(GZ_VY^{b+1})$.
On a nonmaximal class, $\Psi(GZ_VY^j)$ vanishes for $j<b$, and
$\Psi(GZ_VY^b)=\theta$, by the factor identity. Thus the same $G$
works wherever a nonzero row is prescribed. The rank and valuation
conclusions in Section~\ref{z87:kernel-completion} apply with the
parity for the actual $\epsilon$. The extra equation for $s=b+2$
sets $\theta=0$ over $\mathbb Q$.
\end{proof}

The finite check uses a nonzero determinant modulo a denominator-unit
auxiliary prime to prove rational nonsingularity. All scalar equalities
follow from the polynomial identity. For each interval in $\beta=p/N$
we subtract $g-R$ from the ordinary two-kernel feature profile with
its endpoint restrictions and paid Cauchy roots, then take the
pointwise minimum with the inherited estimates. Every normalized
profile is affine between its recorded rational breakpoints.

\paragraph{Independent completions with a prescribed coefficient.}
There is also a rectangular sufficient test that allows a different
monic polynomial on each maximal type. Use the extreme $U,V$ on the
nonmaximal classes. With $h=D_r-U$, $b=D_c-V$, solve the $b$ equations
$\Psi(\chi Z_VY^j)=0$ for $0\le j<b$ by using the $b$ lowest
unknown coefficients. It is sufficient that $h\ge b$ and the
corresponding $b$-row principal cross Gram determinant is nonzero.
If $h\ge b+1$, a nonzero determinant of order $b+1$ allows one to
impose a zero next coefficient as well. Otherwise take the next
coefficient from the rational solution for $\chi$.
For a maximal type $(a,b')$, put $s=D_c-b'+2$. Provided
$0\le s\le D_r-a$, a nonzero order-$s$ cross Gram determinant for
$\Phi(Z_aZ_{b'}Y^{i+j})$ solves the annihilation equations and the
one prescribed next coefficient simultaneously: its right side
prescribes zeros followed by the rational coefficient just obtained.
Thus any collection of maximal types passing this test may be
selected in the one-sided completion theorem. Only determinant
nonvanishing is checked modulo the auxiliary prime.

\paragraph{Separate local systems on selected nonmaximal classes.}
Let $F_M,F_T$ be the excluded maximal and nonmaximal classes,
of sizes $f$ and $v$. Prescribe the common row top
$c^{\alpha_1}P(c^2)$ to vanish on both sets. On each retained type
choose its own monic local polynomial of degree $D_r$ with its actual
row roots. On a maximal type with column root count $b$ require
annihilation through degree $D_c$ and the scalar $\theta$ at the next
degree. On a nonmaximal type require annihilation through degree
$D_c-1$ and the same $\theta$ at degree $D_c$. On excluded types
prescribe the entire leading local feature to be zero.
The complete-feature and zero-class arguments in
Section~\ref{z87:kernel-completion} realize these prescriptions.
The new row dimension modulo the ordinary maximal restriction kernel is
\[
 g=(H_r-f-v)_+-(H_r-w-v)_+.
\]
Indeed the allowed polynomials vanish at $f+v$ distinct square nodes
and their intersection with the old kernel vanishes at $w+v$.
The first-layer pairing vanishes. In the second layer, the identical
scalar gives the same half-field sum as before, because $P$ vanishes
at every excluded node. Consequently the previous $g-R$ valuation
gain holds with this new dimension and its corresponding rank cap.

It is sufficient for the selected types to solve finite rational
linear systems independently. For a row root count $a$ write the
polynomial as $Z_a(Y)(Y^{D_r-a}+\sum_{j<D_r-a}u_j Y^j)$.
Impose all annihilation equations, followed by the specified scalar
equation, using the corresponding local functional. Full row rank
of this rectangular system guarantees a solution for any $\theta$.
One distinguished type can instead determine $\theta$ after the
annihilation equations alone have been solved. If every selected type
can prescribe the scalar, choose $\theta=0$. A nonzero maximal minor
over a denominator-unit auxiliary prime certifies full row rank over
$\mathbb Q$. The discarded classes and their widths enter the
dimension formula before the profile is integrated.

\subsection{Local null spaces and the rank of the first jet}
\label{sec:radical-extension}
Keep the hypotheses of Theorem~\ref{z87:completion-theorem} in the
appropriate parity, including $D_i-2\lfloor m_i/p\rfloor\ge2$.
For a maximal class with row and column root counts $(a,b)$, form the
matrix $A$ of the pairing $\Phi(Z_aY^j Z_bY^i)$ with
$0\le j<D_r-a$ and $0\le i\le D_c-b$. Let $A^+$ include the
additional row $i=D_c+1-b$. Let $l$ be a guaranteed dimension in
$\ker A$, and $z$ a guaranteed dimension in $\ker A^+$, with
$0\le l-z\le1$ and $z\le l$. Put
\[
 L=\sum_{c\in M}l_c,\qquad D=\sum_{c\in M}(l_c-z_c).
\]
These are dimensions for one physical prime. The integral uses the
dimensions divided by $N$, with each class type weighted by its
normalized width.

\begin{proposition}\label{prop:local-radical}
On each selected row type allow a monic solution of the annihilation
equations with its own next coefficient $\theta_c$. Suppose the top
of the row vanishes on excluded types, as above. Let $g$ be its
dimension modulo the ordinary maximal kernel. Fix a reference
$\theta_*$, and let $E$ count the selected nodes at which
$\theta_c\ne\theta_*$. Let $J=0$ for $\theta_*=0$, and otherwise
let $J$ be any of the preceding half-field or dimension bounds on
the full constant-coefficient sum on the space of row polynomials
$\deg P<H_r$ and column polynomials $\deg Q<K_c$. The improvement
over the compatible two-kernel valuation is at least
\[
 g+L-\min\{g+L,K_c,J+E+D\}.
\]
\end{proposition}
\begin{proof}
On each maximal class choose an $l_c$-dimensional subspace of
$\ker A$ containing a $z_c$-dimensional subspace of $\ker A^+$.
Lemma~\ref{lyr:complete-features} realizes these directions with
zero correlated top, zero leading features in every other class,
and the prescribed roots on the chosen class. Their nonzero maximal
restrictions are independent. They are also independent modulo
the ordinary kernel of the $g$ added top directions, since the
latter are distinguished by their top coefficients on $M$.
All $g+L$ added directions are true first row kernels.

For a column in the ordinary maximal kernel the divided leading
feature on a maximal class has degree at most $D_c+1$.
Its coefficient of degree $D_c+1$ is
$T_c'(c)=2c^{1+\alpha_c}V_M'(c^2)Q(c^2)$.
The remaining polynomial is annihilated by each low radical row.
Thus the local rows contribute only through $Q(c^2)$ and their
first-jet functional. At least $z_c$ of their independent directions
give zero, so all low rows together contribute rank at most $D$.
For the added tops the previous half-field computation now has the
weight $\theta_c$ at each retained node, and $P(c^2)=0$ on the
excluded ones. Splitting off the constant $\theta_*$ leaves at most
$E$ rank-one point evaluations. The resulting pairing factors through
the $K_c$ coefficients of $Q$ and has at most $g+L$ independent
rows. Adding its rank bound to the previous second-layer bound
gives the formula, by the same lattice argument as in
Theorem~\ref{z87:completion-theorem}.
\end{proof}

The elementary kernel bound for an $m$-by-$n$ matrix is $(n-m)_+$.
On a maximal class with an anti-invariant functional and even $a,b$,
use powers of $Y+1/2$. Even powers pair only with odd powers, and
vice versa. Its rank is at most
$\min(\lceil n/2\rceil,\lfloor m/2\rfloor)
 +\min(\lfloor n/2\rfloor,\lceil m/2\rceil)$.
Use the corresponding bounds for $m+1$ to obtain $z$.
For $q\le10$ the checker can instead eliminate the exact rational
matrix of the defining functional. The affine system for a monic
row is solved over $\Q$. If its next coefficient is prescribed by the
system, the code substitutes the solution in all equations. If the
coefficient can vary, it verifies both a monic solution with next
coefficient zero and a homogeneous direction with next coefficient one.
For this fixed finite set of rational solutions, reduction excludes
only finitely many physical primes. For larger $q$ a nonzero modular
minor certifies solvability over $\Q$ whenever it is used. The
dimension and parity bounds do not infer rational nullity from a
singular matrix modulo an auxiliary prime.

\subsection{Two-sided low radicals and exact second-layer ranks}
\label{sec:bilateral-rank}
The compatible second-layer bound itself can be sharpened. For a
nonmaximal class let $r_{ij}$ denote an upper bound on the rank of
\[
 (\Psi(Z_a Z_b Y^{u+v}))_{
  0\le u<(D_1-a+i)_+,\ 0\le v<(D_2-b+j)_+},\qquad i,j\in\{0,1\}.
\]
Empty matrices have rank zero. On the two ordinary maximal restriction
kernels the second form is a sum of these local forms, as proved in
Section~\ref{lyr:second-layer}. Removing the top row or column
features costs at most $K_1$ or $K_2$ dimensions. Its rank is therefore
at most the minimum, over the four choices, of
\[
 (1-i)K_1+(1-j)K_2+\sum_{c\in T}r_{ij}(c).
\]
This bound can be used before adding any of the new row spaces.
It is compatible with them because they extend the same ordinary
kernels. Their pairing with restricted ordinary kernels cannot have
larger rank.

For either local functional on $h$ physical nodes its defining formula
is a sum of evaluations and a differentiated polynomial quotient.
After factoring off $Z_a Z_b$, each evaluation has rank at most one
and the nodes with $H(x)=0$ or $Z_a(x)Z_b(x)=0$ give zero. The
quotient part vanishes through total degree $h+2$, since the third
derivative of a quadratic vanishes. Thus in an $n$-by-$m$ moment
matrix the quotient entry vanishes for $u+v<h+3-a-b$. If $t=h+3-a-b$,
the remaining entries have rank at most
$\min(n,m,(n+m-1-t)_+)$ by counting the rows or columns which can
be nonzero. Add the two bounds, cap at $\min(n,m)$, and also apply
the parity bound when available. On the first ten quotient intervals
the replay additionally computes rational ranks by elimination.

Finally take the low class-supported radicals on both sides. Write
$l_{ic},z_{ic},L_i,D_i'$ for the dimensions and sums of
Proposition~\ref{prop:local-radical}, where $i$ designates the side
and $D_i'=\sum_{c\in M}(l_{ic}-z_{ic})$. The two off-diagonal
pairings with the ordinary kernels cost at most
$\min(D_1',K_2)$ and $\min(D_2',K_1)$. The corner between the two
added spaces has rank at most $\sum_{c\in M}\min(l_{1c},l_{2c})$.
Indeed, lifts supported in different maximal classes have zero
order-$p$ pairing. Where both leading reductions vanish, both
factors are divisible by $p$. Where just one is nonzero, the other
has its actual roots on division by $p$ and degree at most its
endpoint degree, since the top polynomial of a low row is identically
zero and so is its derivative. The surviving local factor belongs
to the kernel of the full leading cross Gram matrix on the other
side. Thus only diagonal class blocks can survive in this corner,
and each has the asserted rank bound. The contribution on the
nonmaximal classes is zero because the low rows there have zero
leading features. Counting independent rows and columns and using
the rank subadditivity of the three added blocks proves the gain
\[
 L_1+L_2-\min(D_1',K_2)-\min(D_2',K_1)
 -\sum_{c\in M}\min(l_{1c},l_{2c}).
\]
Take the maximum with zero and the two one-sided gains. All estimates
are lower bounds for the same coefficientwise valuation.

\paragraph{Joining the remaining top coefficients.}
On a nonmaximal root type $(a,b)$, a monic row completion annihilates
all columns of degree below $D_2$ divisible by $Z_b$. A monic column
completion has the transposed property. Require the row polynomial
$Q_1$ describing the ordinary top to vanish at every nonmaximal node
without the first witness, and $Q_2$ to vanish at those without the
second witness. If these sets have sizes $A$ and $C$, the rank
charged for these restrictions is at most
$\min(A,K_1)+\min(C,K_2)$, and the two remaining top dimensions are
$(K_1-A)_+$ and $(K_2-C)_+$. Low features can cancel every leading
maximal restriction, and on a nonmaximal type they can complete the
top by the appropriate monic polynomial independently. After these
basis choices the mixed top--low pairings vanish. Charge the low--low
pairing by the sum of its local ranks.

Where both top completions exist, their pairing has a single value
$\theta_{a,b}$. It is unchanged if either completion is replaced by
another, since the difference is low. For two feasible types $(a,b)$
and $(a',b')$ with $a\ge a'$ and $b\le b'$, pair the monic row of
the first with the monic column of the second. The column is
admissible for the first type and the row for the second type, so
this value equals both $\theta_{a,b}$ and $\theta_{a',b'}$.
Transitivity allows intermediate types whose root counts do not
occur in this particular chamber, provided their two systems have
solutions. In the anti-invariant case with even root counts, a
monic solution can be symmetrized with its reflection, keeping its
leading coefficient and annihilation conditions. If $D_1+D_2$ is
even the paired scalar is therefore zero.

The universal nonmaximal weight in the second-layer pairing cancels
one factor of $V_M$. Apart from a common unit the residual sum is
\[
 \sum_{c\in T}\theta_c c^{5+\epsilon+\alpha_1+\alpha_2}
           V_M(c^2)Q_1(c^2)Q_2(c^2).
\]
Where a top completion was unavailable the product already vanishes.
Subtract any reference scalar. The remaining exceptional nodes each
have rank at most one. The constant part extends over $M$ because
$V_M$ vanishes there. When the displayed exponent of $c$ is even,
the half-field moment sum gives normalized rank at most
$\bigl(w+K_1+K_2-n+O(1)\bigr)_+$ in the unscaled dimensions,
by the same coefficient argument as in the one-sided completion.
The case of zero reference has no constant part. Capping each bound
by the two remaining top dimensions completes the estimate.

\paragraph{A common second pairing for the ordinary and added top rows.}
Another bound counts both types of top rows in one subspace. Orient
the two windows as $r,c$. On $M$ retain monic row solutions annihilating
the full local columns, and on $T$ retain monic row solutions
annihilating the low local columns. Require the correlated top $P$
to vanish on the $f$ nodes admitting no such selected solution. The
top space then has dimension $h=(H_r-f)_+$, including those tops
which already vanish on $M$. It extends the low ordinary kernel by
these $h$ directions and the $L$ independent maximal low radicals
from Proposition~\ref{prop:local-radical}. Its dimension relative
to the original row kernel changes by $h+L-K_r$. On the column side,
start with the ordinary maximal kernel and require $Q_c=0$ at the
$v$ nonmaximal nodes without a monic column completion annihilating
the low rows. This leaves $k=(K_c-v)_+$ column tops. Low features
realize the two sets of monic prescriptions; both kernels are true
first kernels.

The ordinary low--low second pairing costs at most the sum $B_{00}$
of the local ranks on $T$. The low ordinary rows pair trivially with
the completed column tops by construction. The added rows pair
trivially with the low columns, since at $M$ the added rows also
annihilate the divided first column jet when its next coefficient
is zero. Thus the rest of the second form factors through the $k$
coefficients of the restricted $Q_c$. For the global row tops the
universal weights reduce it to the same half-field sum of
$c^{5+\epsilon+\alpha_r+\alpha_c}P(c^2)Q_c(c^2)$ with local
coefficients $\theta_c$. At nonmaximal nodes a missing completion
already forces $P Q_c=0$. On a maximal node a low radical can also
contribute, through just the one evaluation $Q_c(c^2)$.

Choose an arbitrary reference value, prescribing it on all types with
a free scalar. Let $E$ be the union of the active low-radical nodes
and the nodes whose contributing fixed scalar differs from the
reference, counted once. The additional rank $R$ is bounded by
$\min(k,h+L,E+J)$, where $J$ is the full half-field rank bound
for this reference and is zero for a zero reference. Thus the
denominator profile uses the ordinary restriction cost with
second rank $B_{00}+R$ and with the dimension changes
$h+L-K_r$ and $k-K_c$ credited on the two sides. This constructs
one specific pair of subspaces, so the same rank saving is counted
only once.

There is also a bridge between maximal and nonmaximal coefficients.
Let $G_M$ be a monic maximal row solution for roots $(a,b)$ and
$H_T$ a monic nonmaximal column solution for roots $(a',b')$,
in a nonmaximal type with a row solution. If $a\ge a'$ and
$Z_b\mid(Y-\ell)Z_{b'}$, then
\[
 \theta_M=\Phi\bigl(G_M(Y-\ell)H_T\bigr)
          =\Psi(G_M H_T)=\theta_T .
\]
Both outside equalities use monicity and annihilation of the lower
degrees. Within $M$ the earlier nested-space argument uses column
degree $D_c+1$ as the dual degree. All scalar identifications in
this test are rational equalities obtained this way, from a symmetry,
or by exact substitution in the small rational matrices.

\paragraph{Complementary maximal supports for two top polynomials.}
One can retain top directions on both sides. Partition the maximal types
between rows and columns. On a type assigned to side $i$, that side
may use its monic full annihilator and its class-supported low radicals.
On the opposite side impose a zero on its top. If the assigned side
lacks a monic witness, impose a zero on that top as well. On each
nonmaximal type require each retained monic top to annihilate the
opposite low feature space. Denote the forced zero sets by $F_i$,
put $A_i(X)=\prod_{c\in F_i}(X-c^2)$, and write
$P_i=A_iQ_i$ with $h_i=(H_i-|F_i|)_+$ free coefficients.
The original ordinary kernel changes in dimension by $h_i+L_i-K_i$,
where $L_i$ counts the assigned low radicals. The ordinary low--low
pairing again costs $B_{00}$.

The polynomial $W=A_1A_2/V_M$ exists since every maximal node is
forced to be a zero on at least one side. At a maximal row node the
leading column feature vanishes. Its next leading coefficient comes
from differentiating its top. With $X=c^2$,
\[
 P_1(X)P_2'(X)=V_M'(X) W(X)Q_1(X)Q_2(X)
 \qquad (P_2(X)=0,\ V_M(X)=0).
\]
At a column node the transposed identity holds. At an excluded
nonmaximal node the product $P_1P_2$ vanishes. The surviving top
pairing, apart from its common unit, is therefore
\[
 \sum_{c=1}^{n}\theta_c c^{5+\epsilon+\alpha_1+\alpha_2}
              W(c^2)Q_1(c^2)Q_2(c^2).
\]
The local low radicals on opposite sides have disjoint maximal
supports, and their corner is zero by the two-sided low-radical
argument. Each active local row at a maximal row node pairs only
through the one value $Q_2(c^2)$, and transposed local columns only
through $Q_1(c^2)$. At an exceptional top node these and the scalar
correction together still have rank at most one, since the opposite
evaluation is shared.

For a constant reference the rank is at most $\min(h_1,h_2)$, also
at most $(\deg W+h_1+h_2-n+O(1))_+$ in the even-power case and zero
for a zero reference. Add the size of the union of its exceptional
nodes and the active low-radical nodes. Cap the result by the actual
dimensions on both sides. If the numbers of active local row and
column nodes are $d_1,d_2$, one can also cap at
$h_1+\min(d_1,h_2)$ and $h_2+\min(d_2,h_1)$ by partitioning the
rows or columns. These give a second rank bound independent of the
ordinary low--low block. To compare a fixed maximal row scalar for
$(a,b)$ with a fixed maximal column scalar for $(a',b')$, take their
monic solutions $G,H$. If one monic linear factor $Y-s$ fills any
missing roots in both divisibilities
$Z_{a'}\mid(Y-s)Z_a$ and $Z_b\mid(Y-s)Z_{b'}$, evaluate
$\Phi((Y-s)GH)$ in the two orders. The resulting scalars coincide.

\paragraph{Interpolating the finite-field coefficient.}
In the last two constructions, charge once the $d$ active radical
nodes, whose local error and top error share the opposite evaluation.
On the remaining contributing nodes there are $m$ prescribed scalars;
the other contributing scalars can be chosen independently at their
nodes. Write $s=5+\epsilon+\alpha_1+\alpha_2$ and $b=s\bmod 2$.
For the $m$ prescribed nodes interpolate
$S(c^2)=c^b\theta_c$ by a polynomial of degree at most $m-1$
over $\mathbb F_p$. If $m=0$, use $S=0$. On free nodes choose
$\theta_c=c^{-b}S(c^2)$. Distinct half-field indices give distinct
nonzero squares. At nodes on which the product already vanishes
there is no restriction. The residual global term is now the
half-field sum with even weight $c^{s-b}S(c^2)$, including when
the original weight had an odd power. For the complementary
construction its rank is at most
\[
 d+\bigl(\deg W+h_1+h_2+m-n+O(1)\bigr)_+ .
\]
Indeed the coefficient functional on $Q_2$ vanishes on every low
degree whose product with $S W Q_1 X^{(s-b)/2}$ has degree strictly
between zero and $n$. The weight contributes a fixed positive power,
so no constant term intervenes. Truncating by the two original
dimensions remains permissible. In the one-sided construction the
same argument has $H_r+K_c$ in place of $\deg W+h_1+h_2$.
The interpolation polynomial may depend on $p$; only the rank bound
passes to the asymptotic prime integral.

\paragraph{Certifying exceptional local ranks.}
For a fixed type write $R_{ij}$ for the rank over $\mathbb Q$ of its
$(n+i)\times(m+j)$ low moment matrix, $i,j\in\{0,1\}$. The structural
feature count gives an upper bound $U_{ij}$; any nonzero minor over
a good auxiliary field gives a lower bound $L_{ij}$. If
$U_{10}=L_{00}$, the last row of the enlarged matrix is a combination
of the preceding rows over $\mathbb Q$, giving a monic annihilator.
If also $U_{11}=L_{00}$, its next moment is zero: eliminate the last
row and column using the low block, which already has the entire
rank. If $L_{01}>U_{00}$, a low radical changes the next moment by
a nonzero amount, allowing an arbitrary prescribed value. Finally
$U_{01}=L_{00}$ certifies that the next moment is independent of the
particular monic solution. These statements follow by nestedness of
the four matrices. Nonvanishing of a rational minor excludes only
finitely many physical primes for each fixed tuple. The alternative
exact rational calculation clears denominators and uses integer
row operations followed by substitution into the moment equations.

A short explicit recurrence can tighten $U_{ij}$. For a rectangular
$u\times v$ Hankel matrix with moments $s_0,\ldots,s_{u+v-2}$,
suppose $(c_0,\ldots,c_D)\ne0$ satisfies
$\sum_{k=0}^D c_k s_{j+k}=0$ for $0\le j\le u+v-2-D$.
Then the matrix has rank at most $D$. When $D<u$, the $u-D$ shifts
of $\sum_{k=0}^D c_kY^k$ in the row space of polynomials are independent
kernel vectors. When $D\ge u$, the inequality is immediate.
Leading or trailing zero coefficients in this recurrence are allowed.
The calculation proposes coefficients by a discrepancy recurrence
on integer moments and checks every convolution directly over
$\mathbb Z$. Its use as an upper bound requires no assumption about
minimality of the proposed recurrence.

The comparison of fixed scalars may also pass through auxiliary root
spaces. For a given $q,\epsilon$ the functionals $\Phi$ and $\Psi$
are defined on every polynomial, independently of which counts occur
in a physical class. Thus one can insert a space with a neighboring
number of consecutive roots, verify its monic row or column over
$\mathbb Q$, and use the nesting and linear-factor identities above
twice. Only actual square classes contribute dimensions or widths
to the rank count. The auxiliary spaces serve only to prove equalities.

\paragraph{Deep low kernels on the opposite side.}
At a maximal node $c$ assigned to the rows, let $z_c$ be the certified
dimension of the low column space annihilating all rooted row
polynomials through degree $D_r+1$. The correlated column top still
vanishes at $c$. Let $l_c$ be the number of retained local row
directions there, and $u_c\in\{0,1\}$ indicate whether a row
monic completion is retained. At least $(z_c-l_c-u_c)_+$ additional
column directions can be included at no second-rank cost.

Indeed, lift the $z_c$ column features with zero global top and no
leading feature at any other node. For a selected row whose feature
vanishes at $c$, its divided first jet there is rooted and of degree
at most $D_r+1$, so the deep column kills it. At the other maximal
nodes the column has zero leading feature and its divided first jet
has degree at most $D_c$, since its global top is zero. The selected
row features there annihilate all such columns. There is no
nonmaximal contribution. The second pairing with the new column
space therefore factors through the row feature at $c$, a space of
dimension at most $l_c+u_c$. Take the kernel of these functionals
among the $z_c$ columns. This proves the count, including possible
first-order arithmetic corrections when the row has a nonzero
leading feature there. The transposed construction applies at
column-assigned nodes. Added directions on opposite sides have
different leading supports, and the same divided-jet argument kills
their mutual pairing. Thus these gains can be summed over the
maximal classes.

\paragraph{Homogeneous interpolation with two polynomial multipliers.}
Write $d$ for the number of charged active radical nodes and $m$ for
the number of remaining prescribed evaluations in the polynomial
interpolation argument. In the complementary construction set
$E=\deg W+h_1+h_2+m-n$; in the one-sided construction set
$E=H_r+K_c+m-n$. The residual cap in either construction improves to
\[
 d+\max\{0,E/2,E-m\}+O(1).
\]
Here the dimensions and counts are unnormalized finite-prime counts.
For $m=0$ the zero reference already suffices. Otherwise fix an
integer $0\le r<m$. At each of the $m$ prescribed squares $x=c^2$
impose $N(x)=c^{s\bmod2}\theta_c D(x)$, with
$\deg N\le m-1-r$, $\deg D\le r$. These are $m$ homogeneous
linear equations in $m+1$ unknowns. A nonzero solution has $D\ne0$,
since a polynomial of degree less than $m$ cannot vanish at all
$m$ distinct squares unless it is zero.

Factor $D=D_1D_2$ over the physical prime field, and restrict
the two formerly free top polynomials to multiples of $D_1,D_2$.
Write $r'=\deg D$. Charge at most $r'$ dimensions for these restrictions.
At a free contributing node where $D$ vanishes, allocate the entire
power of its linear factor to the side admitting the free scalar.
All other factors may go to either side. The opposite multiplier
is then nonzero at this node. To see that a pole creates no exceptional
rank, recall the definition of a free scalar. There is a homogeneous
low vector with zero monic top, the same annihilations, and next
coefficient one. On the free side take a monic solution with next
coefficient zero, multiplied by $D_i(x)$, and add this low vector
times $c^{-(s\bmod2)}N(x)/D_{3-i}(x)$. Multiply the resulting local
section by the original forced factor and the new free top value.
This realizes the numerator pairing even at a zero of $D_i$.
Away from zeros one may simply prescribe the ratio.
A node with a prescribed scalar satisfies the homogeneous equation
also at a zero of $D$. Forced zeros from the earlier construction
are retained. At charged active nodes the contributions and any
error still share the opposite evaluation, and cost at most $d$.
The formula uses no division at a zero of the opposite multiplier.
In the one-sided construction all free contributing scalars can be
realized on the row side; on a nonmaximal type where the two monic
solutions coexist their scalar is fixed by evaluating their product.

The remaining half-field sum has the previous even weight and
numerator $N$, with the total free degrees reduced by $r'$. Its
terminal cap is $(E-r-r'+O(1))_+$. If a restricted top space is
zero there is no such sum and the charged-node bound alone suffices.
Adding back the dimensions lost gives at most
$d+\max(r,E-r)+O(1)$. Minimizing over $r$ gives the displayed
bound, with an absolute rounding cost absorbed in $O(1)$.
For the one-sided construction the factors recording the forced
row and column zeros are left in the numerator weight: their total
degrees plus the two unrestricted remaining dimensions are at most
$H_r+K_c$ whenever both spaces are nonzero. All local rows used in
lifting numerator values still annihilate the same lower jets, so the
ordinary low--low rank bound continues to apply.

The first jets of the remaining low directions admit a global cap.
Write $l'_c$ for the full low-kernel dimension on the opposite side
and retain
\[
 e_c=(l'_c-l_c-u_c)_+,\qquad e_c^0=(z_c-l_c-u_c)_+
\]
directions and deep directions, respectively. To construct them,
choose lifts of a basis of the row features at $c$, taking the local
row radicals with zero global top. Impose orthogonality to their
second pairings. There are at most $l_c+u_c$ conditions, both
on the full opposite low space and on its deep subspace.
Choose the surviving $e_c$-space to contain $e_c^0$ deep directions.

After subtracting those row lifts, the row has zero leading feature
at $c$. The only term which can survive pairing with an opposite
shallow direction is the degree-$(D_r+1)$ coefficient of the
divided row jet. For a zero feature this is the derivative of the
correlated top coefficient, up to the fixed local factor.
Subtracting a chosen monic lift changes it by a fixed multiple
of the original top value; the zero-top radical lifts do not change
it. Thus this entire block factors through the global row top of
dimension $h_r$. Its rank is at most
$\min(h_r,\sum_c(e_c-e_c^0))$, counting as usual with physical
multiplicity. The transposed block has the identical bound with
$h_c$. The corner between newly added rows and columns still
vanishes: their leading supports are disjoint and their global
tops are zero. Subtract these two rank costs from the additional
dimensions $\sum_c e_c$. One may discard the last nondeep assigned
local direction before doing the count. The calculation tests this
choice when a deep opposite direction replaces the lost dimension.

\paragraph{Exact bordered moment elimination.}
After clearing a common denominator in the moments $s_i$, border the
$m$ by $n$ low Hankel matrix by its next column and row. Pivot only
in the low block. The absence of a residual next-column entry on the
remaining low rows is exactly the solvability condition for a monic
annihilator. A nonzero remaining entry of the bottom row in a free
low column is exactly the existence of a homogeneous low vector
with nonzero next moment. Otherwise, in the consistent case, the
bottom right Schur complement is the unique next coefficient.
Thus one elimination gives both nullities and every fixed jet needed
for scalar comparisons. In fraction-free elimination each trailing
Schur entry is multiplied by the current signed pivot minor. This is
the bordered minor identity, or follows by one further ordinary
elimination step. Division of the final corner by the last pivot
and reversal of the initial moment scaling recover the rational jet.
Skipped pivot columns with a nonzero bottom entry register a free
column direction. No nonzero combination of subsequent pivot rows
can change this test, since they vanish in all skipped columns.
The portable checker uses integer arithmetic and checks every exact
division; the optional GMP evaluator performs the same operations.

\paragraph{A $29\times29$ local example.}
Take $q=12$, $\epsilon=1$, a maximal class with four roots on each
side, and degrees $D_r=32$, $D_c=31$. The low test for a row monic
has $n=32-4=28$ and $m=31-4+1=28$. Its leading determinant is
$704553\pmod{1000003}$, so its rational rank is $28$. In the
enlarged $29\times29$ matrix the centered even and odd dimensions
are $15$ and $14$. The anti-invariant functional vanishes on both
same-parity blocks. Its rank is at most $14+14=28$, hence is exactly
$28$. Eliminating the invertible low block leaves a zero scalar.
Thus the row monic exists and its next moment vanishes exactly.
This example can be replayed separately with
\texttt{python3 experiments/check\_worked\_rank.py}.

A rectangular example at the same $q$ explains the extra arithmetic
saving. Take $\epsilon=0$, a nonmaximal class, $(D_r,D_c)=(30,29)$,
and $(a,b)=(4,2)$. Here $n=26$ and $m=27$. There are more local
moment equations than free coefficients, so a test demanding full
rank on the equation side cannot certify a monic row. The four ranks,
in the order $R_{00},R_{10},R_{01},R_{11}$, are $26,26,26,27$.
For instance the leading $26\times26$ minor is
$447344\pmod{1000003}$. Since the added row leaves the rank
unchanged, a monic solution exists. Its next moment is fixed and
nonzero, as the final rank increase shows. This is the distinction
between full row independence and consistency of the specific monic
system used in the denominator argument.

\subsection{An explicit second-order bound for the small-prime tail}
\label{sec:second-order-tail}
For an integer cutoff $K$, the preceding uniform dimension estimate
$D_i+1-k_i\le 2\delta/\beta+2$ gives a sharper integral if the
prime weight is retained to first order. Put $1/\beta=q+t$,
$0\le t<1$, $e=\lfloor2t\rfloor$, and $y=t-e/2$.
The maximal classes have normalized width $\beta y$.
Dropping the cap $\delta$ on their feature dimension gives
\[
 E(\beta)\le(4+e)\delta+(2\delta+2\beta)y
          =(4+2t)\delta+2\beta y .
\]
\begin{lemma}
The contribution from $0<\beta\le1/K$ is at most
\[
 \frac{5\delta}{K}+\frac1{4K^2}-\frac{\delta}{6(K+1)^2}.
\]
\end{lemma}
\begin{proof}
The derivative of $w(t)=(q+t)^{-2}$ is at most $-2/(q+1)^3$.
Consequently $w(t)+2t/(q+1)^3$ is decreasing. Its covariance
with $t$ on $[0,1]$ is nonpositive, and
$\int_0^1(t-1/2)t\,dt=1/12$, so
\[
 \int_0^1(t-1/2)(q+t)^{-2}\,dt
 \le-1/[6(q+1)^3].
\]
On each half-period the increasing function $y$ has mean $1/4$.
Its covariance with $(q+t)^{-3}$ on each half is nonpositive.
Summing therefore bounds the $2\beta y$ contribution by
$\tfrac12\int_K^\infty x^{-3}\,dx=1/(4K^2)$.
Finally $\sum_{q=K}^\infty(q+1)^{-3}\ge1/[2(K+1)^2]$.
Together with $\int_K^\infty x^{-2}\,dx=1/K$, these give the result.
\end{proof}

\begin{lemma}\label{lem:cubic-prime-tail}
With the same dimension estimate, put $u=K+1/2$. The prime tail is at most
\[
 \frac{5\delta}{K}+\frac1{4K^2}
 -\frac{\delta}{6}\left(\frac1{u^2}+\frac1{u^3}\right)
 -\frac1{24(K+1)^3}.
\]
\end{lemma}
\begin{proof}
For a differentiable function $f$ on $[a,a+h]$, integration by parts gives
\[
 \int_0^h (v-h/2) f(a+v)\,dv
 =\frac12\int_0^h v(h-v) f'(a+v)\,dv .
\]
Take $h=1$, $a=q$ and $f(x)=x^{-2}$. The right side is the
negative of $\int_0^1 v(1-v)(q+v)^{-3}\,dv$. The probability density
$6v(1-v)$ has mean $1/2$. Convexity of $x^{-3}$ bounds the last
integral below by $(q+1/2)^{-3}/6$. The coefficient of this
covariance in the prime tail is $2\delta$. For the sum over $q$,
the trapezoid inequality on each unit interval gives
\[
 \sum_{q=K}^\infty(q+1/2)^{-3}
 \ge \int_K^\infty(x+1/2)^{-3}\,dx+\tfrac12u^{-3}
 =\tfrac12(u^{-2}+u^{-3}).
\]
For the remaining $2\beta y$ term take $h=1/2$, $f(x)=x^{-3}$,
and use the identity on each half of $[q,q+1]$. Since
$f'(x)\le -3/(q+1)^4$, each integral of $(v-1/4)f(a+v)$
is at most $-1/[32(q+1)^4]$. The two halves and the factor $2$
save $1/[8(q+1)^4]$ on that unit interval. Finally
$\sum_{q=K}^\infty(q+1)^{-4}\ge 1/[3(K+1)^3]$ by integration.
Combining both estimates with the mean contributions yields the bound.
\end{proof}

\begin{proposition}\label{prop:phase-prime-tail}
For $H\ge1$ write
\[
 m_H(c)=\min_{\theta\in\mathbb R}\frac1H
               \sum_{j=0}^{H-1}\{\theta+jc\},
 \qquad \{z\}=z-\lfloor z\rfloor .
\]
For either factorial window $s=s_i$ the bound in
Lemma~\ref{lem:cubic-prime-tail} can be decreased by
\[
 \frac{m_H(2s)+m_H(2(s+\delta))}{8(K+H)^2}.
\]
These savings also apply when a window parameter is rational with a
denominator smaller than $H$.
\end{proposition}
\begin{proof}
Put $x=1/\beta$. In a physical square class with normalized coordinate
$z\in(0,1/2)$ the number of window roots is
\[
 k=2\lfloor sx\rfloor+
 \mathbf 1_{z<\{sx\}}+\mathbf 1_{z>1-\{sx\}} .
\]
Thus $k\ge\lfloor 2sx\rfloor$ away from class boundaries, which have
zero area. The degree is $D=\lfloor 2(s+\delta)x\rfloor$. Hence
\[
 D+1-k\le 2\delta x+2-
       \big(1-\{2sx\}+\{2(s+\delta)x\}\big).
\]
The right side is positive. Multiplying it by the normalized maximal
width $y/x$ improves the pointwise tail envelope used above. After the
change of variable, the factor multiplying the parenthesis in its
integral is $y/x^3$.

Group the unit intervals $[q,q+1]$ into consecutive blocks of length
$H$ starting at $K$. For fixed $t\in(0,1)$ the value of $y(t)$ is
the same in each of the $H$ intervals of a block. In that block the
unweighted sum of the parentheses is at least
$H(m_H(2s)+m_H(2(s+\delta)))$. For the first term use
$1-\{v\}\ge\{-v\}$ and $m_H(-c)=m_H(c)$ by reversing the progression.
The weight $x^{-3}$ on block $n\ge0$ is at least
$(K+(n+1)H)^{-3}$. Now $\int_0^1 y(t)\,dt=1/4$ and
\[
 H\sum_{n\ge0}(K+(n+1)H)^{-3}
    \ge \int_{K+H}^{\infty}x^{-3}\,dx=\frac1{2(K+H)^2}.
\]
This proves the stated saving.

For a rational $c=a/b$, order the residues $ja\bmod b$ as
$r_0\le\cdots\le r_{H-1}$ and put $A=\sum_j r_j$. The mean increases
linearly between its jumps and takes its lower value at a jump. Its
minimum is therefore
\[
 m_H(a/b)=\frac1{Hb}\min_{0\le j<H}(A-Hr_j+bj).
\]
For equal residues the minimum uses their first index. This gives a
finite calculation even when $H$ crosses several periods.
\end{proof}

\begin{corollary}\label{cor:root-phase-prime-tail}
If $2\delta K\ge1$, the saving in Proposition~\ref{prop:phase-prime-tail}
may be replaced by
\[
 \frac{2m_H(s)+m_H(2(s+\delta))}{8(K+H)^2}.
\]
\end{corollary}
\begin{proof}
Keep the normalized class coordinate $z$ inside the integral. The two
root counts have the exact expression
$k=\lfloor sx+z\rfloor+\lfloor sx+1-z\rfloor$ almost everywhere.
Consequently the dimension is
\[
 D+1-k=2\delta x+\{sx+z\}+\{sx+1-z\}-\{2(s+\delta)x\}.
\]
The right side is nonnegative under the hypothesis, so taking its
positive part leaves it unchanged. Compared with $2\delta x+2$,
we can now save both $1-\{sx+z\}$ and $1-\{sx+1-z\}$ as well as
$\{2(s+\delta)x\}$. For fixed $t=\{x\}$ the normalized maximal
classes fill $0<z<t$ when $0<t<1/2$, and $1-t<z<1/2$ when
$1/2<t<1$. They therefore occupy the same set of width $y(t)$
in each of the $H$ unit intervals. For any fixed $t,z$ in that set
each of the two root terms has mean at least $m_H(s)$ by the
progression estimate. Integrating in $z$ and $t$ gives the result.
\end{proof}

\begin{corollary}\label{cor:finite-prime-bridge}
Let $S(K)$ denote the right side of Lemma~\ref{lem:cubic-prime-tail}.
For integers $1\le K<J$ the contribution from $1/J<\beta\le1/K$
is at most $S(K)-S(J)$.
\end{corollary}
\begin{proof}
Perform both covariance estimates of the lemma over $K\le q<J$.
The trapezoid inequality now gives
\[
 \sum_{q=K}^{J-1}(q+1/2)^{-3}\ge
 \int_K^J(x+1/2)^{-3}\,dx+
 \tfrac12\big((K+1/2)^{-3}-(J+1/2)^{-3}\big).
\]
For the half-period correction use
$\sum_{q=K}^{J-1}(q+1)^{-4}\ge\int_{K+1}^{J+1}x^{-4}\,dx$.
The mean contributions are $5\delta(K^{-1}-J^{-1})$ and
$(K^{-2}-J^{-2})/4$. Together these are precisely the stated
difference. The proof bounds the finite interval directly, which
is required when later intervals are evaluated by chamber sums.
\end{proof}

On this finite range the root-class correction is available as
well when $2\delta K\ge1$. For a chosen $H$, let $m$ denote the
larger of $2m_H(s)+m_H(2(s+\delta))$ at the two windows. With
$N=\lfloor(J-K)/H\rfloor$, subtract
\[
 \frac{mH}{4}\sum_{j=1}^{N}(K+jH)^{-3}
\]
from $S(K)-S(J)$. Apply Corollary~\ref{cor:root-phase-prime-tail}
block by block, using the least weight in each full block and
$\int_0^1 y(t)\,dt=1/4$. Discard the remaining nonnegative saving.
Taking the best $H$ separately on each gap is permissible because
the integration ranges are disjoint.

\section{An exact discrete Pearson identity}
\label{sec:pearson}

Let $\Bfun(x^j)=B_j$. Extend a functional $\K_X$ to rational functions with
polynomial part and only simple integer poles by
\[
 \K_X(f)=\frac1{24}\Bfun(f''')\quad(f\in\Q[x]),
\]
\[
 \K_X\!\left(\frac1{x-i}\right)=
 \begin{cases}-X+H_i^{(5)},&i\ge0,\\
              -X+H_{-i-1}^{(5)},&i\le-1.
 \end{cases}
\]
Unique partial-fraction decomposition makes this definition unambiguous. We extend it linearly when polynomial or residue coefficients lie in $\C$.
For a Laurent expansion at zero, $[x^4]f$ means the coefficient of $x^4$.

\begin{lemma}\label{lem:K}
For every function in this domain,
\[
 \K_X(f(x+1)-f(x))=[x^4]f(x).
\]
For every even such function,
\[
 \K_X(f)=-\tfrac12[x^4]f.
\]
Moreover, for a rational $h(t)$ with simple poles among $-i^2$,
\[
 \J_X(h(t))=\K_X(x^5h(-x^2)).
\]
\end{lemma}
\begin{proof}
For polynomials the shift formula is the Bernoulli identity
$\Bfun(g(x+1)-g(x))=g'(0)$ applied to $g=f'''$, then divided by $24$.
For $f=1/(x-i)$ with $i\ne0$, both sides equal $-i^{-5}$; at $i=0$ both
are zero. For the even identity, among even polynomial monomials only
$x^4$ contributes, with value $-1/2$. An even simple-pole part is a linear
combination of $1/(x-i)-1/(x+i)$; its functional value is $i^{-5}$ and
its $x^4$ coefficient is $-2i^{-5}$.

The last identity follows on polynomial monomials from the moment formula.
For a pole, use
\[
 \frac{x^5}{i^2-x^2}=-x^3-i^2x-\frac{i^4x}{x^2-i^2},\qquad
 \frac{x}{x^2-i^2}=\frac12\left(\frac1{x-i}+\frac1{x+i}\right).
\]
Its value is exactly $i^4X-i^4H_{i-1}^{(5)}-1/(2i)-1/4$.
Linearity completes all three statements.
\end{proof}

Assume $1\le m_1\le m_2\le N$, and put
\[
 W(x)=x^5\frac{C_{m_1}(-x^2)C_{m_2}(-x^2)}{D_N(-x^2)},\qquad
 \Lfun_X(f)=\K_X(Wf).
\]
The function $W$ is odd, has at most simple integer poles, and vanishes to
order nine at zero. Telescoping consecutive factors gives
\[
 \frac{W(x+1)}{W(x)}=
 \left(\frac{x+1}{x}\right)^8
 \frac{(x+m_1)(x+m_2)(x-N)}
      {(x-m_1+1)(x-m_2+1)(x+N+1)}.
\]
Define the degree-eleven polynomials
\[
 A(x)=(x-1)^8(x-m_1)(x-m_2)(x+N),\qquad
 B(x)=(x+1)^8(x+m_1)(x+m_2)(x-N).
\]
Thus $B(x)=-A(-x)$ and $W(x+1)A(x+1)=W(x)B(x)$.
Applying Lemma~\ref{lem:K} to $WAf$, whose $x^4$ coefficient vanishes, proves
\begin{equation}\label{eq:Pearson}
 \Lfun_X\bigl(B(x)f(x+1)-A(x)f(x)\bigr)=0.
\end{equation}
Also $\Lfun_X$ vanishes on odd polynomials: then $Wf$ is even and has zero
$x^4$ coefficient. Finally
\[
 \Lfun_X(p(-x^2))=
 \J_X\!\left(\frac{C_{m_1}C_{m_2}}{D_N}p(t)\right).
\]
The fixed degree in \eqref{eq:Pearson}, independent of $N$, is the structural
input that the earlier total-positivity argument did not supply.

\section{Every determinant root has multiplicity at most two}
\label{sec:nullity}

\begin{lemma}[Finite Hankel kernel]\label{lem:kernel}
Let $L$ be any linear functional on polynomials over a characteristic-zero
field. Suppose $H=[L(t^{i+j})]_{0\le i,j<r}$ has nullity $m>0$. If $p$ is its
least-degree monic kernel polynomial, of degree $j$, then the leading
$j\times j$ Hankel matrix is nonsingular,
\[
 \ker H=p\,\Q[t]_{\le m-1},\qquad j+m\le r,
\]
and
\[
 L(pt^\ell)=0\qquad(0\le\ell\le r+m-2).
\]
The ground field can replace $\Q$ in the displayed kernel statement.
\end{lemma}
\begin{proof}
If $q$ of degree below $j$ is orthogonal to all degrees below $j$, divide
$t^k=pa_k+b_k$ for $k<r$. Since $\deg(qa_k)\le k-1\le r-2$ and
$\deg b_k<j$, one gets $L(qt^k)=0$ for every $k<r$. Minimality forces $q=0$,
proving nonsingularity. Dividing any kernel polynomial as $pa+b$, and testing
against $t^k$ for $k<j$, similarly gives $b=0$.

Write $T_\ell=L(pt^\ell)$; initially $T_0=\cdots=T_{r-1}=0$.
Only indices through $2r-j-2$ can enter the kernel condition for $pa$.
If all these vanish, every $\deg a\le r-j-1$ is allowed, so $m=r-j$.
Otherwise let $\nu\ge r$ be the first nonzero index. All
$\deg a\le\nu-r$ work. If $\deg a=e\ge\nu-r+1$, choose $k=\nu-e$;
then $0\le k<r$ and $L(pat^k)$ is the leading coefficient of $a$ times
$T_\nu\ne0$. Therefore $m=\nu-r+1$, which proves the claims.
\end{proof}

Define the difference operator
\[
 \mathcal D f=B(x)(f(x+1)-f(x))
             -B(-x)(f(x-1)-f(x)).
\]
It is self-adjoint for the bilinear functional $\Lfun_X$:
\[
 \Lfun_X((\mathcal Df)g)=\Lfun_X(f\mathcal Dg).
\]
Indeed \eqref{eq:Pearson}, applied to $f(x)g(x-1)$, says that the adjoint of
$B(x)T_+$ is $-B(-x)T_-$. The remaining multiplication terms are self-adjoint.
The operator reverses parity and raises degree by at most nine. For the latter,
the possible degree $\deg f+10$ terms cancel because $B$ is monic of odd degree
 eleven.

\begin{lemma}[Divisibility from a polynomial adjoint]\label{lem:remainder}
Suppose $p$ is monic of degree $j$, the leading Hankel matrix $H_j$ is
nonsingular, and $L(pf)=0$ for every $\deg f\le M$. Let $S,S^*$ be polynomial
operators with
\[
 L((Sf)g)=L(fS^*g),
\]
and suppose each raises degree by at most $d$. If $j+d-1\le M$, then
$p\mid Sp$.
\end{lemma}
\begin{proof}
For $j=0$ the assertion is automatic. Otherwise divide $Sp=pV+R$, where
$\deg R<j$ and $\deg V\le d$. For $0\le k<j$,
\[
 L(Rt^k)=L(pS^*(t^k))-L(pVt^k)=0.
\]
Both multipliers have degree at most $j+d-1$. Nonsingularity of $H_j$ gives
$R=0$. This argument does not require $Sp$ to belong to the original finite
Hankel space.
\end{proof}

\begin{theorem}[Uniform nullity]\label{thm:nullity}
If $1\le m_1\le m_2\le N$ and $r>N-m_2+1$, then
\[
 \nullity H_N(X)\le2
\]
for every complex $X$. Every finite zero of $P_N$ therefore has multiplicity
at most two.
\end{theorem}
\begin{proof}
Suppose the nullity is $m\ge3$. Apply Lemma~\ref{lem:kernel} to the functional
$L_X(p)=\J_X(C_{m_1}C_{m_2}p/D_N)$, and put $P(x)=p(-x^2)$.
On the even polynomial space define
\[
 \mathcal Q f=\operatorname{even}[B(x)f(x+1)],\qquad
 \mathcal R f=x\mathcal Df.
\]
Their adjoints for $\Lfun_X$ are
\[
 \mathcal Q^*g=\operatorname{even}[A(x)g(x-1)],\qquad
 \mathcal R^*g=\mathcal D(xg).
\]
The first identity follows from Pearson applied to $f(x)g(x-1)$; the second
follows from self-adjointness of $\mathcal D$. All four operators preserve
evenness and raise $x$-degree by at most ten. For $\mathcal Q$ and
$\mathcal Q^*$, the possible leading degree $2n+11$ is odd and is removed.
For the other two, use the degree-nine gain of $\mathcal D$.

Transporting by $t=-x^2$ therefore gives polynomial operators and adjoints
of degree gain at most five. The kernel lemma supplies $j\le r-m$ and
orthogonality through $M=r+m-2$. Since
\[
 j+4\le r-m+4\le r+m-2\qquad(m\ge3),
\]
Lemma~\ref{lem:remainder} implies
\[
 P\mid\mathcal QP,\qquad P\mid\mathcal RP.
\]
The equality case $m=3$ is included.

At any nonzero root $\alpha$ of $P$, the resulting sum and difference equations
imply $B(\alpha)P(\alpha+1)=0$. A noninteger root could never reach an integer
root of $B$ or zero under repeated addition of one. It would produce infinitely
many distinct roots, a contradiction. Every root of $P$ is therefore an
integer, and every root of the monic $p(t)$ is a nonpositive square $-a^2$.
In particular $p(t)>0$ for $t>0$.

The polynomial $E(t)=\prod_{i=m_2}^N(t+i^2)$ has degree $N-m_2+1<r$, so the
kernel condition gives $L_X(pE)=0$. Exact cancellation instead gives
\[
 L_X(pE)=\J_X(p(t)tC_{m_1}(t))
        =\int_0^\infty p(t)tC_{m_1}(t)\dd\mu(t)>0.
\]
This is independent of $X$ and is a contradiction, even for complex $X$.

Finally, congruence by $G_N^{-1/2}$ converts the real pencil into
$I+(X-\xi)T_N$, with $T_N$ real symmetric. Its finite determinant-root
multiplicity equals the matrix nullity. Constant basis changes and primitive
normalization do not change those multiplicities.
\end{proof}

For comparison, requiring both operator images to lie in the original
finite Hankel kernel gives the weaker bound five. The remainder argument
above removes that unnecessary degree restriction. An alternative direct
argument, using the even part of $B(x)P(x+1)$ and $x$ times its odd part,
gives the bound six. All three bounds suffice for the transfer, but the
strongest proved structural statement is the uniform bound two.

\subsection{A second range for the multiplicity bound}
\label{sec:low-degree-nullity}

\begin{proposition}
For $1\le m_1\le m_2\le N$ and $1\le r\le N-m_1-m_2-1$,
the nullity of the same Hankel pencil is at most two at every complex
$X$. Its finite determinant zeros have multiplicity at most two.
\end{proposition}
\begin{proof}
If the nullity were at least three, the kernel and polynomial-adjoint
arguments in Theorem~\ref{thm:nullity} would again produce a monic
kernel polynomial $p$ of degree $j\le r-3$ all of whose roots are
nonpositive squares. Those arguments precede, and do not use, the
choice of the pole-clearing test polynomial $E$.
The $X$ coefficient of $L_X(p)$ is the sum of the finite residues of
\[
 t^2p(t) C_{m_1}(t) C_{m_2}(t)/D_N(t).
\]
Our hypothesis implies $2+j+m_1+m_2-N\le-2$. The residue at infinity
is zero, so the finite sum vanishes. Therefore
$L_X(p)=L_\xi(p)>0$, since $p$ and the factorial weight are positive on
the open positive half-line. The kernel identity $L_X(p)=0$ gives a
contradiction. The final multiplicity assertion follows by symmetric
congruence as in Theorem~\ref{thm:nullity}.
\end{proof}

This alternative test enlarges the range in which one can use the
bounded-multiplicity transfer. The parameters of Theorem~\ref{thm:main}
use the original pole-clearing test.

\section{A response criterion for the irrationality exponent}
\label{sec:transfer}

\begin{theorem}[Bounded-multiplicity transfer]\label{thm:transfer}
Assume \eqref{eq:degree}--\eqref{eq:gap}, and a uniform $O(1)$ bound for rational
root multiplicities. Suppose, for some $a>0$ and every fixed $C>0$,
\[
 \limsup_{N\to\infty}\frac1{N^2}
 \log\sup_{|h|\le Ce^{-aN}}|\det(I+hT_N)|\le R(a)<\varepsilon.
\]
Then
\[
 \mu(\xi)\le\frac{\lambda a}{\varepsilon-R(a)}.
\]
\end{theorem}
\begin{proof}
Suppose infinitely many rational approximations $x=p/q$ satisfy
$|x-\xi|\le q^{-M}$. Choose integral stages $N\sim c\log q$ with
\[
 \frac{\lambda}{\varepsilon-R(a)}<c<\frac M a.
\]
The fixed stage modulus changes $N$ by $O(1)$ only. Let $m$ be the multiplicity
of $x$ as a root of $P_N$, with $m=0$ if its value is nonzero. The first nonzero
normalized derivative has denominator dividing $q^{\deg P_N-m}$, hence
\[
 \left|\frac{P_N^{(m)}(x)}{m!}\right|\ge q^{-(\deg P_N-m)}.
\]
A Cauchy circle of radius $e^{-aN}$ about $x$ lies within $2e^{-aN}$ of $\xi$.
Using \eqref{eq:det-ratio} gives the upper logarithmic bound
\[
 -\varepsilon N^2+R(a)N^2+amN+o(N^2).
\]
Since $m=O(1)$, comparison contradicts the choice of $c$. Such a $c$ exists
for every $M>\lambda a/(\varepsilon-R(a))$.
\end{proof}

\begin{corollary}[Integrated envelope]\label{cor:integrated}
Suppose a positive contour matrix controlling the determinant has ordered
eigenvalue exponents bounded by a continuous decreasing function $g(t)$ on
$0\le t\le T$, where $0<T<\delta$ and $g(T)>0$. If
\[
 I(T):=\int_0^T g(t)\dd t<\varepsilon,
\]
then $\mu(\xi)<\lambda/T$.
\end{corollary}
\begin{proof}
The response at $a=g(T)$ is at most $I(T)-Ta$. A finite partition of
$[0,T]$ bounds the eigenvalues in fixed index blocks; min-max caps every
remaining eigenvalue by the last block. Refining that finite partition gives
the integral without a uniform asymptotic theorem at moving indices.
Theorem~\ref{thm:transfer} gives
\[
 \mu(\xi)\le\frac{\lambda a}{\varepsilon-I(T)+Ta}<\frac\lambda T.
\]
\end{proof}

\section{A positive contour matrix and polynomial compressions}
\label{sec:contour}

Write $z=t/N^2$. The inherited residue-contour identity, polarized to two
polynomials, represents $M_N$ as an integral of $p(z)q(z)$ on any fixed circle
$|z|=u^2$ with $u>1$. Include the full positive scaling $S_N$ in the contour
measure, so that this is a representation of $M_N$ itself; the same scaling
is included in the physical Gram form $G_N$ and cancels in their ratio.
Choose a real $G_N$-orthonormal polynomial basis. In this section $T_N$ denotes
the matrix of $M_N$ in that basis, which is orthogonally similar to the relative
matrix defined in Section~\ref{sec:setup}. Define $\mathcal A_N$ in the same
basis by taking the absolute contour measure and conjugating the second
polynomial. It is positive definite and real symmetric, the latter by
conjugation symmetry of the circle.

For every principal index set $I$, Andreief's identity and the triangle
inequality give
\[
 |\det(T_N)_I|\le\det(\mathcal A_N)_I.
\]
Consequently, for every complex $h$,
\begin{equation}\label{eq:PSD-domination}
 |\det(I+hT_N)|\le\det(I+|h|\mathcal A_N).
\end{equation}
This is a statement about determinant response; no ordering of the individual
absolute eigenvalues of $T_N$ by those of $\mathcal A_N$ is asserted.

The pointwise weighted-polynomial argument in the preceding report applies
also to this absolute contour integral. To state its rate, set
\[
 V(y)=2\pi|y|+I_1(y)-I_{s_1}(y)-I_{s_2}(y),\qquad
 I_a(y)=\int_0^a\log(y^2+x^2)\dd x.
\]
Let $\eta_b$ be the positive auxiliary measure described in
Section~\ref{sec:modified}; its mass is
\[
 v(b)=\frac{b+1-s_1-s_2-\sqrt{1+b^2}
                   +\sqrt{s_1^2+b^2}+\sqrt{s_2^2+b^2}}2.
\]
Write $b(v)$ for its inverse, $U_b$ for its logarithmic potential, and
$\ell(v)$ for its contact constant $4U_b-V$. On the circle define
\[
 \Phi_v(z)=L(z)+4U_{b(v)}(\sqrt z)-\ell(v),
\]
where
\[
 L(z)=\int_0^{s_1}\log|z+x^2|\dd x+
      \int_0^{s_2}\log|z+x^2|\dd x-
      \int_0^1\log|z+x^2|\dd x.
\]
Its maximum occurs at $z=-u^2$. Put $\Gamma_u(v)=\Phi_v(-u^2)$. The source
mass derivative gives
\begin{equation}\label{eq:Gamma-derivative}
 \frac{\partial\Gamma_u(v)}{\partial v}
       =4\arsinh\frac{u}{b(v)}>0.
\end{equation}
A polynomial of degree below $vN$ has absolute contour norm bounded by
$\exp(N\Gamma_u(v)+o(N))$ times its physical norm. The stronger pointwise
version retains $\Phi_v(z)$ before maximizing around the circle.

Put $y=1+\cos\theta\in[0,2]$ on $z=u^2e^{i\theta}$. The rational part of
$\Phi_v$ has an angular deficit at least
\begin{equation}\label{eq:angular-integral}
 \frac12\int_{s_1+s_2}^1
   \log\left(1+\frac{2u^2x^2}{(u^2-x^2)^2}y\right)\dd x.
\end{equation}
Indeed the integrand is increasing in $x$, so the positive small intervals in
$L$ can be cancelled against a portion of the negative interval. The potential
part is also maximal at the negative point and supplies a further nonnegative
deficit, which may be discarded.

For a partition $s_1+s_2\le\alpha_0<\cdots<\alpha_J=1$, choose either left
endpoints or, using the convexity lemma below, midpoints $c_j$ and set
\[
 K_j=\frac{\alpha_{j+1}-\alpha_j}{2},\qquad
 e_j=\frac{(u^2-c_j^2)^2}{2u^2c_j^2},\qquad
 D(y)=\sum_jK_j\log(1+y/e_j).
\]
Then $D$ is a lower bound for \eqref{eq:angular-integral}.

\begin{lemma}[Midpoint angular bound]\label{lem:midpoint}
For $0\le y\le2$, the function
$f_y(x)=\log(1+2yx^2/(1-x^2)^2)$ is convex on $0\le x<1$.
Thus midpoint quadrature gives a lower bound in \eqref{eq:angular-integral}.
\end{lemma}
\begin{proof}
Apart from a nonnegative factor and a positive denominator, the second
derivative is, with $r=x^2$,
\[
 (1-y/2)(1-r)^2(3r^2+8r+1)
 +(y/2)(3r^4+6r^3+4r^2+2r+1),
\]
which is nonnegative. Jensen's inequality on each interval proves the claim.
\end{proof}

As a first compression choose $p(z)=(z+u^2)^kq(z)$, of codimension $k$.
On the positive axis the squared factor is at least $u^{4k}$, while on the
circle it is $u^{4k}(2y)^k$. The powers of $u$ cancel. With $k/N\to t$,
positive-matrix min-max therefore bounds the $(k+1)$st eigenvalue by the rate
\[
 \Gamma_u(\delta-t)+\max_{0\le y\le2}\{t\log(2y)-D(y)\}.
\]
Using \eqref{eq:Gamma-derivative} gives the explicit envelope
$\gamma-\kappa t+F(t)$, where
$\gamma\ge\Gamma_u(\delta)$ and $\kappa\le4\arsinh(u/b(\delta))$.
The next two sections strengthen the same construction.

\section{Distributed conjugate zeros and an explicit equilibrium}
\label{sec:distributed}

A quadratic factor gives more freedom than repeated zeros at one point. Put
$a=u^2$ and, for $0\le h\le1$, define
\[
 Q_h(z)=z^2+2a(1-h)z+a^2.
\]
On $z\ge0$ it is positive and at least $a^2$. On $|z|=a$, with
$y=1+\cos\theta$, it satisfies
\[
 |Q_h(z)|^2=a^4[2|y-h|]^2.
\]
For $k=2\ell$, take $p(z)=\prod_{j=1}^{\ell}Q_{h_j}(z)\,q(z)$, with
$\deg q<r-k$. This real polynomial subspace has codimension $k$. When
$k/N\to t$, its angular zeros have limiting mass $t/2$.

We allow a nonnegative linear term in the angular deficit:
\begin{equation}\label{eq:Dfinal}
 D(y)=c_\theta y+\sum_{j=1}^J K_j\log(1+y/e_j),
 \qquad c_\theta\ge0,\quad K_j,e_j>0.
\end{equation}
Section~\ref{sec:linear-angular} supplies a certified $c_\theta$. Setting it to
zero is always valid.

\begin{proposition}[Angular equilibrium]\label{prop:angular-eq}
For $0<B_\theta\le1$, set
\[
 S_j=\sqrt{1+B_\theta/e_j},\qquad
 T=\frac{c_\theta B_\theta}{2}+\sum_jK_j(1-S_j^{-1}).
\]
The measure
\[
 \dd\rho(y)=\frac1{2\pi}\sqrt{\frac{B_\theta-y}{y}}
 \left(c_\theta+\sum_j\frac{K_j}{S_j(y+e_j)}\right)\dd y,
 \qquad0<y<B_\theta,
\]
has mass $T/2$. Its potential satisfies
\[
 2U_\rho(y)-D(y)\le\ell_\theta\quad(y\ge0),
\]
with equality on $[0,B_\theta]$. If $B=B(t)$ denotes the endpoint determined
by the mass equation, the resulting optimal angular compression rate is
\begin{equation}\label{eq:Fdist}
 F(t)=t\log\frac{B(t)}2-\frac{c_\theta B(t)}2
             -2\sum_jK_j\log\frac{1+\sqrt{1+B(t)/e_j}}2,
\end{equation}
and $F'(t)=\log(B(t)/2)$.
\end{proposition}
\begin{proof}
The Cauchy transform is
\[
 \frac12\left[c_\theta+\sum_j\frac{K_j}{z+e_j}
 -\sqrt{\frac{z-B_\theta}{z}}
       \left(c_\theta+\sum_j\frac{K_j}{S_j(z+e_j)}\right)\right].
\]
Cancellation at $z=-e_j$, the jump on the interval, and the coefficient at
infinity give respectively analyticity, the stated density, and its mass.
On the support its real part is $D'(y)/2$. For $y>B_\theta$ the derivative of
$2U_\rho-D$ equals
\[
 -\sqrt{\frac{y-B_\theta}{y}}
       \left(c_\theta+\sum_j\frac{K_j}{S_j(y+e_j)}\right)<0.
\]
This proves the global potential inequality. Evaluating its contact constant
at zero gives \eqref{eq:Fdist} after adding the factor $t\log2$ from the
quadratics. The same formula can be checked by differentiating the mass
identity: $F'(t)=\log(B(t)/2)$ and $F(0)=0$.
\end{proof}

\paragraph{From the measure to actual polynomial subspaces.}
For fixed $t>0$, take $\ell=\lceil tN/2\rceil$ and choose one zero in each
of the $\ell$ equal-mass quantile cells of $\rho$. Every cell has mass at
most $1/N$. Its density is
bounded by a constant times $y^{-1/2}$, so an interval of length $h$ has mass
$O(\sqrt h)$. The empirical logarithmic potential satisfies the uniform upper
estimate
\[
 \frac1N\sum_{j=1}^{\ell}\log|y-h_j|
       \le U_\rho(y)+O(\log N/N)\qquad(0\le y\le2).
\]
Here is the relevant error argument. Exclude the cell containing $y$ and its immediate neighbors, at most three
cells in total. If $y$ lies outside the support, use the nearest endpoint
cell and its neighbor instead. Outside the excluded cells, monotonicity of
the logarithm makes the errors telescope. The excluded cells have widths bounded below by a constant times
$N^{-2}$, so the telescoping variation is $O(\log N)$. Inside them, the
layer-cake estimate using
$\rho(I)\le\min(3/N,C\sqrt{|I|})$ bounds the omitted logarithmic integral by
$O(\log N/N)$. The density bound also gives polynomial lower bounds on cell
widths. Integer rounding of $\ell$ changes the mass by $O(1/N)$; scaling the
same potential inequality absorbs that change, since the density bound makes
$U_\rho$ bounded on $[0,2]$. At $t=0$ use the constant polynomial. Thus the total logarithmic
error is $o(N)$, uniformly in angle. This proves an upper polynomial bound,
which is the direction required for min-max; no lower asymptotic for these
angular polynomials is needed.

Consequently the distributed factors replace the one-point angular term by
$F(t)$ in the positive-matrix eigenvalue envelope. Each fixed $t$ is treated
first. A finite partition in $t$ and subsequent refinement give the integrated
response, avoiding a moving-parameter uniformity assumption.

\subsection{An exact primitive for the response}
\label{sec:primitive}

For a selected endpoint $B=B_\theta$, abbreviate
\[
 L_j=\log\frac{1+S_j}{2},\qquad A_\theta=\sum_jK_jL_j,
\]
\[
 H_{ij}=2L_i+2L_j-2\log\frac{S_i+S_j}{2}.
\]
Then the exact integral $J_F(T)=\int_0^T F(t)\dd t$ is
\begin{align}
 J_F(T)={}&\frac{T^2}{2}\log\frac B2
 -\frac{Tc_\theta B}{2}-2TA_\theta
 +\frac{c_\theta^2B^2}{16}\notag\\
 &+\frac{c_\theta}{2}\sum_jK_j\{B-2e_j(S_j-1)\}
 +\frac12\sum_{i,j}K_iK_jH_{ij}.\label{eq:JF}
\end{align}
There is no numerical quadrature in \eqref{eq:JF}. To derive it, integrate
$F'(t)=\log(B(t)/2)$ by parts, then expand $t(B)^2/B$. The pair integral is
\[
 \int_0^B\frac{(1-\sqrt{e_i/(e_i+x)})(1-\sqrt{e_j/(e_j+x)})}{x}\dd x
 =H_{ij}.
\]
Indeed the cross term reduces to
$2\log((\sqrt{1+B/e_i}+\sqrt{1+B/e_j})/2)$. This formula includes $i=j$.

\subsection{A stronger finite angular partition}
\label{sec:gauss}

The midpoint lemma can be strengthened without changing the finite-log form
of $D$. If $\cos\varphi=1-y$, then for $0\le x<1$,
\[
 \log\left(1+\frac{2yx^2}{(1-x^2)^2}\right)
 =2\sum_{n\ge1}\frac{1-\cos(n\varphi)}n x^{2n}.
\]
Every coefficient is nonnegative. Termwise differentiation is valid on
compact subintervals of $[0,1)$, so every derivative is nonnegative there.
In particular two-point Gauss quadrature is a lower bound: its error on a
panel $[\alpha,\beta]$ is
$(\beta-\alpha)^5 f^{(4)}(\zeta)/4320\ge0$.
For the half-integral in \eqref{eq:angular-integral}, each of the two nodes has
weight $(\beta-\alpha)/4$. The exact nodes are
\[
 \frac{\alpha+\beta}{2}\ \pm\
 \frac{\beta-\alpha}{2\sqrt3}.
\]
The checker encloses $1/\sqrt3$ rationally and rounds each node downward;
monotonicity preserves the lower bound. All $K_j$ and $e_j$ used in the
response formula are consequently rational numbers.

\section{The actual physical norm of the zero factors}
\label{sec:modified}

The preceding estimate used only $Q_h(z)\ge a^2$ on the physical axis. Retaining
its dependence on $z$ gives another improvement. We first use repeated linear
factors, then compare the distributed quadratics with them.

For reference, the inherited auxiliary measures on $[-b,b]$ have densities
\[
 \frac{\dd\lambda_b}{\dd x}=\frac1\pi\operatorname{arcosh}\frac b{|x|},
 \qquad
 \frac{\dd\nu_{b,s}}{\dd x}
 =\frac{s\sqrt{b^2-x^2}}{2\pi\sqrt{s^2+b^2}(x^2+s^2)}.
\]
The positive measure is
\[
 \eta_b=\int_0^\infty c(s)\nu_{b,s}\dd s,\qquad
 c(s)=\mathbf1_{s\ge s_1}+\mathbf1_{s\ge s_2}+\mathbf1_{s\le1}\ge1.
\]
In particular $\eta_b\ge\lambda_b/2$.

Fix $u>\max(1,b_0)$, where $b_0=b(\delta)$, and a cap $T_{\max}>0$ satisfying
\begin{equation}\label{eq:physical-guards}
 T_{\max}<\min\left\{\delta,\frac u2,\frac{u^2}{2b_0},
               \frac{u^2}{8\sqrt{4u^2+b_0^2}}\right\}.
\end{equation}
For $0\le t\le T_{\max}$, the exact physical norm of $(z+u^2)^{tN}q(z)$ has
external field $V_t(y)=V(y)-2t\log(y^2+u^2)$. Define $b=b_t$ by
\begin{equation}\label{eq:modified-mass}
 v(b)+\frac{tu}{\sqrt{u^2+b^2}}=\delta,
 \qquad\eta_{b,t}=\eta_b-2t\nu_{b,u}.
\end{equation}
Its mass is $\delta-t$. The endpoint is unique because $v'(b)\ge1/2$ and the
absolute derivative of the added term is at most $t/u<1/2$. It initially lies between
$b_{\min}^{(0)}=b(\delta-T_{\max})$ and $b_0$.

\paragraph{Sharpening the lower support by a finite iteration.}
The mass identity gives a stronger uniform lower bound at negligible
certificate cost. Suppose $\underline b_j\le b_t$ has already been proved
for every $t\le T_{\max}$. Since $q(b)=u/\sqrt{u^2+b^2}$ decreases, choose
an upper rational enclosure
\[
 q_j^+\ge\frac{u}{\sqrt{u^2+\underline b_j^2}}.
\]
Then
\[
 v(b_t)=\delta-tq(b_t)\ge\delta-T_{\max}q_j^+.
\]
A rational lower enclosure $v_{j+1}$ of the right side therefore yields
another valid lower support $\underline b_{j+1}=b(v_{j+1})$, enclosed
in the same directed fashion. Start from $v_0=\delta-T_{\max}$ and perform
four such steps. This is a finite induction; no convergence assertion is
needed. In the following, $v_{\min}$ and $b_{\min}=b(v_{\min})$ denote the
final certified mass and support floors. The original guards involving
$b_0$ are unchanged.

The density of $\eta_b$ is at least $\sqrt{b^2-x^2}/(2\pi b)$, while the
subtracted density is at most
$t\sqrt{b^2-x^2}/(\pi u\sqrt{u^2+b^2})$. Thus
\eqref{eq:physical-guards} ensures positivity. The source Cauchy transforms
give the contact identity $4U_{b,t}-V_t=\ell_{b,t}$ on the support. For $y>b$,
its derivative is
\[
 -2\sqrt{y^2-b^2}\left[
 \int_0^\infty\frac{c(s)s\dd s}{\sqrt{s^2+b^2}(y^2+s^2)}
 -\frac{2tu}{\sqrt{u^2+b^2}(y^2+u^2)}\right].
\]
Keeping only $u\le s\le2u$ in the integral gives the lower bound
$u^2/(\sqrt{4u^2+b^2}(y^2+4u^2))$. The last guard makes the bracket
nonnegative. The weighted-polynomial argument therefore applies to this
modified positive measure, with the same subexponential norm losses.

\subsection{Green functions and a uniform inner-arc bound}

In the squared variable, put
\[
 \phi_b(z)=\frac{2z-b^2+2\sqrt{z(z-b^2)}}{b^2},\qquad
 g_b(z)=\log|\phi_b(z)|,
\]
with the branch satisfying $|\phi_b|>1$ outside $[0,b^2]$. Let
\[
 \tau_b=\left(\frac{u+\sqrt{u^2+b^2}}b\right)^2,
 \qquad
 G_b(z,-u^2)=\log\left|
       \frac{\tau_b\phi_b(z)+1}{\phi_b(z)+\tau_b}\right|.
\]
These are the Green functions with poles at infinity and $-u^2$ respectively.
They are nonnegative away from the interval.

Let $\omega_b$ be the unit arcsine measure on $[-b,b]$, and let $\omega_{b,u}$ be the balayage of the equally weighted points $iu,-iu$ to that interval. Direct subtraction of their explicit densities gives
\[
 2\nu_{b,u}=\omega_{b,u}-\frac{u}{\sqrt{u^2+b^2}}\omega_b.
\]
Consequently, if $c_m=\delta-v(b)=tu/\sqrt{u^2+b^2}$, the modified contour
exponent is exactly
\begin{equation}\label{eq:green-exponent}
 \Phi_{b,t}(z)=\Phi_{v(b)}(z)-2tG_b(z,-u^2)+2c_mg_b(z).
\end{equation}
This identity can also be verified from the displayed conformal functions,
including their constants on the interval.

Choose an angular cutoff $0<Y_c<2$. A basic arcsine estimate gives
$|\phi_b(z)|\ge\tau_b\sqrt{1-Y_c/2}$ on $0\le y\le Y_c$.
The following bound is an intermediate estimate; Section~\ref{sec:support-cancellation} sharpens its angular dependence. For its pushforward unit arcsine measure $\sigma_b$ on $[0,b^2]$, set $\chi(s)=2u^2s/(u^2+s)^2\le1/2$. The chord inequality
\[
 \log(1-\chi y)\ge2\chi\log(1-y/2)
\]
and its exact arcsine average
\[
 \int\chi\dd\sigma_b=\frac{ub^2}{(u^2+b^2)^{3/2}}
\]
give
\[
 |\phi_b(z)|\ge\tau_b S,\quad
 S=(1-Y_c/2)^{C_\sigma},\quad
 C_\sigma=\frac{ub_0^2}{(u^2+b_0^2)^{3/2}},
\]
provided $b_0^2<2u^2$. The average increases with $b$ on this range.
The checker uses an upper enclosure for $C_\sigma$ and a lower enclosure for
$S$. A cubic lower Taylor bound for $e^x$ at negative $x$ suffices; no
floating-point exponential is required.

Let $\tau_0=\tau_{b_0}$ and assume $\tau_0\ge3$, $\tau_0^2S>1$. The inverse
conformal identity
\[
 z+u^2=\frac{b^2}{4}
       \frac{(\phi_b(z)+\tau_b)(\tau_b\phi_b(z)+1)}{\tau_b\phi_b(z)}
\]
implies, uniformly on the inner arc,
\begin{equation}\label{eq:green-bound}
 e^{-G_b(z,-u^2)}\le C_0\sqrt{2y},\qquad
 C_0=\frac{(\tau_0-1)^2\tau_0S}{(\tau_0^2S-1)^2}.
\end{equation}
The expression first decreases in $|\phi_b|$; after substituting
$|\phi_b|\ge\tau_bS$, its logarithmic derivative in $\tau_b$ has numerator
$-S\tau_b^3+3S\tau_b^2-3\tau_b+1<0$ for $\tau_b\ge3$. Since $b\le b_0$,
we have $\tau_b\ge\tau_0$, proving uniformity. In this intermediate estimate, $C_0$ is replaced by a rational upper bound, and that same rational constant is used in both the inner and outer estimates. The cubic bound for $S$ and rational rounding make every subsequent quantity use only square roots, logarithms, and rational operations.

\subsection{A quadratic bound for the mass adjustment}

Write $\kappa(v)=4\arsinh(u/b(v))$. On $[v_{\min},\delta]$, let
\[
 d_{\min}=\frac12\left(1-\frac{b_0}{\sqrt{1+b_0^2}}
 +\frac{b_{\min}}{\sqrt{s_1^2+b_{\min}^2}}
 +\frac{b_{\min}}{\sqrt{s_2^2+b_{\min}^2}}\right)>0,
\]
\[
 L_\kappa=\frac{4u}{b_{\min}\sqrt{u^2+b_{\min}^2}\,d_{\min}},\qquad
 q_{\max}=\frac{u}{\sqrt{u^2+b_{\min}^2}},\qquad
 A_m=\frac{L_\kappa q_{\max}^2}{2}.
\]
Monotonicity of the individual ratios bounds $v'(b)$ below by $d_{\min}$,
so $|\kappa'(v)|\le L_\kappa$. Hence
\begin{align*}
 \Gamma_u(v)+\kappa(v)(\delta-v)-\Gamma_u(\delta)
 &=\int_v^\delta(\kappa(v)-\kappa(w))\dd w\\
 &\le\tfrac12 L_\kappa(\delta-v)^2\le A_m t^2.
\end{align*}
Since $2c_mg_b(z)\le\kappa(v(b))c_m$, equations
\eqref{eq:green-exponent}--\eqref{eq:green-bound} now give the inner-arc rate
\[
 \gamma+ A_mt^2+2t\log C_0+t\log(2y)-D(y),
 \qquad\gamma\ge\Gamma_u(\delta).
\]

\subsection{A certified linear angular deficit}\label{sec:linear-angular}

Let $\nu_v$ be the pushforward of $\eta_{b(v)}$ to $[0,b(v)^2]$. Its angular
potential deficit is
\[
 -\int\log\left(1-\frac{2u^2s}{(u^2+s)^2}y\right)\dd\nu_v(s)
 \ge C_v y.
\]
The arcsine mass derivative gives
\[
 C_v=u\int_0^v\frac{b(w)^2}{(u^2+b(w)^2)^{3/2}}\dd w.
\]
Thus $C_v\ge C_{v_{\min}}$ for all masses used in
\eqref{eq:modified-mass}. A lower enclosure of this latter number is a valid
$c_\theta$ in \eqref{eq:Dfinal}.

There is an elementary closed expression. Put $b=b_{\min}$,
$U=\sqrt{u^2+b^2}$, $A_s=\sqrt{s^2+b^2}$, and
\[
 J_s=\log\frac{A_s+U}{s+u}
       -\frac{u^2}{u^2-s^2}\left(\frac{A_s}{U}-\frac su\right).
\]
Then
\begin{equation}\label{eq:linear-closed}
 C_{v_{\min}}=
 \frac u2\left[\arsinh\frac bu-\frac bU-J_1+J_{s_1}+J_{s_2}\right].
\end{equation}
Indeed $J_s(0)=0$ and
$J_s'(b)=b^3/(\sqrt{s^2+b^2}(u^2+b^2)^{3/2})$.
Inserting the explicit derivative of $v(b)$ proves the identity. All
denominators are nonzero because $u>1>s_i$.

\section{First-moment control and an intermediate response}
\label{sec:combined}

Choose an angular endpoint cap $B_{\max}\le\min(1,Y_c)$. The distributed
quadratic factors satisfy the pointwise physical inequality
\[
 Q_h(z)=(z+a)^2-2ahz\ge(1-h/2)(z+a)^2\qquad(z\ge0).
\]
We now retain the first moment of the angular root distribution instead of
replacing every $h$ by the largest possible endpoint.

\begin{lemma}[First moment and a quadratic physical cost]\label{lem:physical-moment}
Let $T=T(B_{\max})$ for the angular mass function in
Proposition~\ref{prop:angular-eq}, and define
\[
 M_\theta(t)=2\int h\dd\rho_t(h).
\]
For each endpoint $B=B(t)$, this moment is explicitly
\begin{equation}\label{eq:angular-moment}
 M_\theta(t)=\frac{c_\theta B^2}{8}
    +\sum_j\frac{K_jB(S_j-1)}{2S_j(S_j+1)}.
\end{equation}
For every $0\le t\le T$, the squared physical product of distributed
factors, relative to the repeated linear factor, has logarithmic rate at
least $-A_pt^2$, where
\begin{equation}\label{eq:physical-coefficient}
 A_p=\frac{-\log(1-B_{\max}/2)}{B_{\max}}
                 \frac{M_\theta(T)}{T^2}
 \le\frac{-\log(1-B_{\max}/2)}{4T}.
\end{equation}
Larger directed upper bounds for this coefficient remain valid.
\end{lemma}
\begin{proof}
Integrating the density in Proposition~\ref{prop:angular-eq} gives
\eqref{eq:angular-moment}; the substitution $h=B\sin^2\theta$ suffices.
One can also check each term directly. If
$T_j=K_j(1-S_j^{-1})$ and
$M_j=K_je_j(S_j-1)^2/(2S_j)$, then
\[
 \frac{\dd M_j}{\dd B}=\frac B2\frac{\dd T_j}{\dd B}.
\]
The linear terms satisfy the same identity. Therefore
\[
 M_\theta'(t)=\frac{B(t)}2,\qquad
 M_\theta(t)=\frac12\int_0^t B(s)\dd s.
\]
The mass function $T(B)$ is increasing and concave, with $T(0)=0$.
Its inverse $B(t)$ is convex and vanishes at zero, so
$B(s)\le sB(t)/t$ for $0\le s\le t$. It follows that
\[
 M_\theta(t)\le\frac{tB(t)}4,\qquad
 \frac{\dd}{\dd t}\frac{M_\theta(t)}{t^2}
 =\frac{tB(t)/2-2M_\theta(t)}{t^3}\ge0.
\]
In particular $M_\theta(t)\le M_\theta(T)t^2/T^2$.

Concavity of the logarithm gives its chord bound
\[
 \log(1-h/2)\ge\frac h{B_{\max}}\log(1-B_{\max}/2).
\]
The coefficient is negative, so the last moment inequality yields
\[
 2\int\log(1-h/2)\dd\rho_t(h)
 \ge\frac{\log(1-B_{\max}/2)}{B_{\max}}M_\theta(t)
 \ge-A_pt^2.
\]
The bound $M_\theta(T)\le TB_{\max}/4$ gives the final comparison in
\eqref{eq:physical-coefficient}.

For the actual quantile polynomials, the normalized physical product has
logarithm
\[
 \frac2N\sum_j\log(1-h_j/2)
 =2\int\log(1-h/2)\dd\rho_t(h)+O(N^{-1}).
\]
The error estimate follows from bounded variation of this smooth logarithm
on $[0,B_{\max}]$, by summing its oscillations over equal-mass cells.
It is uniform in the physical variable, because the pointwise comparison
already removed that variable. Integer rounding has the same error order.
\end{proof}

Put
\[
 \kappa_*=-2\log C_0,\qquad A_*=A_m+A_p.
\]
The inner-arc compression is bounded by
\begin{equation}\label{eq:final-envelope}
 g(t)=\gamma-\kappa_*t+A_*t^2+F(t).
\end{equation}
For $y\ge Y_c\ge B_{\max}$, every $|y-h_j|\le y$, so the distributed
factor divided by the repeated-zero factor has absolute value at most one.
The outer arc is therefore also bounded by \eqref{eq:final-envelope} if
\begin{equation}\label{eq:outer-check}
 D(Y_c)\ge-2T\log C_0-F(T).
\end{equation}
Both quadratic costs appear in the inner and outer estimates and cancel
from this condition. The right side is increasing in $t\le T$, since
$F'(t)=\log(B(t)/2)$ and $B(t)\le B_{\max}$; checking the endpoint suffices.

The finite guards
\[
 0<T<T_{\max},\qquad B(T)\le B_{\max}\le\min(1,Y_c),
 \qquad\kappa_*>2A_*T_{\max}+\log(B_{\max}/2)
\]
ensure applicability and monotonicity. The checker also verifies $g(T)>0$.
By Corollary~\ref{cor:integrated}, the single decisive integral comparison is
\begin{equation}\label{eq:final-I}
 \boxed{\gamma T-\frac{\kappa_*T^2}{2}
              +\frac{A_*T^3}{3}+J_F(T)<\varepsilon.}
\end{equation}
All quantities in \eqref{eq:outer-check}--\eqref{eq:final-I} are rational
numbers, square roots, and logarithms of explicitly specified positive
numbers. The angular mass $T$ itself is enclosed by directed intervals.
In particular the exact checker computes $A_p$ using upper bounds for
$-\log(1-B_{\max}/2)$ and $M_\theta(T)$, and a positive lower bound for $T$.
It provides a rational upper bound for
$\lambda g(T)/(\varepsilon-I(T)+Tg(T))$ and also the simpler bound
$\lambda/T_{\rm lower}$. It verifies strict inequalities with an explicit
rational choice of the transfer constant $c$.
\section{Coupling the support and Green corrections}\label{sec:support-cancellation}

This section strengthens the earlier combined response without changing the
matrix, its arithmetic normalization, or the multiplicity theorem. We retain
all positivity and support guards of Section~\ref{sec:modified}, including
$u>b_0$, $T_{\max}<u/2$, and the certified interval
$b_{\min}\le b_t\le b_0=b(\delta)$. The physical moment coefficient $A_p$
and positive mass coefficient $A_m$ are those proved above; the angular
measure used to compute $A_p$ will have the effective deficit defined below.

\subsection{A uniform angular Green estimate}

Put $Y=Y_c$, and, for a physical support endpoint $b$, write
\[
 \tau_b=\left(\frac{u+\sqrt{u^2+b^2}}b\right)^2,
 \qquad C_{00}(b)=\frac{\tau_b}{(\tau_b+1)^2}.
\]
Choose $C\ge ub_0^2/(u^2+b_0^2)^{3/2}$ and a rational
$0<S_{\min}\le(1-Y/2)^C$ with $\tau_0^2S_{\min}>1$,
where $\tau_0=\tau_{b_0}\ge3$. The preceding conformal estimate, applied
at the actual support, gives
\[
 e^{-G_b(z,-u^2)}\le\mathcal C(\tau_b,S(y))\sqrt{2y},
 \quad S(y)=(1-y/2)^C,
 \quad\mathcal C(\tau,S)=\frac{(\tau-1)^2\tau S}{(\tau^2S-1)^2}.
\]
Here $C$ is uniform because $b_0<u$ and
$ub^2/(u^2+b^2)^{3/2}$ increases on $0<b<\sqrt2u$.
For $x=-\log S$,
\[
 \frac{\dd}{\dd x}\log\mathcal C(\tau,e^{-x})
 =\frac{\tau^2e^{-x}+1}{\tau^2e^{-x}-1}.
\]
This is bounded above, uniformly for $b\le b_0$ and $y\le Y$, by
$(\tau_0^2S_{\min}+1)/(\tau_0^2S_{\min}-1)$.
Convexity of $-\log(1-y/2)$ therefore proves
\begin{equation}\label{eq:coupled-green-linear}
 e^{-G_b(z,-u^2)}\le C_{00}(b)e^{A_Gy}\sqrt{2y},\qquad
 A_G\ge\frac{\tau_0^2S_{\min}+1}{\tau_0^2S_{\min}-1}
          C\frac{-\log(1-Y/2)}Y.
\end{equation}
The zero-angle constant retains its dependence on $b$.
If the old angular deficit is bounded below by
$c_{\rm phys}y+\sum_jK_j\log(1+y/e_j)$, set
\begin{equation}\label{eq:coupled-effective-deficit}
 c_{\rm eff}=c_{\rm phys}-2A_GT_{\max}\ge0,\qquad
 D_{\rm eff}(y)=c_{\rm eff}y+\sum_jK_j\log(1+y/e_j).
\end{equation}
For every $t\le T_{\max}$, this absorbs the extra Green term $2tA_Gy$.
All charges remain positive, so the angular equilibrium proposition and
its exact primitive apply to $D_{\rm eff}$; denote its constant by
$F_{\rm eff}(t)$ and its endpoint by $B(t)$.

\subsection{Cancellation at the moving physical support}

For fixed $u$, define
\[
 q(v)=\frac{u}{\sqrt{u^2+b(v)^2}},\qquad
 k(v)=4\arsinh\frac{u}{b(v)},\qquad
 k_0(v)=4\log\frac{2\sqrt{u^2+b(v)^2}}{b(v)}.
\]
Then $k_0(v)=-2\log C_{00}(b(v))$. The inherited mass derivative and
direct differentiation give
\begin{equation}\label{eq:coupled-derivatives}
 \Gamma_u'(v)=k(v),\qquad
 k'(v)=-\frac{4u}{b(v)\sqrt{u^2+b(v)^2}\,v'(b(v))},\qquad
 k_0'(v)=q(v)k'(v).
\end{equation}
The modified mass satisfies $v_t+tq(v_t)=\delta$.
Since $v'(b)\ge1/2$ and $|\dd q/\dd b|\le1/u$, the retained guard
$T_{\max}<u/2$ gives $1+tq'(v_t)>0$; also $q'(v)<0$. Hence
\[
 v_t'=-\frac{q(v_t)}{1+tq'(v_t)}<0,\qquad
 0<1+tq'(v_t)\le1.
\]
By the exact modified-exponent identity and
\eqref{eq:coupled-green-linear}, the constant part of the inner upper
bound, before the effective angular deficit and physical moment cost, is
\[
 H(t)=\Gamma_u(v_t)+k(v_t)(\delta-v_t)-t k_0(v_t).
\]
The three identities in \eqref{eq:coupled-derivatives} imply the exact
cancellation
\begin{equation}\label{eq:coupled-H-derivative}
 H'(t)=k'(v_t)v_t'(\delta-v_t)-k_0(v_t)-t k_0'(v_t)v_t'
       =-k_0(v_t).
\end{equation}
Thus a positive quadratic mass cost is unnecessary in the inner estimate.
A negative quadratic correction is available as well. Define
\[
 V_+=\frac12\left(1-\frac{b_{\min}}{\sqrt{1+b_{\min}^2}}
       +\frac{b_0}{\sqrt{s_1^2+b_0^2}}
       +\frac{b_0}{\sqrt{s_2^2+b_0^2}}\right).
\]
Monotonicity of $b/\sqrt{s^2+b^2}$ gives $v'(b)\le V_+$ throughout the
support interval. Consequently
\[
 \frac{\dd}{\dd t}k_0(v_t)
 =\frac{4u^3}{b_t(u^2+b_t^2)^{3/2}v'(b_t)[1+tq'(v_t)]}
 \ge\frac{4u^3}{b_0(u^2+b_0^2)^{3/2}V_+}.
\]
Choose a rational lower bound
\begin{equation}\label{eq:coupled-negative-mass}
 0<A_-\le\frac{2u^3}{b_0(u^2+b_0^2)^{3/2}V_+}.
\end{equation}
Integrating twice, using \eqref{eq:coupled-H-derivative}, proves
$H(t)\le\Gamma_u(\delta)-k_0(\delta)t-A_-t^2$.

\subsection{The final envelope and its outer guard}

Let $\gamma\ge\Gamma_u(\delta)$ and $0<\kappa_0\le k_0(\delta)$ be the
directed bounds in the checker. Write $\beta=A_p-A_-$, which is negative
at the final parameters. The resulting inner envelope is
\begin{equation}\label{eq:coupled-final-envelope}
 \boxed{g_6(t)=\gamma-\kappa_0t+\beta t^2+F_{\rm eff}(t).}
\end{equation}
On the outer arc, the previously proved estimate still has the positive
mass cost: its rate is at most
$\gamma+(A_m+A_p)t^2-D_{\rm old}(y)$. Therefore the sufficient outer
endpoint condition is now
\begin{equation}\label{eq:coupled-outer-check}
 \boxed{D_{\rm eff}(Y)\ge
     \kappa_0T+(A_m+A_-)T^2-F_{\rm eff}(T).}
\end{equation}
Indeed $D_{\rm old}(y)\ge D_{\rm eff}(Y)$ for $y\ge Y$, and the right
side of \eqref{eq:coupled-outer-check}, with $T$ replaced by $t$, increases
because $F_{\rm eff}'(t)=\log(B(t)/2)$. The extra quadratic term in this
outer condition is essential; the earlier cancellation of the inner and
outer positive costs no longer applies.

The final finite guards are
\[
 0<T<T_{\max},\quad B(T)\le B_{\max}\le\min(1,Y),\quad g_6(T)>0,
\]
\begin{equation}\label{eq:coupled-monotonicity}
 -\kappa_0+2\max(\beta,0)T_{\max}+\log(B_{\max}/2)<0.
\end{equation}
The maximum with zero is needed when $\beta<0$. It proves monotonicity
of the full envelope. The unchanged closed angular primitive now gives
\begin{equation}\label{eq:coupled-integral}
 I_6(T)=\gamma T-\frac{\kappa_0T^2}{2}
              +\frac{\beta T^3}{3}
              +\int_0^T F_{\rm eff}(t)\dd t<\varepsilon.
\end{equation}
Putting $a_6=g_6(T)$ and $\mathcal R_6=I_6(T)-Ta_6$, the same bounded
multiplicity transfer yields the upper exponent
$\lambda a_6/(\varepsilon-\mathcal R_6)$. The checker also verifies an
explicit rational transfer constant with both strict margins positive.

Every new numerical direction is explicit: use lower $\tau_0$ and
$S_{\min}$ in the decreasing function $(q+1)/(q-1)$, upper $C$ and
$-\log(1-Y/2)$ to bound $A_G$ upward, and round $c_{\rm eff}$ downward.
Bound $V_+$ upward and $A_-$ downward. An upper rational bound for
$C_{00}(b_0)$ supplies the lower $\kappa_0$; the same constant is used in
the inner and outer calculations. Together with the inherited directed
square-root and logarithm bounds, these operations certify
\eqref{eq:coupled-outer-check}--\eqref{eq:coupled-integral} using rational
intervals alone. The original arithmetic certificate is an independent
input to this analytic calculation.
\section{Shifted contours and an exact window correction}
\label{sc:section}

The response estimate can be sharpened while retaining the same physical
equilibrium measure and the same determinant transfer.  The new contour
is a circle whose centre lies on the negative real axis.  Its leftmost
point is unchanged.  This preserves the zero-angle norm exponent, while
changing the angular logarithmic charges and the Green constant in a
controlled way.  A particular shift also turns the previously discarded
window contribution into an explicit positive linear term.

Put
\begin{equation}
 a=u^2,\qquad \alpha=s_1+s_2,\qquad
 \eta=-\alpha^2,\qquad R=a+\eta,
 \label{sc:geometry}
\end{equation}
and use the contour
\begin{equation}
 z=\eta+R e^{i\theta},\qquad y=1+\cos\theta\in[0,2].
 \label{sc:contour}
\end{equation}
Let $b_0=b(\delta)$.  In addition to the earlier support and response
guards, require
\begin{equation}
 a+2\eta>b_0^2,\qquad
 b_0^2<\frac{2a(a+2\eta)}{a-\eta}.
 \label{sc:guards}
\end{equation}
These are strict rational inequalities after enclosing $b_0$.
The left and right real intersections of the contour are $-a$ and
$a+2\eta$, respectively.  Thus the contour encloses all negative poles
and the physical interval $[0,b_0^2]$.

\subsection{Positive factors and the shifted angular field}
\label{sc:factors}

For $0\le h\le1$, define
\begin{equation}
 Q_h(z)=(z-\eta)^2+2R(1-h)(z-\eta)+R^2
       =(z+a)^2-2Rh(z-\eta).
 \label{sc:quadratic}
\end{equation}
On the physical half-line, put
\[
 \chi_\eta(x)=\frac{2R(x-\eta)}{(x+a)^2}.
\]
The identity
\begin{equation}
 \frac12-\chi_\eta(x)
 =\frac{(x-a-2\eta)^2}{2(x+a)^2}
 \label{sc:physical-factor}
\end{equation}
and $\eta\le0$ give $0\le\chi_\eta(x)\le1/2$ for $x\ge0$.
Consequently $Q_h(x)\ge(1-h/2)(x+a)^2$.  The earlier endpoint-moment
bound for the physical cost therefore applies without modification.
On the shifted circle one has the exact identity
\begin{equation}
 Q_h(\eta+R e^{i\theta})
 =2R^2e^{i\theta}(y-h).
 \label{sc:circle-factor}
\end{equation}
The distributed-root construction and its hard-edge angular equilibrium
retain the same coordinate $y$.

For the rational weight, set
\begin{equation}
 A_\eta(x)=\frac{2R(\eta+x^2)}{(a-x^2)^2},\qquad
 f_y(x)=\log\bigl(1+A_\eta(x)y\bigr).
 \label{sc:rational-field}
\end{equation}
The elementary modulus identity
$|z+x^2|^2/(a-x^2)^2=1+A_\eta(x)y$ shows that its angular deficit equals
\begin{equation}
 D_{\rm rat}(y)=\frac12\left(
 \int_0^1 f_y(x)\,dx-\int_0^{s_1}f_y(x)\,dx
                       -\int_0^{s_2}f_y(x)\,dx\right).
 \label{sc:rational-deficit}
\end{equation}
Since
\[
 \frac{dA_\eta}{d(x^2)}
 =\frac{2R(a+2\eta+x^2)}{(a-x^2)^3}>0,
\]
the function $f_y$ increases on $[0,1]$.  Hence
\begin{equation}
 D_{\rm rat}(y)
 =\frac12\int_\alpha^1 f_y(x)\,dx+W_\eta(y),\qquad
 W_\eta(y)=\frac12\int_0^{s_1}
       \bigl(f_y(s_2+x)-f_y(x)\bigr)\,dx\ge0.
 \label{sc:window-split}
\end{equation}

\begin{lemma}[Positive Gauss lower sums on the shifted contour]
\label{sc:gauss}
For every $0\le y\le2$, the two-node Gauss--Legendre rule on any panel
in $[\alpha,1]$ is a lower bound for the integral of $f_y$.
Rounding both nodes downward preserves that lower bound.
\end{lemma}
\begin{proof}
Write $t=x^2+\eta$.  On the integration interval, $0\le t<R$.
Choose $\phi$ with $\cos\phi=1-y$.  Factoring the numerator gives
\begin{equation}
 f_y(x)=2\sum_{n\ge1}\frac{1-\cos(n\phi)}{n}
             \left(\frac{x^2+\eta}{R}\right)^n.
 \label{sc:positive-series}
\end{equation}
All coefficients are nonnegative.  Each fourth derivative of
$(x^2+\eta)^n$ is nonnegative where $x^2+\eta\ge0$, by the chain rule
with $t'=2x$, $t''=2$, and higher derivatives zero.  Termwise
differentiation on compact subsets proves $f_y^{(4)}\ge0$.
The sign of the Gauss remainder gives the first assertion, and
monotonicity gives the second.
\end{proof}

Thus the positive logarithmic charges have the earlier weights but new
locations:
\begin{equation}
 K_j=\frac{\text{panel width}}4,\qquad
 e_j=\frac{(a-x_j^2)^2}{2R(\eta+x_j^2)}.
 \label{sc:charges}
\end{equation}
At $\eta=-\alpha^2$, the first Gauss node lies strictly above $\alpha$,
so every denominator in \eqref{sc:charges} is positive.

\subsection{An exact linear contribution from the windows}
\label{sc:window}

For $0\le x\le s_1$ one has
$A_\eta(x)\le A_\eta(s_2+x)\le0$.  The two factors in the denominator
below are positive and at most one.  Therefore
\begin{align}
 W_\eta'(y)
 &=\frac12\int_0^{s_1}
 \frac{A_\eta(s_2+x)-A_\eta(x)}
 {(1+A_\eta(s_2+x)y)(1+A_\eta(x)y)}\,dx \notag\\
 &\ge\frac12\int_0^{s_1}
        \bigl(A_\eta(s_2+x)-A_\eta(x)\bigr)\,dx.
 \label{sc:window-derivative}
\end{align}
Since $W_\eta(0)=0$, this proves the global bound
\begin{equation}
 W_\eta(y)\ge c_{\rm win}y\quad(0\le y\le2),\qquad
 c_{\rm win}=\frac{J(\alpha)-J(s_1)-J(s_2)}2>0,
 \label{sc:window-linear}
\end{equation}
where the exact primitive is
\begin{equation}
 J(x)=\frac{R^2x}{a(a-x^2)}
       +\frac{R(\eta-a)}{2u^3}
                   \log\frac{u+x}{u-x}.
 \label{sc:window-primitive}
\end{equation}
Direct differentiation gives $J'=A_\eta$ and $J(0)=0$.
This coefficient needs only rational arithmetic and logarithm enclosures;
no additional window quadrature is required.

\subsection{Physical contribution and its closed primitive}
\label{sc:physical}

Let $\nu_v$ denote the inherited physical equilibrium measure of mass
$v$ on $[0,b(v)^2]$.  Its angular deficit satisfies
\begin{equation}
 -\int\log(1-\chi_\eta(x)y)\,d\nu_v(x)
 \ge \mathcal C_\eta(v)y,
 \qquad \mathcal C_\eta(v)=\int\chi_\eta(x)\,d\nu_v(x).
 \label{sc:physical-deficit}
\end{equation}
In terms of the unshifted coefficient $\mathcal C_0(v)$,
\begin{equation}
 \mathcal C_\eta(v)=\frac Ra\mathcal C_0(v)-2R\eta E_2(v),
 \qquad E_2(v)=\int\frac{d\nu_v(x)}{(a+x)^2}.
 \label{sc:physical-linear}
\end{equation}
The mass derivative of $\nu_v$ is the arcsine probability measure.
Its average of $\chi_\eta$ is
\begin{equation}
 A_{\sigma,\eta}(b)
 =\frac{R\{ab^2-\eta(2a+b^2)\}}
        {u^3(a+b^2)^{3/2}}>0.
 \label{sc:arcsine-coefficient}
\end{equation}
Hence $\mathcal C_\eta$ increases with mass, and the previously certified
mass floor $v_{\min}$ gives the valid coefficient
$\mathcal C_\eta(v_{\min})$.

For exact evaluation put $b=b(v)$, $U=\sqrt{u^2+b^2}$, and
$S_s=\sqrt{s^2+b^2}$.  Then
\begin{align}
 F_0(b)&=\frac{\operatorname{arsinh}(b/u)+b/U}{2u^3},\notag\\
 F_s(b)&=\frac1{2u^3}\left[
       \log\frac{S_s+U}{s+u}
       +\frac{u^2}{u^2-s^2}\left(\frac{S_s}{U}-\frac su\right)
                         \right],\notag\\
 E_2(v)&=\frac{F_0(b)-F_1(b)+F_{s_1}(b)+F_{s_2}(b)}2.
 \label{sc:inverse-square-primitive}
\end{align}
These formulas follow by differentiation using the explicit $v'(b)$
and the arcsine identity
\[
 \int\frac{d\sigma_b(x)}{(a+x)^2}
 =\frac{2a+b^2}{2u^3(a+b^2)^{3/2}}.
\]
Every primitive in \eqref{sc:inverse-square-primitive} vanishes at $b=0$.

\subsection{Green constant and the resulting response}
\label{sc:green}

Let $\varphi_b$ be the exterior map of $[0,b^2]$ and set
\[
 \tau=\left(\frac{u+\sqrt{u^2+b^2}}b\right)^2,
 \qquad \varphi_b(-a)=-\tau.
\]
The inverse-map identity gives
\begin{equation}
 e^{-G_b(z,-a)}
 =\frac{4|z+a|\tau|\varphi_b(z)|}
        {b^2|\tau\varphi_b(z)+1|^2}.
 \label{sc:green-identity}
\end{equation}
Because $|z+a|=R\sqrt{2y}$, the zero-angle constant becomes
\begin{equation}
 C_{00,\eta}(v)=\frac Ra\frac{\tau(v)}{(\tau(v)+1)^2},\qquad
 \kappa_\eta(v)=\kappa_0(v)-2\log\frac Ra.
 \label{sc:green-constant}
\end{equation}
The arcsine logarithmic average gives
\begin{equation}
 \log|\varphi_b(z)|-\log\tau
 \ge A_{\sigma,\eta}(b)\log(1-y/2).
 \label{sc:green-angular}
\end{equation}
The second guard in \eqref{sc:guards} makes
$A_{\sigma,\eta}(b)$ increasing for $0\le b\le b_0$: its derivative
with respect to $b^2$ has the sign of
$a(a+2\eta)-(a-\eta)b^2/2$.  Therefore the new uniform $C_{\max}$ is
$A_{\sigma,\eta}(b_0)$.  The earlier cubic exponential lower bound and
Green logarithmic derivative estimate give the corresponding upper
slope $A_G$.

The factor $R/a$ is independent of the physical mass.  Consequently
$\kappa_\eta'(v)=\kappa_0'(v)$, so the exact support cancellation and
its negative quadratic gain $A_-$ remain valid.  Set
\begin{equation}
 c_{\rm eff}=\mathcal C_\eta(v_{\min})+c_{\rm win}
                     -2A_G T_{\max}>0,
 \qquad
 D_{\rm eff}(y)=c_{\rm eff}y+\sum_jK_j\log(1+y/e_j).
 \label{sc:effective-field}
\end{equation}
Let $F_{\rm eff}$ be the same explicit angular equilibrium value for
these charges and this linear term.  The shifted response envelope is
\begin{equation}
 g_8(t)=\gamma-\kappa_\eta t+(A_p-A_-)t^2+F_{\rm eff}(t).
 \label{sc:response}
\end{equation}
The sufficient outer-arc condition remains
\begin{equation}
 D_{\rm eff}(Y)\ge
 \kappa_\eta T+(A_m+A_-)T^2-F_{\rm eff}(T).
 \label{sc:outer-check}
\end{equation}
Indeed $Q_h(z)/(z+a)^2=(y-h)/y$ on the shifted circle, so the earlier
outer argument dropping the nonnegative Green function is unchanged.
The constants $A_m$, $A_p$, and $A_-$ have their previous meanings.

The positive contour Gram majorant, the closed primitive
$I(T)=\int_0^Tg_8(t)\,dt$, the identity $\mathcal R=I(T)-Tg_8(T)$,
and the bounded-multiplicity irrationality transfer now apply exactly
as before.  The version-eight checker encloses all new quantities with
rational intervals and checks \eqref{sc:guards},
\eqref{sc:effective-field}, and \eqref{sc:outer-check} explicitly.
\section{A larger shift with a retained negative interval}
\label{es:section}

The circle may be shifted slightly farther to the left while retaining
the same response framework.  Choose a rational $L\ge\alpha=s_1+s_2$,
put $\eta=-L^2$, and keep $a=u^2$, $R=a+\eta$ and the strict support and
monotonicity guards of~\eqref{sc:guards}.  All contour, physical-factor,
Green and inverse-square-moment identities remain unchanged.

Define
\[
 A_\eta(x)=\frac{2R(\eta+x^2)}{(a-x^2)^2},\qquad
 f_y(x)=\log(1+A_\eta(x)y).
\]
The exact rational angular deficit splits as
\begin{align}
 D_{\mathrm{rat}}(y)
 &=\frac12\int_L^1 f_y(x)\,dx+W(y)+N(y),
 \label{es:split}\\
 W(y)&=\frac12\int_0^{s_1}
       \bigl(f_y(s_2+x)-f_y(x)\bigr)\,dx,
       \nonumber\\
 N(y)&=\frac12\int_\alpha^L f_y(x)\,dx.
       \nonumber
\end{align}
The positive Gauss argument applies on $[L,1]$.  Both angular
coefficients in the difference defining $W$ are nonpositive, so the
earlier primitive gives the global bound
\begin{equation}
 W(y)\ge C_{\mathrm{win}}y,\qquad
 C_{\mathrm{win}}=
 \frac{J_\eta(\alpha)-J_\eta(s_1)-J_\eta(s_2)}2>0.
 \label{es:positive-window}
\end{equation}
The upper endpoint here is the actual combined window $\alpha$, not
the Gauss starting point $L$.

On $[\alpha,L]$, the coefficient $A_\eta$ is nonpositive and increasing.
The identity
\[
 1+2A_\eta(x)=
 \frac{(a+2\eta+x^2)^2}{(a-x^2)^2}>0
\]
ensures positive denominators throughout $0\le y\le2$.  For
$0\le y\le Y$,
\begin{equation}
 N'(y)\ge E(Y):=
 \frac12\int_\alpha^L
 \frac{A_\eta(x)}{1+A_\eta(x)Y}\,dx,
 \qquad N(y)\ge yE(Y).
 \label{es:inner-negative}
\end{equation}
Also $N'(y)\ge E(2)$ for all $0\le y\le2$.
For fixed $Y\in[0,2]$, the integrand in~\eqref{es:inner-negative}
increases with $x$, since its derivative is
$A_\eta'(x)/(1+A_\eta(x)Y)^2$.  A left rectangular sum therefore gives
a rational lower bound $E^-(Y)\le E(Y)$, despite its negative summands.
Each summand is rounded downward before summation.

Let $C_{\mathrm{phys}}$ be the certified physical linear coefficient
before Green absorption.  Impose the additional guard
\begin{equation}
 C_{\mathrm{phys}}+C_{\mathrm{win}}+E^-(2)>0.
 \label{es:outer-monotonicity}
\end{equation}
The complete physical and rational angular deficit is then increasing
on $[0,2]$: the positive Gauss interval has nonnegative derivative, the
positive window contributes at least $C_{\mathrm{win}}$, the physical
deficit contributes at least $C_{\mathrm{phys}}$, and the negative
interval contributes at least $E^-(2)$.  At the inner cutoff, the
effective linear coefficient is
\begin{equation}
 c_{\mathrm{eff}}=
 C_{\mathrm{phys}}+C_{\mathrm{win}}+E^-(Y)
 -2A_GT_{\max}>0.
 \label{es:effective-slope}
\end{equation}
Thus the existing inner-envelope and outer-endpoint checks still apply.
The varying-slope refinement of Section~\ref{ds:section} is unchanged,
because the retained window coefficients do not depend on the physical
mass or on codimension.

This optional mode is labelled version~8a in the exact checker.  It
records every negative-interval cell and checks
\eqref{es:outer-monotonicity} explicitly.  The original version~8 is
the special case $L=\alpha$, with no negative interval.
\section{Retaining the varying angular slope}
\label{ds:section}

The shifted-circle estimate admits a further improvement without changing
its endpoint threshold or its outer-arc test.  The physical angular
coefficient is larger at smaller codimension, and the Green absorption
cost is smaller there.  Retaining part of this variation reduces the
integrated determinant response.

\subsection{A codimension-dependent linear coefficient}
\label{ds:slope}

Let $c>0$ be the effective linear angular coefficient in the shifted-circle
certificate, and let $T_{\max}$ be its mass cap.  Write the final certified
bootstrap mass floor as
\begin{equation}
 v_{\min}=\delta-q_*T_{\max},\qquad 0<q_*<1.
 \label{ds:mass-floor}
\end{equation}
Here $q_*$ is the upper bound used in the final bootstrap step, so that
$v_t\ge\delta-q_*t$ for $0\le t\le T_{\max}$.  This choice is made from
the stored mass floor; it need not equal a subsequently recomputed
upper bound for $u/\sqrt{u^2+b_{\min}^2}$.

In the notation of Section~\ref{sc:section}, the mass derivative of the
shifted physical coefficient is
\begin{equation}
 A_\sigma(b)=
 \frac{R\{a b^2-\eta(2a+b^2)\}}
 {u^3(a+b^2)^{3/2}}.
 \label{ds:arcsine-derivative}
\end{equation}
The guard in~\eqref{sc:guards} makes this function increasing over the
physical support range.  Let $A_{\min}>0$ be a certified rational lower
bound for $A_\sigma(b_{\min})$, and let $A_G$ be the certified Green
logarithmic slope upper bound.  Since
\[
 v_t-v_{\min}\ge q_*(T_{\max}-t),
\]
the angular field has the stronger linear coefficient
\begin{equation}
 c+h(T_{\max}-t),\qquad
 0<h\le A_{\min}q_*+2A_G.
 \label{ds:h}
\end{equation}
Downward rational rounding of $h$ is permissible.

Choose a rational $T_0$ with $0<T_0\le T$, where $T$ is the mass of the
existing endpoint response.  We use the smaller strengthening
\begin{equation}
 c_t=
 \begin{cases}
 c+h(T_0-t),&0\le t\le T_0,\\
 c,&t\ge T_0.
 \end{cases}
 \label{ds:ct}
\end{equation}
In particular, the endpoint family at $t=T$ is unchanged.

\subsection{Response derivatives at fixed mass}
\label{ds:derivatives}

For a positive linear coefficient $z$, define $B(t,z)$ by inversion of
the increasing concave function
\begin{equation}
 t_z(B)=\frac{zB}{2}+
 \sum_j K_j\left(1-\frac{1}{\sqrt{1+B/e_j}}\right).
 \label{ds:mass-function}
\end{equation}
With $S_j=\sqrt{1+B/e_j}$, the hard-edge equilibrium constant and its
normalized first moment are
\begin{align}
 F(t,z)&=t\log(B/2)-\frac{zB}{2}
       -2\sum_j K_j\log\frac{1+S_j}{2},
       \label{ds:F}\\
 M(t,z)&=\frac{zB^2}{8}
       +\sum_j\frac{K_jB(S_j-1)}{2S_j(S_j+1)}.
       \label{ds:M}
\end{align}
Direct differentiation, while holding $t$ fixed, gives
\begin{equation}
 \frac{\partial F}{\partial z}=-\frac{B(t,z)}2,
 \qquad
 \frac{\partial M}{\partial z}=-\frac{B(t,z)^2}{8}.
 \label{ds:fixed-mass-identities}
\end{equation}
For verification, the derivatives at fixed $z$ are
$F_B=t_z'(B)\log(B/2)$ and
$M_B=(B/2)t_z'(B)$, while
$B_z=-B/(2t_z'(B))$ at fixed mass.
Thus increasing the angular slope decreases both the endpoint constant
and the physical first-moment penalty.  The previously certified
quadratic physical penalty remains valid.

Put $c_0=c+hT_0$.  Since $B(t,z)$ decreases with $z$, integrating the
first identity in~\eqref{ds:fixed-mass-identities} yields
\begin{equation}
 F(t,c_t)\le F(t,c)
       -\frac h2(T_0-t)B(t,c_0),\qquad 0\le t\le T_0.
 \label{ds:F-saving}
\end{equation}
The outer margin also improves: its derivative with respect to $z$ is
$Y-B(t,z)/2>0$, using the existing guard $B\le Y$.
Consequently no new outer-arc numerical test is needed.

\subsection{A valid reduction of the positive-part response}
\label{ds:positive-part}

Let
\begin{equation}
 g_0(t)=\gamma-\kappa t+\beta t^2+F(t,c),
 \qquad a_*=g_0(T).
 \label{ds:old-envelope}
\end{equation}
The determinant response integrates the positive part above $a_*$.
Thus one must check that the proposed subtraction does not lower the
comparison envelope below this threshold.  Define
\begin{equation}
 m=\kappa-\log(B/2)-2\max(\beta,0)T_{\max}>0.
 \label{ds:m}
\end{equation}
The earlier monotonicity argument gives $-g_0'(t)\ge m$ on $[0,T]$.
Require the additional rational guard
\begin{equation}
 \frac{hB}{2}\le m.
 \label{ds:positive-part-guard}
\end{equation}

Suppose $0\le L(t)\le B(t,c_0)$.  Then, for $0\le t\le T_0$,
\[
 \frac h2(T_0-t)L(t)
 \le m(T_0-t)
 \le g_0(t)-g_0(T),
\]
because $B(t,c_0)\le B$.  The comparison envelope obtained from
\eqref{ds:F-saving} therefore remains at least $a_*$ throughout the
range where a saving is claimed.  Both the old integrated exponent $I$
and the old response $\mathcal R=I-Ta_*$ can consequently be reduced by
the certified amount
\begin{equation}
 S=\frac h2\int_0^{T_0}(T_0-t)L(t)\,dt.
 \label{ds:integrated-saving}
\end{equation}
The endpoint threshold remains unchanged.  The resulting transfer bound
is
\begin{equation}
 \mu(\zeta(5))\le
 \frac{\lambda a_*}{\varepsilon-\mathcal R+S}.
 \label{ds:transfer}
\end{equation}

\subsection{Exact rational tangent certificates}
\label{ds:tangents}

The inverse of $t_{c_0}$ is convex.  At any rational $b_j\ge0$,
\begin{equation}
 B(t,c_0)\ge b_j+
 \frac{t-t_{c_0}(b_j)}{t_{c_0}'(b_j)}.
 \label{ds:tangent}
\end{equation}
Suppose rational intervals supply
$0\le f_j\le f_j^+$ for $f_j=t_{c_0}(b_j)$ and
$0<d_j^-\le d_j\le d_j^+$ for $d_j=t_{c_0}'(b_j)$.
For every $t\ge0$, the affine function
\begin{equation}
 \ell_j(t)=\frac{t}{d_j^+}+b_j-\frac{f_j^+}{d_j^-}
 \label{ds:rational-tangent}
\end{equation}
lies below the tangent in~\eqref{ds:tangent}.  Rounding its slope and
intercept downward preserves this assertion.  Hence
\begin{equation}
 L(t)=\max\{0,\ell_1(t),\ldots,\ell_n(t)\}
 \label{ds:L}
\end{equation}
is an admissible lower bound in~\eqref{ds:integrated-saving}.
All intersections of this affine envelope are rational.  Integrating
$(T_0-t)L(t)$ separately on its finitely many intervals uses only rational
arithmetic.  Rounding each nonnegative interval contribution downward
gives a certified lower bound for $S$.

The version-nine checker records the full byte hash of its shifted-circle
source certificate, in version~8 or its extended version~8a, all tangent
intervals, the active affine envelope,
the positive-part guard, and the resulting rational transfer inequality.
The source certificate is replayed separately before applying this
postprocessing step.
\section{A contour crossing the interval of the equilibrium measure}
\label{sec:crossing}

Write $a=u^2$, $\eta=-L^2$, $R=a+\eta$, and $d=a+2\eta$.
The estimates in Sections~\ref{sc:factors} and~\ref{es:section} use the
circle $z=\eta+Re^{i\theta}$ with angular coordinate $y=1+\cos\theta$.
The following proposition allows its right endpoint to lie in the physical
interval and retains the deficit on the full outer arc.

\begin{proposition}\label{prop:crossing}
Assume $u>1$, $\alpha=s_1+s_2<L<1$, $d>0$, and $0<B\le Y<1$.
Use the physical mass floor, the positivity checks and the inner-arc
checks of Section~\ref{es:section}. For $b_{\min}\le b\le b_0$ define
\[
 H(x)=\frac{R\{ax-\eta(2a+x)\}}{u^3(a+x)^{3/2}},
 \qquad x_* = \frac{2a(a+2\eta)}{a-\eta}.
\]
An upper bound on $H$ on $[b_{\min}^2,b_0^2]$ is obtained by evaluating it
at the projection of $x_*$ onto this interval. A lower bound is the
smaller of its endpoint values. The same rules apply to any certified
rational interval enclosing the squared support endpoints.

Let $C$ be the sum of the certified lower bounds for the physical and
positive-window linear coefficients before the Green loss. Write
$A(x)=2R(\eta+x^2)/(a-x^2)^2$, and retain the positive charges $(K_j,e_j)$
from the two-node Gauss rule on $[L,1]$. Set
\[
 N(y)=\frac12\int_\alpha^L\log(1+A(x)y)\,dx,\qquad
 F(y)=Cy+\sum_jK_j\log(1+y/e_j)+N(y).
\]
Let $D_{\rm eff}$ be the certified inner deficit from
Section~\ref{es:section}, including its Green loss.
If
\[
 F(2)\ge D_{\rm eff}(Y)
 \quad\hbox{and}\quad
 D_{\rm eff}(Y)\ge\kappa_\eta T+(A_m+A_-)T^2-F_{\rm eff}(T),
\]
the fixed-slope and varying-slope response bounds hold on this circle.
In the varying-slope bound one uses the lower bound for $H$ specified
above for the mass derivative. For that refinement require the additional
margin $F(2)-D_{\rm eff}(Y)\ge hT_0Y$, where $h,T_0$ have the meanings
in Section~\ref{ds:section}.
\end{proposition}

\begin{proof}
In the scaled residue identity the integrand for a pair of test polynomials
is a constant times
\[
 z^2 p(z)q(z)
 \frac{\prod_{j=0}^{m_1-1}(z+(j/N)^2)
       \prod_{j=0}^{m_2-1}(z+(j/N)^2)}
      {\prod_{j=1}^{N}(z+(j/N)^2)}\,dz.
\]
Its finite poles lie in $[-1,0]$. The left endpoint is $-a<-1$ and the
right endpoint is $d>0$. The residue theorem therefore gives the same
bilinear form on this circle. The rational factors in the angular potential
have a fixed positive distance from their pole and zero locations.

For the weighted polynomial estimate, fix the mass and take the weighted
supremum on $[0,b^2]$. Outside the interval the weighted maximum principle
applies as in Section~\ref{sec:contour}. On the interval the corresponding
inequality follows directly from that same supremum and the equilibrium
contact identity. The arcsine potential is continuous up to the interval,
and both Green functions extend there with value zero. Thus the inequality
holds on the entire circle, including $z=d$ when $d\le b^2$.
On $0\le y\le Y$ we have $\Re z=-a+Ry<0$, so the inner-arc exterior-map
calculation uses exterior points throughout. The physical quadratic-factor
inequality follows from~\eqref{sc:physical-factor} for every $x\ge0$.
The arcsine coefficient is $H(b^2)$ by~\eqref{sc:arcsine-coefficient}.
Its derivative has the sign of $2a(a+2\eta)-(a-\eta)x$. This proves
both extremum rules and supplies the bounds for the inner Green loss and
the varying physical slope.

For the outer arc, drop the nonnegative finite-pole Green function.
The physical deficit dominates its linear lower bound for $0\le y<2$
by~\eqref{sc:physical-deficit}. At $y=2$ the same inequality follows
by the one-sided limit. The identity $1-2\chi_\eta(x)
=(x-d)^2/(x+a)^2$ gives an integrable logarithmic singularity at $x=d$
for the stated equilibrium densities. The positive-window and Gauss
estimates give $F(y)$ as a minorant of the complete raw angular deficit.
Each logarithm in $F$ is concave in $y$. Its argument in the negative
integral stays positive because
$1+2A(x)=(d+x^2)^2/(a-x^2)^2>0$. Consequently $F$ is concave on
$[Y,2]$ and its minimum is bounded below by its smaller endpoint value.
At $Y$, the left rectangular rule in~\eqref{es:inner-negative} and the
nonnegative Green loss give $F(Y)\ge D_{\rm eff}(Y)$. The assumed
right endpoint inequality proves the same bound throughout the outer arc.
The right endpoint can be certified directly since
\[
 N(2)=\int_\alpha^L \log\frac{d+x^2}{a-x^2}\,dx .
\]
The integrand increases on this interval, so each left rectangle is a
lower bound. These estimates establish the usual outer comparison at
every allowed codimension. For $c_t$ in~\eqref{ds:ct}, the required
inner cutoff value can increase by at most $hT_0Y$. The additional
right-endpoint margin absorbs this increase. At the left endpoint the
improved physical and Green bounds of Section~\ref{ds:section} apply
directly. The outer comparison improves with $c_t$ because its derivative
with respect to $c_t$ is $Y-B(t,c_t)/2>0$. The determinant and transfer arguments now
apply to this one fixed contour.
\end{proof}

\section{Curvature of the physical and rational angular field}
\label{sec:curvature}

On a strongly shifted circle the negative interval $[\alpha,L]$ affects
the linear estimate appreciably. Its curvature can instead be paid for by
the physical measure. Let $C_*$ be a positive lower bound for
$\int\chi_\eta\,d\nu_v$ over the allowed mass interval. Let $D$ be
the sum of the physical deficit, the positive window $W$ in
\eqref{es:split}, and the negative interval $N$ in
Proposition~\ref{prop:crossing}. The Gauss charges are kept separately.

\begin{lemma}\label{lem:curvature}
Suppose $J_-$ bounds
$\tfrac12\int_\alpha^L(A(x)/(1+A(x)Y))^2\,dx$ from above. If
$0<d_2\le (C_*^2/\delta-J_-)/2$, then, for $0\le y\le Y$,
\[
 D(y)\ge \left(\int\chi_\eta\,d\nu_v+C_{\rm win}
            +\tfrac12\int_\alpha^L A(x)\,dx\right)y+d_2 y^2.
\]
\end{lemma}
\begin{proof}
The second derivative of the physical term is
$\int \chi_\eta^2/(1-\chi_\eta y)^2\,d\nu_v
\ge C_*^2/\delta$ by Cauchy--Schwarz and $v\le\delta$.
On $[\alpha,L]$ the absolute value of $A/(1+Ay)$ increases with $y$,
so $N''(y)\ge-J_-$. In each difference defining $W$, the two
coefficients satisfy $A(x)\le A(x+s_2)\le0$. The second derivative of
the difference is nonnegative because the square of $A/(1+Ay)$
decreases as $A$ increases to zero. Hence $D''(y)\ge2d_2$.
At zero all terms vanish. The derivative there is the coefficient
displayed in the lemma. Integrating twice proves the result.
The function $A$ increases on $[\alpha,L]$. Left endpoints on a
rational partition thus bound its integral from below and the integral
defining $J_-$ from above.
\end{proof}

For the outer comparison in Proposition~\ref{prop:crossing}, a concave
minorant on $[Y,2]$ is obtained by replacing the physical part by
$Cy$ and retaining $N(y)$ and the log charges. Its value at $Y$ alone
may be slightly below the improved linear cutoff. The missing physical
surplus can be retained on the whole outer interval. Indeed the physical
deficit minus its tangent at zero is increasing, and for $y\ge Y$ it
is at least $C_*^2Y^2/(2\delta)$. Adding this constant to the earlier
concave minorant gives a valid outer minorant. At the left endpoint it
absorbs the negative curvature of $N$ using the same $J_-$, even if
the positive window is replaced by its tangent. Its right endpoint
is at least the previously certified bound. In particular the improved
linear comparison remains valid.

\begin{proposition}\label{prop:curved-transfer}
Let $c_t$ be the varying coefficient of~\eqref{ds:ct}, $B$ the original
angular support endpoint, and $c_0=c+hT_0$. Put
$z_*=c_0+d_2B$. For $0\le t\le t_1<T_0$, choose any step function
$0\le b(t)\le B(t,z_*)$. Suppose the right-endpoint margin after the
varying-slope check exceeds $d_2BY$, and
\[
 (m-hB/2)(T_0-t_1)\ge d_2B^2/4.
\]
The previous response can be reduced further by
$\tfrac14d_2\int_0^{t_1}b(t)^2\,dt$ at its original threshold.
\end{proposition}
\begin{proof}
Fix $t$ and write $b=b(t)$, $x=b/2$. The inequality
$d_2y^2\ge2d_2xy-d_2x^2$ allows the larger linear coefficient
$c_t+2d_2x$ at an additive cost $d_2x^2$ in the inner envelope.
Since $c_t+2d_2x\le z_*$, the fixed-mass identity~\eqref{ds:fixed-mass-identities}
reduces its equilibrium constant by at least
$d_2xB(t,z_*)\ge d_2xb$. The physical first-moment penalty decreases
with the slope. After the additive cost, the saving is at least
$d_2b^2/4$.

At $Y$ the outer minorant just proved includes $d_2Y^2$. This covers
the cost of increasing the linear cutoff by $d_2bY$ because $b\le B\le Y$.
The extra right-endpoint margin covers the same increase at $2$.
The outer comparison then follows from concavity after retaining the
constant physical surplus. Its endpoint test improves when the linear
slope increases, since its derivative is $Y-B(t,z)/2>0$.
The positive-part margin in Section~\ref{ds:positive-part} is at
least $(m-hB/2)(T_0-t)$ after the varying-slope saving. The guard at
$t_1$ permits the new subtraction throughout $[0,t_1]$. For larger
$t$ use the preceding estimate. Integration gives the claim.

For a rational support grid $b_j$ the mass upper enclosures at $c_0$
give those at $z_*$ by adding $d_2Bb_j/2$. If an upper enclosure at
$b_j$ is at most $t$, monotonicity gives $B(t,z_*)\ge b_j$.
Thus the step integral is a finite sum of rational lower rectangles.
\end{proof}

\subsection{Moving past the origin after pole cancellation}
\label{sec:past-origin}

\begin{proposition}\label{prop:past-origin}
Proposition~\ref{prop:crossing} and the curvature bound continue to hold
with $-s_2^2/2<d<0$. The contour is then disjoint from the positive
physical interval.
\end{proposition}
\begin{proof}
The second factorial numerator cancels all denominator factors indexed
by $1\le j<m_2$. The remaining poles are at $-(j/N)^2$ with $j\ge m_2$;
they lie strictly inside the circle for all sufficiently large $N$.
Passing the origin thus preserves the residue identity. In the continuum
rational potential one may likewise cancel the common interval $[0,s_2]$
before taking limits. Remaining numerator logarithms with a zero on
the circle admit uniform upper bounds by truncation from below, followed
by the limit in the truncation level; their logarithmic singularities
are integrable. The pole interval stays a positive distance away.

To verify the paired-window inequality, put $v=x^2$, $w=(s_2+x)^2$
for $0\le x\le s_1$. In the variable $v$ the derivative of the angular
coefficient equals $2R(d+v)/(a-v)^3$. The factors $d+v$ and
$(a-v)^{-3}$ are both increasing. Their covariance under uniform measure
on $[v,w]$ is nonnegative, as follows by integrating the product of
their differences over two independent points. Moreover the average
of the first factor is $d+(v+w)/2\ge d+s_2^2/2>0$. Integration of
the derivative therefore gives $A(s_2+x)\ge A(x)$. Both coefficients
are nonpositive, which is all that the paired-window slope and curvature
arguments require. At $y=2$ the denominator logarithm may diverge
at one point; the integral inequality follows by its one-sided limit.
On $[s_1+s_2,L]$ the derivative of $A$ is positive and $d+x^2>0$,
so the negative-interval bounds and endpoint logarithm still apply.
The physical-factor identity is algebraic for either sign of $d$.
The arcsine extrema use the same projected critical point, now to the
left of positive arguments. All inner, outer, and positive-part checks
are then exactly those in the preceding propositions.
\end{proof}

A second treatment of the signed windows uses the whole available
interval $-s_2^2<d<0$. After common-factor cancellation the signed
angular deficit, apart from the positive Gauss interval, is exactly
\[
 D_w(y)=\frac12\int_{s_2}^L\log(1+A(x)y)\,dx
       -\frac12\int_0^{s_1}\log(1+A(x)y)\,dx .
\]
Its initial slope is $[J(L)-J(s_2)-J(s_1)]/2$, where $J'=A$ and
$J(0)=0$. The second summand is convex. On the inner arc the negative
curvature from the first summand is bounded by
$\frac12\int_{s_2}^L[A(x)/(1+A(x)Y)]^2\,dx$. Thus Lemma~\ref{lem:curvature}
holds with this constant and the displayed initial slope. For the
outer minorant replace the second summand by its tangent
$-J(s_1)y/2$, retaining the first summand and the physical surplus
in Section~\ref{sec:curvature}. The resulting function is concave.
On $[s_2,L]$ both $A(x)$ and $(d+x^2)/(a-x^2)$ are increasing
because $d+s_2^2>0$. Left rectangles therefore certify the curvature
upper bound and the right-endpoint logarithmic lower bound respectively.
The inner left endpoint follows by integrating the same bounds on the
second derivatives. These formulas also apply for $d\ge0$.

\subsection{A contour selecting only the outer poles}
\label{sec:split-residues}
Let $d=-v_*^2$, with $s_2<v_*<1$. In the scaled residue formula a
circle with right endpoint $d$ selects only the poles indexed by
$j/N>v_*$. We add the remaining residues as discrete masses.

\begin{proposition}\label{prop:split-residues}
Define $g(x)=x(\log|x|-1)$ for $x\ne0$ and set $g(0)=0$.
For $s_2\le t\le v_*$ put
\begin{align*}
 H(t)={}&g(s_1-t)+g(s_1+t)+g(s_2-t)+g(s_2+t)\\
 &-g(1-t)-g(1+t)+4U_{b(\delta)}(it)-\ell(\delta).
\end{align*}
If $\sup H<a_*$, the omitted residues do not increase the limiting
determinant response at threshold $a_*$. The contour part obeys the
preceding estimates when the negative integral and outer-arc comparison
are made as described below.
\end{proposition}
\begin{proof}
Set $v_N=(\lfloor Nv_*\rfloor+1/2)/N$ and choose the circle with
endpoints $-u^2$ and $-v_N^2$. This avoids all poles. The residue
identity expresses the original coefficient form as the sum of its
contour integral and the signed residues at $s_2\le j/N<v_N$.
Use the same physical orthonormal basis for both. On the disjoint
union of the circle and the omitted nodes, the absolute measure gives
a positive matrix $A_C+A_D$. The principal-minor Andreief argument of
\eqref{eq:PSD-domination} works on that union, including the atoms.

At $t=j/N$ the logarithm of each absolute rational residue weight,
after the common scaling is removed, is at most $N$ times the six
$g$-terms displayed above, with error $O(\log N)$ uniformly in $j$.
Indeed, factor each difference of squares into $(j-k)(j+k)/N^2$
and omit the zero denominator factor $k=j$. Logarithms of the
resulting factorials satisfy the uniform Stirling estimate
$\log m!=m\log m-m+O(\log(N+2))$ for $0\le m\le 2N$.
The endpoint shifts, the removed denominator factor, and $t^4$
contribute $O(\log N)$. The physical supremum and Markov comparison
from Section~\ref{sec:contour} bound the squared evaluation by
$\exp\{N(4U_{b(\delta)}(it)-\ell(\delta))+o(N)\}$ times the
physical squared norm in the same scaling. The error is uniform on
this compact interval; its polynomial factor comes from the physical
norm comparison. Summing at most $N$ nodes gives
$\|A_D\|\le\exp\{N(\sup H+o(1))\}$. For $|h|\le Ce^{-a_*N}$,
the determinant of $I+|h|(A_C+A_D)$ exceeds that of $I+|h|A_C$
by a factor of at most $(1+|h|\|A_D\|)^r$. Its logarithm is $o(1)$.

For completeness, on the varying circle the Riemann sums of the pole
logarithms converge to their limiting integrals uniformly after division
by $N$. Near the crossing, $|z+(k/N)^2|$ is bounded below by
$|(k/N)^2-v_N^2|$ for the indices in a fixed neighborhood of $v_*$:
the nearest point on the circle to those nodes is its right endpoint.
The half-step spacing gives a lower bound proportional to
$(1+|k-\lfloor Nv_*\rfloor|)/N$ there. Truncating $\log|z+x^2|$ at
$\log\eta$ therefore costs at most
$O(\eta(1+|\log\eta|)+\log N/N)$ in the normalized sums, by summing
$\log(N/j)$ for $j=O(\eta N)$. The limiting integral obeys the
same truncation estimate. Away from this neighborhood the logarithms
are continuous and the remaining distances have a uniform positive
lower bound. Fixed truncation gives uniform ordinary Riemann-sum
convergence. Letting $\eta\downarrow0$ proves the assertion. The circle
parameters converge, and the angular polynomial factors use the actual
radius and center at each stage. Their continuous physical factors and
their exact angular squared-modulus identity have the same limits.

It remains to check the outer minorant. Use the cancelled-window
formula for $D_w$ in Section~\ref{sec:past-origin}. The coefficient
$A$ has its minimum at $x=v_*$, with $A(v_*)=-1/2$. On each panel,
the minimum over its two endpoints and any contained critical point
therefore supplies the inner curvature bound and the lower logarithm
at a fixed intermediate $Y<y_0<2$. At $y=2$ the negative interval
contributes
\[
 [g(x-v_*)+g(x+v_*)+g(u-x)-g(u+x)]_{x=s_2}^{L}.
\]
This is valid across $v_*$ by integration of the logarithmic singularity.
The physical deficit minus its linear tangent dominates
$J(y)=-\delta\log(1-C_*y/\delta)-C_*y$ by Jensen's inequality,
where $C_*$ is the mass-uniform physical slope lower bound.
Indeed the normalized surplus increases with the mean coefficient,
and for a fixed first moment decreases with the mass. On $[Y,y_0]$
use the earlier concave minorant and retain the uniform physical surplus
at $Y$ when checking the left endpoint. On $[y_0,2]$ add the tangent
$J(y_0)+(y-y_0)J'(y_0)$ to the ordinary concave minorant. It suffices
to check the middle and right values exceed the common cutoff by the
specified slope margins. The strict inequalities persist for all
sufficiently large $N$. The limiting eigenvalue and transfer arguments
then apply.
\end{proof}

For interval evaluation of $H$, the physical term is bounded at the
right endpoint of each cell, since it equals
$V(0)+4\int_0^\delta\operatorname{arsinh}(t/b(v))\,dv$ and increases in
$t$. Bound each $g$ from its endpoint values and any of $-1,0,1$ lying
in that cell's argument interval. This handles endpoints where an
individual logarithm diverges.

\paragraph{Physical factors when the whole circle is negative at its right endpoint.}
If $d<0$ then $\chi_\eta(x)=2R(x-\eta)/(x+a)^2$ decreases for $x\ge0$,
since its derivative has the sign of $d-x$. Consequently on the
entire physical axis we may replace $1/2$ by $c_*=-2R\eta/a^2$
in the comparison of $Q_h(x)$ with $(x+a)^2$. The endpoint first-moment
penalty uses $-\log(1-c_* B)$ in place of $-\log(1-B/2)$.
For the arcsine Green estimate, concavity gives
$\log(1-\chi y)\ge (\chi/c_*)\log(1-c_*y)$ for $0\le\chi\le c_*$.
The exterior-map modulus changes by one half the arcsine average of
this logarithm, by its logarithmic potential formula. Thus in the
Green bound one replaces $\log(1-y/2)$ by
$\log(1-c_*y)/(2c_*)$. The remaining mass and inverse-support identities
are unaltered.

For the cancelled window the inner estimate can retain an additional
curvature contribution. Put $A_+=\max\{A(0),A(s_1)\}<0$.
The only possible interior extremum of $A$ is a minimum, so the
second derivative of $-\frac12\int_0^{s_1}\log(1+A(x)y)\,dx$
is at least $s_1 A_+^2/2$ throughout the inner arc. This amount
can be added to $C_*^2/\delta$ in Lemma~\ref{lem:curvature}.
On the outer arc the same window minus its tangent at zero is
increasing. Retaining its surplus at $Y$ therefore also justifies
the strengthened left endpoint of the outer comparison.

More precisely, on each strip $l\le y\le h\le Y$ extend the physical
coefficient distribution by an atom at zero to total mass $\delta$.
The function $x\mapsto x^2/(1-lx)^2$ is increasing and convex on
$0\le lx<1$. Jensen's inequality therefore gives the uniform physical
second derivative bound $C_*^2/[\delta(1-lC_*/\delta)^2]$.
Bound the negative curvature on the same strip at $h$. The smallest
resulting lower bound for the total second derivative over the strips
applies throughout $[0,Y]$. At the outer left endpoint one retains
the actual physical and positive-window surpluses at $Y$; both are
nondecreasing on the outer arc.

\paragraph{Checking separated angular intervals.}
Set $q_+=-J(s_1)/s_1$ and $C=C_*$. The two convex terms of the
cancelled-window deficit have respective initial slopes at least $C$
and $s_1q_+/2$. Subtract their actual initial tangents. The sum of
their surpluses is at least
\[
 J_+(y)=-\delta\log(1-Cy/\delta)-Cy
         -\tfrac{s_1}{2}\{\log(1-q_+y)+q_+y\}.
\]
This follows from Jensen's inequality for the logarithmic surplus.
Moreover, on an inner panel $[l,h]$ their joint second derivative
is bounded below by
\[
 C^2/[\delta(1-Cl/\delta)^2]+s_1q_+^2/[2(1-q_+l)^2].
\]
Indeed $z\mapsto[z/(1-lz)]^2$ is increasing and convex. For the
physical measure, Jensen may be applied after padding by a zero
coefficient to reach mass $\delta$. From this lower bound subtract
the panelwise rational upper bound on the negative integral's
absolute curvature at $h$. Denote half the result by $c_j$.
For equal-width consecutive inner panels a uniform quadratic lower
bound uses $d_2=\min_{k\ge1}k^{-1}\sum_{j=1}^k c_j$.
The average of the piecewise-constant $c_j$ over $[0,y]$ takes its
minimum at a panel endpoint or has its initial constant value.
Hence integrating twice gives $D(y)-D'(0)y\ge d_2 y^2$ even if a
late $c_j$ is negative. The same calculation applies to the surrogate
that retains $J_+$ exactly and uses the certified slope lower bounds.

On each outer subinterval add a tangent to $J_+$ to the ordinary
concave minorant. Its two endpoint values bound the entire subinterval
from below. On the first interval take the tangent at $Y$ so that the
inner comparison provides that endpoint. Every other endpoint is
checked using rational lower rectangles for the negative logarithm or
its closed endpoint primitive at $2$. The minimum margin includes
the subsequent varying-slope and curvature corrections.

\paragraph{Separating the inner Green cutoff from the outer comparison.}
Let $0<B<Y_1<Y<2$ and let $A_1,A_2$ bound the logarithmic Green loss
per unit $y$ on $[0,Y_1]$ and $[0,Y]$ by the same construction, with
$A_2\ge A_1>0$. Use $A_1$ in the effective linear coefficient.
On $[Y_1,Y]$ the cost that is missing from this estimate is at most
$2(A_2-A_1)T_{\max}y$: the proved full-arc upper estimate is $2A_2ty$
at mass $t$, while $2A_1ty$ is already charged. This also holds after
the mass-dependent slope correction.
Set $v=(B+Y_1)/2$, and write $S_j=\sqrt{1+B/e_j}$.
A sufficient test to absorb the extra cost is
\[
 (1-v/Y_1)\sqrt{1-B/v}
 \left(c+\sum_j\frac{K_j}{S_j(Y+e_j)}\right)
 \ge 2(A_2-A_1)T_{\max}.
\]
Indeed the angular equilibrium at any mass $t\le T$ and linear coefficient
at least $c$ has endpoint $B_t\le B$. Outside that endpoint its
Euler--Lagrange deficit has derivative
\[
 \sqrt{1-B_t/y}\left(c_t+\sum_j
 \frac{K_j}{\sqrt{1+B_t/e_j}(y+e_j)}\right).
\]
This is the derivative computed in Proposition~\ref{prop:angular-eq}
with its sign reversed. Integrating from $v$ to $y\in[Y_1,Y]$ gives
at least $y$ times the left side of the sufficient test. The argument
holds with the larger slopes used by the curvature tangent as well.
Thus the middle arc satisfies the same equilibrium bound as the
inner arc, and the original piecewise comparison applies from $Y$
onward.

\subsection{The physical cost as a harmonic average}
\label{sec:harmonic-physical}
For a circle lying strictly to the left of the physical interval, the
cost of the conjugate zeros can be bounded using harmonic measure.
Fix a physical mass $v$ in the admissible range, write $I=[0,b(v)^2]$,
and denote the harmonic extension of a continuous boundary function by
$\mathcal H_b$. It is the bounded solution on the exterior including
infinity. For the potential and Green-function conventions see also \cite{saff2010}. Let $g_b$ be the Green function with pole at infinity and
$G_b$ the finite-pole Green function.
\begin{lemma}\label{lem:harmonic-zeros}
Let $p$ have degree at most $\delta N$ with a prescribed collection
of exterior zeros $\zeta_j$ counted with multiplicity. For
$T_p=\max_{x\in I}|p(x)|e^{-NV(\sqrt x)/2}$ one has in the exterior
\[
 \log(|p(z)|^2/T_p^2)\le N(2U_{\nu_v}(z)-\ell(v))
       +2N(\delta-v)g_b(z)-2\sum_jG_b(z,\zeta_j).
\]
If $Q_h$ has the two conjugate roots $\zeta_h,\bar\zeta_h$ on
the circle, then
\begin{align*}
 -G_b(z,\zeta_h)-G_b(z,\bar\zeta_h)+2G_b(z,-a)
 ={}&\log\left|Q_h(z)/(z+a)^2\right|\\
 &-\mathcal H_b[\log(1-\chi_\eta h)](z).
\end{align*}
\end{lemma}
\begin{proof}
Subtract the potential and infinity Green terms from $\log|p|$ and
add the prescribed finite Green functions. The result is subharmonic because the singularities of the selected zeros cancel. At infinity its
logarithmic coefficient is $\deg p-\delta N\le0$, so it has a
subharmonic extension there. The boundary contact identity and the
maximum principle give the inequality. For the equality cancel the
identical logarithmic singularities on both sides. Their difference
is bounded at infinity and harmonic throughout the exterior. On the
boundary the finite Green functions vanish and
$Q_h(x)/(x+a)^2=1-\chi_\eta(x)h$. Uniqueness of the bounded harmonic
extension proves the assertion.
\end{proof}

Apply the lemma with the same $v=v_t$ as in the support cancellation
identity. On the circle the exact quotient in the logarithm still
has modulus $|y-h|/y$. The physical comparison therefore costs only
the harmonic average of $-\log(1-\chi_\eta h)$ at the contour point.
The auxiliary term $2(\delta-v_t)g_b(z)$ is bounded by its value at
$-a$ since every distance from $z$ to a nonnegative point is largest
at that endpoint. Thus the constant and moving-support estimates are
the same ones used in the transfer above.

For $d=-v_*^2$ let $c_*=-2R\eta/a^2$ and let $C_{\max}$ bound the
arcsine average of $\chi_\eta$ on the entire physical mass range.
Each point $z$ of the contour has $|\varphi_b(z)|\ge
|\varphi_{b_0}(-v_*^2)|$: the arcsine potential of $I$ gives the
first bound at fixed $b$, and the explicit map at the negative
endpoint gives the bound for $b\le b_0$. In the exterior disk
the Poisson density relative to uniform angle is at least
$(|\varphi_b(z)|-1)/(|\varphi_b(z)|+1)$. Under the inverse map
uniform angle pushes forward to the arcsine measure. Hence for any
$p_0\le v_*/\sqrt{v_*^2+b_0^2}$,
\[
 \mathcal H_b[\chi_\eta](z)
 \le c_*-p_0(c_*-C_{\max})=:c_\omega .
\]
Indeed apply the density lower bound to $c_*-\chi_\eta\ge0$.
Convexity first in $\chi$ and then in $h$ gives
\[
 -\log(1-\chi h)
 \le \frac{\chi}{c_*}\frac{h}{B}[-\log(1-c_*B)]
 \qquad(0\le\chi\le c_*,\ 0\le h\le B).
\]
Consequently the physical first-moment coefficient in the fixed-slope
response can be multiplied by $c_\omega/c_*$. This bound is uniform
on both arcs, and the moment continues to decrease under a slope
increase. The outer-arc test cancels the physical coefficient in the
same manner as before.

\paragraph{A uniform Poisson integral.}
A sharper alternative uses the angular distribution itself. Write
$d=a+2\eta=-v_*^2<0$, $R=a+\eta$, and $x=c(1-\cos\theta)/2$ with
$c=b^2$. Then
\[
 \chi_\eta(x)=2R\left((a+x)^{-1}-R(a+x)^{-2}\right),\qquad
 \chi_\eta'(x)=2R(d-x)/(a+x)^3<0 .
\]
For Poisson parameter $\rho$, the convolution of this even, decreasing
boundary function with $P_\rho$ is largest at $\theta=0$. Indeed, by
layering their positive superlevel sets, this reduces to intersecting
two circular arcs, whose overlap is maximal when their centers coincide.
The centered integral is increasing in $\rho$: the $P_\rho$ mass of
$[-t,t]$ is $(2/\pi)\arctan(((1+\rho)/(1-\rho))\tan(t/2))$ for
$0<t<\pi$. It decreases as $c$ increases by pointwise comparison.
It therefore suffices to take $c=b_{\min}^2$ and
$\rho=b_0^2/(\sqrt{v_*^2+b_0^2}+v_*)^2$, rounding these bounds outward.
If $D=\sqrt{a(a+c)}$, $t=c/(\sqrt{a+c}+\sqrt a)^2$, $y=\rho t$ and
$H=(1+y)/(1-y)$, the two centered averages are
\[
 K_1=H/D,\qquad
 K_2=\frac{(2a+c)H}{2D^3}+\frac{2y}{D^2(1-y)^2}.
\]
The first formula follows by multiplying the Fourier coefficients of
$(a+c(1-\cos\theta)/2)^{-1}$ by $\rho^{|n|}$. The second is
$-\partial_a K_1$ at fixed $c,\rho$, using $\partial_a t=-t/D$.
Thus $2R(K_1-RK_2)$ is a uniform upper bound for
$\mathcal H_b[\chi_\eta]$ in the moment estimate.

\paragraph{A secant bound from the exact Green function.}
Continue to assume $a+2\eta<0$. The branch of
$\sqrt{(z-c)/z}$ with positive real part maps the slit exterior to
the right half-plane, taking $-a$ to $\sqrt{(a+c)/a}$. Its Green
function therefore gives, for $c=b^2$,
\[
 e^{-G_b(z,-a)}
 =\frac{c|z+a|}{|\sqrt{a(c-z)}+\sqrt{(a+c)(-z)}|^2}.
\]
The square roots in this expression have positive real parts on the
circle. Normalize the denominator by its value $4a(a+c)$ at $-a$.
In terms of $y=1+\cos\theta$ the resulting function is
\[
 D_c(y)=\tfrac14\left(A+B+2\sqrt{(AB+E)/2}\right)
\]
where $A=\sqrt{1-\chi_\eta(0)y}$, $B=\sqrt{1-\chi_\eta(c)y}$
and $E=1-R[2(a-R)+c]y/[a(a+c)]$. Here $AB$ is a geometric mean
of two positive affine functions and is concave. The square root
is increasing and concave, so $D_c$ is positive and concave on
$[0,2]$. In particular $\log D_c$ is concave and $D_c(0)=1$.
Both $\chi_\eta(c)$ and $R[2(a-R)+c]/[a(a+c)]$ decrease with $c$,
because $2R-a<0$. Thus, for $0\le y\le Y$,
\[
 e^{-G_b(z,-a)}\le C_{00,\eta}(b)\sqrt{2y}\exp(A_Gy),
 \qquad A_G= -\log D_{b_{\min}^2}(Y)/Y.
\]
This bound applies at both cutoffs in the separated angular estimate.

\subsection{Using the same zero subspace on the outer arc and the atoms}
Fix a rational inner cutoff $Y>B$. For $y\ge Y$ every angular factor
has $|y-h|\le y$. The exact Green denominator bound at $Y=2$ gives
a single constant $A_{\mathrm{all}}$ for the whole circle. Before
subtracting the actual angular field $D_{\mathrm{raw}}(y)$, the angular
cost at mass $t$ is therefore at most
$2t A_{\mathrm{all}}y+t\log(2y)$. As
$F_c'(t)=\log(B_c(t)/2)$, the difference of this cost and $F_c(t)$
is increasing for $y\ge Y$. It suffices to verify the uniform slack
\[
 S=\inf_{Y\le y\le2}
 \{D_{\mathrm{raw}}(y)-2T A_{\mathrm{all}}y-T\log(2y)+F_c(T)\}>0.
\]
For the varying slope and curvature the effective linear coefficient
increases by at most $hT_0+d_2 B$. Its equilibrium constant falls by
at most $(hT_0+d_2 B)B/2$ by the fixed-mass derivative identity.
The physical moment penalty can always be bounded by its old upper
bound. The replay conservatively reserves $(hT_0+d_2 B)Y$ from $S$.
The outer test now includes both endpoints of every subinterval.
On each use the midpoint tangent of the convex Jensen surplus, retain
the concave charge and window logs, subtract the linear Green term,
and use the upper endpoint for $\log(2y)$. Concavity of the resulting
minorant makes its endpoint bounds valid on the entire subinterval.

The excluded atoms can also be estimated on precisely the subspace
used on the contour. Put $c=b(\delta)^2$ and let $w=j/N$ be an
excluded node. For a prescribed zero $\zeta$ with angular coordinate
$h\le B$, one has $|w^2+\zeta|\le a-w^2$: the squared distance is
$(a-w^2)^2+2R(w^2+\eta)h$ and $w^2\le v_*^2<-\eta$.
The exact Green identity with left endpoint $-w^2$ gives
\[
 G_b(-w^2,\zeta)\ge
 \log\frac{[w\sqrt{c+x_0}+\sqrt{w^2+c}\sqrt{x_0}]^2}
                 {c(a-w^2)},\qquad x_0=a-RB>0 .
\]
Indeed $\Re\sqrt z\ge\sqrt{\Re z}$ for $\Re z\ge0$ on the principal
branch. Use the selected-zero inequality at $v=\delta$ on the atoms:
their contribution on this subspace has exponent at most $H(w)-2tG$,
where $G$ is the displayed lower bound. Each real interval is checked
using its upper bound for $H$ and its lower bound for $G$.

Specifically, choose $m_0>0$ below the fixed-envelope monotonicity
margin after subtracting $hB/2$, let $T_0\le T$ be the certified mass
lower bound, and set $t_c=T_0-d_2B^2/(4m_0)>0$. The pointwise
upper envelope used in the transfer calculation (including its
assigned response savings) is at least
$a_*+\max\{0,m_0(T_0-t)-d_2B^2/4\}$ on $[0,T_0]$.
This follows from the derivative bound for the fixed envelope, the
varying-slope saving bounded by $h(T_0-t)B/2$, and the curvature
saving bounded by $d_2 B^2/4$. In the end interval the assigned
saving is clipped as in Proposition~\ref{prop:curved-transfer}.
It therefore suffices on each real interval to check
$H<a_*+m_0t_c$ and $H-2t_cG<a_*$. Linear interpolation handles
$[0,t_c]$, and monotonicity of $H-2tG$ handles $t\ge t_c$.
For $t\in[T_0,T]$ only the fixed bound is used and is at least its
endpoint $a_*$. At $t=T$ the remaining eigenvalues are bounded by
min--max. These are estimates on the same subspace for the sum of
the two positive forms. Adding their two norms costs only a factor
two in the pointwise exponential estimate.

\paragraph{Higher angular moments and the physical variance.}
The inner arc satisfies $\Re z\le-a+RY$, so the Poisson cost there
can use $v_*^2=a-RY$ in the distance estimate. For the outer estimate
reserve $(A_{p,\mathrm{all}}-A_{p,\mathrm{in}})_+T_{\max}^2$ from
its slack, using the original valid whole-circle quadratic coefficient.
This allows the smaller coefficient in the common envelope.

For an angular equilibrium measure $\rho$ supported on $[0,b]$,
the density after multiplication by $h$ is proportional to
$\sqrt{h(b-h)}$ times a positive decreasing function by
Proposition~\ref{prop:angular-eq}. Hence if $M=2\int h\,d\rho$,
\[
 2\int h^n\,d\rho/M
 \le b^{n-1}\frac{4}{4^n(n+1)}\binom{2n}{n}
 \qquad (n\ge1)
\]
by reversed covariance with the increasing function $h^{n-1}$
and the elementary beta integral. The physical harmonic measure has
$\mathcal H[\chi^n]\le c_*^{n-1}c_\omega$. Applying the positive
series for $-\log(1-\chi h)$ gives a physical quadratic coefficient
at most $M(T)c_\omega P(c_*B)/T^2$, where
\[
 P(x)=\sum_{n\ge1}\frac{4}{n(n+1)4^n}\binom{2n}{n}x^{n-1}.
\]
After eight terms its tail is at most $x^8/[9(1-x)]$.

For the physical curvature let $V_2$ bound below the unnormalized
variance of $\chi$ for the physical mass padded by zeros up to
$\delta$. For $f(x)=x^2/(1-lx)^2$, $f''\ge2$, hence Jensen applied
to $f(x)-x^2$ adds $V_2$ to the second-derivative bound.
The variance is $(2\delta)^{-1}\iint(\chi-\chi')^2$. To bound it
uniformly use the submeasure $\nu_{v_{\min}}$, obtained by integrating
the unit arcsine mass derivative from zero to $v_{\min}$. Split
this mass interval into twelve equal parts and the arcsine angle
$\theta\in[0,\pi]$ into sixteen equal parts. On each rectangle
$x=b(v)^2(1-\cos\theta)/2$ has certified endpoints, hence so does
the decreasing $\chi(x)$. For each pair of disjoint coefficient
intervals, include the square of their distance times the product
of their masses divided by $\delta$, taking the pair once. Discarding
the other nonnegative contributions proves the stated lower bound.

\paragraph{Keeping the complete quadratic external field.}
On the inner arc one may keep $d_2y^2$ in the equilibrium problem.
For slope $c$ and quadratic coefficient $d\ge0$, the hard-edge
density is that in Proposition~\ref{prop:angular-eq} with the linear
part of the bracket replaced by $c+dB+2dy$. Indeed the Cauchy
transform equals one half of the difference between $D'(z)$ and this
full bracket (including the rational terms) multiplied by
$\sqrt{1-B/z}$. Pole cancellation, the jump, and the coefficient at
infinity give the density and mass. In particular
\[
 t_{c,d}(B)=t_{c,0}(B)+3dB^2/4,\qquad
 F_{c,d}(t)=t\log(B/2)-cB/2-3dB^2/8-2\sum_j K_j\log((1+S_j)/2).
\]
To see the second identity, differentiate the density with respect
to $B$. Its quadratic contribution is
$3dB/(4\pi\sqrt{y(B-y)})$. The other terms give the arcsine
mass derivative as before. Differentiation of the contact constant
at the origin now gives $\partial_t F=\log(B/2)$ including its
$t\log 2$ normalization. The displayed formula has this derivative
and value zero at $t=0$. At fixed mass it gives
$\partial_d F=-3B^2/8$. The first moment still satisfies $M'(t)=B/2$.
Adding a nonnegative quadratic coefficient shrinks $B(t)$ and $M(t)$.
When higher angular moments are used, the replay also checks
$\sum_j K_j/[S_j(B+e_j)^2]>2d_2$ at the original largest endpoint.
The full bracket is then decreasing throughout every smaller support,
as required for the higher-moment inequality.

For $c_t\le c_0$ and $0\le d\le d_2$, the inverse mass endpoint is
at least $B(t,c_0,d_2)$. Thus at each selected mass the contour
saving beyond the varying slope is at least $3d_2 b(t)^2/8$ for
any step function $0\le b(t)\le B(t,c_0,d_2)$. A certified upper
bound on the mass at a test endpoint $b$ is the old $c_0$ mass upper
bound plus $3d_2b^2/4$. Stop the saving at or before
$T_0-3d_2B^2/(8m_0)$ to preserve the positive-part threshold,
and use this larger reserve in the atomic interpolation test.
The previous outer reserve $d_2BY$ exceeds $3d_2B^2/8$.

\paragraph{Integrating the physical cost at its own mass.}
Let $H=c_\omega P(c_*B)$ use the inner-arc series bound and let
$\beta_p$ be the upper coefficient used for the original physical
cost. The more accurate upper bound at mass $t$ is
$H M_{c,0}(t)$. Increasing the linear and quadratic coefficients
only reduces this moment. At a test endpoint $b$ one has
\[
 M_{c,0}(t_{c,0}(b))
 = cb^2/8+\sum_j K_j b(S_j-1)/[2S_j(S_j+1)],\quad
 S_j=\sqrt{1+b/e_j}.
\]
Partition $[0,B]$ into equal endpoint increments. Use a rational
lower bound on $t_{c,0}(b)$ as the right edge of a mass cell and
an upper bound on this displayed moment there. If the cell starts
at $l$, then $(\beta_p l^2-H M_{\rm up})_+$ is a uniform saving
over that cell. The outer test additionally reserves
$\beta_pT_{\max}^2$. The same reserve is added to $3d_2B^2/8$
both in the positive-part cutoff for this extra saving and in the
atomic interpolation test. Thus every credited cell is covered by
the outer, atomic and transfer estimates at the original threshold.

A stronger cell integral follows from the identity $M'(t)=B(t)/2$.
The endpoint increases with mass, so $M$ is convex. If $u_l,u_h$
bound its values from above at the two rational mass edges $l,h$,
their linear interpolant bounds $M$ from above on the whole cell.
The endpoint estimates just used give such bounds also on the
truncated final cell, by monotonicity, and $u_0=0$. Hence the
certified saving on the cell can be increased to
\[
(h-l)\left(\beta_p(l^2+lh+h^2)/3-H(u_l+u_h)/2\right)_+ .
\]
Indeed the integral of the positive part dominates both zero and
the integral of its argument. The endpoint and transfer reserves
are unchanged.

The quadratic saving admits an analogous integration between inverse
mass samples. Write $U_j$ for the upper mass enclosure at $b_j$
including $3d_2b_j^2/4$, and $D_j$ for an upper bound on the
derivative of the linear-field mass there, increased by $3d_2B/2$.
Each charge contributes a decreasing derivative to the linear-field
mass. Thus $D_j$ bounds the derivative of the full mass throughout
$[b_j,B]$. For $t\ge U_j$ the inverse endpoint is at least
\[
 \min\{B,\ b_j+(t-U_j)/D_j\}.
\]
Indeed the mass at the displayed endpoint is at most $t$, by
integration of the derivative bound. Integrating the square of this
capped affine function on each mass cell gives a rational lower
bound in the complete quadratic saving. This estimate also applies
to the last truncated mass cell.

\paragraph{A separable sum for the angular pair energy.}
For a quadrature charge set $K_i>0$ put $S_i=\sqrt{1+B/e_i}$,
$r_i=(S_i-1)/(S_i+1)$ and $K=\sum_i K_i$. Direct multiplication gives
\[
 \log\frac{S_i+S_j}{2}
 =\log\frac{1+S_i}{2}+\log\frac{1+S_j}{2}+\log(1-r_i r_j).
\]
Consequently the ordered weighted sum of these logarithms is
\[
 2K\sum_iK_i\log\frac{1+S_i}{2}
 -\sum_{k\ge1}\frac{1}{k}\left(\sum_iK_i r_i^k\right)^2.
\]
Here $0\le r_i<1$. Truncating the nonnegative series after $m$ terms
omits at most $K^2r_{\max}^{2m+2}/[(m+1)(1-r_{\max}^2)]$.
This identity permits high-order panel calculations with work linear
in the number of charges for each additional series term.

\paragraph{A cubic gain from the moving physical endpoint.}
Let $b_0=b(\delta)$, $b_-=b(v_{\min})$, $a=u^2$, and write $v'(b)$
for the derivative of the physical mass. The constant $H(t)$ preceding
the angular integral has $H'(t)=-k_0(v_t)$ and
$v_t+tq(v_t)=\delta$, $q(v)=u/\sqrt{a+b(v)^2}$.
In particular,
\[
 -H''(t)=\frac{4u^3}{b_t(a+b_t^2)^{3/2}v'(b_t)[1+tq'(v_t)]}.
\]
Use certified $0<V_-\le v'(b)\le V_+$ on $[b_-,b_0]$.
The explicit mass formula gives
$v''(b)\ge -1/[2(1+b_-^2)^{3/2}]$. Therefore for
$f(b)=1/[b(a+b^2)^{3/2}v'(b)]$,
\[
 -\partial_b\log f(b)\ge
 W=1/b_0+3b_-/(a+b_0^2)-1/[2(1+b_-^2)^{3/2}V_-]>0 .
\]
The support equation gives
$-b_t'\ge S=u/[\sqrt{a+b_0^2}V_+]$ and
$-q'(v_t)\ge Q=ub_-/[(a+b_0^2)^{3/2}V_+]$.
Thus $f(b_t)\ge f(b_0)(1+WS t)$ and
$(1+tq'(v_t))^{-1}\ge1+Qt$. For
$A\le 2u^3 f(b_0)$ and $C\le A(WS+Q)/3$, integrating twice yields
\[
 H(t)\le H(0)-k_0(\delta)t-A t^2-C t^3 .
\]
The replay credits only the increment $D=A-A_-\ge0$ over the
previous quadratic bound. It checks that
$D T_{\max}^2+C T_{\max}^3$ fits in the remaining outer and atomic
margins and adds that quantity to the positive-part cutoff reserve.
On the retained interval $[0,\tau]$ the extra integral saving is
$D\tau^3/3+C\tau^4/4$.

\paragraph{Cell integration of the physical second derivative.}
The same implicit equation gives a direct bound on $-H''$ on each
mass interval $[l,h]$. Enclose $b_h$ from below by $x$ and $b_l$
from above by $y$ using rational bisection. The residual bracket is
checked at both ends and its derivative
$v'(b)-tu b/(a+b^2)^{3/2}$ is positive throughout the support range.
Since $b_t$ decreases, a lower bound for $-H''(t)$ on the cell is
\[
 J=\frac{4u^3}
 {y(a+y^2)^{3/2}\{V-lu x/(a+y^2)^{3/2}\}},
 \qquad
 V=\tfrac12\left(1-\frac{x}{\sqrt{1+x^2}}
       +\frac{y}{\sqrt{s_1^2+y^2}}+\frac{y}{\sqrt{s_2^2+y^2}}\right).
\]
The cubic proof established $-H''(t)\ge 2A+6Ct$ pointwise.
One may therefore add $g=(J-2A-6Ch)_+$ to that lower bound on
the cell. For the extra saving reserve
$U=\max g\, T_{\max}^2/2$ in both the outer and atomic estimates.
Shorten the previous cubic cutoff by at least $U/m_0$, where $m_0$
is the positive monotonicity margin. If $z$ is the shortened
cutoff, the additional integrated saving from the cell equals
\[
 \frac g6\left((z-l)_+^3-(z-h)_+^3\right).
\]
This follows by integrating the Hessian twice to obtain the gain
in the physical constant and once more over the retained masses.

\paragraph{Integration over the full slope interval.}
Fix $t\le T_0$ and $c\le z\le c_t=c+h(T_0-t)\le c_0=c+hT_0$.
Choose an endpoint sample $b$ with an upper mass enclosure
$t_{c_0}(b)\le t$, and let $D\ge t_{c_0}'(b)$.
For every such $z$, $B(t,z)\ge B(t,c_0)\ge b$.
The derivative $t_z'(B)$ decreases with $B$ and increases with
$z$, so $-\partial_z B=B/(2t_z'(B))\ge b/(2D)$.
Integration first from $z$ to $c_0$ and then from $c$ to $c_t$
gives
\[
 F(t,c)-F(t,c_t)\ge
 \tfrac h2(T_0-t) B(t,c_0)+
 \frac{b h^2}{8D}(T_0^2-t^2).
\]
On each mass cell the final term has an elementary rational
integral. The first term already supplied the tangent estimate.
Both together are bounded by the full slope difference, itself
at most $h(T_0-t)B/2$ since $B(t,z)\le B(T,c)=B$.
Thus the same positive-part and outer reserves cover the refinement.
The quadratic-field comparison begins at $c_t$, after this entire
linear-field difference has been taken.

\paragraph{Physical moments in the strengthened angular field.}
The first moment of the full quadratic equilibrium obeys
$M(t,c,d)=\frac12\int_0^t B(s,c,d)\,ds$. At fixed mass,
implicit differentiation in the mass equation and the change of
variable $s=t_{c,d}(x)$ give
\[
 -\partial_c M=B(t,c,d)^2/8,\qquad
 -\partial_d M=B(t,c,d)^3/8 .
\]
For example $-\partial_d B=3B^2/(4t'_{c,d}(B))$, and the Jacobian
in the change of variable cancels the denominator. On a mass cell
$[l,h]$, choose $b_*$ with $t_{c_0,d_2}(b_*)\le l$.
Integrating the two identities along the path from $(c,0)$ to
$(c_t,d_2)$ reduces the upper bound on the mean moment by at least
\[
 \{h_0(T_0-h)b_*^2+d_2 b_*^3\}/8,
 \qquad h_0=(c_0-c)/T_0 .
\]
Every intermediate endpoint is at least $b_*$. In the same cell
the original upper sample $b_j$ bounds every actual endpoint from
above. Positivity of the higher-moment series permits the sharper
factor $P(c_*b_j)$ there. Multiplying the revised nonnegative
upper mean by this cell factor and by the upper harmonic coefficient
gives the new physical cost. The entire saving is still at most
$\beta_p t^2$ pointwise before cell integration, since the actual
physical cost is nonnegative; the existing reserve therefore suffices.

The same centered Poisson comparison applies to each power
$\chi_\eta^n$, which is positive and decreasing in the physical coordinate.
Set $c=b_{\min}^2$ (rounded down), take the outward upper bound
$\bar\rho$ for the inner-arc Poisson parameter, and partition
$[0,\pi]$ into 128 equal angles. The centered integrand
$P_{\bar\rho}(\theta)\chi_\eta(c(1-\cos\theta)/2)^n$ decreases
on that interval, so the normalized left-endpoint sum is an upper
bound for its Poisson average. Dyadic half-angle radicals enclose
every cosine. For each $1\le n\le8$ let $U_n$ be the smaller of
this bound and $c_\omega c_*^{n-1}$. In the physical series on a
cell of endpoint bound $b_j$, the coefficient of
$4\binom{2n}{n} b_j^{n-1}/[4^n n(n+1)]$ can be replaced by $U_n$.
All terms are positive, and the same geometric estimate handles
the remaining tail with $x=c_*b_j$.

The centered averages also have exact radical enclosures. With
the fixed outward values of $c,\bar\rho$ just used, expand at $z=0$
\[
 D(z)=\sqrt{(a-z)(a+c-z)},\quad
 Y(z)=\bar\rho\,[2a-2z+c-2D(z)]/c,\quad
 K(z)=\frac{1+Y(z)}{(1-Y(z))D(z)} .
\]
The formula for $K_1$ identifies $K(z)$ with the centered Poisson
average of $(a+x-z)^{-1}$ near zero. Its coefficient $K_j=[z^j]K(z)$
is the average of $(a+x)^{-j-1}$ by the uniformly convergent
geometric series for $|z|<a$. Hence the average of $\chi_\eta^k$ is
\[
 (2R)^k\sum_{j=0}^k {k\choose j}(-R)^j K_{k+j-1}.
\]
The replay obtains the first sixteen coefficients by squaring $D$
and multiplying out $(1-Y)DK=1+Y$, with outward rounding also on
the alternating sum. This gives a second upper bound for each
$U_k$ in the cell integral.

\paragraph{Integrating the higher equilibrium moments.}
Let $M_n(t,c,d)=2\int y^n\,d\rho$ and
$A_n=4^{-n}{2n\choose n}$. The arcsine mass derivative used above
has $n$th moment $A_n B^n$ with this normalization, hence
\[
 M_n= A_n\int_0^t B(s,c,d)^n\,ds,\quad
 -\partial_c M_n=\frac{n A_n B^{n+1}}{2(n+1)},\qquad
 -\partial_d M_n=\frac{3n A_n B^{n+2}}{4(n+2)} .
\]
Both field derivatives follow from the same change of variable as
for $M_1$. At the base field $(c,0)$, on each cell $[l,h]$ set
$\Delta_n=A_n b_j^n(h-l)$. An upper bound $u$ on $M_n(l,c,0)$
therefore gives $u+\Delta_n$ at the right edge and
$u+\Delta_n/2$ for the cell mean. Carry the right-edge bound
to the next cell. For the strengthened field subtract from the mean
the two displayed derivative gains with $B$ replaced by the lower
endpoint $b_*$ and the field increments replaced by $h_0(T_0-h)$
and $d_2$. Take the smaller of this bound and the earlier
beta-integral bound for each $n$. The finite sum with coefficients
$U_n/n$ and the same positive tail then bounds the mean physical
cost. All power moments are compared on the same cell and the
positive-part estimate uses the total cost once.

For the base field $(c,0)$ the mass derivative $t'_c(B)>0$ is
decreasing, so its inverse $B(t,c,0)$ is convex. In a cell
$[l,h]$, let $L$ and $H$ be upper bounds on the respective
endpoints. At $t=l+s(h-l)$ the endpoint is at most
$(1-s)L+sH$. For an upper bound $u$ on the initial moment,
the right-edge increment and the increment of the cell mean are
therefore at most, respectively,
\[
 \frac{A_n(h-l)}{n+1}\sum_{k=0}^n L^{n-k}H^k,
 \qquad
 \frac{A_n(h-l)}{(n+1)(n+2)}
 \sum_{k=0}^n(n-k+1)L^{n-k}H^k .
\]
These are the integrals of the chord power against $ds$ and
$(1-s)\,ds$, by the integral of each Bernstein monomial. The
estimate remains valid in the truncated last cell. For $n=1$
the result may also be compared with the secant estimate from
the exact base-field moment at both edges.

\paragraph{Integrating the quadratic saving across both field coefficients.}
Write $d=d_2$ and let $b=b_j$, $D=D_j$, $U=U_j$ be the samples
used for the capped affine lower bound on the endpoint. Fix
$t\ge U$ in this cell and $0\le x\le d$. The mass identity gives
\[
 t_{c_t,x}(b)=t_{c_0,d}(b)-h_0 t b/2-3(d-x)b^2/4,
 \qquad c_0-c_t=h_0t .
\]
The endpoint for $(c_t,x)$ lies in $[b,B]$. Its mass derivative
on this interval is at most $D$. Thus the endpoint is at least
$b+(t-U)/D+At+G(d-x)$, with $A=h_0 b/(2D)$, $G=3b^2/(4D)$.
This also bounds it below by the original capped affine function
plus the last two terms. Since the capped function is at least
$b$, squaring and averaging in $x$ adds at least
\[
 2b(At+Gd/2)+A^2t^2+AGdt+G^2d^2/3
\]
to its squared lower bound. Integrate this polynomial over the mass
cell and multiply by $3d/8$, using $-\partial_x F=3B(t,c_t,x)^2/8$.
This estimates the same quadratic field difference on the same
mass interval. Its pointwise upper reserve remains $3dB^2/8$.

The support used by the physical harmonic estimate equals $[0,b_t^2]$
at angular mass $t$. The preceding centered Poisson comparison is
decreasing in $c$ and increasing in the upper bound for $\rho$.
For $t\in[l,h]$, the implicit support equation bounds $b_t$
between a lower enclosure of $b_h$ and an upper enclosure of
$b_l$. Choose $c$ as the square of the former and use the latter
in $\rho=b^2/(\sqrt{a-RY+b^2}+\sqrt{a-RY})^2$.
The same resolvent then bounds each centered moment throughout
this cell. The computation bisects the physical support on a second
mass grid; each angular cell uses the grid edges immediately outside
it. Taking the minimum with the uniform bound term by term preserves
the earlier estimates and their reserves.

\paragraph{Physical power moments for the curvature bound.}
For the submeasure $\nu_v$ at the certified floor $v$, write
$J_k=\langle(a+x)^{-k}\rangle$ for the arcsine mean on $[0,c]$.
Put $D^2=a(a+c)$ and $w=a+c/2$. The generating function
$\sum_{k\ge 1}J_k z^{k-1}=(D^2-2wz+z^2)^{-1/2}$ gives
$J_1=D^{-1}$ and
$kD^2J_{k+1}=(2k-1)wJ_k-(k-1)J_{k-1}$.
Thus the arcsine mean of $\chi^n$ follows from the same alternating
sum as in the centered Poisson formula. Since $\chi$ decreases
in $x$, the mean of each positive power decreases in $c$.
Integrate the arcsine mass derivative up to $v$ using endpoints
of 128 equal mass intervals. Right endpoints give lower bounds
$L_n$ for $\int \chi^n\,d\nu_v$; left endpoints give an upper bound
$U_1$ for the first moment. An uncertainty in $c$ contributes
at most $2R|\Delta c|/a^2$ to the first mean since
$|\chi'(x)|=2R(x-a-2\eta)/(a+x)^3\le 2R/a^2$.

The unnormalized variance of the full padded measure is at least
the variance of its submeasure $\nu_v$: minimize the integral of
$(\chi-z)^2$ over constants $z$ before and after restricting the
integral. The latter variance is at least $L_2-U_1^2/v$.
In a curvature strip with left edge $l\ge0$, positivity also gives
the direct physical second-derivative bound
$\sum_{k=0}^4(k+1)l^k L_{k+2}$, by expanding
$\chi^2/(1-l\chi)^2$ on that submeasure. Each strip uses the
larger of this bound and the variance-enhanced Jensen bound.

\paragraph{A quadratic minorant from stripwise second derivatives.}
Let $q(s)$ be the stepwise lower bound for $D''/2$, with the
curvature-strip width rescaled to one. Subtract $D(0)+D'(0)y$.
At every positive $y$ the coefficient of a valid quadratic
minorant is bounded below by
$2y^{-2}\int_0^y(y-s)q(s)\,ds$. In the strip $j\le y\le j+1$
this expression is $q_j+P/y+Q/y^2$, where
$P=2(\int_0^j q(s)\,ds-jq_j)$ and
$Q=j^2q_j-2\int_0^j s q(s)\,ds$. Its minimum is attained at
an edge or at $-2Q/P$ when that point belongs to the open strip;
its limit at zero is $q_0$. All tests are rational. The resulting
minimum gives $D(y)\ge D(0)+D'(0)y+d_2y^2$ on the whole inner
arc. This pointwise field comparison suffices for the full quadratic
equilibrium. The separate joint outer and atomic inequalities are
reevaluated with the resulting $d_2$.

The full quadratic field permits a separate coefficient on the near
arc. Let $d$ be a quadratic minorant on $[0,Y]$ and $d_1\ge d$
one on $[0,Y_1]$, where $B<Y_1<Y$. Let $s>0$ be the certified
middle-arc Euler--Lagrange deficit per unit $y$ remaining after the
Green correction. Choose $d'\le\min(d_1,d+s/Y)$ with $d'\ge d$.
Use the equilibrium for $d'$ on the whole inner arc. On $[Y_1,Y]$
the excess $(d'-d)y^2$ is at most $sy$. The bound on the available
deficit used for the linear problem applies to this equilibrium:
its endpoint decreases and its exterior derivative has the same
positive bracket with the added terms $d'B_t+2d'y$. Thus the
Green correction and the quadratic excess together fit the
middle-arc deficit. Beyond $Y$ the joint outer and atomic tests
are performed afresh with $d'$, as is the condition that the
equilibrium density bracket decreases on its support.

There is also a direct lower series for the physical field together
with the two signed raw windows. Write $q_+$ for the lower bound
on the mean of $-A$ on $[0,s_1]$, and let $N_n$ be an upper
sum for $\tfrac12\int_{s_2}^L(-A)^n\,dx$, using an upper bound
$A_*$ for $-A$ on the negative window. After subtracting the
actual initial tangent, the quotient by $y^2$ is at least
\[
 \sum_{n=2}^m \frac{L_n+s_1q_+^n/2-N_n}{n}y^{n-2}
       -\frac{N_{m+1}y^{m-1}}{(m+1)(1-A_*y)} .
\]
This follows from the logarithm series, Jensen's inequality on
the positive window, and $N_{m+1+k}\le A_*^kN_{m+1}$ for the
underlying integrals. The replay uses the same pointwise coefficient
majorants in all upper sums. The expansion is uniform up to $Y$
because $A_*Y<1$. The signed polynomial is bounded below on each
of 64 subintervals using the appropriate endpoint for each monomial,
and the negative tail uses the right endpoint. The least of these
bounds can improve either the near or the full quadratic coefficient.

\paragraph{Closed Stieltjes transform of the physical submeasure.}
Fix its mass $v$ and corresponding radius $b$. For $x>0$ write
$u_x=\sqrt{x}$, $w_x=\sqrt{x+b^2}$ and
$J_r(x)=\log[(\sqrt{r^2+b^2}+w_x)/(r+u_x)]$. Its Stieltjes
transform has the following expression:
\[
 H(x)=\int\frac{d\nu_v(t)}{x+t}
 =\frac{\log[(b+w_x)/u_x]-J_1(x)+J_{s_1}(x)+J_{s_2}(x)}
        {2u_x}.
\]
Indeed the arcsine mass derivative at radius $z$ contributes
$1/(u_x\sqrt{x+z^2})$ per unit mass. Integrate this against
$\{1-z/\sqrt{1+z^2}+z/\sqrt{s_1^2+z^2}
+z/\sqrt{s_2^2+z^2}\}\,dz/2$ from $0$ to $b$. The first
primitive is the inverse hyperbolic sine; the other three follow
by differentiating $\log(\sqrt{r^2+z^2}+\sqrt{x+z^2})$.

Consequently $[z^j]H(a-z)=\int(a+t)^{-j-1}\,d\nu_v(t)$.
The binomial expansion of $\chi^n$ gives its full moment by a
finite alternating sum of these coefficients. The computation
uses the outward enclosure of the fixed radius $b$, the relations
$U^2=a-z$, $W^2=a+b^2-z$, and $X(\log X)'=X'$ to enclose
twenty Taylor coefficients. Only the constant logarithms require
transcendental evaluation. This gives ten moment lower bounds and
an upper bound for the first moment, allowing both the variance
estimate and the signed-field expansion to use the same submeasure.

\paragraph{Green cutoffs on individual mass intervals.}
On $l\le t\le h$ choose $b$ with $h\le t_{c,0}(b)$ and
$b_*\le B(l,c_0,d')$. For $b<Y_s<Y_1$, the radical Green
denominator gives a slope $A_s\le A_1$ up to $Y_s$. In that
range a further linear field $g(t)y$, $g(t)=2t(A_1-A_s)$, can
be retained. Check $2(A_1-A_s)\le h_0$, so $c_t+g(t)\le c_0$.
The new equilibrium endpoint is at most $b$ and at least $b_*$.
Writing $w=(b+Y_s)/2$, require
\[
 (1-w/Y_s)\sqrt{1-b/w}
 \left(c+\sum_j\frac{K_j}{S_j(B)(Y+e_j)}\right)
 \ge 2(A_2-A_s)h+(d'-d)Y .
\]
The exterior derivative bound of the preceding argument now applies
from $w$ to every $y\in[Y_s,Y]$. This pays the full middle-arc
Green discrepancy and any localized quadratic excess there.
The outer and atomic estimates use the same upper support $B$;
the additional change in the contact constant is at most
$h(A_1-A_s)B$. The replay checks this against the remaining margins
after the physical reserves. It also checks that, with this quantity
added to all earlier reserves, the positive-part margin at $h$ stays
positive. Increasing the slope decreases all angular moments, so all
previous physical cost bounds remain valid. At fixed $t,d'$ the
change of the contact constant is at least $g(t)b_*/2$ in absolute
value, by integrating $-\partial_c F=B(t,c,d')/2$. Thus this cell
contributes an additional saving $(A_1-A_s)b_*(h^2-l^2)/2$.
Cells without a successful comparison keep the preceding estimate.

The signed-window moments themselves have elementary enclosures.
Set $K_j=\int_l^r(a-x^2)^{-j}\,dx$ for either interval. Then
\[
 K_1=\frac1{2u}\log\frac{(u+r)(u-l)}{(u-r)(u+l)},\qquad
 2(k-1)aK_k=\left[\frac{x}{(a-x^2)^{k-1}}\right]_l^r+(2k-3)K_{k-1}.
\]
Differentiation of the endpoint term proves the recurrence.
Since $-A(x)=2R((a-x^2)^{-1}-R(a-x^2)^{-2})$, the binomial
theorem integrates every required power on both windows. Lower
enclosures on the positive window and upper enclosures on the
negative window give the signed coefficients. In the tail one
may take $A_*=1/2$, since $(a-x^2-2R)^2\ge0$.

For an upper endpoint $b$, the exterior derivative permits the full
factor $y^{-1}\int_b^y\sqrt{1-b/x}\,dx$. It is nondecreasing for
$y\ge b$: extend the integrand by zero to $[0,b]$ and average a
nondecreasing function on $[0,y]$. Its value at the proposed lower
cutoff is bounded below by 64 equal-width left rectangles, including
the vanishing first rectangle. This gives a second rigorous test for
both the common middle-arc comparison and its mass-cell version.

On the same mass cells one may sharpen the physical initial slope.
Let $P(v)=\int\chi\,d\nu_v$, and write $A_{\min}$ for the
previous lower bound on $P'(v)$. In $[l,h]$, an upper bound $q_h$
for $u/\sqrt{a+b_h^2}$, capped by the bootstrap bound $q_*$,
gives $v_t\ge\delta-q_ht$. Integrate a nonnegative lower bound
for $P'(v)-A_{\min}$ up to $\delta-q_hh$ from $v_{\min}$, the
mass floor, to obtain $E$. The arcsine derivative formula bounds
this integrand on 32 rational mass intervals. Since
$v_{\min}=\delta-q_*T_{\max}$, the available slope increment over
$c_t$ and the previously used Green increment is at least
\[
 E+A_{\min}(q_*-q_h)l+
       (A_{\min}q_*+2A_1)(T_{\max}-T_0).
\]
Here one also uses $h_0\le A_{\min}q_*+2A_1$.
Cap the nonnegative improvement by $(h_0-2A_1)l$ to keep the
total slope below $c_0$. A constant increment $p$ within this
cap gives an additional saving $p b_*(h-l)/2$. Its pointwise cost
in the upper reserve is at most $pB/2$ and is checked jointly
with the Green increment on this same cell.

\section{The earlier reference certificate}
\label{sec:certificate}

The rational windows in this reference calculation are
\[
 s_1=\CertSOne,\qquad s_2=\CertSTwo,\qquad\delta=\CertDelta.
\]
The inherited base calculation uses cutoff $\CertBaseCutoff$ and
$\CertCells$ exact affine cells. On the refined prime bands, the
correction checker takes the lower envelope of the inherited profile
and all applicable bounds in Sections~\ref{lyr:section},
\ref{lyrx:section}, \ref{sec:odd-max-transition}, and~\ref{conf:section}. It checks
the floor values, root-count orderings, eligibility inequalities, and
every crossing of the affine candidates at exact rational endpoints.
The correction has $\CertCorrectionCells$ checked cells.

The general second-layer refinements are used for $q\le\CertArithmeticQmax$.
Retained-root, equal-even/one-odd third-layer, and both-odd third-layer
refinements are restricted respectively to
$q\le\CertMixedQmax$, $q\le\CertThirdQmax$, and $q\le\CertOddQmax$.
Saturated odd flags are used for $q\le\CertSaturatedQmax$.
The central-root two-layer counts are used for
$q\le\CertCentralSecondQmax$, and the adjacent-quotient odd,
mixed, and vanishing-cross-constraint refinements for
$q\le\CertCrossedQmax$. The determinant profiles use the
restricted-kernel corollaries; the stronger exact-rank statements
retain their stated finite guards.
The two unconditional containment-core families are used through
$q\le\CertCoreQmax$. The all-parity confluent family is used through
$q\le\CertConfluentQmax$ and has $\CertConfluentWitnesses$ distinct
fixed rational low-matrix witnesses in the final calculation.
Only the floor version of the odd-head extension is needed by the
production profile. The optional ceiling theorem is retained as a proved
extension, but is not used in the reported numerical bound.
The upper odd maximal-root transition theorem is used through
$q\le\CertUpperOddQmax$. Its $\CertUpperOddWitnesses$ contributing tuples
have the additional cross-matrix witness and the required nonzero
$\theta,\kappa$ witnesses. Both saturated and safely unsaturated
branches are restricted to $u_2<w\le u_1$ with $k\ge0$.
The same finite cutoff includes $\CertOuterUpperWitnesses$ outer upper
tuples, $\CertPaidLowerWitnesses$ one-sided lower tuples, and
$\CertOuterLowerWitnesses$ outer lower Cauchy tuples. Their respective
guards are both heads saturated with $w>u_1$; the paid one-sided lower
transition with $\tau_1<w\le\tau_2$; and the Cauchy-degenerate outer
lower transition with $w>\tau_2$ and $H_1+H_2\le w$. The last family
uses the equal-coefficient flag lemma, including singular leading odd
restrictions. Each retains $k\ge0$ and its stated local matrix tests.
All other bands, the small-prime tail, and the large-prime estimate retain
the inherited bounds. Every fixed rational matrix that is actually used
has an exact modular nonsingularity witness in the JSON certificate.
At the final windows, $1/(s_2-s_1)<\CertCommonRootCutoff$.
Both the common-quotient
geometry and the adjacent-quotient geometry above require
$s_2-s_1<\beta$. Hence their proved central-root and third-layer
candidates have no eligible bands with $q\ge\CertCommonRootCutoff$; the cutoff $128$
already covers their possible prime bands. The retained-root Cauchy
family is not subject to this restriction.

The exact integral is rounded conservatively before transfer. The
\emph{applied} saving in the arithmetic upper rate is
\[
 \eta_{\rm applied}=\CertAppliedSavingExact.
\]
The difference between the unrounded saving and this applied saving is
nonnegative and at most $10^{-18}$. Thus the exact final gap is
\[
 \varepsilon=\CertGapExact,
 \qquad C_{\rm ar}^{\rm up}=\CertArithmeticExact.
\]
The analytic lower rate is unchanged by the arithmetic correction.

The spectral response uses the shifted circle of
Sections~\ref{sc:section} and~\ref{es:section},
with
\[
 u=\CertU,\qquad B_{\max}=\CertAngularB,
 \qquad Y_c=\CertYcut,\qquad T_{\max}=\CertTcap.
\]
The actual combined window is $\alpha=s_1+s_2$. The positive
Gauss interval starts at $L=\CertAlphaZero\ge\alpha$.
There are $\CertPanels$ panels with edges
\[
 x_j=1-(1-L)(1-j/J)^3,\qquad J=\CertPanels.
\]
Two downward-rounded Gauss nodes on each panel give
$\CertCharges$ positive angular charges. The centre is
$\eta=-L^2$ and the radius is $u^2+\eta$. The negative interval
$[\alpha,L]$ is retained using $\CertExtendedPanels$ downward-rounded
left-sum cells. Its contribution enters both the effective inner slope
and the full outer-monotonicity guard. All contour, support,
positivity, and outer-arc guards are checked with rational intervals.
No fitted spectral values enter the certificate.
The varying-slope argument of Section~\ref{ds:section} then reduces the
integrated response by at least $\CertDynamicSaving$. Its certificate
uses $\CertDynamicTangents$ rational tangents, slope improvement
$h\ge\CertDynamicH$, and positive-part guard margin at least
$\CertDynamicGuard$. The shifted-circle bound before this final
postprocessing is at most $\CertPreviousResponseMuUpper$.

\begin{center}
\begin{tabular}{@{}ll@{}}
\toprule
Quantity & Certified directed bound\\
\midrule
Degree rate $\lambda$ & $\CertLambda$\\
Arithmetic upper rate & $C_{\rm ar}\le\CertArithmetic$\\
Analytic lower rate & $C_{\rm an}\ge\CertAnalytic$\\
Positive gap & $\varepsilon\ge\CertGap$\\
Zero-angle norm rate & $\gamma\le\CertGamma$\\
Shifted Green decrease & $\kappa_\eta\ge\CertKappa$\\
Combined quadratic coefficient & $\beta\le\CertBeta$\\
Effective angular coefficient & $c_{\rm eff}\ge\CertLinear$\\
Negative mass gain & $A_-\ge\CertMassGain$\\
Physical comparison cost & $A_p\le\CertPhysicalBeta$\\
\bottomrule
\end{tabular}
\end{center}

Four support iterations give $b_{\min}\ge\CertBMinLower$.
The positive mass cost used on the outer arc satisfies
$A_m\le\CertMassBeta$. Write $I_9$ for the integrated exponent
after the valid varying-slope subtraction; its endpoint threshold
$a_8=g_8(T)$ is unchanged. The exact calculation gives
\[
 \CertTLower\le T\le\CertTUpper,
\]
\[
 \CertILower\le I_9(T)\le\CertIUpper<\varepsilon,
\]
\[
 \varepsilon-I_9(T)\ge\CertISlack>0.
\]
The outer condition \eqref{sc:outer-check} has margin at least
$\CertOuterSlack>0$. With $a_8=g_8(T)$ and
$\mathcal R_9=I_9(T)-Tg_8(T)$, the final enclosures are
\[
 \CertThresholdLower\le a_8\le\CertThresholdUpper,
\]
\[
 \mathcal R_9\le\CertResponseUpper.
\]
Consequently,
\begin{equation}\label{eq:final-number}
 \mu(\zeta(5))\le\frac{\lambda a_{8,\rm up}}
                           {\varepsilon-\mathcal R_{9,\rm up}}
 =\CertMuExact<\CertClaim.
\end{equation}
The exact fraction lies between
$\CertMuLower$ and $\CertMuDetailed$.
For a strict application of Theorem~\ref{thm:transfer}, the checker
also supplies a rational $c>0$ and verifies
\[
 \varepsilon-\mathcal R_{9,\rm up}-\lambda/c>0,
 \qquad \CertClaim-a_{8,\rm up}c>0.
\]
The exact $c$, both strict margins, and every intermediate enclosure
are retained in the response certificate.

\subsection{A short complete certificate for integer exponent 90}

A separate companion uses only the general second-layer Cauchy bound
and the inherited relative-norm transfer. Its rational parameters are
\[
 s_1=\frac{29163}{250000},\qquad
 s_2=\frac{15899}{125000},\qquad
 \delta=\frac{1117219}{1000000}.
\]
At arithmetic cutoff $1/256$, its inherited base has $3837$ cells and
its correction has $2120$ checked cells. Its exact gap is
\[
 \varepsilon_{\rm short}=\frac{949231736337453}{12500000000000000},
 \qquad
 \gamma_{\rm norm}\le\frac{1955707855932751}{250000000000000}.
\]
With $\lambda=109101/125000$, this proves
\[
 \frac{\lambda\gamma_{\rm norm}}{\varepsilon_{\rm short}}
 =\frac{71123227596706355617}{791026446947877500}
 <89.912578613686761<89.92<90.
\]
The companion does not use the third arithmetic layer, distributed
zeros, the shifted contour, or the root-multiplicity theorem. Indeed,
when $|h|<\exp(-\gamma' N)$ with $\gamma'>\gamma_{\rm norm}$, the
relative operator norm is less than one, so $P_N(\xi+h)\ne0$.
The integer-polynomial lower bound then gives the inherited norm
transfer. Its checker supplies a rational transfer constant and two
strict positive margins. This gives a compact complete finite replay
of the integer-$90$ consequence.

\section{Replaying the reference configuration}
\label{sec:replay}

The production chain has four stages:
\begin{enumerate}[leftmargin=*]
\item Recompute the inherited finite arithmetic and analytic base at
the final rational windows.
\item Recompute the lower envelope of the proved local arithmetic
profiles, including every fixed modular witness actually used, and
integrate its saving with exact rational arithmetic.
\item Recompute the shifted-circle response, retained negative interval,
and both strict transfer inequalities using the corrected base bytes
as a hash-linked input.
\item Recompute the varying-slope response saving from the hash-linked
shifted-circle certificate and verify the final strict transfer.
\end{enumerate}

All production calculations use the Python standard library. Square
roots are enclosed using integer square roots. Logarithms use an
explicit range-reduced series and remainder. The final evaluator encloses
this positive series on the fixed integer grid $2^{-112}$; directed
integer division controls every operation and the geometric remainder.
Local moments are evaluated using the exact reciprocal-series identity
for polynomial quotients, and centered parity splits each confluent
low determinant into two transposed blocks. The Hankel block determinant
uses a monic orthogonal-polynomial recurrence when all leading pivots
are nonzero, and ordinary modular Gaussian elimination otherwise.
These accelerations change the evaluation method, not the mathematical
functional or the certified matrix.
Running with Python
assertions disabled is rejected. The launcher checks source and
certificate hashes, parameter consistency, directed rounding, and
every mathematical field of the fresh outputs against the frozen
certificates.

From the extracted package directory, run
\begin{verbatim}
python3 run_all.py --simple
python3 run_all.py --response-only
python3 run_all.py
\end{verbatim}
The first command recomputes the short integer-$90$ companion.
The second recomputes the arithmetic correction and response while
using the hash-verified inherited base; its output explicitly states
that this run did not recompute the inherited base integral. The
third recomputes every finite stage. Add \texttt{--output replay.json}
to save complete replay records. Earlier configurations, proofs, and
experiments are retained separately in the archive.

\paragraph{Mathematical review.}
The internal reviews covered the Newton and Bernoulli precision,
maximal-class cancellation, complete low-feature isomorphism, correlated
top flags, Cauchy moments and finite offsets, retained central roots,
first-jet identity, maximal Schur correction, both-odd boundary flags,
restricted second kernels, adjacent-quotient coefficient cancellation,
and passage to the $N^2$ rate. This continuation additionally reviews the
containment-core singular/invertible dichotomy, canonical maximal
radicals, the exact Christoffel identity, the first-order unit coefficient,
confluent residues, crossed flags giving true first kernels, and the
all-parity odd-head constraints. A separate pointwise implementation audit
compares the new rates with the exact production profile.
The upper odd transition proof additionally checks the maximal cross
pivots, polynomial interpolation of the missing column directions,
the bounded full-field coefficient remainder, and the transformed boundary
form in the unsaturated case.
The shifted-contour geometry, positive
Gauss remainder, retained negative interval, window primitive, physical
moment primitive, Green monotonicity, exact support cancellation, and
varying-slope response reduction were checked separately.
The written asymptotic arguments remain mathematical inputs to the
finite replay; a successful replay is not a formal proof of them.

\paragraph{Original-matrix diagnostics.}
An accelerated exact constructor was first compared entrywise with
the retained original constructor. Subsequent Smith calculations
were used to select possible theorem extensions. One useful check takes
\[
 (p,N,m_1,m_2,r)=(251,1530,179,197,1707).
\]
The finite both-odd theorem predicts first rank $339$, second kernel
dimension $60$, and third rank at most $6$. The original matrix gives
valuation counts
\[
 (n_0,n_1,n_2,n_{\ge3})=(339,1308,6,54),
\]
matching all three counts exactly. A second example with
$(m_1,m_2,r)=(152,158,1746)$ has $k_f=12<w=24$ but
$k_f+H_1=27\ge w$, and therefore tests the relaxed guard. It gives
\[
 (n_0,n_1,n_2,n_{\ge3})=(351,1362,18,15),
\]
again matching the finite theorem. A further test at
$(p,N,m_1,m_2,r)=(101,617,57,61,704)$ has
$k_f=4$, $w=11$, and $(H_1,H_2)=(3,7)$. It fails the former
boundary-surjectivity guard, while its first three Smith nullities
sum to $567$, exactly the conservative determinant lower bound.

For the mixed adjacent-quotient family,
$(p,N,m_1,m_2,r)=(101,814,88,111,860)$ has first rank $102$,
second kernel dimension $45$, and third rank $4$. The chosen
left and right second-kernel subspaces have dimensions $41$ and
$45$, with zero third restriction, so their valuation lower bound
is $758+41+45=844$. These diagnostics check the finite mechanisms;
they are not substituted for the uniform proofs.

For the new leading containment-core bound, original-matrix calculations at
\[
 (p,N,m_1,m_2,r)=(101,821,97,105,916),
\]
and at $(151,1227,144,157,1369)$ and $(211,1714,202,220,1912)$
give first ranks $238$, $346$, and $471$, respectively. Each equals
the finite bound of Section~\ref{lyrx:section}, including its interpolation
offset. The all-parity example in Section~\ref{conf:all-parities}
likewise attains the stated first-two-nullity lower bound. These checks
distinguish a rate improvement from bounded finite offsets.

\section{Arithmetic gains in the reference configuration}

The preceding release gave $\mu(\zeta(5))<88.58905$. The largest new
saving comes from retaining maximal-class radicals that an earlier
second-layer construction discarded. Reflection about $Y=-1/2$ gives
canonical even radicals, and the exact Christoffel relation between the
maximal and nonmaximal functionals makes their mixed blocks vanish.
The remaining first jets and local top coefficients form one global
residue pairing with double poles at the maximal classes. Its rank is
controlled by a degree deficit. This is the confluent arithmetic theorem
of Section~\ref{conf:section}.

Root counts need not be constant on the maximal classes. Crossed flags
force a top coefficient to vanish both when its own root space is too
small and when the opposite root space contains unwanted directions.
These two conditions make the chosen spaces actual first kernels of the
original matrix, a necessary hypothesis for the second-layer valuation
bound. Truncating odd endpoint heads then reduces to even effective
spaces. The discarded opposite heads impose explicit Cauchy constraints;
keeping all four caps on their rank extends the saving to every parity.

The containment-core argument gives a smaller additional improvement in
both the first and second local layers. If the fixed core matrix is
singular, its rank defect already gives the desired estimate. If it is
invertible, root cancellation on the exceptional classes reduces the
remaining top pairing to ordinary interpolation. No rational matrix
nonvanishing hypothesis is needed for this dichotomy. The common mixed
head and capped crossed-head extensions remove further restrictions
from the earlier valuation profiles.

Further steps cross maximal-root transitions in the
upper odd bands. An odd column interpolation space fills the missing
maximal column features. Its first-jet correction cancels the exceptional
boundary contribution; the remaining full-field term has bounded rank.
In the unsaturated case a shorter interpolation polynomial makes that
term vanish exactly. Both cases reduce to a boundary form on the remaining
common-root classes, and yield the third-layer estimates of
Section~\ref{sec:odd-max-transition}.

The outer upper extension fills both missing maximal feature spaces and
uses an exact projection identity on the larger local matrix. Omitting
the free zero-class directions leaves a bounded finite loss and preserves
the limiting rate. In the lower transition, the enlarged leading rank
must be paid explicitly. A one-sided version retains the same boundary
mechanism with this cost.

In the final outer lower family, a Cauchy identity makes all four leading,
mixed, and third boundary blocks carry the same scalar $\kappa$.
Let $a$ be the boundary dimension, $h$ the dimension of each head flag,
and $\ell$ their pairing rank. The residual odd equations cost at most
$h-\ell$ dimensions on either even side, while the third Schur form has
rank at most $a-2h+\ell$. In the valuation bound these unknown ranks
cancel. This gives the explicit rate $4\delta+(e+3)w-\beta$ on the
stated guarded region, without assuming the heads span the boundary.

The arithmetic envelope is minimized over these proved candidates and
integrated exactly. The rational windows are then retuned jointly with
the spectral response. The final analytic proof uses one fixed contour.
A proposed intermediate patching of bounds from different contours was
withdrawn: the associated positive majorant matrices differ, so pointwise
spectral envelopes cannot simply be mixed. No such patching enters the
final certificate. The inherited shifted-circle and varying-slope proofs
remain the analytic inputs; faster exact logarithm and moment evaluation
make the larger finite replay practical.

Optional ceiling enlargements and forced-zero-head variants produced no
additional gain at the tested windows. Original-matrix surveys also found
several apparent residuals that stayed bounded as the prime increased.
They do not yield an $N^2$ arithmetic saving. These observations guide the
search, but are not exclusion theorems or claims of global optimality.
The quantitative extension remains an ordinary proof with exact finite
certificates; it is not Lean verified and supplies no explicit denominator
threshold $q_0$.

\paragraph{The next integer benchmark.}
With the final windows and analytic envelope held fixed, an additional
arithmetic-rate saving strictly greater than
\[
 \frac{\lambda a_{8,\rm up}}{\CertNextInteger}
 +\mathcal R_{9,\rm up}-\varepsilon
 \;\le\;\CertNextIntegerGap
\]
would suffice for an exponent below $\CertNextInteger$. This is a conditional
optimization benchmark, not an achieved result.

The extra retained-root mechanism has a simple scope-limited tail bound.
For matched unpaid Cauchy candidates $E_j$ and retained-root candidates
$E_j-C_j$, one has $0\le C_j\le t\le\beta/2$. Taking lower envelopes,
including any common competing bounds, therefore changes the profile
by at most $\beta/2$. Beyond $q=192$, its total possible extra-root
contribution is bounded above by
\[
 \int_0^{1/193}\frac{\beta}{2}\,d\beta
 =\frac{1}{148996}<0.000006712.
\]
This compares with an envelope already containing the corresponding
unpaid Cauchy candidates; savings from extending that unpaid family are
separate. At the frozen analytic response the displayed tail bound is
smaller than the gap needed for $\CertNextInteger$. It does not exclude a better
parameter regime or other arithmetic mechanisms. Finding another local
rank saving on intervals of prime ratios remains a more substantial
optimization objective than extending this extra-root cutoff alone.

\section{Numerical proof of the current bound}
\label{sec:current-cert}
For the windows in Theorem~\ref{thm:main} the exact degree is
$N-m_2+1$, so $\lambda=\frac{87095923}{100000000}$. The determinant is evaluated along
multiples of the common denominator of the window parameters.
The arithmetic profile for the new rational windows is integrated afresh on disjoint intervals $1/(j+1)<\beta<1/j$ through $j=77056$. To the resulting upper sum we add the exact large-prime area and the analytic small-prime bound at $1/77057$. The finite-interval form of the cubic small-prime estimate supplies the integral upper bound for $43137\le j\le 65536$. The later prime bands are included individually. On each complete block in the indicated gap, the two floor functions counting physical roots give the additional phase saving proved above. The number of disjoint bands using the local radical calculation is 2, and the number using the selected-class calculation is 1. Their corrections are measured against the compatible profile before integration. The selected integration routines include homogeneous interpolation of the two top multipliers, the rectangular rank comparisons, the opposite-side deep kernels and identification of equal local coefficients. The corresponding local rank arguments appear above; each interval uses the strongest applicable pointwise inequality. Section~\ref{sec:second-order-tail} supplies the second-order correction for the unintegrated prime tail. The third-order correction in Lemma~\ref{lem:cubic-prime-tail} further reduces the same tail. The two exact root-class floors supply the additional rational saving of Corollary~\ref{cor:root-phase-prime-tail}. For this rational window the threshold and determinant response depend only on the contour and the determinant. The sharpened arithmetic rate can therefore be substituted at the same threshold. The supplement recomputes both estimates at the stated parameters.

The resulting rate bounds are
\begin{align*}
 C_{\rm an}&\ge \frac{51621616696872997397}{10000000000000000000}, & C_{\rm ar}&\le \frac{5085415449269262704417}{1000000000000000000000},\\
 \varepsilon&=\frac{76746220418037035283}{1000000000000000000000}.
\end{align*}
In the spectral calculation $u=\frac{511129}{500000}$, $L=\frac{189721}{200000}$ and
$B=\frac{5919716361}{50000000000}$. We take $Y=\frac{398200017}{2000000000}$ and $T_{\max}=\frac{1238596617}{125000000000}$.
The positive interval is divided into 144 panels with two Gauss nodes
on each. The negative interval is divided into 3072 panels.
The circle selects the outer poles. The omitted atoms and the contour are estimated on a common polynomial subspace as in Section~\ref{sec:harmonic-physical}. The atomic interpolation test covers 128 real intervals and has a uniform slack of at least $\frac{72902624265465633037}{2000000000000000000000}$. The outer arc is bounded with the Green denominator on the whole circle. The inner second-derivative estimate uses 16 equal intervals. The outer comparison uses 17 intervals, with a separate tangent bound on each. The quadratic coefficient is $\frac{38095965257043936151}{10000000000000000000000}$ after the middle-arc comparison. The harmonic estimate of Section~\ref{sec:harmonic-physical} gives $\mathcal H_b[\chi_\eta]\le 0.192522218007$ throughout the contour. On the inner arc the Poisson bound is $0.190658945962$. Eight angular moments bound the physical cost. The physical variance is at least $0.000692552972879$. The curvature estimate uses 10 physical power moments integrated against the arcsine mass derivative. These moments come from the closed Stieltjes transform of the physical submeasure. The inner equilibrium measure retains the complete quadratic field. The separate mass integral for the physical penalty saves a further $\frac{2820375434064140531282853}{500000000000000000000000000000000}$ in the response. The last integral uses the convex secant bound for the first moment on each mass cell. The cubic correction for the moving physical endpoint contributes a further $\frac{131557768073018123787089}{100000000000000000000000000000000}$ to the integrated response saving. Accounting for endpoint motion across the slope interval saves $\frac{8446485887275886911758491}{2000000000000000000000000000000000}$ in the response. Each physical moment uses the local support bound and the derivatives with respect to both angular field coefficients. The centered Poisson resolvent and the separate equilibrium moment integrals bound the first 8 terms of the physical cost series. The calculation integrates powers of the endpoint chord on each mass interval. The physical support for these Poisson averages is enclosed on 32 separate mass intervals. The quadratic saving integrates the endpoint bound as both field coefficients vary. The mass-dependent Green and physical slope estimates together save $\frac{177998658674232208178321}{125000000000000000000000000000000}$ in the response. Integrating the physical second derivative on 64 mass intervals saves another $\frac{5660590359283847309}{62500000000000000000000000000000}$ in the response. The exact Green denominator gives the inner exponential slope $A_G\le 0.096352993179$ at $Y_1=\frac{1799}{15000}$. The certificate uses the explicit centered Poisson integral there. The following bounds use outward rounding.
\begin{center}
\begin{tabular}{@{}ll@{}}\toprule
Quantity & Certified bound \\ \midrule
Outer chord margin, after varying the slope & $>0.001055$ \\
Threshold $a_*$ & $\le 7.6876187143654598741988$ \\
Determinant response $\mathcal R(a_*)$ & $\le 0.0005712160662811348734$ \\
Integrated eigenvalue bound & $\le 0.0767458536910915696699$ \\
Number of inverse-mass tangents & 1025 \\ \bottomrule
\end{tabular}
\end{center}
The angular-field bounds and the outer-arc checks described above control both angular regions.
The combined angular refinements lower the fixed-slope response by at least $\frac{733444506556413319059111287}{2000000000000000000000000000000000}$. The complete quadratic field contributes at least $\frac{6166897135375542456625629}{100000000000000000000000000000000}$.
The exact response upper bound used in the quotient equals the upper bound from the fixed-slope certificate minus the rational response saving. The decimals in the table have been rounded outward for display.
Applying Theorem~\ref{thm:transfer} gives
\[
 \mu(\zeta(5))\le\frac{\lambda a_{*,\mathrm{up}}}
 {\varepsilon-\mathcal R_{\mathrm{up}}}
 =\frac{4463734983998220580418722476616000}{50783336234503933606439686370429} < 87.897631683469558 .
\]
\begin{center}
\begin{tabular}{@{}lc@{}}\toprule
Successive certified estimates in this project & Strict upper bound \\ \midrule
Preceding six-hour calculation & $87.925687270635509$ \\
Support-crossing circle with the previous arithmetic & $87.925292030199392$ \\
Curvature estimate with refined prime bands & $87.925159833785116$ \\
Selecting contour and a separate residue bound & $87.924297304244208$ \\
Selected local completions and a common contour subspace & $87.918141483721804$ \\
Joint count of restricted rows and columns & $87.908950341025563$ \\
Strongest registered certificate at the three-hour mark & $87.898063774316436$ \\
Current certificate & $87.897631683469558$ \\ \bottomrule
\end{tabular}
\end{center}
At this threshold and degree ratio an extra arithmetic saving exceeding
$0.000785943648214$ would suffice for $\mu(\zeta(5))<87$.

\paragraph{Reproducibility.}
The supplement contains the rational window parameters, contour recipe,
arithmetic inputs, and both spectral certificates. From its top directory run
\begin{center}\texttt{python3 replay.py}\end{center}
The replay uses Python's standard library and checks the mathematical outputs
against the saved fractions. The finite calculation is interpreted using the
ordinary proofs in this article and the qualitative determinant construction
cited in Section~\ref{sec:setup}. The quantitative bound has no Lean
certification.
To regenerate the selected prime bands from their defining equations before
assembling the arithmetic bound, run \texttt{python3 replay.py --full-bands}.
The two worked moment systems can also be checked separately with
\texttt{python3 experiments/check\_worked\_rank.py}.

\paragraph{Independent finite checks.}
Of the 1493 selected interval files, 277 had also been regenerated individually from the rational chamber equations when this article was assembled. The supplement identifies these files by their SHA-256 digests.
For the reference factorial window the independent rational rank computation covered $q=12$ (189 matrices), $q=20$ (181 matrices). It constructed the moment entries by polynomial division and harmonic residues, then compared exact matrix ranks with both the modular lower bounds and the structural upper bounds. Each interval certificate also reproduced its saved affine integrals and witness hashes.
All 370 of these rank intervals were singletons.
At a reference contour, independent 65-digit integration in the physical radius and angle checked all 10 moments from the closed Stieltjes transform. Each rational lower bound was within $10^{-24}$ of the numerical integral. The proof uses the rational enclosures.
For 77 finite parameter choices, at primes 61, 79, 101, direct Smith calculations at both $X=0$ and $X=1$ met the two-layer lower bound from the complementary-kernel argument.
The homogeneous interpolation estimate was also checked against 70 of these finite pencils, using the two previously computed Smith specializations in each case.
\paragraph{Other determinant families tested.} Requiring opposite weighted evaluations at neighboring poles cancels pairs of residues and lowers the slope rank. In finite examples the extra denominators of the constrained bases offset this benefit. A positive modification of the Fourier weight also suppresses the leading mode, but the normalized determinant lost its decay in the finite tests. Neither experiment supplies an asymptotic exponent; the numerical proof above uses the unconstrained two-window determinant.

\begin{thebibliography}{9}\raggedright
\bibitem{ng} David Niedbala Giraudin, \emph{A new upper bound for the irrationality
exponent of $\zeta(5)$}, arXiv:2611.01708v1, 1 November 2026.
\url{https://arxiv.org/abs/2611.01708}.
\bibitem{independent} Zeta Irrationality Project, \emph{Independent factorial-window
proof}, source snapshot \texttt{zeta5\_journal.tex}; included unchanged in the
accompanying replay package. Its exact provenance is recorded in the manifest.
\bibitem{previous} Zeta Irrationality Project, \emph{Optimizing an irrationality-exponent
bound for $\zeta(5)$: factorial-window determinants, local rank refinements, and exact
certificates}, 17 November 2026; preceding quantitative report and replay package.
\bibitem{saff2010} E. B. Saff, \emph{Logarithmic potential theory with applications to approximation theory}, Surveys in Approximation Theory (2010). \url{https://arxiv.org/abs/1010.3760}.
\end{thebibliography}
\end{document}
